%%%
%%% Y
%%%
\section*{Y}\addcontentsline{toc}{section}{Y}\addcontentsline{loh}{figure}{\#\#\#\#\#\#\#\# Y}

%%%%%%%%%% 丫 %%%%%%%%%%
\subsection*{丫}\addcontentsline{loh}{figure}{丫 \dpy{ya1}}

\begin{EntryWithPhonetic}{丫}{ya1}{3}{⼁}
  \definition{s.}{bifurcação; forquilha | Dialeto: menina}
\end{EntryWithPhonetic}

\begin{EntryWithPhonetic}{丫头}{ya1tou5}{3,5}{⼁,⼤}[HSK 7-9]
  \definition[个]{s.}{garota; menina; o penteado com coque usado por crianças antigamente, com os cabelos espetados em ambos os lados da cabeça; usado metaforicamente para se referir a meninas | empregada doméstica; criada}
  \synonymref{姑娘}{gu1niang5}
  \synonymref{女儿}{nv3'er2}
  \synonymref{女孩}{nv3hai2}
\end{EntryWithPhonetic}

%%%%%%%%%% 压 %%%%%%%%%%
\subsection*{压}\addcontentsline{loh}{figure}{压 \dpy{ya1}}

\begin{EntryWithPhonetic}{压}{ya1}{6}{⼚}[HSK 3]
  \definition{v.}{pressionar; empurrar para baixo; segurar; pesar | acalmar emoções agitadas ou situações ruins; tranquilizar | intimidar; reprimir; exercer pressão sobre; usar poder, posição ou padrões morais para coagir ou restringir as pessoas, impedindo"-as de se expressar, decidir ou se desenvolver livremente | aproximar"-se; estar chegando perto | arquivar; deixar de lado | pressionar; metáfora para uma grande carga emocional e psicológica | superar; ultrapassar; voz, capacidade e presença mais fortes do que os outros | apostar em um determinado resultado ao jogar | pressionar; força na superfície de contato do objeto}
  \seeref{ya4}
\end{EntryWithPhonetic}

\begin{EntryWithPhonetic}{压倒}{ya1dao3}{6,10}{⼚,⼈}[HSK 7-9]
  \definition{v.}{dominar; subjugar; prevalecer sobre}
  \antonymref{反弹}{fan3tan2}
\end{EntryWithPhonetic}

\begin{EntryWithPhonetic}{压力}{ya1li4}{6,2}{⼚,⼒}[HSK 3]
  \definition[份,个]{s.}{pressão; força atuando perpendicularmente à superfície de um objeto | pressão; força esmagadora; metáfora para a força que coage e intimida as pessoas (principalmente nos aspectos espirituais e psicológicos) | tensão; fardo; os encargos econômicos, psicológicos e espirituais impostos pelo mundo exterior}
\end{EntryWithPhonetic}

\begin{EntryWithPhonetic}{压迫}{ya1po4}{6,8}{⼚,⾡}[HSK 6]
  \definition{v.}{oprimir; reprimir; confiar no poder para suprimir e forçar | contrair; uma força externa comprime uma parte de um organismo}
\end{EntryWithPhonetic}

\begin{EntryWithPhonetic}{压岁钱}{ya1sui4qian2}{6,6,10}{⼚,⼭,⾦}
  \definition{s.}{dinheiro da sorte | dinheiro dado às crianças como presente no Ano Novo Chinês}
\end{EntryWithPhonetic}

\begin{EntryWithPhonetic}{压碎}{ya1sui4}{6,13}{⼚,⽯}
  \definition{v.}{esmagar em pedaços}
\end{EntryWithPhonetic}

\begin{EntryWithPhonetic}{压缩}{ya1suo1}{6,14}{⼚,⽷}[HSK 7-9]
  \definition{v.}{comprimir; aplicar pressão; reduzir o volume | reduzir; condensar; cortar; reduzir (pessoal, financiamento, duração, etc.)}
  \synonymref{紧缩}{jin3suo1}
  \synonymref{浓缩}{nong2suo1}
  \synonymref{收缩}{shou1suo1}
  \antonymref{发展}{fa1zhan3}
  \antonymref{拉开}{la1/kai5}
  \antonymref{延伸}{yan2shen1}
\end{EntryWithPhonetic}

\begin{EntryWithPhonetic}{压抑}{ya1yi4}{6,7}{⼚,⼿}[HSK 7-9]
  \definition{adj.}{reprimido; deprimido; sentir"-se desconfortável porque as próprias emoções e sentimentos não podem ser liberados}
  \definition{v.}{deprimir; restringir; inibir; conter; reprimir as emoções e a força, impedindo que sejam plenamente expressas ou utilizadas}
  \seealsoref{压制}{ya1zhi4}
  \synonymref{克制}{ke4zhi4}
  \synonymref{压迫}{ya1po4}
  \synonymref{压制}{ya1zhi4}
  \synonymref{抑制}{yi4zhi4}
  \antonymref{发挥}{fa1hui1}
  \antonymref{轻松}{qing1song1}
\end{EntryWithPhonetic}

\begin{EntryWithPhonetic}{压韵}{ya1yun4}{6,13}{⼚,⾳}
  \variantof{押韵}
\end{EntryWithPhonetic}

\begin{EntryWithPhonetic}{压制}{ya1zhi4}{6,8}{⼚,⼑}[HSK 7-9]
  \definition{v.}{suprimir; sufocar; inibir | Militar: neutralizar (fogo inimigo por bombardeio maciço, etc.) | pressionar; fazer algo pressionando; fabricar utilizando métodos de pressão}
  \seealsoref{压迫}{ya1po4}
  \seealsoref{压抑}{ya1yi4}
  \synonymref{逼迫}{bi1po4}
  \synonymref{克制}{ke4zhi4}
  \synonymref{强迫}{qiang3po4}
  \synonymref{压迫}{ya1po4}
  \synonymref{压抑}{ya1yi4}
  \synonymref{要挟}{yao1xie2}
  \synonymref{抑制}{yi4zhi4}
  \antonymref{鼓励}{gu3li4}
  \antonymref{激励}{ji1li4}
  \antonymref{推广}{tui1guang3}
\end{EntryWithPhonetic}

%%%%%%%%%% 呀 %%%%%%%%%%
\subsection*{呀}\addcontentsline{loh}{figure}{呀 \dpy{ya1}}

\begin{EntryWithPhonetic}{呀}{ya1}{7}{⼝}
  \definition{interj.}{``Ah!''; ``Oh!''; interjeição que expressa surpresa}
  \definition{s.}{Onomatopéia: rangido; guincho (o som de uma porta)}
  \seeref{ya5}
\end{EntryWithPhonetic}

%%%%%%%%%% 押 %%%%%%%%%%
\subsection*{押}\addcontentsline{loh}{figure}{押 \dpy{ya1}}

\begin{EntryWithPhonetic}{押}{ya1}{8}{⼿}[HSK 7-9]
  \definition*{s.}{Sobrenome: Ya}
  \definition{s.}{assinatura; marca em vez de assinatura; nome assinado ou símbolo desenhado}
  \definition{v.}{dar como garantia; hipotecar; penhorar | deter; levar sob custódia | escoltar | assinar (um documento, contrato, etc.); colocar sua assinatura (ou marcar no lugar da assinatura)}
\end{EntryWithPhonetic}

\begin{EntryWithPhonetic}{押后}{ya1hou4}{8,6}{⼿,⼝}
  \definition{v.}{encerrar | adiar}
\end{EntryWithPhonetic}

\begin{EntryWithPhonetic}{押金}{ya1jin1}{8,8}{⼿,⾦}[HSK 5]
  \definition[笔,份,些]{s.}{caução; sinal; depósito; dinheiro como garantia}
\end{EntryWithPhonetic}

\begin{EntryWithPhonetic}{押送}{ya1song4}{8,9}{⼿,⾡}
  \definition{v.}{enviar sob escolta | transportar um detido}
\end{EntryWithPhonetic}

\begin{EntryWithPhonetic}{押运}{ya1yun4}{8,7}{⼿,⾡}
  \definition{v.}{escoltar sob guarda | escoltar (bens ou fundos)}
\end{EntryWithPhonetic}

\begin{EntryWithPhonetic}{押韵}{ya1yun4}{8,13}{⼿,⾳}
  \definition{v.}{rimar}
\end{EntryWithPhonetic}

\begin{EntryWithPhonetic}{押注}{ya1zhu4}{8,8}{⼿,⽔}
  \definition{v.}{apostar}
\end{EntryWithPhonetic}

\begin{EntryWithPhonetic}{押租}{ya1zu1}{8,10}{⼿,⽲}
  \definition{s.}{depósito de aluguel}
\end{EntryWithPhonetic}

%%%%%%%%%% 哑 %%%%%%%%%%
\subsection*{哑}\addcontentsline{loh}{figure}{哑 \dpy{ya1}}

\begin{EntryWithPhonetic}{哑}{ya1}{9}{⼝}
  \definition{interj.}{interjeição que expressa surpresa, o mesmo que 呀}
  \seeref{ya3}
  \seealsoref{呀}{ya1}
\end{EntryWithPhonetic}

%%%%%%%%%% 鸦 %%%%%%%%%%
\subsection*{鸦}\addcontentsline{loh}{figure}{鸦 \dpy{ya1}}

\begin{EntryWithPhonetic}{鸦}{ya1}{9}{⿃}
  \definition[只]{s.}{corvo}
\end{EntryWithPhonetic}

\begin{EntryWithPhonetic}{鸦雀无声}{ya1que4-wu2sheng1}{9,11,4,7}{⿃,⾫,⽆,⼠}[HSK 7-9]
  \definition{expr.}{``Nem mesmo um corvo ou pardal pode ser ouvido.''; silêncio absoluto; nenhum som podia ser ouvido; descreve um estado de extremo silêncio; também pode ser uma metáfora para uma situação em que ninguém se atreve a falar ou todos permanecem em silêncio}
  \synonymref{无话可说}{wu2hua4-ke3shuo1}
  \antonymref{沸沸扬扬}{fei4fei4yang2yang2}
  \antonymref{七嘴八舌}{qi1zui3-ba1she2}
\end{EntryWithPhonetic}

%%%%%%%%%% 鸭 %%%%%%%%%%
\subsection*{鸭}\addcontentsline{loh}{figure}{鸭 \dpy{ya1}}

\begin{EntryWithPhonetic}{鸭}{ya1}{10}{⿃}
  \definition[只]{s.}{pato | Gíria: prostituto}
\end{EntryWithPhonetic}

\begin{EntryWithPhonetic}{鸭子}{ya1zi5}{10,3}{⿃,⼦}[HSK 5]
  \definition[只,群]{s.}{pato | Gíria: prostituto}
\end{EntryWithPhonetic}

%%%%%%%%%% 雅 %%%%%%%%%%
\subsection*{雅}\addcontentsline{loh}{figure}{雅 \dpy{ya1}}

\begin{EntryWithPhonetic}{雅}{ya1}{12}{⾫}
  \variantof{鸦}
\end{EntryWithPhonetic}

%%%%%%%%%% 牙 %%%%%%%%%%
\subsection*{牙}\addcontentsline{loh}{figure}{牙 \dpy{ya2}}

\begin{EntryWithPhonetic}{牙}{ya2}{4}{⽛}[HSK 4][Kangxi 92]
  \definition*{s.}{Sobrenome: Ya}
  \definition[颗,排]{s.}{dente | marfim | algo semelhante a um dente}
\end{EntryWithPhonetic}

\begin{EntryWithPhonetic}{牙齿}{ya2chi3}{4,8}{⽛,⿒}[HSK 7-9]
  \definition[颗,排,口]{adv.}{dente; o termo geral para dentes}
  \definition[颗]{s.}{dente}
\end{EntryWithPhonetic}

\begin{EntryWithPhonetic}{牙膏}{ya2gao1}{4,14}{⽛,⾁}[HSK 7-9]
  \definition[支,管,盒]{s.}{pasta de dente; a pasta de dentes usada para escovar os dentes é feita de glicerina, pó dental, cola branca em pó, água, sacarina, amido, etc., e é embalada em um recipiente flexível de metal ou plástico}
\end{EntryWithPhonetic}

\begin{EntryWithPhonetic}{牙行}{ya2hang2}{4,6}{⽛,⾏}
  \definition{s.}{corretor | \emph{broker}}
\end{EntryWithPhonetic}

\begin{EntryWithPhonetic}{牙刷}{ya2shua1}{4,8}{⽛,⼑}[HSK 4]
  \definition[把,个,支]{s.}{escova de dentes}
\end{EntryWithPhonetic}

\begin{EntryWithPhonetic}{牙线}{ya2xian4}{4,8}{⽛,⽷}
  \definition[条]{s.}{fio dental}
\end{EntryWithPhonetic}

\begin{EntryWithPhonetic}{牙医}{ya2yi1}{4,7}{⽛,⼖}
  \definition{s.}{dentista}
\end{EntryWithPhonetic}

%%%%%%%%%% 芽 %%%%%%%%%%
\subsection*{芽}\addcontentsline{loh}{figure}{芽 \dpy{ya2}}

\begin{EntryWithPhonetic}{芽}{ya2}{7}{⾋}[HSK 7-9]
  \definition[颗,枚,个]{s.}{broto; rebento; germe; botão | algo que se parece com um broto}
\end{EntryWithPhonetic}

%%%%%%%%%% 崖 %%%%%%%%%%
\subsection*{崖}\addcontentsline{loh}{figure}{崖 \dpy{ya2}}

\begin{EntryWithPhonetic}{崖}{ya2}{11}{⼭}
  \definition{s.}{precipício | penhasco}
\end{EntryWithPhonetic}

%%%%%%%%%% 哑 %%%%%%%%%%
\subsection*{哑}\addcontentsline{loh}{figure}{哑 \dpy{ya3}}

\begin{EntryWithPhonetic}{哑}{ya3}{9}{⼝}[HSK 7-9]
  \definition{adj.}{mudo | rouco; rouco | (um projétil, bomba, etc.) não explodir | mudo; sem som}
  \definition{s.}{mudo; pessoas que são incapazes de emitir sons ou falar}
  \seeref{ya1}
\end{EntryWithPhonetic}

%%%%%%%%%% 雅 %%%%%%%%%%
\subsection*{雅}\addcontentsline{loh}{figure}{雅 \dpy{ya3}}

\begin{EntryWithPhonetic}{雅}{ya3}{12}{⾫}
  \definition*{s.}{Sobrenome: Ya}
  \definition{adj.}{refinado; elegante; nobre; não vulgar | adequado; correto; padrão; em conformidade}
  \definition{adv.}{frequentemente; geralmente  | muito; de fato; extremamente}
  \definition{s.}{honoríficos, usados para expressar sentimentos e ações em relação à outra parte | uma seção do Livro das Canções composta por hinos dinásticos; um tipo de poema do Livro das Canções | amizade; conhecimento}
  \antonymref{俗}{su2}
\end{EntryWithPhonetic}

%%%%%%%%%% 亚 %%%%%%%%%%
\subsection*{亚}\addcontentsline{loh}{figure}{亚 \dpy{ya4}}

\begin{EntryWithPhonetic}{亚}{ya4}{6}{⼆}
  \definition*{s.}{Ásia, abreviação de 亚洲 | Sobrenome: Ya}
  \definition{adj.}{inferior | abaixo do padrão | (química) de menor valência atômica}
  \definition{pref.}{sub-}
  \seealsoref{亚洲}{ya4zhou1}
\end{EntryWithPhonetic}

\begin{EntryWithPhonetic}{亚军}{ya4jun1}{6,6}{⼆,⼍}[HSK 5]
  \definition[个]{s.}{segundo lugar; vice"-campeão; medalhista de prata}
\end{EntryWithPhonetic}

\begin{EntryWithPhonetic}{亚热带}{ya4re4dai4}{6,10,9}{⼆,⽕,⼱}
  \definition{s.}{zona ou clima subtropical; subtropical; semitropical}
\end{EntryWithPhonetic}

\begin{EntryWithPhonetic}{亚细亚洲}{ya4xi4ya4zhou1}{6,8,6,9}{⼆,⽷,⼆,⽔}
  \definition*{s.}{Ásia}
\end{EntryWithPhonetic}

\begin{EntryWithPhonetic}{亚运会}{ya4yun4hui4}{6,7,6}{⼆,⾡,⼈}[HSK 4]
  \definition*{s.}{Jogos Asiáticos}
\end{EntryWithPhonetic}

\begin{EntryWithPhonetic}{亚洲}{ya4zhou1}{6,9}{⼆,⽔}
  \definition*{s.}{Ásia, abreviação de 亚细亚洲}
  \seealsoref{亚细亚洲}{ya4xi4ya4zhou1}
\end{EntryWithPhonetic}

\begin{EntryWithPhonetic}{亚洲人}{ya4zhou1ren2}{6,9,2}{⼆,⽔,⼈}
  \definition{s.}{asiático | pessoa ou povo da Ásia}
\end{EntryWithPhonetic}

%%%%%%%%%% 压 %%%%%%%%%%
\subsection*{压}\addcontentsline{loh}{figure}{压 \dpy{ya4}}

\begin{EntryWithPhonetic}{压}{ya4}{6}{⼚}
  \definition{adv.}{fundamentalmente; nunca (usado principalmente em frases negativas)}
  \seeref{ya1}
  \seealsoref{压根儿}{ya4gen1r5}
\end{EntryWithPhonetic}

\begin{EntryWithPhonetic}{压根儿}{ya4gen1r5}{6,10,2}{⼚,⽊,⼉}
  \definition{adv.}{(geralmente no negativo) nunca; fundamentalmente}
\end{EntryWithPhonetic}

%%%%%%%%%% 呀 %%%%%%%%%%
\subsection*{呀}\addcontentsline{loh}{figure}{呀 \dpy{ya5}}

\begin{EntryWithPhonetic}{呀}{ya5}{7}{⼝}[HSK 4]
  \definition{part.}{usado no lugar de 啊 quando a palavra anterior termina com o som a, e, i, o ou ü}
  \seeref{ya1}
  \seealsoref{啊}{a5}
\end{EntryWithPhonetic}

%%%%%%%%%% 咽 %%%%%%%%%%
\subsection*{咽}\addcontentsline{loh}{figure}{咽 \dpy{yan1}}

\begin{EntryWithPhonetic}{咽}{yan1}{9}{⼝}
  \definition{s.}{deglutição; garganta; faringe}
  \seeref{yan4}
  \seeref{ye4}
\end{EntryWithPhonetic}

\begin{EntryWithPhonetic}{咽喉}{yan1hou2}{9,12}{⼝,⼝}[HSK 7-9]
  \definition{pref.}{bronco-}
  \definition{s.}{Anatomia: faringe e laringe; garganta | Figurativo: passagem estratégica (ou vital); entroncamento (ou ligação) fundamental; passagem de transporte perigosa}
  \synonymref{喉咙}{hou2long2}
  \synonymref{嗓子}{sang3zi5}
\end{EntryWithPhonetic}

%%%%%%%%%% 殷 %%%%%%%%%%
\subsection*{殷}\addcontentsline{loh}{figure}{殷 \dpy{yan1}}

\begin{EntryWithPhonetic}{殷}{yan1}{10}{⽎}
  \definition{adj.}{vermelho"-escuro}
  \seeref{yin1}
  \seeref{yin3}
\end{EntryWithPhonetic}

%%%%%%%%%% 烟 %%%%%%%%%%
\subsection*{烟}\addcontentsline{loh}{figure}{烟 \dpy{yan1}}

\begin{EntryWithPhonetic}{烟}{yan1}{10}{⽕}[HSK 3]
  \definition[股,支,根,盒,包]{s.}{fumaça; gás produzido pela combustão de materiais, misturado com pequenas partículas não completamente queimadas | névoa; neblina | tabaco; planta de tabaco | fumo; cigarro; termo geral para cigarros, charutos, etc. | ópio | fuligem; fumaça de carvão}
  \definition{v.}{ficar irritado com a fumaça (os olhos lacrimejam ou não conseguem abrir)}
\end{EntryWithPhonetic}

\begin{EntryWithPhonetic}{烟草}{yan1cao3}{10,9}{⽕,⾋}
  \definition{s.}{tabaco}
\end{EntryWithPhonetic}

\begin{EntryWithPhonetic}{烟囱}{yan1cong1}{10,7}{⽕,⼞}[HSK 7-9]
  \definition[根]{s.}{chaminé; uma estrutura que proporciona ventilação para gases de combustão quentes ou fumaça provenientes de uma caldeira, fogão, forno ou lareira}
  \synonymref{突}{tu1}
\end{EntryWithPhonetic}

\begin{EntryWithPhonetic}{烟花}{yan1hua1}{10,7}{⽕,⾋}[HSK 6]
  \definition[场,朵]{s.}{fogos de artifício; uma coisa que emite faíscas de várias cores quando exposta à observação | prostituta; antigamente, referia"-se a algo relacionado à prostituição}
\end{EntryWithPhonetic}

\begin{EntryWithPhonetic}{烟火}{yan1huo3}{10,4}{⽕,⽕}[HSK 7-9]
  \definition{s.}{fumaça e fogo | comida cozida | descendentes; incenso oferecido aos ancestrais; metaforicamente referindo"-se aos descendentes | guerra; fumaça e chamas da guerra; uma metáfora para a guerra}
  \seeref{yan1huo5}
\end{EntryWithPhonetic}

\begin{EntryWithPhonetic}{烟火}{yan1huo5}{10,4}{⽕,⽕}
  \definition{s.}{fogos de artifício}
  \synonymref{烟花}{yan1hua1}
  \synonymref{焰火}{yan4huo3}
\end{EntryWithPhonetic}

\begin{EntryWithPhonetic}{烟头}{yan1tou2}{10,5}{⽕,⼤}
  \definition[根]{s.}{bituca de cigarro}
\end{EntryWithPhonetic}

\begin{EntryWithPhonetic}{烟叶}{yan1ye4}{10,5}{⽕,⼝}
  \definition{s.}{folha de tabaco}
\end{EntryWithPhonetic}

\begin{EntryWithPhonetic}{烟雨}{yan1yu3}{10,8}{⽕,⾬}
  \definition{s.}{chuvisco | garoa}
\end{EntryWithPhonetic}

%%%%%%%%%% 淹 %%%%%%%%%%
\subsection*{淹}\addcontentsline{loh}{figure}{淹 \dpy{yan1}}

\begin{EntryWithPhonetic}{淹}{yan1}{11}{⽔}[HSK 7-9]
  \definition{adj.}{amplo | atrasado; alongado}
  \definition{v.}{inundar; submergir; alagar | sentir formigamento devido ao suor; ter a pele irritada ou com coceira devido ao suor, etc.}
\end{EntryWithPhonetic}

%%%%%%%%%% 燕 %%%%%%%%%%
\subsection*{燕}\addcontentsline{loh}{figure}{燕 \dpy{yan1}}

\begin{EntryWithPhonetic}{燕}{yan1}{16}{⽕}
  \definition*{s.}{Hebei do Norte | Yan, um estado vassalo de Zhou nas atuais províncias de Hebei e Liaoning | Sobrenome: Yan}
  \seeref{yan4}
\end{EntryWithPhonetic}

%%%%%%%%%% 延 %%%%%%%%%%
\subsection*{延}\addcontentsline{loh}{figure}{延 \dpy{yan2}}

\begin{EntryWithPhonetic}{延}{yan2}{6}{⼵}[HSK 7-9]
  \definition*{s.}{Sobrenome: Yan}
  \definition{v.}{prolongar; estender; alongar | adiar; atrasar | envolver (um professor, conselheiro, etc.); enviar para; convidar}
\end{EntryWithPhonetic}

\begin{EntryWithPhonetic}{延长}{yan2chang2}{6,4}{⼵,⾧}[HSK 4]
  \definition{v.}{estender; prolongar; alongar; aumentar o tempo, a distância ou a duração de algo específico}
\end{EntryWithPhonetic}

\begin{EntryWithPhonetic}{延缓}{yan2huan3}{6,12}{⼵,⽶}[HSK 7-9]
  \definition{v.}{atrasar; adiar; postergar; protelar}
  \synonymref{减速}{jian3/su4}
  \antonymref{加速}{jia1su4}
  \antonymref{提前}{ti2qian2}
\end{EntryWithPhonetic}

\begin{EntryWithPhonetic}{延期}{yan2/qi1}{6,12}{⼵,⽉}[HSK 4]
  \definition{v.+compl.}{atrasar; adiar; postergar}
\end{EntryWithPhonetic}

\begin{EntryWithPhonetic}{延伸}{yan2shen1}{6,7}{⼵,⼈}[HSK 5]
  \definition{v.}{estender; esticar; alongar; estender"-se}
\end{EntryWithPhonetic}

\begin{EntryWithPhonetic}{延误}{yan2wu4}{6,9}{⼵,⾔}[HSK 7-9]
  \definition{v.}{falhar; incorrer em prejuízo devido a atraso; atrasar; adiar}
  \synonymref{耽搁}{dan1ge5}
  \synonymref{耽误}{dan1wu5}
  \synonymref{拖延}{tuo1yan2}
\end{EntryWithPhonetic}

\begin{EntryWithPhonetic}{延续}{yan2xu4}{6,11}{⼵,⽷}[HSK 4]
  \definition{v.}{durar; continuar; prosseguir; continuar como antes; prolongar}
\end{EntryWithPhonetic}

%%%%%%%%%% 严 %%%%%%%%%%
\subsection*{严}\addcontentsline{loh}{figure}{严 \dpy{yan2}}

\begin{EntryWithPhonetic}{严}{yan2}{7}{⼀}[HSK 4]
  \definition*{s.}{Sobrenome: Yan}
  \definition{adj.}{apertado; próximo | rigoroso; severo; duro; áspero; rigoroso; austero | severo; extremo; difícil}
  \definition{s.}{pai; refere"-se ao pai}
  \antonymref{宽}{kuan1}
\end{EntryWithPhonetic}

\begin{EntryWithPhonetic}{严格}{yan2ge2}{7,10}{⼀,⽊}[HSK 4]
  \definition{adj.}{rígido; estrito; rigoroso; muito consciente e meticuloso na implementação de sistemas e no domínio de padrões}
  \definition{v.}{tornar (sistemas, provisões, etc.) rigorosos}
  \synonymref{严谨}{yan2jin3}
  \synonymref{严峻}{yan2jun4}
  \synonymref{严厉}{yan2li4}
  \synonymref{严肃}{yan2su4}
  \synonymref{用心}{yong4 xin1}
  \antonymref{放松}{fang4song1}
  \antonymref{宽容}{kuan1rong2}
  \antonymref{马虎}{ma3hu5}
  \antonymref{松弛}{song1chi2}
\end{EntryWithPhonetic}

\begin{EntryWithPhonetic}{严谨}{yan2jin3}{7,13}{⼀,⾔}[HSK 7-9]
  \definition{adj.}{estrito; rigoroso; cuidadoso e preciso; rigoroso e cauteloso | compacto; bem"-feito; completo; sem omissões}
\end{EntryWithPhonetic}

\begin{EntryWithPhonetic}{严禁}{yan2jin4}{7,13}{⼀,⽰}[HSK 7-9]
  \definition{v.}{proibir estritamente}
  \synonymref{不准}{bu4 zhun3}
  \synonymref{禁止}{jin4zhi3}
\end{EntryWithPhonetic}

\begin{EntryWithPhonetic}{严峻}{yan2jun4}{7,10}{⼀,⼭}[HSK 7-9]
  \definition{adj.}{severo; rigoroso; implacável; carrancudo | severo; sério; difícil; crítico; desafiador}
  \seealsoref{严重}{yan2zhong4}
  \synonymref{恶化}{e4hua4}
  \synonymref{严格}{yan2ge2}
  \synonymref{严厉}{yan2li4}
  \synonymref{严肃}{yan2su4}
  \synonymref{严重}{yan2zhong4}
  \antonymref{慈祥}{ci2xiang2}
  \antonymref{和气}{he2qi5}
\end{EntryWithPhonetic}

\begin{EntryWithPhonetic}{严厉}{yan2li4}{7,5}{⼀,⼚}[HSK 5]
  \definition{adj.}{severo; rigoroso; as palavras e atitudes de crítica ou punição são muito sérias e severas}
  \synonymref{强势}{qiang2shi4}
  \synonymref{严格}{yan2ge2}
  \synonymref{严峻}{yan2jun4}
  \synonymref{严肃}{yan2su4}
  \antonymref{慈祥}{ci2xiang2}
  \antonymref{和蔼}{he2'ai3}
  \antonymref{宽容}{kuan1rong2}
  \antonymref{平和}{ping2he2}
  \antonymref{温和}{wen1he2}
\end{EntryWithPhonetic}

\begin{EntryWithPhonetic}{严密}{yan2mi4}{7,11}{⼀,⼧}[HSK 7-9]
  \definition{adj.}{apertado; fechado; tudo está intimamente conectado, sem lacunas | bem"-unido; unido; completo; sem omissões}
  \definition{v.}{apertar}
  \synonymref{紧密}{jin3mi4}
  \synonymref{精密}{jing1mi4}
  \synonymref{精细}{jing1xi4}
  \synonymref{严谨}{yan2jin3}
  \antonymref{马虎}{ma3hu5}
  \antonymref{疏忽}{shu1hu5}
  \antonymref{松弛}{song1chi2}
  \antonymref{随便}{sui2/bian4}
\end{EntryWithPhonetic}

\begin{EntryWithPhonetic}{严肃}{yan2su4}{7,8}{⼀,⾀}[HSK 5]
  \definition{adj.}{sério; solene; sincero; (expressão, atmosfera, etc.) faz as pessoas se sentirem admiradas e desconfortáveis | sóbrio; grave; sério; sincero}
  \definition{v.}{aplicar rigorosamente; fazer algo sério}
  \synonymref{古板}{gu3ban3}
  \synonymref{稳重}{wen3zhong4}
  \synonymref{严格}{yan2ge2}
  \synonymref{严峻}{yan2jun4}
  \synonymref{严厉}{yan2li4}
  \antonymref{好笑}{hao3xiao4}
  \antonymref{胡闹}{hu2nao4}
  \antonymref{活泼}{huo2po5}
  \antonymref{轻松}{qing1song1}
  \antonymref{随便}{sui2/bian4}
  \antonymref{随和}{sui2he5}
  \antonymref{戏谑}{xi4xue4}
\end{EntryWithPhonetic}

\begin{EntryWithPhonetic}{严重}{yan2zhong4}{7,9}{⼀,⾥}[HSK 4]
  \definition{adj.}{sério; grave; crítico; severo}
  \synonymref{惨重}{can3zhong4}
  \synonymref{恶化}{e4hua4}
  \synonymref{首要}{shou3yao4}
  \synonymref{危机}{wei1ji1}
  \synonymref{危急}{wei1ji2}
  \synonymref{严峻}{yan2jun4}
  \synonymref{要紧}{yao4jin3}
  \synonymref{重要}{zhong4yao4}
  \synonymref{主要}{zhu3yao4}
  \antonymref{平常}{ping2chang2}
  \antonymref{轻微}{qing1wei1}
  \antonymref{稍微}{shao1wei1}
\end{EntryWithPhonetic}

\begin{EntryWithPhonetic}{严重打伤}{yan2zhong4 da3 shang1}{7,9,5,6}{⼀,⾥,⼿,⼈}
  \definition{s.}{gravemente ferido}
\end{EntryWithPhonetic}

\begin{EntryWithPhonetic}{严重地}{yan2zhong4 di4}{7,9,6}{⼀,⾥,⼟}
  \definition{adv.}{seriamente | gravemente}
\end{EntryWithPhonetic}

\begin{EntryWithPhonetic}{严重关切}{yan2zhong4guan1qie4}{7,9,6,4}{⼀,⾥,⼋,⼑}
  \definition{s.}{preocupação séria}
\end{EntryWithPhonetic}

\begin{EntryWithPhonetic}{严重后果}{yan2zhong4hou4guo3}{7,9,6,8}{⼀,⾥,⼝,⽊}
  \definition{s.}{consequências sérias | repercursões graves}
\end{EntryWithPhonetic}

\begin{EntryWithPhonetic}{严重破坏}{yan2zhong4 po4huai4}{7,9,10,7}{⼀,⾥,⽯,⼟}
  \definition{s.}{destruição grave}
\end{EntryWithPhonetic}

\begin{EntryWithPhonetic}{严重伤害}{yan2zhong4 shang1hai4}{7,9,6,10}{⼀,⾥,⼈,⼧}
  \definition{s.}{ferimento grave; lesão grave}
\end{EntryWithPhonetic}

\begin{EntryWithPhonetic}{严重危害}{yan2zhong4 wei1hai4}{7,9,6,10}{⼀,⾥,⼙,⼧}
  \definition{s.}{perigo crítico | dano grave}
\end{EntryWithPhonetic}

\begin{EntryWithPhonetic}{严重问题}{yan2zhong4 wen4ti2}{7,9,6,15}{⼀,⾥,⾨,⾴}
  \definition{s.}{problema sério}
\end{EntryWithPhonetic}

\begin{EntryWithPhonetic}{严重性}{yan2zhong4xing4}{7,9,8}{⼀,⾥,⼼}
  \definition{s.}{seriedade | gravidade}
\end{EntryWithPhonetic}

%%%%%%%%%% 言 %%%%%%%%%%
\subsection*{言}\addcontentsline{loh}{figure}{言 \dpy{yan2}}

\begin{EntryWithPhonetic}{言}{yan2}{7}{⾔}[Kangxi 149]
  \definition*{s.}{Sobrenome: Yan}
  \definition{s.}{palavra; discurso; o que foi dito | palavra; caracter; uma frase ou palavra chinesa}
  \definition{v.}{dizer; falar}
\end{EntryWithPhonetic}

\begin{EntryWithPhonetic}{言辞}{yan2ci2}{7,13}{⾔,⾟}[HSK 7-9]
  \definition{s.}{as palavras de alguém; o que alguém diz; a escolha de palavras}
  \synonymref{推辞}{tui1ci2}
  \synonymref{言论}{yan2lun4}
\end{EntryWithPhonetic}

\begin{EntryWithPhonetic}{言论}{yan2lun4}{7,6}{⾔,⾔}[HSK 7-9]
  \definition{s.}{discurso; pontos de vista políticos; opinião sobre assuntos públicos; expressão das próprias convicções políticas; discussões sobre política e assuntos públicos em geral}
  \synonymref{言辞}{yan2ci2}
  \synonymref{舆论}{yu2lun4}
  \antonymref{默默}{mo4mo4}
\end{EntryWithPhonetic}

\begin{EntryWithPhonetic}{言行}{yan2xing2}{7,6}{⾔,⾏}[HSK 7-9]
  \definition{s.}{palavras e ações; declarações e atos}
\end{EntryWithPhonetic}

\begin{EntryWithPhonetic}{言语}{yan2yu3}{7,9}{⾔,⾔}[HSK 5]
  \definition{s.}{verbal; fala; linguagem falada; conversa; palavras}
\end{EntryWithPhonetic}

%%%%%%%%%% 岩 %%%%%%%%%%
\subsection*{岩}\addcontentsline{loh}{figure}{岩 \dpy{yan2}}

\begin{EntryWithPhonetic}{岩}{yan2}{8}{⼭}
  \definition*{s.}{Sobrenome: Yan}
  \definition{s.}{rocha; pedra | penhasco; rochedo; picos de montanhas formados por rochas salientes | caverna (nas colinas)}
\end{EntryWithPhonetic}

\begin{EntryWithPhonetic}{岩石}{yan2shi2}{8,5}{⼭,⽯}[HSK 7-9]
  \definition[块]{s.}{rocha; pedra; as rochas ígneas são agregados compostos por um ou mais minerais e constituem o principal material que forma a crosta terrestre; de acordo com sua formação, podem ser divididas em três categorias: rochas ígneas, rochas sedimentares e rochas metamórficas}
  \antonymref{石头}{shi2tou5}
\end{EntryWithPhonetic}

%%%%%%%%%% 沿 %%%%%%%%%%
\subsection*{沿}\addcontentsline{loh}{figure}{沿 \dpy{yan2}}

\begin{EntryWithPhonetic}{沿}{yan2}{8}{⽔}[HSK 6]
  \definition{prep.}{ao longo}
  \definition{s.}{beira; borda; acabamento}
  \definition{v.}{seguir (uma tradição, padrão, etc.) | enfeitar (com fita, faixa, etc.)}
\end{EntryWithPhonetic}

\begin{EntryWithPhonetic}{沿岸}{yan2'an4}{8,8}{⽔,⼭}[HSK 7-9]
  \definition{s.}{ao longo da margem (ou costa); litorâneo; ribeirinho | costeiro; litorâneo; áreas próximas a rios, lagos e ao mar}
\end{EntryWithPhonetic}

\begin{EntryWithPhonetic}{沿海}{yan2hai3}{8,10}{⽔,⽔}[HSK 6]
  \definition{s.}{costa; litoral; área ou região ao longo da costa}
\end{EntryWithPhonetic}

\begin{EntryWithPhonetic}{沿途}{yan2tu2}{8,10}{⽔,⾡}[HSK 7-9]
  \definition{adv.}{a caminho; ao longo de uma viagem | ao longo do caminho (estrada)}
\end{EntryWithPhonetic}

\begin{EntryWithPhonetic}{沿线}{yan2xian4}{8,8}{⽔,⽷}[HSK 7-9]
  \definition{s.}{ao longo da linha (isto é, uma linha ferroviária, rodoviária, aérea ou marítima); locais ao longo de ferrovias, rodovias ou rotas aéreas}
\end{EntryWithPhonetic}

\begin{EntryWithPhonetic}{沿着}{yan2zhe5}{8,11}{⽔,⽬}[HSK 6]
  \definition{prep.}{ao longo (de uma determinada rota)}
\end{EntryWithPhonetic}

%%%%%%%%%% 炎 %%%%%%%%%%
\subsection*{炎}\addcontentsline{loh}{figure}{炎 \dpy{yan2}}

\begin{EntryWithPhonetic}{炎}{yan2}{8}{⽕}
  \definition{adj.}{escaldante; ardente}
  \definition{s.}{inflamação | poder; influência}
\end{EntryWithPhonetic}

\begin{EntryWithPhonetic}{炎热}{yan2re4}{8,10}{⽕,⽕}[HSK 7-9]
  \definition{adj.}{ardente; escaldante; extremamente quente}
  \synonymref{火热}{huo3re4}
  \synonymref{闷热}{men1re4}
  \antonymref{冰冷}{bing1leng3}
  \antonymref{寒冷}{han2leng3}
  \antonymref{凉爽}{liang2shuang3}
  \antonymref{清凉}{qing1liang2}
\end{EntryWithPhonetic}

\begin{EntryWithPhonetic}{炎症}{yan2zheng4}{8,10}{⽕,⽧}[HSK 7-9]
  \definition{s.}{inflamação}
\end{EntryWithPhonetic}

%%%%%%%%%% 研 %%%%%%%%%%
\subsection*{研}\addcontentsline{loh}{figure}{研 \dpy{yan2}}

\begin{EntryWithPhonetic}{研}{yan2}{9}{⽯}
  \definition{s.}{(abreviação)  pesquisador adjunto, 副研}
  \definition{v.}{moer; esmerilhar; triturar; pulverizar | estudar; pesquisar}
  \seealsoref{副研}{fu4yan2}
\end{EntryWithPhonetic}

\begin{EntryWithPhonetic}{研发}{yan2fa1}{9,5}{⽯,⼜}[HSK 6]
  \definition{s.}{pesquisa e desenvolvimento; P\&D}
  \definition{v.}{pesquisar e/ou desenvolver}
\end{EntryWithPhonetic}

\begin{EntryWithPhonetic}{研究}{yan2jiu1}{9,7}{⽯,⽳}[HSK 4]
  \definition{v.}{estudar; pesquisar | discutir; considerar}
\end{EntryWithPhonetic}

\begin{EntryWithPhonetic}{研究生}{yan2jiu1sheng1}{9,7,5}{⽯,⽳,⽣}[HSK 4]
  \definition[位,名,个,些]{s.}{pós"-graduando; estudante de pós"-graduação}
\end{EntryWithPhonetic}

\begin{EntryWithPhonetic}{研究所}{yan2jiu1suo3}{9,7,8}{⽯,⽳,⼾}[HSK 5]
  \definition[家,个]{s.}{instituto de pesquisa; instituição de pesquisa científica envolvida em pesquisas em um determinado campo}
\end{EntryWithPhonetic}

\begin{EntryWithPhonetic}{研讨}{yan2tao3}{9,5}{⽯,⾔}[HSK 7-9]
  \definition{v.}{discutir; deliberar}
  \synonymref{考虑}{kao3lv4}
  \synonymref{商讨}{shang1tao3}
  \synonymref{商量}{shang1liang5}
  \synonymref{探讨}{tan4tao3}
  \synonymref{研究}{yan2jiu1}
\end{EntryWithPhonetic}

\begin{EntryWithPhonetic}{研制}{yan2zhi4}{9,8}{⽯,⼑}[HSK 4]
  \definition{v.}{desenvolver; fabricar; produzir | triturar; (medicina chinesa) moer}
\end{EntryWithPhonetic}

%%%%%%%%%% 盐 %%%%%%%%%%
\subsection*{盐}\addcontentsline{loh}{figure}{盐 \dpy{yan2}}

\begin{EntryWithPhonetic}{盐}{yan2}{10}{⽫}[HSK 4]
  \definition[袋,勺,把,包,粒]{s.}{sal (de cozinha) | Química: sal (produto formado pela neutralização de um ácido por uma base)}
\end{EntryWithPhonetic}

%%%%%%%%%% 阎 %%%%%%%%%%
\subsection*{阎}\addcontentsline{loh}{figure}{阎 \dpy{yan2}}

\begin{EntryWithPhonetic}{阎}{yan2}{11}{⾨}
  \definition*{s.}{Sobrenome: Yan}
  \definition{s.}{Arcaico: o portão de uma viela; entrada para uma viela}
\end{EntryWithPhonetic}

\begin{EntryWithPhonetic}{阎王}{yan2wang5}{11,4}{⾨,⽟}[HSK 7-9]
  \definition*{s.}{Yama; Rei do Inferno}
  \definition{s.}{uma pessoa extremamente cruel e violenta; uma metáfora para uma pessoa extremamente rigorosa ou cruel}
\end{EntryWithPhonetic}

%%%%%%%%%% 颜 %%%%%%%%%%
\subsection*{颜}\addcontentsline{loh}{figure}{颜 \dpy{yan2}}

\begin{EntryWithPhonetic}{颜}{yan2}{15}{⾴}
  \definition*{s.}{Sobrenome: Yan}
  \definition{s.}{rosto; semblante; expressão facial | rosto; prestígio; dignidade | cor}
\end{EntryWithPhonetic}

\begin{EntryWithPhonetic}{颜色}{yan2se4}{15,6}{⾴,⾊}[HSK 2]
  \definition[个,种]{s.}{cor; a sensação visual de um objeto é uma impressão diferente produzida pelas diferentes quantidades de luz absorvidas e refletidas pelo objeto | tez; semblante; aparência; geralmente se refere à aparência de uma garota | olhar severo no rosto como um aviso; um olhar ou ação que faz os outros parecerem particularmente ferozes | a expressão mostrada no rosto}
\end{EntryWithPhonetic}

%%%%%%%%%% 广 %%%%%%%%%%
\subsection*{广}\addcontentsline{loh}{figure}{广 \dpy{yan3}}

\begin{EntryWithPhonetic}{广}{yan3}{3}{⼴}[Kangxi 53]
  \definition[家]{s.}{casa ou edifício construído contra ou ao longo da encosta de uma montanha ou penhasco}
  \seeref{an1}
  \seeref{guang3}
\end{EntryWithPhonetic}

%%%%%%%%%% 衍 %%%%%%%%%%
\subsection*{衍}\addcontentsline{loh}{figure}{衍 \dpy{yan3}}

\begin{EntryWithPhonetic}{衍}{yan3}{9}{⾏}
  \definition{adj.}{redundante; supérfluo}
  \definition{s.}{Literário: planície baixa; terreno baixo e plano | Literário: pântano; brejo; charco}
  \definition{v.}{Literário: amplificar; desenvolver; disseminar; executar; dar plena expressão a | Literário: ser redundante; ser supérfluo}
\end{EntryWithPhonetic}

\begin{EntryWithPhonetic}{衍生}{yan3sheng1}{9,5}{⾏,⽣}[HSK 7-9]
  \definition{v.}{derivar; um composto mais simples pode ter seus átomos ou grupos de átomos substituídos por outros átomos ou grupos de átomos para formar um composto mais complexo | evoluir; desenvolver; reproduzir e prosperar}
  \synonymref{繁殖}{fan2zhi2}
  \synonymref{掌握}{zhang3wo4}
\end{EntryWithPhonetic}

%%%%%%%%%% 掩 %%%%%%%%%%
\subsection*{掩}\addcontentsline{loh}{figure}{掩 \dpy{yan3}}

\begin{EntryWithPhonetic}{掩}{yan3}{11}{⼿}
  \definition{v.}{cobrir; esconder | fechar; encerrar | Dialeto: ficar espremido (ou beliscado) ao fechar uma porta, tampa, etc. | atacar de surpresa; emboscar e atacar}
\end{EntryWithPhonetic}

\begin{EntryWithPhonetic}{掩盖}{yan3gai4}{11,11}{⼿,⽫}[HSK 7-9]
  \definition{v.}{cobrir; cobrir com um cobertor | ocultar; encobrir}
  \seealsoref{遮盖}{zhe1gai4}
  \synonymref{包围}{bao1wei2}
  \synonymref{保护}{bao3hu4}
  \synonymref{覆盖}{fu4gai4}
  \synonymref{掩护}{yan3hu4}
  \synonymref{掩饰}{yan3shi4}
  \synonymref{隐藏}{yin3cang2}
  \synonymref{隐瞒}{yin3man2}
  \synonymref{遮盖}{zhe1gai4}
  \antonymref{暴露}{bao4lu4}
  \antonymref{揭露}{jie1lu4}
  \antonymref{揭示}{jie1shi4}
  \antonymref{揭晓}{jie1xiao3}
  \antonymref{坦白}{tan3bai2}
\end{EntryWithPhonetic}

\begin{EntryWithPhonetic}{掩护}{yan3hu4}{11,7}{⼿,⼿}[HSK 7-9]
  \definition[面]{s.}{escudo}
  \definition{v.}{proteger; blindar; cobrir}
  \synonymref{保护}{bao3hu4}
  \synonymref{掩盖}{yan3gai4}
  \synonymref{掩饰}{yan3shi4}
  \antonymref{伤害}{shang1hai4}
\end{EntryWithPhonetic}

\begin{EntryWithPhonetic}{掩饰}{yan3shi4}{11,8}{⼿,⾷}[HSK 7-9]
  \definition{v.}{mascarar; ocultar; encobrir; disfarçar; tentar esconder os pensamentos, sentimentos e emoções dos outros}
  \synonymref{掩盖}{yan3gai4}
  \synonymref{掩护}{yan3hu4}
  \synonymref{隐瞒}{yin3man2}
  \synonymref{遮盖}{zhe1gai4}
  \synonymref{装饰}{zhuang1shi4}
  \antonymref{暴露}{bao4lu4}
  \antonymref{表白}{biao3bai2}
  \antonymref{衬托}{chen4tuo1}
  \antonymref{揭露}{jie1lu4}
  \antonymref{流露}{liu2lu4}
\end{EntryWithPhonetic}

%%%%%%%%%% 眼 %%%%%%%%%%
\subsection*{眼}\addcontentsline{loh}{figure}{眼 \dpy{yan3}}

\begin{EntryWithPhonetic}{眼}{yan3}{11}{⽬}[HSK 2]
  \definition{clas.}{usado para grandes coisas ocas: poços, fogões, panelas, etc.}
  \definition[双,只]{s.}{olho; o órgão visual dos humanos ou animais | abertura; pequeno furo; pequeno buraco | ponto"-chave; refere"-se ao ponto"-chave das coisas | armadilha; um termo do jogo Go que se refere a um espaço vazio cercado pelas peças de um jogador, onde o outro jogador não pode colocar uma peça, a menos que haja circunstâncias especiais | uma batida sem acento na música tradicional chinesa}
\end{EntryWithPhonetic}

\begin{EntryWithPhonetic}{眼柄}{yan3bing3}{11,9}{⽬,⽊}
  \definition{s.}{pedúnculo ocular (de crustáceo, etc.)}
\end{EntryWithPhonetic}

\begin{EntryWithPhonetic}{眼袋}{yan3dai4}{11,11}{⽬,⾐}
  \definition{s.}{inchaço sob os olhos}
\end{EntryWithPhonetic}

\begin{EntryWithPhonetic}{眼光}{yan3guang1}{11,6}{⽬,⼉}[HSK 5]
  \definition{s.}{olho; visão | visão; percepção; previsão; capacidade de observar e identificar coisas | vista; ponto de vista}
\end{EntryWithPhonetic}

\begin{EntryWithPhonetic}{眼红}{yan3hong2}{11,6}{⽬,⽷}[HSK 7-9]
  \definition{adj.}{furioso; irritado; com os olhos ardendo de fúria}
  \definition{v.}{cobiçar; ter inveja; sentir ciúmes; ter olhos verdes; ser extremamente invejoso e ciumento da fama, fortuna ou bens materiais alheios, e também desejar muito obtê-los}
  \synonymref{忌妒}{ji4du4}
  \antonymref{羡慕}{xian4mu4}
\end{EntryWithPhonetic}

\begin{EntryWithPhonetic}{眼花缭乱}{yan3hua1liao2luan4}{11,7,15,7}{⽬,⾋,⽷,⼄}
  \definition{v.}{ficar deslumbrado | deslumbrar}
\end{EntryWithPhonetic}

\begin{EntryWithPhonetic}{眼界}{yan3jie4}{11,9}{⽬,⽥}[HSK 7-9]
  \definition{s.}{campo de visão; campo visual; perspectiva; a abrangência das coisas vistas, referindo"-se metaforicamente à amplitude do conhecimento}
  \synonymref{见识}{jian4shi5}
  \synonymref{视野}{shi4ye3}
  \synonymref{眼光}{yan3guang1}
  \synonymref{阅历}{yue4li4}
\end{EntryWithPhonetic}

\begin{EntryWithPhonetic}{眼镜}{yan3jing4}{11,16}{⽬,⾦}[HSK 4]
  \definition[副]{s.}{óculos; óculos de grau; lentes usadas nos olhos para melhorar a visão ou proteger os olhos, feitas de vidro ou cristal incolor ou colorido}
\end{EntryWithPhonetic}

\begin{EntryWithPhonetic}{眼睛}{yan3jing5}{11,13}{⽬,⽬}[HSK 2]
  \definition[双,只]{s.}{olho(s)}
\end{EntryWithPhonetic}

\begin{EntryWithPhonetic}{眼看}{yan3kan4}{11,9}{⽬,⽬}[HSK 6]
  \definition{adv.}{em breve; em um momento; imediatamente}
  \definition{v.}{observar impotentemente; olhar passivamente; observar (o que está acontecendo)}
\end{EntryWithPhonetic}

\begin{EntryWithPhonetic}{眼泪}{yan3lei4}{11,8}{⽬,⽔}[HSK 4]
  \definition[滴,行]{s.}{lágrimas; termo genérico para lágrimas; fluido incolor e transparente secretado pelas glândulas lacrimais no olho, que serve para proteger o olho}
\end{EntryWithPhonetic}

\begin{EntryWithPhonetic}{眼里}{yan3li3}{11,7}{⽬,⾥}[HSK 4]
  \definition{s.}{aos olhos de alguém; na opinião (ou visão) de alguém}
\end{EntryWithPhonetic}

\begin{EntryWithPhonetic}{眼前}{yan3qian2}{11,9}{⽬,⼑}[HSK 3]
  \definition{adv.}{agora; (no) momento}
  \definition{s.}{diante dos olhos; diante de | agora; (no) momento}
\end{EntryWithPhonetic}

\begin{EntryWithPhonetic}{眼色}{yan3se4}{11,6}{⽬,⾊}[HSK 7-9]
  \definition{s.}{olhar significativo; insinuação transmitida com os olhos; o contato visual serve para expressar significado ou emoção aos outros | capacidade de fazer o que bem entender; refere"-se à capacidade de fazer julgamentos rápidos com base na situação e de tomar medidas adequadas com flexibilidade}
  \seealsoref{眼神}{yan3shen2}
  \seealsoref{眼神}{yan3shen5}
\end{EntryWithPhonetic}

\begin{EntryWithPhonetic}{眼神}{yan3shen2}{11,9}{⽬,⽰}[HSK 7-9]
  \definition{s.}{olhar; expressão nos olhos; a expressão transmitida pelos olhos | visão; vista}
  \seeref{yan3shen5}
  \seealsoref{眼色}{yan3se4}
  \synonymref{目光}{mu4guang1}
\end{EntryWithPhonetic}

\begin{EntryWithPhonetic}{眼神}{yan3shen5}{11,9}{⽬,⽰}
  \definition{s.}{olhar; expressão nos olhos; contato visual}
  \seeref{yan3shen2}
  \seealsoref{眼色}{yan3se4}
  \synonymref{目光}{mu4guang1}
\end{EntryWithPhonetic}

\begin{EntryWithPhonetic}{眼下}{yan3xia4}{11,3}{⽬,⼀}[HSK 7-9]
  \definition{s.}{no momento; atualmente; agora}
  \seealsoref{眼前}{yan3qian2}
\end{EntryWithPhonetic}

\begin{EntryWithPhonetic}{眼证}{yan3zheng4}{11,7}{⽬,⾔}
  \definition{s.}{testemunha ocular}
\end{EntryWithPhonetic}

%%%%%%%%%% 演 %%%%%%%%%%
\subsection*{演}\addcontentsline{loh}{figure}{演 \dpy{yan3}}

\begin{EntryWithPhonetic}{演}{yan3}{14}{⽔}[HSK 3]
  \definition{v.}{desenvolver; evoluir | deduzir; elaborar | exercitar; praticar | representar; atuar; encenar | desempenhar}
\end{EntryWithPhonetic}

\begin{EntryWithPhonetic}{演变}{yan3bian4}{14,8}{⽔,⼜}[HSK 7-9]
  \definition{v.}{evoluir; desenvolver; desenvolvimento e mudanças (frequentemente referindo"-se àquelas que têm longa duração)}
  \synonymref{变迁}{bian4qian1}
  \antonymref{静止}{jing4zhi3}
\end{EntryWithPhonetic}

\begin{EntryWithPhonetic}{演播室}{yan3bo1shi4}{14,15,9}{⽔,⼿,⼧}[HSK 7-9]
  \definition{s.}{estúdio de transmissão}
\end{EntryWithPhonetic}

\begin{EntryWithPhonetic}{演唱}{yan3chang4}{14,11}{⽔,⼝}[HSK 3]
  \definition{v.}{cantar em uma performance; apresentar canções, óperas, peças teatrais, etc.}
\end{EntryWithPhonetic}

\begin{EntryWithPhonetic}{演唱会}{yan3chang4hui4}{14,11,6}{⽔,⼝,⼈}[HSK 3]
  \definition[个,场]{s.}{recital vocal; concerto vocal; uma forma de apresentação centrada no canto, acompanhada por movimentos de dança simples}
\end{EntryWithPhonetic}

\begin{EntryWithPhonetic}{演出}{yan3chu1}{14,5}{⽔,⼐}[HSK 3]
  \definition[场,次]{s.}{show; concerto; performance}
  \definition{v.}{apresentar; representar; fazer um show; apresentar peças teatrais, danças, artes cênicas, acrobacias, etc. para o público apreciar}
\end{EntryWithPhonetic}

\begin{EntryWithPhonetic}{演歌}{yan3ge5}{14,14}{⽔,⽋}[HSK 7-9]
  \definition*{s.}{Enka (gênero japonês de baladas sentimentais)}
\end{EntryWithPhonetic}

\begin{EntryWithPhonetic}{演技}{yan3ji4}{14,7}{⽔,⼿}[HSK 7-9]
  \definition{s.}{atuação; habilidades de atuação; habilidades de performance; as habilidades de atuação referem"-se à capacidade de um ator de criar personagens usando diversas técnicas e métodos}
\end{EntryWithPhonetic}

\begin{EntryWithPhonetic}{演讲}{yan3jiang3}{14,6}{⽔,⾔}[HSK 4]
  \definition[场,次]{s.}{palestra; discurso; ato ou a atividade de apresentar ou expressar ideias, opiniões ou informações oralmente em público ou diante de um público}
  \definition{v.}{dar uma palestra; fazer um discurso; informar o público sobre uma determinada área de conhecimento ou opinião sobre um determinado assunto}
\end{EntryWithPhonetic}

\begin{EntryWithPhonetic}{演练}{yan3lian4}{14,8}{⽔,⽷}[HSK 7-9]
  \definition{v.}{realizar um exercício ou prática}
  \synonymref{练习}{lian4xi2}
  \synonymref{训练}{xun4lian4}
\end{EntryWithPhonetic}

\begin{EntryWithPhonetic}{演示}{yan3shi4}{14,5}{⽔,⽰}[HSK 7-9]
  \definition{v.}{mostrar; demonstrar; mostrar o processo de desenvolvimento e mudança das coisas usando experimentos, objetos físicos ou gráficos}
  \synonymref{示范}{shi4fan4}
\end{EntryWithPhonetic}

\begin{EntryWithPhonetic}{演说}{yan3shuo1}{14,9}{⽔,⾔}[HSK 7-9]
  \definition{s.}{fala; discurso; opiniões expressas}
  \definition{v.}{proferir um discurso; fazer um pronunciamento; discursar}
  \seealsoref{演讲}{yan3jiang3}
  \synonymref{演讲}{yan3jiang3}
\end{EntryWithPhonetic}

\begin{EntryWithPhonetic}{演习}{yan3xi2}{14,3}{⽔,⼄}[HSK 7-9]
  \definition{v.}{manobrar; exercitar; treinar; praticar; realizar exercícios de campo de acordo com o plano estabelecido (utilizado principalmente em aplicações militares)}
  \synonymref{练习}{lian4xi2}
  \synonymref{实习}{shi2xi2}
  \antonymref{实践}{shi2jian4}
\end{EntryWithPhonetic}

\begin{EntryWithPhonetic}{演戏}{yan3/xi4}{14,6}{⽔,⼽}[HSK 7-9]
  \definition{v.+compl.}{encenar uma peça; atuar em uma peça; apresentar peças de teatro e óperas | representar; fingir; uma metáfora para enganar as pessoas}
\end{EntryWithPhonetic}

\begin{EntryWithPhonetic}{演艺圈}{yan3yi4quan1}{14,4,11}{⽔,⾋,⼞}[HSK 7-9]
  \definition{s.}{indústria do entretenimento; \emph{show business}; \emph{show biz}}
\end{EntryWithPhonetic}

\begin{EntryWithPhonetic}{演绎}{yan3yi4}{14,8}{⽔,⽷}[HSK 7-9]
  \definition{v.}{elaborar; descrever detalhadamente; imaginar e desenvolver com base em certos princípios | elaborar; expor sobre; apresentar coisas e conceitos abstratos; expressar | deduzir; um método de pensar sobre problemas, derivando conclusões de princípios gerais para casos específicos}
  \synonymref{表演}{biao3yan3}
  \synonymref{推理}{tui1li3}
  \antonymref{归纳}{gui1na4}
\end{EntryWithPhonetic}

\begin{EntryWithPhonetic}{演员}{yan3yuan2}{14,7}{⽔,⼝}[HSK 3]
  \definition[个,位,名]{s.}{ator; artista; pessoas que participam de apresentações teatrais, cinematográficas, de dança, de artes cênicas, de acrobacias, etc.}
\end{EntryWithPhonetic}

\begin{EntryWithPhonetic}{演奏}{yan3zou4}{14,9}{⽔,⼤}[HSK 6]
  \definition{v.}{tocar um instrumento musical; fazer uma apresentação instrumental}
\end{EntryWithPhonetic}

%%%%%%%%%% 厌 %%%%%%%%%%
\subsection*{厌}\addcontentsline{loh}{figure}{厌 \dpy{yan4}}

\begin{EntryWithPhonetic}{厌}{yan4}{6}{⼚}
  \definition{v.}{estar farto de; estar entediado com; estar cansado de | ficar satisfeito; satisfazer | ter nojo de; detestar}
  \synonymref{恶}{wu4}
  \antonymref{喜}{xi3}
\end{EntryWithPhonetic}

\begin{EntryWithPhonetic}{厌烦}{yan4fan2}{6,10}{⼚,⽕}[HSK 7-9]
  \definition{v.}{estar farto de; estar de saco cheio de}
  \synonymref{讨厌}{tao3/yan4}
  \synonymref{厌倦}{yan4juan4}
  \antonymref{爱好}{ai4hao4}
  \antonymref{可爱}{ke3'ai4}
  \antonymref{耐心}{nai4xin1}
  \antonymref{喜欢}{xi3huan5}
  \antonymref{兴趣}{xing4qu4}
\end{EntryWithPhonetic}

\begin{EntryWithPhonetic}{厌倦}{yan4juan4}{6,10}{⼚,⼈}[HSK 7-9]
  \definition{v.}{estar cansado de; estar farto de; perder o interesse em uma atividade e faltar vontade de continuá"-la}
  \synonymref{讨厌}{tao3/yan4}
  \synonymref{厌烦}{yan4fan2}
  \antonymref{痴迷}{chi1mi2}
  \antonymref{迷恋}{mi2lian4}
\end{EntryWithPhonetic}

%%%%%%%%%% 咽 %%%%%%%%%%
\subsection*{咽}\addcontentsline{loh}{figure}{咽 \dpy{yan4}}

\begin{EntryWithPhonetic}{咽}{yan4}{9}{⼝}[HSK 7-9]
  \definition{v.}{engolir; permitir que alimentos ou outras substâncias presentes na boca passem pela faringe e cheguem ao esôfago | engolir (metaforicamente); essa metáfora descreve o ato de suprimir ou não expressar as próprias palavras ou a raiva}
  \seeref{yan1}
  \seeref{ye4}
  \antonymref{吐}{tu3}
\end{EntryWithPhonetic}

%%%%%%%%%% 宴 %%%%%%%%%%
\subsection*{宴}\addcontentsline{loh}{figure}{宴 \dpy{yan4}}

\begin{EntryWithPhonetic}{宴}{yan4}{10}{⼧}
  \definition{adj.}{tranquilo e confortável}
  \definition[个,场]{s.}{festa; banquete | facilidade e conforto; felicidade; lazer}
  \definition{v.}{entreter em um banquete; festejar}
\end{EntryWithPhonetic}

\begin{EntryWithPhonetic}{宴会}{yan4hui4}{10,6}{⼧,⼈}[HSK 6]
  \definition[个,场,次]{s.}{festa; banquete; jantar; uma reunião onde convidados e anfitriões bebem e comem juntos (referindo"-se a uma ocasião mais solene)}
\end{EntryWithPhonetic}

%%%%%%%%%% 艳 %%%%%%%%%%
\subsection*{艳}\addcontentsline{loh}{figure}{艳 \dpy{yan4}}

\begin{EntryWithPhonetic}{艳}{yan4}{10}{⾊}
  \definition{adj.}{magnífico; belo; colorido; chamativo; brilhante | amoroso}
  \definition{v.}{Litarário: admirar; invejar}
  \antonymref{素}{su4}
\end{EntryWithPhonetic}

\begin{EntryWithPhonetic}{艳丽}{yan4li4}{10,7}{⾊,⼀}[HSK 7-9]
  \definition{adj.}{maravilhoso; magnífico; de cores vivas e belo | lindo; bela (escrita)}
  \synonymref{灿烂}{can4lan4}
  \synonymref{华丽}{hua2li4}
  \synonymref{美丽}{mei3li4}
  \synonymref{倾城}{qing1cheng2}
  \synonymref{鲜艳}{xian1yan4}
  \synonymref{秀丽}{xiu4li4}
  \antonymref{丑陋}{chou3lou4}
  \antonymref{难看}{nan2kan4}
  \antonymref{朴素}{pu3su4}
\end{EntryWithPhonetic}

%%%%%%%%%% 验 %%%%%%%%%%
\subsection*{验}\addcontentsline{loh}{figure}{验 \dpy{yan4}}

\begin{EntryWithPhonetic}{验}{yan4}{10}{⾺}[HSK 7-9]
  \definition{s.}{efeito pretendido; resultados esperados | eficácia; efetividade}
  \definition{v.}{examinar; verificar; testar; confirmar | provar ser eficaz; produzir o resultado esperado | examinar; verificar; testar; investigar | provar ser eficaz; produzir o resultado esperado; alcançar os resultados esperados | para verificar; isso foi confirmado por meio da experiência prática e de outros meios}
\end{EntryWithPhonetic}

\begin{EntryWithPhonetic}{验收}{yan4shou1}{10,6}{⾺,⽁}[HSK 7-9]
  \definition{v.}{verificar e aceitar; os itens são inspecionados de acordo com determinados padrões antes de serem aceitos}
  \synonymref{检验}{jian3yan4}
\end{EntryWithPhonetic}

\begin{EntryWithPhonetic}{验证}{yan4zheng4}{10,7}{⾺,⾔}[HSK 7-9]
  \definition{v.}{verificar; provar; isso foi confirmado por meio de experimentos e verificado por meio de testes}
  \synonymref{论证}{lun4zheng4}
\end{EntryWithPhonetic}

%%%%%%%%%% 焰 %%%%%%%%%%
\subsection*{焰}\addcontentsline{loh}{figure}{焰 \dpy{yan4}}

\begin{EntryWithPhonetic}{焰}{yan4}{12}{⽕}
  \definition{adj.}{arrogante; metáfora para poder e prestígio}
  \definition{s.}{chama; labareda}
\end{EntryWithPhonetic}

\begin{EntryWithPhonetic}{焰火}{yan4huo3}{12,4}{⽕,⽕}[HSK 7-9]
  \definition[种]{s.}{fogos de artifício}
  \synonymref{火花}{huo3hua1}
  \synonymref{烟火}{yan1huo5}
\end{EntryWithPhonetic}

%%%%%%%%%% 燕 %%%%%%%%%%
\subsection*{燕}\addcontentsline{loh}{figure}{燕 \dpy{yan4}}

\begin{EntryWithPhonetic}{燕}{yan4}{16}{⽕}
  \definition{s.}{andorinha (uma família de pássaros) | conforto; desfruta | festa; banquete; jantar}
  \seeref{yan1}
\end{EntryWithPhonetic}

\begin{EntryWithPhonetic}{燕子}{yan4zi5}{16,3}{⽕,⼦}[HSK 7-9]
  \definition[只,窝]{s.}{andorinha}
\end{EntryWithPhonetic}

%%%%%%%%%% 扬 %%%%%%%%%%
\subsection*{扬}\addcontentsline{loh}{figure}{扬 \dpy{yang2}}

\begin{EntryWithPhonetic}{扬}{yang2}{6}{⼿}[HSK 7-9]
  \definition*{s.}{Yangzhou, abreviação de 扬州 | Sobrenome: Yang}
  \definition{v.}{levantar | separar e espalhar; peneirar | espalhar; fazer conhecido}
  \seealsoref{扬州}{yang2zhou1}
  \antonymref{避}{bi4}
  \antonymref{抑}{yi4}
\end{EntryWithPhonetic}

\begin{EntryWithPhonetic}{扬雄}{yang2xiong2}{6,12}{⼿,⾫}
  \definition*{s.}{Yang Xiong (53 AC-18 DC), estudioso, poeta e lexicógrafo, autor do primeiro dicionário de dialeto chinês 方言}
  \seealsoref{方言}{fang1yan2}
\end{EntryWithPhonetic}

\begin{EntryWithPhonetic}{扬州}{yang2zhou1}{6,6}{⼿,⼮}
  \definition*{s.}{Yangzhou, uma cidade na província de Jiangsu}
\end{EntryWithPhonetic}

%%%%%%%%%% 羊 %%%%%%%%%%
\subsection*{羊}\addcontentsline{loh}{figure}{羊 \dpy{yang2}}

\begin{EntryWithPhonetic}{羊}{yang2}{6}{⽺}[HSK 3][Kangxi 123]
  \definition*{s.}{Sobrenome: Yang}
  \definition[只,头,群]{s.}{carneiro; ovelha; bode; cabra; antílope}
\end{EntryWithPhonetic}

%%%%%%%%%% 阳 %%%%%%%%%%
\subsection*{阳}\addcontentsline{loh}{figure}{阳 \dpy{yang2}}

\begin{EntryWithPhonetic}{阳}{yang2}{6}{⾩}
  \definition*{s.}{Sobrenome: Yang}
  \definition{adj.}{em relevo | aberto; evidente; revelado | pertencente a este mundo; preocupado com os seres vivos; superstição se refere a coisas que pertencem aos vivos e ao mundo | positivo; carregado positivamente}
  \definition{s.}{o Sol; luz solar | ao sul de uma colina ou ao norte de um rio | Yang (o princípio masculino ou positivo da natureza); na filosofia chinesa antiga, refere"-se a um dos dois opostos que permeiam todas as coisas no universo (o outro lado é Yin) | masculino; refere"-se aos órgãos genitais masculinos ou à função reprodutiva | varanda; refere"-se ao lugar onde o sol brilha}
  \seealsoref{阴阳}{yin1yang2}
  \antonymref{阴}{yin1}
\end{EntryWithPhonetic}

\begin{EntryWithPhonetic}{阳光}{yang2guang1}{6,6}{⾩,⼉}[HSK 3]
  \definition{adj.}{alegre; otimista; personalidade positiva e alegre; cheio de vitalidade juvenil | aberto; transparente; público; conduzido sob supervisão pública}
  \definition[缕,束,道]{s.}{luz do sol; raio de sol}
\end{EntryWithPhonetic}

\begin{EntryWithPhonetic}{阳台}{yang2tai2}{6,5}{⾩,⼝}[HSK 4]
  \definition[个,块,处]{s.}{varanda; terraço; sacada; pequeno terraço do edifício com grades para se refrescar, tomar sol ou olhar o horizonte}
\end{EntryWithPhonetic}

\begin{EntryWithPhonetic}{阳性}{yang2xing4}{6,8}{⾩,⼼}[HSK 7-9]
  \definition{s.}{positivo; um resultado positivo é uma forma de representar os resultados de um teste ou exame laboratorial realizado durante o diagnóstico de uma doença; um resultado positivo indica a presença de um patógeno no organismo ou uma reação alérgica a um medicamento | gênero masculino; em algumas línguas, os substantivos (assim como os pronomes e adjetivos) são classificados como femininos, masculinos ou femininos, masculinos e neutros}
  \seealsoref{性}{xing4}
  \antonymref{阴性}{yin1xing4}
\end{EntryWithPhonetic}

%%%%%%%%%% 杨 %%%%%%%%%%
\subsection*{杨}\addcontentsline{loh}{figure}{杨 \dpy{yang2}}

\begin{EntryWithPhonetic}{杨}{yang2}{7}{⽊}
  \definition*{s.}{Sobrenome: Yang}
  \definition[颗]{s.}{álamo (\emph{Populus}); o álamo é uma árvore decídua com folhas alternadas, ovais ou ovais"-lanceoladas, e amentos; existem muitas variedades, incluindo o álamo"-prateado, o álamo"-branco e o álamo"-de"-folhas"-pequenas}
\end{EntryWithPhonetic}

\begin{EntryWithPhonetic}{杨树}{yang2shu4}{7,9}{⽊,⽊}[HSK 7-9]
  \definition[棵,株]{s.}{álamo; diversas árvores do gênero \emph{Populus}}
\end{EntryWithPhonetic}

%%%%%%%%%% 洋 %%%%%%%%%%
\subsection*{洋}\addcontentsline{loh}{figure}{洋 \dpy{yang2}}

\begin{EntryWithPhonetic}{洋}{yang2}{9}{⽔}[HSK 6]
  \definition*{s.}{Sobrenome: Yang}
  \definition{adj.}{vasto; rico; transbordante | estrangeiro (especialmente ocidental) | moderno}
  \definition[个,片]{s.}{oceano | moeda de prata}
  \seealsoref{土}{tu3}
  \antonymref{土}{tu3}
\end{EntryWithPhonetic}

\begin{EntryWithPhonetic}{洋葱}{yang2cong1}{9,12}{⽔,⾋}
  \definition{s.}{cebola}
\end{EntryWithPhonetic}

\begin{EntryWithPhonetic}{洋溢}{yang2yi4}{9,13}{⽔,⽔}[HSK 7-9]
  \definition{v.}{transbordar de; estar repleto de}
  \seealsoref{充满}{chong1man3}
  \antonymref{隐藏}{yin3cang2}
\end{EntryWithPhonetic}

%%%%%%%%%% 仰 %%%%%%%%%%
\subsection*{仰}\addcontentsline{loh}{figure}{仰 \dpy{yang3}}

\begin{EntryWithPhonetic}{仰}{yang3}{6}{⼈}[HSK 6]
  \definition*{s.}{Sobrenome: Yang}
  \definition{v.}{levantar | virar para cima | admirar; respeitar | confiar em; depender de}
  \antonymref{俯}{fu3}
\end{EntryWithPhonetic}

%%%%%%%%%% 养 %%%%%%%%%%
\subsection*{养}\addcontentsline{loh}{figure}{养 \dpy{yang3}}

\begin{EntryWithPhonetic}{养}{yang3}{9}{⼋}[HSK 2]
  \definition*{s.}{Sobrenome: Yang}
  \definition{adj.}{adotivo; órfão; adotado; não biológico}
  \definition{s.}{qualidade; (caráter moral, desempenho acadêmico, etc.) boas qualidades}
  \definition{v.}{apoiar; prover; fornecer dinheiro e materiais necessários para viver | aumentar; manter; crescer; alimentar os animais e cuidar de suas vidas para que possam crescer | dar à luz | formar; adquirir; cultivar | descansar; curar; convalescer; recuperar a saúde | manter; manter em bom estado | deixar (o cabelo) crescer | ajudar; apoiar | cultivar (plantações ou flores)}
\end{EntryWithPhonetic}

\begin{EntryWithPhonetic}{养成}{yang3cheng2}{9,6}{⼋,⼽}[HSK 4]
  \definition{v.}{cultivar; desenvolver; cultivar para formar; nutrir para crescer}
\end{EntryWithPhonetic}

\begin{EntryWithPhonetic}{养分}{yang3fen4}{9,4}{⼋,⼑}
  \definition{s.}{nutriente}
\end{EntryWithPhonetic}

\begin{EntryWithPhonetic}{养活}{yang3huo5}{9,9}{⼋,⽔}[HSK 7-9]
  \definition{v.}{apoiar; prover para; providenciar despesas de subsistência ou itens de primeira necessidade para garantir a sobrevivência | criar animais}
  \synonymref{成活}{cheng2huo2}
  \synonymref{赡养}{shan4yang3}
\end{EntryWithPhonetic}

\begin{EntryWithPhonetic}{养老}{yang3 lao3}{9,6}{⼋,⽼}[HSK 6]
  \definition{v.}{prover assistência aos idosos (geralmente os pais) | viver a vida na aposentadoria; refere"-se ao idoso que descansa em casa}
\end{EntryWithPhonetic}

\begin{EntryWithPhonetic}{养老金}{yang3lao3jin1}{9,6,8}{⼋,⽼,⾦}[HSK 7-9]
  \definition{s.}{pensão; gratificação; aposentadorias para funcionários mais velhos | pensão de velhice}
\end{EntryWithPhonetic}

\begin{EntryWithPhonetic}{养老院}{yang3lao3yuan4}{9,6,9}{⼋,⽼,⾩}[HSK 7-9]
  \definition{s.}{lar de idosos; residência para a terceira idade; as casas de repouso são instituições administradas pelo governo ou por coletivos que acolhem idosos que vivem sozinhos}
\end{EntryWithPhonetic}

\begin{EntryWithPhonetic}{养料}{yang3liao4}{9,10}{⼋,⽃}
  \definition{s.}{nutrição}
\end{EntryWithPhonetic}

\begin{EntryWithPhonetic}{养生}{yang3sheng1}{9,5}{⼋,⽣}[HSK 7-9]
  \definition{v.}{cuidar da vida; conservar as energias vitais; preservar a saúde; manter"-se saudável; manter sua saúde e aumente sua vitalidade}
\end{EntryWithPhonetic}

\begin{EntryWithPhonetic}{养殖}{yang3zhi2}{9,12}{⼋,⽍}[HSK 7-9]
  \definition{v.}{cultivar; criar; reproduzir; cultivar e reproduzir (plantas e animais aquáticos)}
\end{EntryWithPhonetic}

%%%%%%%%%% 氧 %%%%%%%%%%
\subsection*{氧}\addcontentsline{loh}{figure}{氧 \dpy{yang3}}

\begin{EntryWithPhonetic}{氧}{yang3}{10}{⽓}[HSK 7-9]
  \definition{s.}{oxigênio (O)}
\end{EntryWithPhonetic}

\begin{EntryWithPhonetic}{氧气}{yang3qi4}{10,4}{⽓,⽓}[HSK 6]
  \definition{s.}{oxigênio (O); gás oxigênio}
\end{EntryWithPhonetic}

%%%%%%%%%% 痒 %%%%%%%%%%
\subsection*{痒}\addcontentsline{loh}{figure}{痒 \dpy{yang3}}

\begin{EntryWithPhonetic}{痒}{yang3}{11}{⽧}[HSK 7-9]
  \definition{v.}{coçar; fazer cócegas | estar ansioso (se coçando) para fazer algo; estar ansioso para tentar}
\end{EntryWithPhonetic}

%%%%%%%%%% 样 %%%%%%%%%%
\subsection*{样}\addcontentsline{loh}{figure}{样 \dpy{yang4}}

\begin{EntryWithPhonetic}{样}{yang4}{10}{⽊}[HSK 6]
  \definition{clas.}{usado para tipos de coisas}[这里有四样东西。===Há quatro coisas aqui.]
  \definition[个]{s.}{aparência; aspecto;  forma; aparência; a forma do objeto | modelo; amostra; padrão; coisas usadas como padrões | ar; maneira; aparência; a aparência ou expressão de uma pessoa | tendência; probabilidade; a situação ou tendência das coisas}
\end{EntryWithPhonetic}

\begin{EntryWithPhonetic}{样本}{yang4ben3}{10,5}{⽊,⽊}[HSK 7-9]
  \definition{s.}{amostra; espécime; folhas impressas ou recortadas com imagens de produtos, ou cadernos de tecido, usados para fins publicitários | livro de amostras; trechos de publicações são usados para fins publicitários ou para solicitar \emph{feedback}}
\end{EntryWithPhonetic}

\begin{EntryWithPhonetic}{样品}{yang4pin3}{10,9}{⽊,⼝}[HSK 7-9]
  \definition[件,个]{s.}{amostra; espécime; itens de amostra (frequentemente usados para promoção de produtos ou teste de materiais)}
  \synonymref{样本}{yang4ben3}
\end{EntryWithPhonetic}

\begin{EntryWithPhonetic}{样儿}{yang4r5}{10,2}{⽊,⼉}
  \definition{s.}{aparência | forma | modelo}
  \seealsoref{样子}{yang4zi5}
\end{EntryWithPhonetic}

\begin{EntryWithPhonetic}{样样}{yang4yang4}{10,10}{⽊,⽊}
  \definition{adv.}{todos os tipos}
\end{EntryWithPhonetic}

\begin{EntryWithPhonetic}{样章}{yang4zhang1}{10,11}{⽊,⾳}
  \definition{s.}{capítulo de amostra}
\end{EntryWithPhonetic}

\begin{EntryWithPhonetic}{样子}{yang4zi5}{10,3}{⽊,⼦}[HSK 2]
  \definition[个,种,副]{s.}{forma; aparência; estilo | ar; maneira; modalidade; estado | tendência; probabilidade; usado com 看 e 照 para expressar uma estimativa de uma tendência | modelo; amostra; padrão; uma pessoa ou coisa que pode ser usada como um padrão para as pessoas verificarem, seguirem ou aprenderem com ela}
  \seealsoref{看}{kan4}
  \seealsoref{样儿}{yang4r5}
  \seealsoref{照}{zhao4}
\end{EntryWithPhonetic}

%%%%%%%%%% 约 %%%%%%%%%%
\subsection*{约}\addcontentsline{loh}{figure}{约 \dpy{yao1}}

\begin{EntryWithPhonetic}{约}{yao1}{6}{⽷}
  \definition{adj.}{econômico; frugal | simples; breve | indistinto}
  \definition{adv.}{cerca de; ao redor; aproximadamente}
  \definition{s.}{pacto; acordo; nomeação; coisa prometida}
  \definition{v.}{marcar uma consulta; organizar | perguntar ou convidar com antecedência | restringir; conter | reduzir (fração aproximada)}
  \seeref{yue1}
\end{EntryWithPhonetic}

%%%%%%%%%% 妖 %%%%%%%%%%
\subsection*{妖}\addcontentsline{loh}{figure}{妖 \dpy{yao1}}

\begin{EntryWithPhonetic}{妖}{yao1}{7}{⼥}
  \definition{adj.}{maligno e fraudulento | sedutor; encantador | paquerador}
  \definition[个,只]{s.}{\emph{goblin}; demônio; espírito maligno}
\end{EntryWithPhonetic}

\begin{EntryWithPhonetic}{妖怪}{yao1guai4}{7,8}{⼥,⼼}[HSK 7-9]
  \definition[个]{s.}{monstro; \emph{goblin}; duende; demônio; os monstros das lendas e histórias geralmente têm uma aparência estranha ou aterrorizante e podem prejudicar os outros de alguma forma}
  \synonymref{魔鬼}{mo2gui3}
\end{EntryWithPhonetic}

%%%%%%%%%% 祅 %%%%%%%%%%
\subsection*{祅}\addcontentsline{loh}{figure}{祅 \dpy{yao1}}

\begin{EntryWithPhonetic}{祅}{yao1}{8}{⽰}
  \definition{s.}{espírito maligno | \emph{goblin} | bruxaria}
  \variantof{妖}
\end{EntryWithPhonetic}

%%%%%%%%%% 要 %%%%%%%%%%
\subsection*{要}\addcontentsline{loh}{figure}{要 \dpy{yao1}}

\begin{EntryWithPhonetic}{要}{yao1}{9}{⾑}
  \definition*{s.}{Sobrenome: Yao}
  \definition{v.}{exigir; pedir; requerer; solicitar; buscar; insistir com base em algo em que se apoia | forçar; coagir; ameaçar}
  \seeref{yao4}
\end{EntryWithPhonetic}

\begin{EntryWithPhonetic}{要求}{yao1qiu2}{9,7}{⾑,⽔}[HSK 2]
  \definition[个,点]{s.}{exigência; demanda; reivindicação; desejos ou condições específicas propostas}
  \definition{v.}{pedir; exigir; exigir; reivindicar; apresentar desejos ou condições específicas, esperando que sejam satisfeitos ou realizados}
\end{EntryWithPhonetic}

\begin{EntryWithPhonetic}{要挟}{yao1xie2}{9,9}{⾑,⼿}
  \definition{v.}{chantagear | ameaçar}
\end{EntryWithPhonetic}

%%%%%%%%%% 腰 %%%%%%%%%%
\subsection*{腰}\addcontentsline{loh}{figure}{腰 \dpy{yao1}}

\begin{EntryWithPhonetic}{腰}{yao1}{13}{⾁}[HSK 4]
  \definition*{s.}{Sobrenome: Yao}
  \definition[个,尺]{s.}{cintura; região lombar | cós | bolso | parte do meio das coisas | lombo}
\end{EntryWithPhonetic}

\begin{EntryWithPhonetic}{腰包}{yao1bao1}{13,5}{⾁,⼓}
  \definition{s.}{pochete | bolso}
\end{EntryWithPhonetic}

\begin{EntryWithPhonetic}{腰椎}{yao1zhui1}{13,12}{⾁,⽊}
  \definition{s.}{vértebra lombar (espinha dorsal inferior)}
\end{EntryWithPhonetic}

%%%%%%%%%% 邀 %%%%%%%%%%
\subsection*{邀}\addcontentsline{loh}{figure}{邀 \dpy{yao1}}

\begin{EntryWithPhonetic}{邀}{yao1}{16}{⾡}[HSK 7-9]
  \definition{v.}{convidar | solicitar; procurar; requerer | parar; conter; interceptar}
  \seealsoref{请}{qing3}
  \seealsoref{邀请}{yao1qing3}
\end{EntryWithPhonetic}

\begin{EntryWithPhonetic}{邀请}{yao1qing3}{16,10}{⾡,⾔}[HSK 5]
  \definition[份,个]{s.}{convite}
  \definition{v.}{convidar; solicitar; convidar pessoas para irem à sua casa ou a um local combinado}
\end{EntryWithPhonetic}

%%%%%%%%%% 尧 %%%%%%%%%%
\subsection*{尧}\addcontentsline{loh}{figure}{尧 \dpy{yao2}}

\begin{EntryWithPhonetic}{尧}{yao2}{6}{⼪}
  \definition*{s.}{Yao, um monarca lendário da China antiga | Sobrenome: Yao}
\end{EntryWithPhonetic}

%%%%%%%%%% 陶 %%%%%%%%%%
\subsection*{陶}\addcontentsline{loh}{figure}{陶 \dpy{yao2}}

\begin{EntryWithPhonetic}{陶}{yao2}{10}{⾩}
  \definition*{s.}{Cao Dao (nome de uma pessoa antiga); Gaotao: Um nome próprio dos tempos antigos}[皋陶是舜的臣子。===Gao Yao era um ministro de Shun.]
  \seeref{tao2}
\end{EntryWithPhonetic}

%%%%%%%%%% 窑 %%%%%%%%%%
\subsection*{窑}\addcontentsline{loh}{figure}{窑 \dpy{yao2}}

\begin{EntryWithPhonetic}{窑}{yao2}{11}{⽳}[HSK 7-9]
  \definition[个]{s.}{forno | mina de carvão; mina de carvão produzida utilizando métodos tradicionais | habitação em caverna | Dialeto: olaria}
\end{EntryWithPhonetic}

%%%%%%%%%% 谣 %%%%%%%%%%
\subsection*{谣}\addcontentsline{loh}{figure}{谣 \dpy{yao2}}

\begin{EntryWithPhonetic}{谣}{yao2}{12}{⾔}
  \definition*{s.}{Sobrenome: Yao}
  \definition[派,个]{s.}{balada; rima; canções folclóricas | boato; rumor}
\end{EntryWithPhonetic}

\begin{EntryWithPhonetic}{谣言}{yao2yan2}{12,7}{⾔,⾔}[HSK 7-9]
  \definition[个,种]{s.}{rumor; boato; fofoca; alegação infundada}
  \seealsoref{留言}{liu2yan2}
  \synonymref{谎言}{huang3yan2}
  \antonymref{实话}{shi2hua4}
\end{EntryWithPhonetic}

%%%%%%%%%% 摇 %%%%%%%%%%
\subsection*{摇}\addcontentsline{loh}{figure}{摇 \dpy{yao2}}

\begin{EntryWithPhonetic}{摇}{yao2}{13}{⼿}[HSK 4]
  \definition{v.}{chacoalhar; ondular; balançar; fazer com que um objeto se mova para frente e para trás | agitar algo | sacudir; chacoalhar; agitar algo para que se mova}
\end{EntryWithPhonetic}

\begin{EntryWithPhonetic}{摇摆}{yao2bai3}{13,13}{⼿,⼿}[HSK 7-9]
  \definition{v.}{balançar; oscilar; empinar; mover ou alternar entre direções opostas | vacilar; oscilar}
  \seealsoref{摇晃}{yao2huang4}
  \synonymref{摆动}{bai3dong4}
  \synonymref{动摇}{dong4yao2}
  \synonymref{摇晃}{yao2huang4}
\end{EntryWithPhonetic}

\begin{EntryWithPhonetic}{摇滚}{yao2gun3}{13,13}{⼿,⽔}[HSK 7-9]
  \definition{s.}{\emph{rock'n roll}; a música \emph{rock}, um gênero popular no Ocidente, é caracterizada por seu som rico e ritmo forte}
\end{EntryWithPhonetic}

\begin{EntryWithPhonetic}{摇晃}{yao2huang4}{13,10}{⼿,⽇}[HSK 7-9]
  \definition{v.}{balançar; oscilar; tremer; estremecer | agitar; sacudir}
  \seealsoref{摇摆}{yao2bai3}
  \synonymref{摆动}{bai3dong4}
  \synonymref{动摇}{dong4yao2}
  \synonymref{晃荡}{huang4dang5}
  \synonymref{摇摆}{yao2bai3}
  \antonymref{安稳}{an1wen3}
  \antonymref{固定}{gu4ding4}
  \antonymref{静止}{jing4zhi3}
  \antonymref{平衡}{ping2heng2}
  \antonymref{平稳}{ping2wen3}
  \antonymref{稳定}{wen3ding4}
  \antonymref{稳固}{wen3gu4}
\end{EntryWithPhonetic}

\begin{EntryWithPhonetic}{摇篮}{yao2lan2}{13,16}{⼿,⽵}[HSK 7-9]
  \definition{s.}{berço; enxoval para bebê em formato de cesta, que pode balançar de um lado para o outro | o berço de algo; uma metáfora para o lugar onde o talento floresce ou a origem de coisas importantes}
\end{EntryWithPhonetic}

\begin{EntryWithPhonetic}{摇头}{yao2/tou2}{13,5}{⼿,⼤}[HSK 5]
  \definition{v.+compl.}{sacudir; balançar a cabeça; balançar a cabeça para a esquerda e para a direita, indicando negação, desacordo ou impedimento}
\end{EntryWithPhonetic}

\begin{EntryWithPhonetic}{摇摇欲坠}{yao2yao2-yu4zhui4}{13,13,11,7}{⼿,⼿,⽋,⼟}[HSK 7-9]
  \definition{expr.}{cambaleante; em ruínas; à beira do colapso}
  \antonymref{根深蒂固}{gen1shen1-di4gu4}
\end{EntryWithPhonetic}

%%%%%%%%%% 遥 %%%%%%%%%%
\subsection*{遥}\addcontentsline{loh}{figure}{遥 \dpy{yao2}}

\begin{EntryWithPhonetic}{遥}{yao2}{13}{⾡}
  \definition{adj.}{distante; remoto; longe}
\end{EntryWithPhonetic}

\begin{EntryWithPhonetic}{遥控}{yao2kong4}{13,11}{⾡,⼿}[HSK 7-9]
  \definition{s.}{controle remoto}
  \definition{v.}{controlar remotamente; dirigir operações de um local remoto; controlar dispositivos e máquinas a uma certa distância usando tecnologia | manter o controle remoto; metáfora para controlar ou comandar à distância}
  \synonymref{操控}{cao1kong4}
\end{EntryWithPhonetic}

\begin{EntryWithPhonetic}{遥远}{yao2yuan3}{13,7}{⾡,⾡}[HSK 7-9]
  \definition{adj.}{distante; remoto; (distância no tempo ou no espaço) muito longe}
  \seealsoref{远}{yuan3}
  \synonymref{边远}{bian1yuan3}
  \synonymref{长远}{chang2yuan3}
  \synonymref{日后}{ri4hou4}
  \synonymref{远处}{yuan3chu4}
  \antonymref{附近}{fu4jin4}
  \antonymref{临近}{lin2jin4}
  \antonymref{眼前}{yan3qian2}
\end{EntryWithPhonetic}

%%%%%%%%%% 咬 %%%%%%%%%%
\subsection*{咬}\addcontentsline{loh}{figure}{咬 \dpy{yao3}}

\begin{EntryWithPhonetic}{咬}{yao3}{9}{⼝}[HSK 5]
  \definition{v.}{morder; estalar; pressionar os dentes superiores e inferiores com força | latir | agarrar; morder | incriminar outra pessoa (geralmente inocente) quando culpada ou interrogada | pronunciar; articular; pronunciar corretamente | corroer (metais); irritar (a pele) | ser minucioso (com relação ao uso de palavras) | aproximar"-se de; pressionar em direção a; avançar sobre}
\end{EntryWithPhonetic}

%%%%%%%%%% 药 %%%%%%%%%%
\subsection*{药}\addcontentsline{loh}{figure}{药 \dpy{yao4}}

\begin{EntryWithPhonetic}{药}{yao4}{9}{⾋}[HSK 2]
  \definition*{s.}{Sobrenome: Yao}
  \definition[片,粒,颗,瓶,服]{s.}{droga; loção; remédio; medicamento; substâncias que podem prevenir e tratar doenças, pragas ou melhorar funções corporais | certos produtos químicos com efeitos específicos}
  \definition{v.}{curar com remédios; tomar remédios para tratar doenças | matar com veneno; envenenar}
\end{EntryWithPhonetic}

\begin{EntryWithPhonetic}{药补}{yao4bu3}{9,7}{⾋,⾐}
  \definition{s.}{suplemento dietético medicinal que ajuda a melhorar a saúde}
\end{EntryWithPhonetic}

\begin{EntryWithPhonetic}{药材}{yao4cai2}{9,7}{⾋,⽊}[HSK 7-9]
  \definition{s.}{materiais medicinais; drogas brutas | ingrediente medicinal; ervas medicinais chinesas que foram cortadas e processadas para decocção}
\end{EntryWithPhonetic}

\begin{EntryWithPhonetic}{药典}{yao4dian3}{9,8}{⾋,⼋}
  \definition{s.}{farmacopéia}
\end{EntryWithPhonetic}

\begin{EntryWithPhonetic}{药店}{yao4dian4}{9,8}{⾋,⼴}[HSK 2]
  \definition[家]{s.}{farmácia; drogaria; lojas que vendem medicamentos}
\end{EntryWithPhonetic}

\begin{EntryWithPhonetic}{药方}{yao4fang1}{9,4}{⾋,⽅}[HSK 7-9]
  \definition[个,些,种]{s.}{prescrição; os nomes, dosagens e usos de vários medicamentos combinados para tratar uma determinada doença | receita médica; o papel com a receita médica escrita nele}
  \synonymref{药品}{yao4pin3}
\end{EntryWithPhonetic}

\begin{EntryWithPhonetic}{药房}{yao4fang2}{9,8}{⾋,⼾}
  \definition{s.}{farmácia | drogaria}
\end{EntryWithPhonetic}

\begin{EntryWithPhonetic}{药罐}{yao4guan4}{9,23}{⾋,⽸}
  \definition{s.}{frasco de remédio; pote de remédio}
\end{EntryWithPhonetic}

\begin{EntryWithPhonetic}{药片}{yao4pian4}{9,4}{⾋,⽚}[HSK 2]
  \definition[颗,片]{s.}{pílula; comprimido; preparações em comprimidos}
\end{EntryWithPhonetic}

\begin{EntryWithPhonetic}{药品}{yao4pin3}{9,9}{⾋,⼝}[HSK 6]
  \definition[个,些,种,类,批]{s.}{medicamentos e reagentes químicos; um termo geral para vários medicamentos e reagentes químicos}
\end{EntryWithPhonetic}

\begin{EntryWithPhonetic}{药签}{yao4qian1}{9,13}{⾋,⽵}
  \definition{s.}{cotonete médico}
\end{EntryWithPhonetic}

\begin{EntryWithPhonetic}{药膳}{yao4shan4}{9,16}{⾋,⾁}
  \definition{s.}{alimentos medicamentosos; alimentos cozidos com ervas medicinais | cozinha medicinal}
\end{EntryWithPhonetic}

\begin{EntryWithPhonetic}{药水}{yao4shui3}{9,4}{⾋,⽔}[HSK 2]
  \definition*{s.}{Yaksu na Coreia do Norte, perto da fronteira com Liaoning e a província de Jilin}
  \definition{s.}{medicamento líquido; líquido medicinal | loção | remédio engarrafado | medicamento em forma líquida}
\end{EntryWithPhonetic}

\begin{EntryWithPhonetic}{药丸}{yao4wan2}{9,3}{⾋,⼂}
  \definition[粒]{s.}{pílula}
\end{EntryWithPhonetic}

\begin{EntryWithPhonetic}{药物}{yao4wu4}{9,8}{⾋,⽜}[HSK 4]
  \definition[种]{s.}{droga; medicamento; remédio; substâncias que controlam doenças, pragas, etc.}
\end{EntryWithPhonetic}

%%%%%%%%%% 要 %%%%%%%%%%
\subsection*{要}\addcontentsline{loh}{figure}{要 \dpy{yao4}}

\begin{EntryWithPhonetic}{要}{yao4}{9}{⾑}[HSK 1,4]
  \definition{adj.}{importante; essencial}
  \definition{conj.}{suponha; no caso; se, indicando um relacionamento hipotético | ou; ou\dots ou\dots}
  \definition{s.}{ponto principal; manchete; conteúdo importante}
  \definition{v.}{querer; desejar; pensar | querer; pedir; deseja; querer obter; querer manter | recuperar algo; dizer a alguém para guardar algo para você ou devolver | pedir (ou querer) que alguém faça algo; pedir a alguém para fazer algo, quando usado para conseguir que alguém faça algo, tem um tom de comando e pode ser indelicado | precisar; tomar; pegar | deve; deveria; é necessário (imperativo, essencial) que\dots | estar indo para | querer; ter um desejo por; expressar determinação ou desejo de fazer algo | poder; dever;  indica uma estimativa, usada para comparação}
  \seeref{yao1}
  \seealsoref{要是}{yao4shi5}
\end{EntryWithPhonetic}

\begin{EntryWithPhonetic}{要不是}{yao4bu2shi4}{9,4,9}{⾑,⼀,⽇}[HSK 7-9]
  \definition{conj.}{se não fosse por; as orações de ligação, usadas na primeira oração, indicam uma relação hipotética, equivalente a 如果不是}
\end{EntryWithPhonetic}

\begin{EntryWithPhonetic}{要不}{yao4bu4}{9,4}{⾑,⼀}[HSK 7-9]
  \definition{conj.}{ou então; caso contrário; se você não fizer isso (haverá um resultado ruim) | usado para propor educadamente; usado para fazer uma sugestão educadamente | ou; se você não fizer isso, faça aquilo}
  \seealsoref{不然}{bu4ran2}
  \synonymref{要么}{yao4me5}
\end{EntryWithPhonetic}

\begin{EntryWithPhonetic}{要不然}{yao4bu5ran2}{9,4,12}{⾑,⼀,⽕}[HSK 6]
  \definition{conj.}{caso contrário; ou então; se você não fizer isso (haverá um resultado ruim) | ou então; usado entre duas frases em um relacionamento de escolha; significa escolher uma entre as duas; equivalente a 要不}
  \seealsoref{要不}{yao4bu4}
\end{EntryWithPhonetic}

\begin{EntryWithPhonetic}{要点}{yao4dian3}{9,9}{⾑,⽕}[HSK 7-9]
  \definition[个]{s.}{pontos principais; ponto"-chave; conteúdo importante; partes principais | fortaleza; posição"-chave; ponto estratégico; fortalezas importantes}
  \synonymref{亮点}{liang4dian3}
  \synonymref{重点}{zhong4dian3}
\end{EntryWithPhonetic}

\begin{EntryWithPhonetic}{要害}{yao4hai4}{9,10}{⾑,⼧}[HSK 7-9]
  \definition{s.}{parte vital; ponto crucial; vulnerabilidade; partes fatais do corpo | ponto estratégico; metaforicamente, refere"-se a uma parte importante ou a um local de importância militar}
  \synonymref{把柄}{ba3bing3}
  \synonymref{关键}{guan1jian4}
\end{EntryWithPhonetic}

\begin{EntryWithPhonetic}{要好}{yao4hao3}{9,6}{⾑,⼥}[HSK 6]
  \definition{adj.}{estar em bons termos; ser amigos próximos; relacionamento harmonioso | estar ansioso para melhorar a si mesmo; esforçar"-se para progredir | ansioso para melhorar a si mesmo; esforçar"-se para progredir}
\end{EntryWithPhonetic}

\begin{EntryWithPhonetic}{要谎}{yao4huang3}{9,11}{⾑,⾔}
  \definition{v.}{pedir um preço enorme (como primeiro passo de negociação)}
\end{EntryWithPhonetic}

\begin{EntryWithPhonetic}{要紧}{yao4jin3}{9,10}{⾑,⽷}[HSK 7-9]
  \definition{adj.}{importante; essencial | crítico; sério}
  \definition{adv.}{com pressa; imediatamente;}
  \synonymref{急迫}{ji2po4}
  \synonymref{紧急}{jin3ji2}
  \synonymref{紧迫}{jin3po4}
  \synonymref{迫切}{po4qie4}
  \synonymref{首要}{shou3yao4}
  \synonymref{危机}{wei1ji1}
  \synonymref{严重}{yan2zhong4}
  \synonymref{重要}{zhong4yao4}
  \synonymref{主要}{zhu3yao4}
\end{EntryWithPhonetic}

\begin{EntryWithPhonetic}{要领}{yao4ling3}{9,11}{⾑,⾴}[HSK 7-9]
  \definition{s.}{pontos principais; essenciais; essência; pontos"-chave}
  \synonymref{办法}{ban4fa3}
  \synonymref{措施}{cuo4shi1}
  \synonymref{方法}{fang1fa3}
  \synonymref{手段}{shou3duan4}
  \synonymref{手腕}{shou3wan4}
\end{EntryWithPhonetic}

\begin{EntryWithPhonetic}{要么}{yao4me5}{9,3}{⾑,⼃}[HSK 6]
  \definition{conj.}{ou; ou\dots ou\dots; indica uma escolha entre duas situações ou dois desejos}
  \seealsoref{要么……要么……}{yao4me5 yao4me5}
\end{EntryWithPhonetic}

\begin{EntryWithPhonetic}{要么……要么……}{yao4me5 yao4me5}{9,3,9,3}{⾑,⼃,⾑,⼃}
  \definition{conj.}{ou\dots ou\dots}
  \seealsoref{要么}{yao4me5}
\end{EntryWithPhonetic}

\begin{EntryWithPhonetic}{要命}{yao4/ming4}{9,8}{⾑,⼝}[HSK 7-9]
  \definition{adv.}{extremamente; terrivelmente; em grau extremo}
  \definition{v.}{matar; levar alguém à morte | ser um incômodo (usado para reclamar de algo); causar sérios problemas a alguém (quando em estado de ansiedade ou queixa)}
\end{EntryWithPhonetic}

\begin{EntryWithPhonetic}{要强}{yao4qiang2}{9,12}{⾑,⼸}[HSK 7-9]
  \definition{adj.}{ambicioso; ávido por se destacar; empreendedor; altamente competitivo e relutante em ficar para trás dos outros}
  \synonymref{当然}{dang1ran2}
\end{EntryWithPhonetic}

\begin{EntryWithPhonetic}{要是}{yao4shi5}{9,9}{⾑,⽇}[HSK 3]
  \definition{conj.}{se; suponha; no caso de; conecta frases, expressa uma relação hipotética, equivalente a 如果, e pode ser usado em conjunto com 的话}
  \seealsoref{的话}{de5hua4}
  \seealsoref{如果}{ru2guo3}
\end{EntryWithPhonetic}

\begin{EntryWithPhonetic}{要是……的话}{yao4shi5 de5hua4}{9,9,8,8}{⾑,⽇,⽩,⾔}
  \definition{conj.}{se for assim\dots}
\end{EntryWithPhonetic}

\begin{EntryWithPhonetic}{要死}{yao4si3}{9,6}{⾑,⽍}
  \definition{adv.}{extremamente | muito}
\end{EntryWithPhonetic}

\begin{EntryWithPhonetic}{要素}{yao4su4}{9,10}{⾑,⽷}[HSK 6]
  \definition[个]{s.}{fator essencial; elemento"-chave; os elementos essenciais que compõem as coisas}
\end{EntryWithPhonetic}

\begin{EntryWithPhonetic}{要义}{yao4yi4}{9,3}{⾑,⼂}
  \definition{s.}{resumo | o essencial}
\end{EntryWithPhonetic}

%%%%%%%%%% 钥 %%%%%%%%%%
\subsection*{钥}\addcontentsline{loh}{figure}{钥 \dpy{yao4}}

\begin{EntryWithPhonetic}{钥}{yao4}{9}{⾦}
  \definition{s.}{chave}
\end{EntryWithPhonetic}

\begin{EntryWithPhonetic}{钥匙}{yao4shi5}{9,11}{⾦,⼔}[HSK 7-9]
  \definition[把,串,个]{s.}{chave; esta é uma ferramenta usada para destrancar fechaduras; algumas fechaduras só podem ser trancadas com ela}
\end{EntryWithPhonetic}

\begin{EntryWithPhonetic}{钥匙洞孔}{yao4shi5dong4kong3}{9,11,9,4}{⾦,⼔,⽔,⼦}
  \definition{s.}{buraco da fechadura}
\end{EntryWithPhonetic}

\begin{EntryWithPhonetic}{钥匙卡}{yao4shi5ka3}{9,11,5}{⾦,⼔,⼘}
  \definition{s.}{cartão de acesso | cartão"-chave; cartão magnético}
\end{EntryWithPhonetic}

\begin{EntryWithPhonetic}{钥匙孔}{yao4shi5kong3}{9,11,4}{⾦,⼔,⼦}
  \definition{s.}{buraco da fechadura}
\end{EntryWithPhonetic}

\begin{EntryWithPhonetic}{钥匙圈}{yao4shi5quan1}{9,11,11}{⾦,⼔,⼞}
  \definition{s.}{chaveiro}
\end{EntryWithPhonetic}

%%%%%%%%%% 耀 %%%%%%%%%%
\subsection*{耀}\addcontentsline{loh}{figure}{耀 \dpy{yao4}}

\begin{EntryWithPhonetic}{耀}{yao4}{20}{⽻}
  \definition{s.}{brilho; raios de luz | honra; crédito | brilho; orgulho; glória | brilho; luz; radiância, esplendor}
  \definition{v.}{brilhar; iluminar; deslumbrar | vangloriar"-se de; louvar | gabar"-se | ser deslumbrante; ser ofuscante}
  \synonymref{荣}{rong2}
\end{EntryWithPhonetic}

\begin{EntryWithPhonetic}{耀眼}{yao4yan3}{20,11}{⽻,⽬}[HSK 7-9]
  \definition{adj.}{(luz) deslumbrante; ofuscante; a luz era muito forte e deixava as pessoas com tontura | deslumbrante; ofuscante; descrever algo como deslumbrante e chamativo}
  \synonymref{闪耀}{shan3yao4}
  \synonymref{醒目}{xing3mu4}
\end{EntryWithPhonetic}

%%%%%%%%%% 椰 %%%%%%%%%%
\subsection*{椰}\addcontentsline{loh}{figure}{椰 \dpy{ye1}}

\begin{EntryWithPhonetic}{椰}{ye1}{12}{⽊}
  \definition[只,棵]{s.}{coqueiro; coco}
\end{EntryWithPhonetic}

\begin{EntryWithPhonetic}{椰汁}{ye1zhi1}{12,5}{⽊,⽔}
  \definition{s.}{água de coco}
\end{EntryWithPhonetic}

\begin{EntryWithPhonetic}{椰子}{ye1zi5}{12,3}{⽊,⼦}[HSK 7-9]
  \definition[个,棵,片]{s.}{palmeira de coco; coqueiro; essa espécie de árvore cresce em regiões tropicais e subtropicais; possui tronco reto, sem galhos, e folhas que crescem no topo | coco; o fruto desta planta tem formato oval, com casca externa marrom"-amarelada e polpa branca e suculenta; pode ser prensado para extração de óleo ou usado para fazer bebidas}
\end{EntryWithPhonetic}

%%%%%%%%%% 噎 %%%%%%%%%%
\subsection*{噎}\addcontentsline{loh}{figure}{噎 \dpy{ye1}}

\begin{EntryWithPhonetic}{噎}{ye1}{15}{⼝}
  \definition{v.}{engasgar | sufocar}
\end{EntryWithPhonetic}

%%%%%%%%%% 爷 %%%%%%%%%%
\subsection*{爷}\addcontentsline{loh}{figure}{爷 \dpy{ye2}}

\begin{EntryWithPhonetic}{爷}{ye2}{6}{⽗}
  \definition[个,位,名,些]{s.}{(dialeto) pai | (dialeto) avô | (uma forma respeitosa de se dirigir a um homem idoso) tio | (uma forma de se dirigir a um oficial ou homem rico) senhor; mestre; lorde; o antigo nome para burocratas, pessoas ricas, etc. | deus; forma de tratamento de um adorador para um deus}
\end{EntryWithPhonetic}

\begin{EntryWithPhonetic}{爷爷}{ye2ye5}{6,6}{⽗,⽗}[HSK 1]
  \definition[个,位]{s.}{avô (paterno)}
\end{EntryWithPhonetic}

%%%%%%%%%% 也 %%%%%%%%%%
\subsection*{也}\addcontentsline{loh}{figure}{也 \dpy{ye3}}

\begin{EntryWithPhonetic}{也}{ye3}{3}{⼄}[HSK 1]
  \definition*{s.}{Sobrenome: Ye}
  \definition{adv.}{também; igualmente; assim como; da mesma forma; usado em frases simples, implica que é igual a outra coisa | assim como (expressar ênfase) | (expressar que as consequências são as mesmas) | também (expressar ufemismo; expressar um tom diplomático) | usado em frases compostas paralelas, indica que duas ou mais coisas têm algo em comum (pode ser usado em todas as frases ou apenas na última frase)}
  \definition{part.}{usado no meio de uma frase, destacando um elemento da frase sobre o qual deve ser feita uma afirmação | usado no final de uma frase, indicando uma explicação ou um julgamento; usado no final da frase, indica tom afirmativo e também pode reforçar o tom interrogativo, exclamativo ou imperativo}
  \synonymref{还}{hai2}
  \synonymref{亦}{yi4}
\end{EntryWithPhonetic}

\begin{EntryWithPhonetic}{也好}{ye3hao3}{3,6}{⼄,⼥}[HSK 5]
  \definition{part.}{pode não ser uma má ideia; também pode ser | (reduplicado) se\dots ou\dots; não importa se | pode não ser uma má ideia | se\dots ou\dots; usado em conjunto, significa que não está condicionado a uma determinada situação}
\end{EntryWithPhonetic}

\begin{EntryWithPhonetic}{也就是}{ye3jiu4shi4}{3,12,9}{⼄,⼪,⽇}
  \definition{adv.}{i.e., isso é | ou seja}
\end{EntryWithPhonetic}

\begin{EntryWithPhonetic}{也就是说}{ye3jiu4shi4shuo1}{3,12,9,9}{⼄,⼪,⽇,⾔}[HSK 7-9]
  \definition{adv.}{em outras palavras | então | isto é; ou seja | por isso; assim}
\end{EntryWithPhonetic}

\begin{EntryWithPhonetic}{也许}{ye3xu3}{3,6}{⼄,⾔}[HSK 2]
  \definition{adv.}{talvez; provavelmente; estou com medo; para expressar incerteza; para expressar uma alta probabilidade}
  \synonymref{或许}{huo4xu3}
  \synonymref{或者}{huo4zhe3}
\end{EntryWithPhonetic}

\begin{EntryWithPhonetic}{也有今天}{ye3you3jin1tian1}{3,6,4,4}{⼄,⽉,⼈,⼤}
  \definition{expr.}{obter apenas o que merece | todo cachorro tem seu dia | obter a sua parte (coisas boas ou ruins) | servir alguém bem}
\end{EntryWithPhonetic}

%%%%%%%%%% 野 %%%%%%%%%%
\subsection*{野}\addcontentsline{loh}{figure}{野 \dpy{ye3}}

\begin{EntryWithPhonetic}{野}{ye3}{11}{⾥}[HSK 6]
  \definition*{s.}{Sobrenome: Ye}
  \definition{adj.}{(de plantas ou animais) selvagem; incultivado; não domesticado; indomável (opp. 家) | rude; áspero | desenfreado; abandonado; indisciplinado | ilícito; sem licença}
  \definition{s.}{espaço aberto; o aberto | limite; fronteira | não está no poder; fora do cargo}
  \seealsoref{家}{jia1}
\end{EntryWithPhonetic}

\begin{EntryWithPhonetic}{野餐}{ye3can1}{11,16}{⾥,⾷}[HSK 7-9]
  \definition[个,次]{s.}{piquenique; diversas opções de comida para refeições ao ar livre}
  \definition{v.}{fazer um piquenique; ir a um piquenique; comer em meio à natureza selvagem}
\end{EntryWithPhonetic}

\begin{EntryWithPhonetic}{野炊}{ye3chui1}{11,8}{⾥,⽕}[HSK 7-9]
  \definition{v.}{cozinhar ao ar livre; cozinhar sobre uma fogueira na natureza selvagem}
\end{EntryWithPhonetic}

\begin{EntryWithPhonetic}{野蛮}{ye3man2}{11,12}{⾥,⾍}[HSK 7-9]
  \definition{adj.}{selvagem; incivilizado; inculto | cruel; brutal; bárbaro; desumano}
  \synonymref{残忍}{can2ren3}
  \synonymref{粗暴}{cu1bao4}
  \synonymref{粗鲁}{cu1lu3}
  \synonymref{凶恶}{xiong1'e4}
  \antonymref{仁慈}{ren2ci2}
  \antonymref{温柔}{wen1rou2}
  \antonymref{文明}{wen2ming2}
\end{EntryWithPhonetic}

\begin{EntryWithPhonetic}{野生}{ye3sheng1}{11,5}{⾥,⽣}[HSK 6]
  \definition{adj.}{selvagem; não cultivado; não domesticado}
\end{EntryWithPhonetic}

\begin{EntryWithPhonetic}{野兽}{ye3shou4}{11,11}{⾥,⼋}[HSK 7-9]
  \definition[头]{s.}{besta; animal selvagem; mamíferos que não sejam animais de criação}
  \antonymref{美女}{mei3nv3}
\end{EntryWithPhonetic}

\begin{EntryWithPhonetic}{野外}{ye3wai4}{11,5}{⾥,⼣}[HSK 7-9]
  \definition{s.}{campo; ar livre; campo aberto; lugares distantes de áreas residenciais}
  \synonymref{郊外}{jiao1wai4}
  \synonymref{旷野}{kuang4ye3}
\end{EntryWithPhonetic}

\begin{EntryWithPhonetic}{野心}{ye3xin1}{11,4}{⾥,⼼}[HSK 7-9]
  \definition{s.}{carreirismo; ambição desmedida; plano insensato; um desejo enorme e excessivo por território, poder ou fama}
  \synonymref{打算}{da3suan5}
  \synonymref{计划}{ji4hua4}
  \synonymref{盘算}{pan2suan5}
  \synonymref{企图}{qi3tu2}
  \synonymref{妄想}{wang4xiang3}
  \synonymref{希望}{xi1wang4}
  \synonymref{阴谋}{yin1mou2}
\end{EntryWithPhonetic}

\begin{EntryWithPhonetic}{野营}{ye3ying2}{11,11}{⾥,⾋}[HSK 7-9]
  \definition{v.}{acampar; ficar em um acampamento temporário sem abrigo; acampar e praticar certas atividades na natureza}
\end{EntryWithPhonetic}

%%%%%%%%%% 业 %%%%%%%%%%
\subsection*{业}\addcontentsline{loh}{figure}{业 \dpy{ye4}}

\begin{EntryWithPhonetic}{业}{ye4}{5}{⼀}[HSK 7-9]
  \definition*{s.}{Sobrenome: Ye}
  \definition{adv.}{já; indica que a ação foi concluída, equivalente a 已经}
  \definition{s.}{comércio; indústria; ramo de negócios | emprego; ocupação; profissão | curso de estudo | causa; empreendimento | propriedade | carma; o budismo se refere a todas as ações, palavras e pensamentos humanos como carma, que são chamados de carma corporal, carma da fala e carma mental; o carma inclui aspectos bons e ruins, geralmente referindo"-se ao destino ou ao pecado}
  \definition{v.}{envolver"-se em; exercer uma determinada profissão}
  \seealsoref{已经}{yi3jing1}
\end{EntryWithPhonetic}

\begin{EntryWithPhonetic}{业绩}{ye4ji4}{5,11}{⼀,⽷}[HSK 7-9]
  \definition[项]{s.}{conquista excepcional; feito exemplar; os méritos alcançados e os empreendimentos concluídos; as principais realizações}
  \synonymref{分数}{fen1shu4}
  \synonymref{功绩}{gong1ji4}
  \synonymref{事迹}{shi4ji4}
\end{EntryWithPhonetic}

\begin{EntryWithPhonetic}{业务}{ye4wu4}{5,5}{⼀,⼒}[HSK 5]
  \definition[项,笔,个,类,种]{s.}{negócios; trabalho vocacional; trabalho profissional}
  \synonymref{交易}{jiao1yi4}
  \synonymref{生意}{sheng1yi5}
  \synonymref{营业}{ying2ye4}
\end{EntryWithPhonetic}

\begin{EntryWithPhonetic}{业余}{ye4yu2}{5,7}{⼀,⼈}[HSK 4]
  \definition{adj.}{tempo livre; depois do expediente; fora do horário de trabalho | amador; não profissional}
  \antonymref{职业}{zhi2ye4}
  \antonymref{专业}{zhuan1ye4}
\end{EntryWithPhonetic}

%%%%%%%%%% 叶 %%%%%%%%%%
\subsection*{叶}\addcontentsline{loh}{figure}{叶 \dpy{ye4}}

\begin{EntryWithPhonetic}{叶}{ye4}{5}{⼝}
  \definition*{s.}{Sobrenome: Ye}
  \definition[枝]{s.}{folha; folhagem | coisa parecida com uma folha | página; folha | parte de um período histórico; segmentos de período mais longos | lóbulo; lóbulos do cérebro, pulmões e fígado}
\end{EntryWithPhonetic}

\begin{EntryWithPhonetic}{叶子}{ye4zi5}{5,3}{⼝,⼦}[HSK 4]
  \definition[片]{s.}{folha; termo genérico para as folhas de uma planta}
\end{EntryWithPhonetic}

%%%%%%%%%% 页 %%%%%%%%%%
\subsection*{页}\addcontentsline{loh}{figure}{页 \dpy{ye4}}

\begin{EntryWithPhonetic}{页}{ye4}{6}{⾴}[HSK 1][Kangxi 181]
  \definition{clas.}{página; folha de papel; lâmina; antigamente, referia"-se a uma folha de um livro encadernado; atualmente, refere"-se a uma das faces de um livro impresso em ambos os lados}
  \definition{s.}{página; folha de papel; folhas soltas de um livro}
\end{EntryWithPhonetic}

%%%%%%%%%% 夜 %%%%%%%%%%
\subsection*{夜}\addcontentsline{loh}{figure}{夜 \dpy{ye4}}

\begin{EntryWithPhonetic}{夜}{ye4}{8}{⼣}[HSK 2]
  \definition{s.}{noite; tarde; noturno; o período do anoitecer ao amanhecer; em meteorologia, refere"-se especificamente ao período das 20h do dia atual às 8h do dia seguinte}
  \antonymref{日}{ri4}
  \antonymref{昼}{zhou4}
\end{EntryWithPhonetic}

\begin{EntryWithPhonetic}{夜班}{ye4ban1}{8,10}{⼣,⽟}[HSK 7-9]
  \definition[个]{s.}{turno da noite; turnos noturnos}
\end{EntryWithPhonetic}

\begin{EntryWithPhonetic}{夜场}{ye4chang3}{8,6}{⼣,⼟}
  \definition{s.}{show noturno (em um teatro, etc.) | local de entretenimento noturno (bar, boate, discoteca, etc.)}
\end{EntryWithPhonetic}

\begin{EntryWithPhonetic}{夜店}{ye4dian4}{8,8}{⼣,⼴}
  \definition{s.}{boate | \emph{nightclub}}
\end{EntryWithPhonetic}

\begin{EntryWithPhonetic}{夜间}{ye4jian5}{8,7}{⼣,⾨}[HSK 5]
  \definition{s.}{noite; à noite; noturno; durante a noite}
\end{EntryWithPhonetic}

\begin{EntryWithPhonetic}{夜里}{ye4li5}{8,7}{⼣,⾥}[HSK 2]
  \definition{s.}{noturno; à noite; o período do anoitecer ao amanhecer}
\end{EntryWithPhonetic}

\begin{EntryWithPhonetic}{夜幕}{ye4mu4}{8,13}{⼣,⼱}
  \definition{s.}{escuridão crescente; cortina da noite; o véu da noite}
\end{EntryWithPhonetic}

\begin{EntryWithPhonetic}{夜鸟}{ye4niao3}{8,5}{⼣,⿃}
  \definition{s.}{ave noturna}
\end{EntryWithPhonetic}

\begin{EntryWithPhonetic}{夜深人静}{ye4shen1ren2jing4}{8,11,2,14}{⼣,⽔,⼈,⾭}
  \definition{expr.}{``Na calada da noite.''; ``No silêncio da noite.''}
\end{EntryWithPhonetic}

\begin{EntryWithPhonetic}{夜生活}{ye4sheng1huo2}{8,5,9}{⼣,⽣,⽔}
  \definition{s.}{vida noturna}
\end{EntryWithPhonetic}

\begin{EntryWithPhonetic}{夜市}{ye4shi4}{8,5}{⼣,⼱}[HSK 7-9]
  \definition{s.}{feira noturna; mercado noturno; mercados que funcionam à noite | negócios à noite; atividades comerciais noturnas}
\end{EntryWithPhonetic}

\begin{EntryWithPhonetic}{夜晚}{ye4wan3}{8,11}{⼣,⽇}[HSK 7-9]
  \definition[个]{s.}{noite; entardecer}
  \seealsoref{夜里}{ye4li5}
  \seealsoref{夜间}{ye4jian5}
\end{EntryWithPhonetic}

\begin{EntryWithPhonetic}{夜校}{ye4xiao4}{8,10}{⼣,⽊}[HSK 7-9]
  \definition{s.}{escola noturna (ou vespertina)}
\end{EntryWithPhonetic}

\begin{EntryWithPhonetic}{夜夜}{ye4ye4}{8,8}{⼣,⼣}
  \definition{adv.}{toda noite}
\end{EntryWithPhonetic}

\begin{EntryWithPhonetic}{夜以继日}{ye4yi3ji4ri4}{8,4,10,4}{⼣,⼈,⽷,⽇}[HSK 7-9]
  \definition{expr.}{24 horas por dia; dia e noite; prolongar o dia incluindo a noite; dia após dia; estender o dia incluindo a noite}
  \synonymref{废寝忘食}{fei4qin3-wang4shi2}
  \antonymref{无所事事}{wu2suo3shi4shi4}
\end{EntryWithPhonetic}

\begin{EntryWithPhonetic}{夜总会}{ye4zong3hui4}{8,9,6}{⼣,⼼,⼈}[HSK 7-9]
  \definition{s.}{boate; cabaré; discoteca; casa noturna}
\end{EntryWithPhonetic}

%%%%%%%%%% 咽 %%%%%%%%%%
\subsection*{咽}\addcontentsline{loh}{figure}{咽 \dpy{ye4}}

\begin{EntryWithPhonetic}{咽}{ye4}{9}{⼝}
  \definition{v.}{engasgar (ao chorar)}
  \seeref{yan1}
  \seeref{yan4}
\end{EntryWithPhonetic}

%%%%%%%%%% 液 %%%%%%%%%%
\subsection*{液}\addcontentsline{loh}{figure}{液 \dpy{ye4}}

\begin{EntryWithPhonetic}{液}{ye4}{11}{⽔}
  \definition{s.}{líquido; fluido; suco}
\end{EntryWithPhonetic}

\begin{EntryWithPhonetic}{液晶}{ye4jing1}{11,12}{⽔,⽇}[HSK 7-9]
  \definition[块,个,层]{s.}{cristal líquido; os cristais líquidos podem ser usados em materiais de exibição na indústria eletrônica, bem como em testes não destrutivos e diagnósticos médicos}
\end{EntryWithPhonetic}

\begin{EntryWithPhonetic}{液体}{ye4ti3}{11,7}{⽔,⼈}[HSK 7-9]
  \definition[种,毫升]{adj.}{fluido; líquido; licor; substâncias com volume definido, mas sem forma definida, que podem fluir, são líquidas à temperatura ambiente; óleo, água, álcool e mercúrio são exemplos de líquidos}
  \definition{s.}{líquido}
\end{EntryWithPhonetic}

%%%%%%%%%% 一 %%%%%%%%%%
\subsection*{一}\addcontentsline{loh}{figure}{一 \dpy{yi1}}

\begin{EntryWithPhonetic}{一}{yi1}[(\dpy{yao1})]{1}{⼀}[HSK 1][Kangxi 1]
  \definition{adv.}{uma vez; assim que; indica que duas ações ocorreram em um intervalo de tempo muito curto, uma terminando e a outra começando imediatamente em seguida | indica que primeiro se realiza uma ação e, em seguida, o resultado dessa ação  | indica uma ação única, indicando que a ação é muito curta ou apenas uma tentativa}
  \definition{num.}{um; 1 | pronunciado como \dpy{yao1} quando dito número a número | igual; refere"-se ao mesmo ou igual | inteiro; todo; por toda parte | exclusivo ou único | refere"-se a algo específico | também; caso contrário; referindo"-se a outro ou mais um}
  \definition{part.}{antes de certas palavras para dar ênfase}
  \definition{prep.}{cada; por; toda vez}
  \definition{s.}{uma nota da escala em Gongchepu (工尺谱), correspondente ao 17 na notação musical numerada}
  \seeref{yi2}
  \seeref{yi4}
  \seealsoref{工尺谱}{gong1 che3 pu3}
  \synonymref{一阵}{yi2zhen4}
  \antonymref{多次}{duo1ci4}
\end{EntryWithPhonetic}

\begin{EntryWithPhonetic}{一把手}{yi1ba3shou3}{1,7,4}{⼀,⼿,⼿}[HSK 7-9]
  \definition{s.}{parceiro; participante; membro; mão; um dos participantes | uma boa mão; pessoa capaz | chefe; líder; primeiro em; pessoa número um; pessoa principal responsável}
\end{EntryWithPhonetic}

\begin{EntryWithPhonetic}{一个劲儿}{yi1ge5jin4r5}{1,3,7,2}{⼀,⼈,⼒,⼉}[HSK 7-9]
  \definition{adv.}{um após o outro; continuamente; persistentemente}
  \synonymref{一鼓作气}{yi4gu3-zuo4qi4}
\end{EntryWithPhonetic}

\begin{EntryWithPhonetic}{一会儿……一会儿……}{yi1hui4r5 yi1hui4r5}{1,6,2,1,6,2}{⼀,⼈,⼉,⼀,⼈,⼉}
  \definition{adv.}{um tempo\dots um tempo\dots}
\end{EntryWithPhonetic}

\begin{EntryWithPhonetic}{一……就……}{yi1 jiu4}{1,12}{⼀,⼪}
  \definition{expr.}{logo que |  uma vez que}
\end{EntryWithPhonetic}

\begin{EntryWithPhonetic}{一揽子}{yi1lan3zi5}{1,12,3}{⼀,⼿,⼦}[HSK 7-9]
  \definition{adj.}{inclusivo; não discriminatório; tratamento indiscriminado ou não seletivo de todas as coisas}
\end{EntryWithPhonetic}

\begin{EntryWithPhonetic}{一无所有}{yi1wu2suo3you3}{1,4,8,6}{⼀,⽆,⼾,⽉}[HSK 7-9]
  \definition{expr.}{completamente desprovido; sem nada para esfregar; absolutamente nada; isso pode se referir a dinheiro, mas também a conquistas ou conhecimento}
  \antonymref{比比皆是}{bi3bi3-jie1shi4}
  \antonymref{一应俱全}{yi4ying1-ju4quan2}
  \antonymref{应有尽有}{ying1you3-jin4you3}
\end{EntryWithPhonetic}

\begin{EntryWithPhonetic}{一线}{yi1xian4}{1,8}{⼀,⽷}[HSK 7-9]
  \definition{s.}{frente de guerra; linha de frente | vanguarda (da produção, pesquisa, etc.) | raio; brilho | um fio; um raio de; um vislumbre de}
\end{EntryWithPhonetic}

\begin{EntryWithPhonetic}{一心一意}{yi1xin1-yi1yi4}{1,4,1,13}{⼀,⼼,⼀,⼼}[HSK 7-9]
  \definition{expr.}{dedicar"-se inteiramente a; concentrar"-se totalmente em; empenhar"-se em; um pensamento, uma intenção; isso se refere a um foco singular, desprovido de outros pensamentos}
  \synonymref{诚心诚意}{cheng2xin1-cheng2yi4}
  \synonymref{聚精会神}{ju4jing1-hui4shen2}
  \synonymref{全心全意}{quan2xin1-quan2yi4}
  \antonymref{东张西望}{dong1zhang1-xi1wang4}
\end{EntryWithPhonetic}

\begin{EntryWithPhonetic}{一一}{yi1yi1}{1,1}{⼀,⼀}[HSK 7-9]
  \definition{adv.}{um por um; um após o outro}
\end{EntryWithPhonetic}

%%%%%%%%%% 伊 %%%%%%%%%%
\subsection*{伊}\addcontentsline{loh}{figure}{伊 \dpy{yi1}}

\begin{EntryWithPhonetic}{伊}{yi1}{6}{⼈}
  \definition*{s.}{Iraque, abreviação de 伊拉克 | Irã,abreviação de  伊朗 | Sobrenome: Yi}
  \definition{part.}{(chinês clássico) partícula introdutória sem significado específico}
  \definition{pron.}{Literário: pronome de terceira pessoa do singular (ele ou ela) | pronome de segunda pessoa do singular (você) | disso (precedendo um substantivo)}
  \seealsoref{伊拉克}{yi1la1ke4}
  \seealsoref{伊朗}{yi1lang3}
\end{EntryWithPhonetic}

\begin{EntryWithPhonetic}{伊拉克}{yi1la1ke4}{6,8,7}{⼈,⼿,⼗}
  \definition*{s.}{Iraque}
\end{EntryWithPhonetic}

\begin{EntryWithPhonetic}{伊朗}{yi1lang3}{6,10}{⼈,⽉}
  \definition*{s.}{Irã}
\end{EntryWithPhonetic}

\begin{EntryWithPhonetic}{伊马姆}{yi1ma3mu3}{6,3,8}{⼈,⾺,⼥}
  \definition*{s.}{Iman}
  \seealsoref{伊玛目}{yi1ma3mu4}
  \seealsoref{伊曼}{yi1man4}
  \seealsoref{伊斯兰}{yi1si1lan2}
\end{EntryWithPhonetic}

\begin{EntryWithPhonetic}{伊玛目}{yi1ma3mu4}{6,7,5}{⼈,⽟,⽬}
  \definition*{s.}{Empréstimo linguístico: Iman (Islã)}
  \seealsoref{伊马姆}{yi1ma3mu3}
  \seealsoref{伊曼}{yi1man4}
  \seealsoref{伊斯兰}{yi1si1lan2}
\end{EntryWithPhonetic}

\begin{EntryWithPhonetic}{伊曼}{yi1man4}{6,11}{⼈,⽈}
  \definition*{s.}{Iman (nome próprio)}
  \seealsoref{伊马姆}{yi1ma3mu3}
  \seealsoref{伊玛目}{yi1ma3mu4}
  \seealsoref{伊斯兰}{yi1si1lan2}
\end{EntryWithPhonetic}

\begin{EntryWithPhonetic}{伊斯兰}{yi1si1lan2}{6,12,5}{⼈,⽄,⼋}
  \definition*{s.}{Islã}
  \seealsoref{伊马姆}{yi1ma3mu3}
  \seealsoref{伊玛目}{yi1ma3mu4}
  \seealsoref{伊曼}{yi1man4}
\end{EntryWithPhonetic}

\begin{EntryWithPhonetic}{伊斯兰教}{yi1si1lan2jiao4}{6,12,5,11}{⼈,⽄,⼋,⽁}[HSK 7-9]
  \definition{s.}{islã; islamismo; uma das principais religiões do mundo, fundada pelo árabe Maomé no início do século VII d.C., floresceu na Ásia Ocidental e no Norte da África; foi introduzida na China durante a Dinastia Tang e também é conhecida na China como religião islâmica}
\end{EntryWithPhonetic}

%%%%%%%%%% 衣 %%%%%%%%%%
\subsection*{衣}\addcontentsline{loh}{figure}{衣 \dpy{yi1}}

\begin{EntryWithPhonetic}{衣}{yi1}{6}{⾐}
  \definition[件]{s.}{roupa}
  \seeref{yi4}
\end{EntryWithPhonetic}

\begin{EntryWithPhonetic}{衣服}{yi1fu5}{6,8}{⾐,⽉}[HSK 1]
  \definition[套,件]{s.}{roupas; vestuário; algo que se veste para cobrir o corpo e se proteger do frio}
\end{EntryWithPhonetic}

\begin{EntryWithPhonetic}{衣柜}{yi1gui4}{6,8}{⾐,⽊}
  \definition[个]{s.}{armário | guarda"-roupa}
\end{EntryWithPhonetic}

\begin{EntryWithPhonetic}{衣甲}{yi1jia3}{6,5}{⾐,⽥}
  \definition{s.}{armadura}
\end{EntryWithPhonetic}

\begin{EntryWithPhonetic}{衣架}{yi1jia4}{6,9}{⾐,⽊}[HSK 3]
  \definition[个,副,组]{s.}{cabideiro; móvel para pendurar roupas | estatura; figura; refere"-se ao tipo físico de uma pessoa; estrutura corporal}
\end{EntryWithPhonetic}

\begin{EntryWithPhonetic}{衣食住行}{yi1-shi2-zhu4-xing2}{6,9,7,6}{⾐,⾷,⼈,⾏}[HSK 7-9]
  \definition{expr.}{``Vestuário, alimentação, alojamento e transporte.''; itens essenciais do dia a dia; necessidades básicas da vida; vestuário, alimentação, moradia e transporte; referem"-se às necessidades básicas de vida}
\end{EntryWithPhonetic}

%%%%%%%%%% 医 %%%%%%%%%%
\subsection*{医}\addcontentsline{loh}{figure}{医 \dpy{yi1}}

\begin{EntryWithPhonetic}{医}{yi1}{7}{⼖}
  \definition*{s.}{Sobrenome: Yi}
  \definition{s.}{médico | medicina; ciência médica}
  \definition{v.}{curar; tratar}
\end{EntryWithPhonetic}

\begin{EntryWithPhonetic}{医疗}{yi1liao2}{7,7}{⼖,⽧}[HSK 4]
  \definition{s.}{tratamento médico; tratamento de doenças}
\end{EntryWithPhonetic}

\begin{EntryWithPhonetic}{医生}{yi1sheng1}{7,5}{⼖,⽣}[HSK 1]
  \definition[位,个,名]{s.}{médico; clínico; pessoa que possui conhecimentos médicos e cuja profissão é tratar doenças}
\end{EntryWithPhonetic}

\begin{EntryWithPhonetic}{医务}{yi1wu4}{7,5}{⼖,⼒}[HSK 7-9]
  \definition{pref.}{utilizado como modificador com o significado de 'médico' para formar compostos que indiquem pessoal, serviços ou assuntos de saúde}[医务工作===Trabalho médico]
  \definition{s.}{equipe médica | assuntos médicos}
\end{EntryWithPhonetic}

\begin{EntryWithPhonetic}{医学}{yi1xue2}{7,8}{⼖,⼦}[HSK 4]
  \definition{s.}{medicina; iatrologia; ciência médica; ciência da prevenção e do tratamento de doenças e da proteção e promoção da saúde humana}
\end{EntryWithPhonetic}

\begin{EntryWithPhonetic}{医药}{yi1yao4}{7,9}{⼖,⾋}[HSK 6]
  \definition{s.}{medicina | médico | cuidados médicos e medicamentos | medicamento (droga) | farmacêutica}
\end{EntryWithPhonetic}

\begin{EntryWithPhonetic}{医院}{yi1yuan4}{7,9}{⼖,⾩}[HSK 1]
  \definition[家,所,个]{s.}{hospital; instituições que tratam e cuidam de pacientes, e também realizam exames de saúde, prevenção de doenças, etc.}
\end{EntryWithPhonetic}

%%%%%%%%%% 依 %%%%%%%%%%
\subsection*{依}\addcontentsline{loh}{figure}{依 \dpy{yi1}}

\begin{EntryWithPhonetic}{依}{yi1}{8}{⼈}[HSK 7-9]
  \definition*{s.}{Sobrenome: Yi}
  \definition{prep.}{de acordo com; à luz de; julgando por | 依 + {objeto} + 看/说,\dots; de acordo com a perspectiva/opinião da pessoa/objeto que toma a decisão}[依医生说，你需要休息。===Segundo o médico, você precisa descansar.]
  \definition{v.}{depender de; ser dependente de; confiar em | cumprir; ouvir; ceder a | inclinar"-se; descansar sobre (ou contra)}
  \seealsoref{看}{kan4}
  \seealsoref{说}{shuo1}
  \synonymref{挨}{ai1}
\end{EntryWithPhonetic}

\begin{EntryWithPhonetic}{依次}{yi1ci4}{8,6}{⼈,⽋}[HSK 6]
  \definition{adv.}{sucessivamente; na ordem correta; em ordem}
\end{EntryWithPhonetic}

\begin{EntryWithPhonetic}{依法}{yi1fa3}{8,8}{⼈,⽔}[HSK 5]
  \definition{adv.}{e acordo com regras (ou métodos) fixas | de acordo com a lei; por força da lei; em conformidade com as disposições legais}
\end{EntryWithPhonetic}

\begin{EntryWithPhonetic}{依旧}{yi1jiu4}{8,5}{⼈,⽇}[HSK 5]
  \definition{adv.}{ainda; como antes; como sempre}
\end{EntryWithPhonetic}

\begin{EntryWithPhonetic}{依据}{yi1ju4}{8,11}{⼈,⼿}[HSK 5]
  \definition{prep.}{julgando por; de acordo com; à luz de; com base em; de acordo com; introduzir algo que possa servir como premissa ou base}
  \definition[个]{s.}{base; evidência; fundamento; base para tomar uma decisão ou realizar uma ação}
  \definition{v.}{basear"-se em; confiar em; depdender de; usar algo como premissa ou base}
\end{EntryWithPhonetic}

\begin{EntryWithPhonetic}{依靠}{yi1kao4}{8,15}{⼈,⾮}[HSK 4]
  \definition{s.}{apoio; suporte; algo em que se apoiar; alguém ou algo em quem você pode confiar}
  \definition{v.}{depender de; confiar em (alguém ou alguma coisa para atingir um determinado objetivo)}
\end{EntryWithPhonetic}

\begin{EntryWithPhonetic}{依赖}{yi1lai4}{8,13}{⼈,⾙}[HSK 6]
  \definition{v.}{confiar em; ser dependente de; ser completamente dependente e inseparável | depender de; ser mutuamente dependentes e inseparáveis}
\end{EntryWithPhonetic}

\begin{EntryWithPhonetic}{依然}{yi1ran2}{8,12}{⼈,⽕}[HSK 4]
  \definition{adv.}{ainda; como antes}
  \definition{v.}{estar quieto; estar como antes; estar como o original, sem alterações}
\end{EntryWithPhonetic}

\begin{EntryWithPhonetic}{依托}{yi1tuo1}{8,6}{⼈,⼿}[HSK 7-9]
  \definition{s.}{suporte; apoio; reforço; pessoas ou coisas em que você pode confiar}
  \definition{v.}{confiar em algo (para apoio, etc.); depender de}
  \synonymref{依靠}{yi1kao4}
\end{EntryWithPhonetic}

\begin{EntryWithPhonetic}{依偎}{yi1wei1}{8,11}{⼈,⼈}
  \definition{v.}{aninhar"-se | aconchegar"-se}
\end{EntryWithPhonetic}

\begin{EntryWithPhonetic}{依依不舍}{yi1yi1-bu4she3}{8,8,4,8}{⼈,⼈,⼀,⾆}[HSK 7-9]
  \definition{expr.}{relutar em se separar; um sentimento de relutância em se separar de alguém; não suportar a ideia de partir (separar"-se); não conseguir se desvencilhar de\dots; sentimentos de pesar pela separação; algo que pesa na alma; ter um grande (forte; invencível) apego a\dots; incapaz de se separar um do outro; incapaz de se desvencilhar de; relutante em se separar de (algo)}
  \synonymref{恋恋不舍}{lian4lian4-bu4she4}
\end{EntryWithPhonetic}

\begin{EntryWithPhonetic}{依照}{yi1zhao4}{8,13}{⼈,⽕}[HSK 5]
  \definition{prep.}{de acordo com; à luz de; introduzir certos padrões para os eventos, o que equivale a 按照}
  \definition{v.}{seguir (com base em algo)}
  \seealsoref{按照}{an4zhao4}
\end{EntryWithPhonetic}

%%%%%%%%%% 毉 %%%%%%%%%%
\subsection*{毉}\addcontentsline{loh}{figure}{毉 \dpy{yi1}}

\begin{EntryWithPhonetic}{毉}{yi1}{18}{⼖}
  \variantof{医}
\end{EntryWithPhonetic}

%%%%%%%%%% 一 %%%%%%%%%%
\subsection*{一}\addcontentsline{loh}{figure}{一 \dpy{yi2}}

\begin{EntryWithPhonetic}{一}{yi2}[(一 $+$ 4º tom)]{1}{⼀}[Kangxi 1]
  \definition{num.}{um; 1 | um (artigo)}
  \seeref{yi1}
  \seeref{yi4}
\end{EntryWithPhonetic}

\begin{EntryWithPhonetic}{一半}{yi2ban4}{1,5}{⼀,⼗}[HSK 1]
  \definition{num.}{metade; em parte; uma metade}
  \antonymref{全部}{quan2bu4}
\end{EntryWithPhonetic}

\begin{EntryWithPhonetic}{一辈子}{yi2bei4zi5}{1,12,3}{⼀,⾞,⼦}[HSK 5]
  \definition{s.}{uma vida inteira; vida inteira; toda a vida; durante toda a vida; enquanto se vive; todo o tempo entre o nascimento e a morte}
  \synonymref{一生}{yi4sheng1}
\end{EntryWithPhonetic}

\begin{EntryWithPhonetic}{一不小心}{yi2 bu4 xiao3xin1}{1,4,3,4}{⼀,⼀,⼩,⼼}[HSK 7-9]
  \definition{expr.}{acidentalmente}
\end{EntryWithPhonetic}

\begin{EntryWithPhonetic}{一部分}{yi2bu4fen5}{1,10,4}{⼀,⾢,⼑}[HSK 2]
  \definition{adj.}{parcial}
  \definition{adv.}{parcialmente}
  \definition{num.}{parte; porção; seção; fração}
  \synonymref{个别}{ge4bie2}
  \synonymref{局部}{ju2bu4}
  \synonymref{局限}{ju2xian4}
  \synonymref{一方面}{yi4fang1mian4}
\end{EntryWithPhonetic}

\begin{EntryWithPhonetic}{一刹那}{yi2cha4na4}{1,8,6}{⼀,⼑,⾢}[HSK 7-9]
  \definition{s.}{num instante; num relâmpago; num piscar de olhos | daqui a pouco}
  \seealsoref{一会儿}{yi2hui4r5}
  \synonymref{一瞬间}{yi2shun4jian1}
\end{EntryWithPhonetic}

\begin{EntryWithPhonetic}{一次性}{yi2ci4xing4}{1,6,8}{⼀,⽋,⼼}[HSK 6]
  \definition{adj.}{único; uso único; descartável (produtos); apenas uma vez, sem necessidade ou necessidade de fazer novamente}
\end{EntryWithPhonetic}

\begin{EntryWithPhonetic}{一大早}{yi2da4zao3}{1,3,6}{⼀,⼤,⽇}[HSK 7-9]
  \definition{s.}{de manhã cedo; ao amanhecer}
\end{EntryWithPhonetic}

\begin{EntryWithPhonetic}{一代}{yi2dai4}{1,5}{⼀,⼈}[HSK 6]
  \definition{s.}{uma dinastia | era; época atual | vida; geração; toda a vida de uma pessoa}
\end{EntryWithPhonetic}

\begin{EntryWithPhonetic}{一带}{yi2dai4}{1,9}{⼀,⼱}[HSK 5]
  \definition{s.}{a área em torno de um determinado local; refere"-se a um determinado local e suas proximidades}
\end{EntryWithPhonetic}

\begin{EntryWithPhonetic}{一旦}{yi2dan4}{1,5}{⼀,⽇}[HSK 5]
  \definition{adv.}{uma vez; no caso; agora que | de repente; uma vez}
  \definition{s.}{em um único dia; em um tempo muito curto;}
  \synonymref{万一}{wan4yi1}
\end{EntryWithPhonetic}

\begin{EntryWithPhonetic}{一道}{yi2dao4}{1,12}{⼀,⾡}[HSK 6]
  \definition{adv.}{juntos; lado a lado; junto com}
  \synonymref{一起}{yi4qi3}
\end{EntryWithPhonetic}

\begin{EntryWithPhonetic}{一定}{yi2ding4}{1,8}{⼀,⼧}[HSK 2]
  \definition{adj.}{certo; particular; tendo um certo nível de especificidade; (objeto, situação) determinado em um ou mais | devido; certo; sempre foi assim, não vai mudar | fixo; especificado; há requisitos claros quanto à maneira, método, quantidade, etc.}
  \definition{adv.}{certamente; necessariamente; expressando determinação ou certeza | certamente; indica especulação ou avaliação de que um evento ou situação definitivamente acontecerá ou realmente existirá}
  \synonymref{必定}{bi4ding4}
  \synonymref{必然}{bi4ran2}
  \synonymref{必需}{bi4xu1}
  \synonymref{肯定}{ken3ding4}
  \synonymref{确定}{que4ding4}
  \synonymref{特定}{te4ding4}
  \synonymref{指定}{zhi3ding4}
  \antonymref{不必}{bu2bi4}
  \antonymref{不定}{bu2ding4}
  \antonymref{大概}{da4gai4}
  \antonymref{大致}{da4zhi4}
  \antonymref{或许}{huo4xu3}
  \antonymref{可能}{ke3neng2}
  \antonymref{莫非}{mo4fei1}
  \antonymref{倘若}{tang3ruo4}
  \antonymref{未必}{wei4bi4}
  \antonymref{也许}{ye3xu3}
\end{EntryWithPhonetic}

\begin{EntryWithPhonetic}{一动不动}{yi2dong4-bu2dong4}{1,6,4,6}{⼀,⼒,⼀,⼒}[HSK 7-9]
  \definition{expr.}{imóvel}[他一动不动地站着。===Ele permaneceu imóvel.]
  \antonymref{变幻莫测}{bian4huan4-mo4ce4}
\end{EntryWithPhonetic}

\begin{EntryWithPhonetic}{一度}{yi2du4}{1,9}{⼀,⼴}[HSK 7-9]
  \definition{adv.}{em uma ocasião; uma vez ou por um período de tempo | uma vez; em determinado momento; por um tempo}
  \seealsoref{一向}{yi2xiang4}
  \synonymref{不停}{bu4ting2}
  \synonymref{一贯}{yi2guan4}
  \synonymref{一味}{yi2wei4}
  \synonymref{一向}{yi2xiang4}
\end{EntryWithPhonetic}

\begin{EntryWithPhonetic}{一概}{yi2gai4}{1,13}{⼀,⽊}[HSK 7-9]
  \definition{adv.}{totalmente; todos e quaisquer; sem exceção; em geral, sem exceção}
  \seealsoref{一律}{yi2lv4}
  \synonymref{绝对}{jue2dui4}
  \synonymref{美满}{mei3man3}
  \synonymref{全部}{quan2bu4}
  \synonymref{十足}{shi2zu2}
  \synonymref{统统}{tong3tong3}
  \synonymref{完全}{wan2quan2}
  \synonymref{万万}{wan4wan4}
  \synonymref{一律}{yi2lv4}
  \synonymref{一切}{yi2qie4}
  \synonymref{一致}{yi2zhi4}
\end{EntryWithPhonetic}

\begin{EntryWithPhonetic}{一概而论}{yi2gai4'er2lun4}{1,13,6,6}{⼀,⽊,⽽,⾔}[HSK 7-9]
  \definition{expr.}{generalizar; agrupar assuntos diferentes; tratar ou processar as coisas usando o mesmo padrão (frequentemente usado em negação); tratar assuntos diferentes como se fossem iguais (geralmente na forma negativa)}
  \synonymref{相提并论}{xiang1ti2-bing4lun4}
\end{EntryWithPhonetic}

\begin{EntryWithPhonetic}{一个样}{yi2ge5yang4}{1,3,10}{⼀,⼈,⽊}
  \definition{s.}{o mesmo}
  \seealsoref{一样}{yi2yang4}
\end{EntryWithPhonetic}

\begin{EntryWithPhonetic}{一共}{yi2gong4}{1,6}{⼀,⼋}[HSK 2]
  \definition{adv.}{completamente; em tudo; no todo}
  \synonymref{全部}{quan2bu4}
  \synonymref{全面}{quan2mian4}
  \synonymref{全体}{quan2ti3}
  \synonymref{所有}{suo3you3}
  \synonymref{统统}{tong3tong3}
  \synonymref{一切}{yi2qie4}
  \synonymref{整个}{zheng3ge4}
  \synonymref{总共}{zong3gong4}
\end{EntryWithPhonetic}

\begin{EntryWithPhonetic}{一贯}{yi2guan4}{1,8}{⼀,⾙}[HSK 6]
  \definition{adj./adv.}{do começo ao fim; inabalável; consistente; persistente; o tempo todo}
  \synonymref{不断}{bu2duan4}
  \synonymref{从来}{cong2lai2}
  \synonymref{通常}{tong1chang2}
  \synonymref{向来}{xiang4lai2}
  \synonymref{一度}{yi2du4}
  \synonymref{一向}{yi2xiang4}
  \synonymref{一直}{yi4zhi2}
  \antonymref{偶尔}{ou3'er3}
  \antonymref{偶然}{ou3ran2}
\end{EntryWithPhonetic}

\begin{EntryWithPhonetic}{一晃}{yi2huang4}{1,10}{⼀,⽇}[HSK 7-9]
  \definition{adv.}{num piscar de olhos; rapidamente; descreve a rapidez com que o tempo passa (implicando que algo acontece sem que você sequer perceba)}
  \seeref{yi4huang3}
  \synonymref{毕竟}{bi4jing4}
\end{EntryWithPhonetic}

\begin{EntryWithPhonetic}{一会儿}{yi2hui4r5}{1,6,2}{⼀,⼈,⼉}[HSK 1,2]
  \definition{adv.}{agora\dots agora\dots; um momento\dots o próximo\dots; usado antes de dois antônimos, indica a alternância de situações}
  \definition{s.}{um pouquinho de tempo; muito pouco tempo}
  \synonymref{不一会儿}{bu4 yi2hui4r5}
  \synonymref{紧接着}{jin3 jie1zhe5}
  \synonymref{一下子}{yi2xia4zi5}
\end{EntryWithPhonetic}

\begin{EntryWithPhonetic}{一技之长}{yi2ji4zhi1chang2}{1,7,3,4}{⼀,⼿,⼂,⾧}[HSK 7-9]
  \definition{expr.}{proficiência em uma área específica (ou campo); habilidade profissional; especialidade; habilidade em uma área especializada}
  \antonymref{一事无成}{yi2shi4-wu2cheng2}
\end{EntryWithPhonetic}

\begin{EntryWithPhonetic}{一句话}{yi2ju4hua4}{1,5,8}{⼀,⼝,⾔}[HSK 5]
  \definition{s.}{em resumo; em uma palavra; expressar um conteúdo complexo de forma sucinta | trabalho fácil; fácil de fazer; descrever uma tarefa ou trabalho como muito simples e fácil de realizar}
\end{EntryWithPhonetic}

\begin{EntryWithPhonetic}{一块}{yi2kuai4}{1,7}{⼀,⼟}
  \definition{adv.}{(principalmente mandarim) juntos}
  \synonymref{一路}{yi2lu4}
  \synonymref{一齐}{yi4qi2}
  \synonymref{一起}{yi4qi3}
  \synonymref{一同}{yi4tong2}
  \antonymref{各自}{ge4zi4}
\end{EntryWithPhonetic}

\begin{EntryWithPhonetic}{一块儿}{yi2kuai4r5}{1,7,2}{⼀,⼟,⼉}[HSK 1]
  \definition{adv.}{juntos; em conjunto}
  \definition{s.}{no mesmo lugar; no mesmo local}
  \synonymref{一起}{yi4qi3}
\end{EntryWithPhonetic}

\begin{EntryWithPhonetic}{一路}{yi2lu4}{1,13}{⼀,⾜}[HSK 5]
  \definition{adv.}{o tempo todo; persistentemente; continuamente | juntos; sem parar; continuamente}
  \definition{s.}{o mesmo caminho; a mesma rota; ao longo de toda a viagem, ao longo do caminho | do mesmo tipo; da mesma categoria}
\end{EntryWithPhonetic}

\begin{EntryWithPhonetic}{一路平安}{yi2lu4-ping2'an1}{1,13,5,6}{⼀,⾜,⼲,⼧}[HSK 2]
  \definition{expr.}{Boa viagem!; Tenha uma boa viagem!}
  \definition{v.}{ter uma viagem agradável}
  \synonymref{一路顺风}{yi2lu4shun4feng1}
  \synonymref{一帆风顺}{yi4fan1-feng1shun4}
\end{EntryWithPhonetic}

\begin{EntryWithPhonetic}{一路上}{yi2lu4shang4}{1,13,3}{⼀,⾜,⼀}[HSK 6]
  \definition{s.}{ao longo do caminho; todo o caminho}
\end{EntryWithPhonetic}

\begin{EntryWithPhonetic}{一路顺风}{yi2lu4shun4feng1}{1,13,9,4}{⼀,⾜,⾴,⾵}[HSK 2]
  \definition{expr.}{ter uma viagem agradável; toda a viagem foi segura e tranquila; é uma metáfora para cada etapa do processo de lidar com algo que ocorre sem problemas | Tenha uma boa viagem!; Boa viagem!}
  \synonymref{一路平安}{yi2lu4-ping2'an1}
  \synonymref{一帆风顺}{yi4fan1-feng1shun4}
\end{EntryWithPhonetic}

\begin{EntryWithPhonetic}{一律}{yi2lv4}{1,9}{⼀,⼻}[HSK 4]
  \definition{adj.}{igual; semelhante; uniforme; parecido; idêntico}
  \definition{adv.}{todos; tudo; sem exceção; enfatiza que todos devem ser assim, sem exceção, e é usado principalmente em regulamentos ou requisitos}
  \synonymref{绝对}{jue2dui4}
  \synonymref{同等}{tong2deng3}
  \synonymref{完全}{wan2quan2}
  \synonymref{一概}{yi2gai4}
  \synonymref{一样}{yi2yang4}
  \synonymref{一致}{yi2zhi4}
  \synonymref{整齐}{zheng3qi2}
  \antonymref{不同}{bu4tong2}
  \antonymref{例外}{li4wai4}
\end{EntryWithPhonetic}

\begin{EntryWithPhonetic}{一面}{yi2mian4}{1,9}{⼀,⾯}[HSK 7-9]
  \definition{adv.}{ao mesmo tempo; simultaneamente}
  \definition{s.}{um lado | um aspecto}
  \synonymref{部分}{bu4fen5}
  \synonymref{个别}{ge4bie2}
  \synonymref{个人}{ge4ren2}
  \synonymref{一边}{yi4bian1}
  \antonymref{全体}{quan2ti3}
  \antonymref{四处}{si4chu4}
\end{EntryWithPhonetic}

\begin{EntryWithPhonetic}{一目了然}{yi2mu4-liao3ran2}{1,5,2,12}{⼀,⽬,⼅,⽕}[HSK 7-9]
  \definition{expr.}{claro à primeira vista; óbvio à primeira vista; isso pode ser visto claramente à primeira vista}
  \antonymref{错综复杂}{cuo4zong1-fu4za2}
  \antonymref{眼花缭乱}{yan3hua1liao2luan4}
\end{EntryWithPhonetic}

\begin{EntryWithPhonetic}{一切}{yi2qie4}{1,4}{⼀,⼑}[HSK 3]
  \definition{pron.}{tudo; todo; todas as coisas}
  \synonymref{一共}{yi2gong4}
\end{EntryWithPhonetic}

\begin{EntryWithPhonetic}{一事无成}{yi2shi4-wu2cheng2}{1,8,4,6}{⼀,⼅,⽆,⼽}[HSK 7-9]
  \definition{expr.}{não realizar nada; não chegar a lugar nenhum; ser um fracasso total; não ter conquistado nada}
\end{EntryWithPhonetic}

\begin{EntryWithPhonetic}{一瞬间}{yi2shun4jian1}{1,17,7}{⼀,⽬,⾨}[HSK 7-9]
  \definition{s.}{instante; fração de segundo; um instante; num piscar de olhos, o tempo passa muito depressa; esta expressão é usada para descrever um período de tempo extremamente curto}
  \synonymref{一刹那}{yi2cha4na4}
  \antonymref{一辈子}{yi2bei4zi5}
\end{EntryWithPhonetic}

\begin{EntryWithPhonetic}{一味}{yi2wei4}{1,8}{⼀,⼝}[HSK 7-9]
  \definition{adv.}{persistentemente; teimosamente; cegamente; invariavelmente}
  \synonymref{一度}{yi2du4}
  \synonymref{一向}{yi2xiang4}
\end{EntryWithPhonetic}

\begin{EntryWithPhonetic}{一系列}{yi2xi4lie4}{1,7,6}{⼀,⽷,⼑}[HSK 7-9]
  \definition{adj.}{série de; muitos relacionados; um após o outro}
  \synonymref{一连串}{yi4lian2chuan4}
\end{EntryWithPhonetic}

\begin{EntryWithPhonetic}{一下}{yi2xia4}{1,3}{⼀,⼀}
  \definition{adv.}{em um curto tempo | rapidamente}
  \synonymref{统计}{tong3ji4}
\end{EntryWithPhonetic}

\begin{EntryWithPhonetic}{一下儿}{yi2xia4r5}{1,3,2}{⼀,⼀,⼉}[HSK 1,5]
  \definition{s.}{um tempo; um momento}
\end{EntryWithPhonetic}

\begin{EntryWithPhonetic}{一下子}{yi2xia4zi5}{1,3,3}{⼀,⼀,⼦}[HSK 5]
  \definition{adv.}{tudo de uma vez; de repente; em pouco tempo; em um curto espaço de tempo}
  \synonymref{一会儿}{yi2hui4r5}
\end{EntryWithPhonetic}

\begin{EntryWithPhonetic}{一向}{yi2xiang4}{1,6}{⼀,⼝}[HSK 5]
  \definition{adv.}{desde o início; indica do passado até o presente}
  \synonymref{不断}{bu2duan4}
  \synonymref{从来}{cong2lai2}
  \synonymref{历来}{li4lai2}
  \synonymref{向来}{xiang4lai2}
  \synonymref{一度}{yi2du4}
  \synonymref{一贯}{yi2guan4}
  \synonymref{一味}{yi2wei4}
  \synonymref{一直}{yi4zhi2}
  \antonymref{偶尔}{ou3'er3}
  \antonymref{有时}{you3shi2}
\end{EntryWithPhonetic}

\begin{EntryWithPhonetic}{一样}{yi2yang4}{1,10}{⼀,⽊}[HSK 1]
  \definition{adj.}{o mesmo; tão\dots quanto\dots; igualmente; semelhante}
  \definition{part.}{na mesma medida; anexado a verbos ou palavras nominais, indica uma comparação ou semelhança, equivalente a 似的}
  \seealsoref{似的}{shi4de5}
  \synonymref{雷同}{lei2tong2}
  \synonymref{通常}{tong1chang2}
  \synonymref{同样}{tong2yang4}
  \synonymref{相似}{xiang1si4}
  \synonymref{相通}{xiang1tong1}
  \synonymref{相同}{xiang1tong2}
  \synonymref{一律}{yi2lv4}
  \antonymref{不同}{bu4tong2}
\end{EntryWithPhonetic}

\begin{EntryWithPhonetic}{一再}{yi2zai4}{1,6}{⼀,⼌}[HSK 4]
  \definition{adv.}{repetidamente; de novo e de novo; repetidas vezes; uma e outra vez}
  \synonymref{常常}{chang2chang2}
  \synonymref{反复}{fan3fu4}
  \synonymref{屡次}{lv3ci4}
  \synonymref{频繁}{pin2fan2}
  \synonymref{频频}{pin2pin2}
  \synonymref{再三}{zai4san1}
  \antonymref{不再}{bu2zai4}
\end{EntryWithPhonetic}

\begin{EntryWithPhonetic}{一战}{yi2zhan4}{1,9}{⼀,⼽}
  \definition*{s.}{Primeira Guerra Mundial}
\end{EntryWithPhonetic}

\begin{EntryWithPhonetic}{一阵}{yi2zhen4}{1,6}{⼀,⾩}[HSK 7-9]
  \definition{s.}{um repique; uma explosão; um feitiço; refere"-se ao período de tempo durante o qual uma ação ou situação dura}
  \synonymref{一番}{yi4 fan1}
\end{EntryWithPhonetic}

\begin{EntryWithPhonetic}{一致}{yi2zhi4}{1,10}{⼀,⾄}[HSK 4]
  \definition{adj.}{equado; idêntico; uniforme; unânime; nenhuma diferença (de opinião ou ação)}
  \definition{adv.}{juntos; em conjunto}
  \synonymref{一概}{yi2gai4}
  \synonymref{一律}{yi2lv4}
  \antonymref{分歧}{fen1qi2}
  \antonymref{相反}{xiang1fan3}
\end{EntryWithPhonetic}

%%%%%%%%%% 仪 %%%%%%%%%%
\subsection*{仪}\addcontentsline{loh}{figure}{仪 \dpy{yi2}}

\begin{EntryWithPhonetic}{仪}{yi2}{5}{⼈}
  \definition*{s.}{Sobrenome: Yi}
  \definition{s.}{aparência; porte | cerimônia; rito | presente; dádiva | aparelho; instrumento}
  \definition{v.}{(literário) admirar; ansiar por}
\end{EntryWithPhonetic}

\begin{EntryWithPhonetic}{仪表}{yi2biao3}{5,8}{⼈,⾐}[HSK 7-9]
  \definition[个,种]{s.}{aparência, referindo"-se a uma boa aparência; porte (incluindo aparência, postura, atitude, etc. da aparência de uma pessoa) | dispositivo de medição; instrumento (para medir temperatura, pressão, eletricidade, pressão arterial, etc.)}
  \synonymref{风采}{feng1cai3}
  \synonymref{风度}{feng1du4}
  \synonymref{风范}{feng1fan4}
  \synonymref{面貌}{mian4mao4}
  \synonymref{容貌}{rong2mao4}
\end{EntryWithPhonetic}

\begin{EntryWithPhonetic}{仪器}{yi2qi4}{5,16}{⼈,⼝}[HSK 6]
  \definition[台]{s.}{aparelho; instrumento; ferramentas ou equipamentos utilizados para observação, medição, inspeção, etc. em pesquisas ou experimentos científicos; geralmente, são relativamente precisos e padronizados}
  \synonymref{器械}{qi4xie4}
  \synonymref{设备}{she4bei4}
\end{EntryWithPhonetic}

\begin{EntryWithPhonetic}{仪式}{yi2shi4}{5,6}{⼈,⼷}[HSK 6]
  \definition{s.}{rito; cerimônia; procedimento e formato da cerimônia}
  \synonymref{典礼}{dian3li3}
  \synonymref{庆典}{qing4dian3}
\end{EntryWithPhonetic}

%%%%%%%%%% 怡 %%%%%%%%%%
\subsection*{怡}\addcontentsline{loh}{figure}{怡 \dpy{yi2}}

\begin{EntryWithPhonetic}{怡}{yi2}{8}{⼼}
  \definition*{s.}{Sobrenome: Yi}
  \definition{adj.}{Literário: feliz; alegre; animado}
\end{EntryWithPhonetic}

\begin{EntryWithPhonetic}{怡然自得}{yi2ran2-zi4de2}{8,12,6,11}{⼼,⽕,⾃,⼻}[HSK 7-9]
  \definition{expr.}{feliz e contente; ser feliz e estar satisfeito consigo mesmo; sentir uma sensação de felicidade plena}
\end{EntryWithPhonetic}

%%%%%%%%%% 姨 %%%%%%%%%%
\subsection*{姨}\addcontentsline{loh}{figure}{姨 \dpy{yi2}}

\begin{EntryWithPhonetic}{姨}{yi2}{9}{⼥}[HSK 7-9]
  \definition[个,位,名,些]{s.}{irmã da mãe; tia materna; tia | irmã da esposa; cunhada}
\end{EntryWithPhonetic}

\begin{EntryWithPhonetic}{姨父}{yi2fu5}{9,4}{⼥,⽗}
  \definition{s.}{marido da irmã da mãe ou da tia materna; tio}
  \antonymref{阿姨}{a1yi2}
\end{EntryWithPhonetic}

\begin{EntryWithPhonetic}{姨妈}{yi2ma1}{9,6}{⼥,⼥}
  \definition[个,位,名,些]{s.}{tia; tia materna; irmã da mãe; tia (referindo"-se a uma mulher casada)}
  \synonymref{阿姨}{a1yi2}
  \synonymref{姨娘}{yi2niang2}
\end{EntryWithPhonetic}

\begin{EntryWithPhonetic}{姨娘}{yi2niang2}{9,10}{⼥,⼥}
  \definition{s.}{Obsoleto: forma de tratamento para a concubina do pai | Dialeto: (casada) tia materna; tia}
  \synonymref{阿姨}{a1yi2}
  \synonymref{大妈}{da4ma1}
  \synonymref{姨妈}{yi2ma1}
\end{EntryWithPhonetic}

%%%%%%%%%% 移 %%%%%%%%%%
\subsection*{移}\addcontentsline{loh}{figure}{移 \dpy{yi2}}

\begin{EntryWithPhonetic}{移}{yi2}{11}{⽲}[HSK 4]
  \definition*{s.}{Sobrenome: Yi}
  \definition{v.}{mover; remover; deslocar; mudar | mudar; alterar}
\end{EntryWithPhonetic}

\begin{EntryWithPhonetic}{移动}{yi2dong4}{11,6}{⽲,⼒}[HSK 4]
  \definition{v.}{deslocar; mover; mudar}
\end{EntryWithPhonetic}

\begin{EntryWithPhonetic}{移交}{yi2jiao1}{11,6}{⽲,⼇}[HSK 7-9]
  \definition{v.}{entregar; transferir; disponibilizar sob a custódia de alguém; transferir pessoas ou coisas para as partes relevantes | passar o cargo para um sucessor; a pessoa que originalmente era responsável pela gestão transferiu as responsabilidades para a pessoa que a assumiu antes de deixar a empresa}
  \synonymref{吩咐}{fen1fu4}
  \synonymref{交接}{jiao1jie1}
  \synonymref{嘱咐}{zhu3fu5}
\end{EntryWithPhonetic}

\begin{EntryWithPhonetic}{移民}{yi2min2}{11,5}{⽲,⽒}[HSK 4]
  \definition[个,批]{s.}{emigrante; migrantes; aqueles que se mudam para um país ou estado estrangeiro para se estabelecer}
  \definition{v.}{migrar; imigrar}
\end{EntryWithPhonetic}

\begin{EntryWithPhonetic}{移植}{yi2zhi2}{11,12}{⽲,⽊}[HSK 7-9]
  \definition{v.}{adaptar; transplantar; arrancar as mudas semeadas no canteiro ou no arrozal, ou desenterar com a terra e plantar no campo | enxerto; transplante; o transplante consiste em fixar um tecido ou órgão de um organismo a uma parte defeituosa de outro organismo, permitindo que se cure gradualmente, como no caso de transplantes de córnea, transplantes de pele, transplantes ósseos e transplantes de vasos sanguíneos}
  \synonymref{种植}{zhong4zhi2}
\end{EntryWithPhonetic}

%%%%%%%%%% 遗 %%%%%%%%%%
\subsection*{遗}\addcontentsline{loh}{figure}{遗 \dpy{yi2}}

\begin{EntryWithPhonetic}{遗}{yi2}{12}{⾡}
  \definition*{s.}{Sobrenome: Yi}
  \definition{s.}{descarga involuntária de urina, etc. | algo perdido}
  \definition{v.}{perder | omitir | deixar para trás; guardar; não dar | deixar para trás após a morte; legar; transmitir}
\end{EntryWithPhonetic}

\begin{EntryWithPhonetic}{遗案}{yi2'an4}{12,10}{⾡,⽊}
  \definition{s.}{(lei) caso não resolvido}
\end{EntryWithPhonetic}

\begin{EntryWithPhonetic}{遗产}{yi2chan3}{12,6}{⾡,⼇}[HSK 4]
  \definition[笔,份]{s.}{legado; herança; patrimônio; propriedade deixada pelo falecido | patrimônio; riqueza cultural ou riqueza material transmitida pela história}
\end{EntryWithPhonetic}

\begin{EntryWithPhonetic}{遗传}{yi2chuan2}{12,6}{⾡,⼈}[HSK 4]
  \definition{v.}{herdar, descender, transmitir, passar adiante}
\end{EntryWithPhonetic}

\begin{EntryWithPhonetic}{遗骸}{yi2hai2}{12,15}{⾡,⾻}
  \definition{v.}{restos mortais}
\end{EntryWithPhonetic}

\begin{EntryWithPhonetic}{遗憾}{yi2han4}{12,16}{⾡,⼼}[HSK 6]
  \definition{adj.}{triste; arrependido; contrito; sentir pena de situações que estão fora de controle ou são insatisfatórias}
  \definition{s.}{pena; arrependimento; sentindo pena que os desejos não se realizaram}
\end{EntryWithPhonetic}

\begin{EntryWithPhonetic}{遗迹}{yi2ji4}{12,9}{⾡,⾡}
  \definition{s.}{vestígio histórico; sítio; vestígio; traço; ruína; vestígios deixados por tempos antigos ou eras passadas}
\end{EntryWithPhonetic}

\begin{EntryWithPhonetic}{遗留}{yi2liu2}{12,10}{⾡,⽥}[HSK 7-9]
  \definition{v.}{deixar para trás; passar adiante; coisas ou problemas do passado permanecem e ainda existem hoje}
  \synonymref{残留}{can2liu2}
  \synonymref{遗传}{yi2chuan2}
  \synonymref{遗落}{yi2luo4}
\end{EntryWithPhonetic}

\begin{EntryWithPhonetic}{遗落}{yi2luo4}{12,12}{⾡,⾋}
  \definition{v.}{esquecer | deixar para trás (inadvertidamente) | deixar de fora | omitir}
\end{EntryWithPhonetic}

\begin{EntryWithPhonetic}{遗男}{yi2nan2}{12,7}{⾡,⽥}
  \definition{s.}{órfão | filho póstumo}
\end{EntryWithPhonetic}

\begin{EntryWithPhonetic}{遗弃}{yi2qi4}{12,7}{⾡,⼶}[HSK 7-9]
  \definition{v.}{desertar; abandonar; renegar; descartar; dar o fora | abandonar; deixar para trás; abandonar parentes pelos quais se tem a obrigação legal ou moral de sustentar ou prover}
  \synonymref{丢掉}{diu1diao4}
  \synonymref{放弃}{fang4qi4}
  \synonymref{抛弃}{pao1qi4}
  \synonymref{扔掉}{reng1diao4}
  \antonymref{抚养}{fu3yang3}
  \antonymref{肩负}{jian1fu4}
  \antonymref{赡养}{shan4yang3}
  \antonymref{收养}{shou1yang3}
  \antonymref{饲养}{si4yang3}
  \antonymref{托付}{tuo1fu4}
  \antonymref{寻找}{xun2zhao3}
\end{EntryWithPhonetic}

\begin{EntryWithPhonetic}{遗体}{yi2ti3}{12,7}{⾡,⼈}[HSK 7-9]
  \definition{s.}{restos mortais; ossos | restos de animais mortos ou plantas secas}
  \synonymref{尸体}{shi1ti3}
\end{EntryWithPhonetic}

\begin{EntryWithPhonetic}{遗忘}{yi2wang4}{12,7}{⾡,⼼}[HSK 7-9]
  \definition{v.}{(alguém) esquecer; (lgo) ter um lapso de memória}
\end{EntryWithPhonetic}

\begin{EntryWithPhonetic}{遗物}{yi2wu4}{12,8}{⾡,⽜}[HSK 7-9]
  \definition{s.}{coisas deixadas pelo falecido | relíquia | restos mortais; remanescente; monumento; coisas deixadas pelos antigos ou pelos mortos}
  \synonymref{遗产}{yi2chan3}
\end{EntryWithPhonetic}

\begin{EntryWithPhonetic}{遗愿}{yi2yuan4}{12,14}{⾡,⽕}[HSK 7-9]
  \definition{s.}{desejo não realizado do falecido; últimos desejos do falecido; pedido}
  \synonymref{遗嘱}{yi2zhu3}
\end{EntryWithPhonetic}

\begin{EntryWithPhonetic}{遗址}{yi2zhi3}{12,7}{⾡,⼟}[HSK 7-9]
  \definition[处]{s.}{ruínas; sítios históricos; os sítios originais de antigas cidades e edifícios que há muito se perderam ou foram destruídos}
  \synonymref{遗产}{yi2chan3}
  \synonymref{遗迹}{yi2ji4}
\end{EntryWithPhonetic}

\begin{EntryWithPhonetic}{遗嘱}{yi2zhu3}{12,15}{⾡,⼝}[HSK 7-9]
  \definition{s.}{testamento; última vontade; últimas palavras; instruções escritas ou orais, feitas antes da morte ou na véspera do falecimento, referentes aos procedimentos para os assuntos da pessoa após a morte}
  \definition{v.}{fazer um testamento; deixar as últimas palavras}
  \synonymref{遗愿}{yi2yuan4}
\end{EntryWithPhonetic}

%%%%%%%%%% 颐 %%%%%%%%%%
\subsection*{颐}\addcontentsline{loh}{figure}{颐 \dpy{yi2}}

\begin{EntryWithPhonetic}{颐}{yi2}{13}{⾴}
  \definition{s.}{bochecha}
  \definition{v.}{manter"-se em forma; cuidar de si mesmo}
\end{EntryWithPhonetic}

\begin{EntryWithPhonetic}{颐和园}{yi2he2yuan2}{13,8,7}{⾴,⼝,⼞}
  \definition*{s.}{Palácio de Verão}
\end{EntryWithPhonetic}

%%%%%%%%%% 疑 %%%%%%%%%%
\subsection*{疑}\addcontentsline{loh}{figure}{疑 \dpy{yi2}}

\begin{EntryWithPhonetic}{疑}{yi2}{14}{⽦}
  \definition{adj.}{duvidoso; incerto}
  \definition{v.}{duvidar; desacreditar; suspeitar}
\end{EntryWithPhonetic}

\begin{EntryWithPhonetic}{疑点}{yi2dian3}{14,9}{⽦,⽕}[HSK 7-9]
  \definition{s.}{pontos suspeitos; ponto duvidoso; ponto questionável; ponto incerto ou pouco claro; dúvida | áreas de dúvida; áreas de incerteza}
\end{EntryWithPhonetic}

\begin{EntryWithPhonetic}{疑惑}{yi2huo4}{14,12}{⽦,⼼}[HSK 7-9]
  \definition[个]{s.}{dúvida; coisas que não se entende}
  \definition{v.}{sentir"-se inseguro; não estar convencido; não entender; não conseguir descobrir, não conseguir compreender}
  \seealsoref{疑问}{yi2wen4}
  \synonymref{怀疑}{huai2yi2}
  \synonymref{可疑}{ke3yi2}
  \synonymref{困惑}{kun4huo4}
  \synonymref{困扰}{kun4rao3}
  \synonymref{迷惑}{mi2huo4}
  \synonymref{奇怪}{qi2guai4}
  \synonymref{嫌疑}{xian2yi2}
  \antonymref{断定}{duan4ding4}
  \antonymref{确信}{que4xin4}
  \antonymref{无疑}{wu2yi2}
  \antonymref{相信}{xiang1xin4}
  \antonymref{信任}{xin4ren4}
  \antonymref{信用}{xin4yong4}
\end{EntryWithPhonetic}

\begin{EntryWithPhonetic}{疑虑}{yi2lv4}{14,10}{⽦,⾌}[HSK 7-9]
  \definition{s.}{dúvidas; receios; preocupações decorrentes da dúvida}
  \definition{v.}{ter dúvidas/reclamações/preocupações; demonstrar desconfiança e suspeita}
\end{EntryWithPhonetic}

\begin{EntryWithPhonetic}{疑问}{yi2wen4}{14,6}{⽦,⾨}[HSK 4]
  \definition[个,些]{s.}{dúvida; consulta; pergunta; questionamento; coisas que não podem ser determinadas ou explicadas}
\end{EntryWithPhonetic}

%%%%%%%%%% 乙 %%%%%%%%%%
\subsection*{乙}\addcontentsline{loh}{figure}{乙 \dpy{yi3}}

\begin{EntryWithPhonetic}{乙}{yi3}{1}{⼄}[HSK 5][Kangxi 5]
  \definition*{s.}{Sobrenome: Yi}
  \definition{num.}{segundo}
  \definition{s.}{o segundo lugar do Tian Gan | uma nota da escala em Gongchepu (工尺谱); nível superior na música tradicional chinesa}
  \seealsoref{工尺谱}{gong1 che3 pu3}
\end{EntryWithPhonetic}

%%%%%%%%%% 已 %%%%%%%%%%
\subsection*{已}\addcontentsline{loh}{figure}{已 \dpy{yi3}}

\begin{EntryWithPhonetic}{已}{yi3}{3}{⼰}[HSK 3]
  \definition{adv.}{já | posteriormente; mais tarde; depois de algum tempo | demasiadamente; excessivamente}
  \definition{v.}{terminar; parar; cessar}
\end{EntryWithPhonetic}

\begin{EntryWithPhonetic}{已故}{yi3gu4}{3,9}{⼰,⽁}
  \definition{adj.}{falecido}
\end{EntryWithPhonetic}

\begin{EntryWithPhonetic}{已婚}{yi3hun1}{3,11}{⼰,⼥}
  \definition{adj.}{casado}
\end{EntryWithPhonetic}

\begin{EntryWithPhonetic}{已经}{yi3jing1}{3,8}{⼰,⽷}[HSK 2]
  \definition{adv.}{já; indica que uma ação ou mudança foi concluída ou atingiu um determinado nível}
\end{EntryWithPhonetic}

\begin{EntryWithPhonetic}{已久}{yi3jiu3}{3,3}{⼰,⼃}
  \definition{adv.}{já faz muito tempo}
\end{EntryWithPhonetic}

\begin{EntryWithPhonetic}{已灭}{yi3mie4}{3,5}{⼰,⽕}
  \definition{adj.}{extinto}
\end{EntryWithPhonetic}

\begin{EntryWithPhonetic}{已然}{yi3ran2}{3,12}{⼰,⽕}
  \definition{adv.}{já | já ser assim}
\end{EntryWithPhonetic}

\begin{EntryWithPhonetic}{已知}{yi3zhi1}{3,8}{⼰,⽮}
  \definition{v.}{conhecido (ter ciência)}
\end{EntryWithPhonetic}

%%%%%%%%%% 以 %%%%%%%%%%
\subsection*{以}\addcontentsline{loh}{figure}{以 \dpy{yi3}}

\begin{EntryWithPhonetic}{以}{yi3}{4}{⼈}[HSK 7-9]
  \definition*{s.}{Sobrenome: Yi}
  \definition{conj.}{e; bem como; o mesmo que 而 | de modo a; a fim de; usado no início da frase seguinte para expressar o propósito de atingir um determinado objetivo}
  \definition{prep.}{por; com | de acordo com | por causa de; porque | em (uma data fixa); em (um certo momento) | colocado antes de palavras posicionais simples, indica os limites de tempo, posição e quantidade}
  \seealsoref{而}{er2}
  \seealsoref{用}{yong4}
\end{EntryWithPhonetic}

\begin{EntryWithPhonetic}{以便}{yi3bian4}{4,9}{⼈,⼈}[HSK 5]
  \definition{conj.}{para que; de modo que; a fim de; com o objetivo de; para o propósito de; usado no início da segunda parte da frase, indica que o objetivo mencionado na segunda parte será facilmente alcançado}
  \seealsoref{便于}{bian4yu2}
  \seealsoref{好}{hao3}
  \synonymref{促进}{cu4jin4}
  \synonymref{促使}{cu4shi3}
  \synonymref{推动}{tui1 dong4}
\end{EntryWithPhonetic}

\begin{EntryWithPhonetic}{以此}{yi3ci3}{4,6}{⼈,⽌}
  \definition{conj.}{por este motivo; por esta razão; por isso}
  \synonymref{乃至}{nai3zhi4}
  \synonymref{甚至}{shen4zhi4}
  \synonymref{以致}{yi3zhi4}
\end{EntryWithPhonetic}

\begin{EntryWithPhonetic}{以后}{yi3hou4}{4,6}{⼈,⼝}[HSK 2]
  \definition{s.}{depois; mais tarde; após; daqui em diante}
  \synonymref{此后}{ci3hou4}
  \synonymref{从此}{cong2ci3}
  \synonymref{后来}{hou4lai2}
  \synonymref{今后}{jin1hou4}
  \synonymref{往后}{wang3hou4}
  \synonymref{以来}{yi3lai2}
  \synonymref{之后}{zhi1hou4}
  \antonymref{刚才}{gang1cai2}
  \antonymref{往日}{wang3ri4}
  \antonymref{以前}{yi3qian2}
\end{EntryWithPhonetic}

\begin{EntryWithPhonetic}{以及}{yi3ji2}{4,3}{⼈,⼃}[HSK 4]
  \definition{conj.}{assim como; juntamente como; bem como; também}
  \synonymref{还有}{hai2you3}
  \synonymref{乃至}{nai3zhi4}
\end{EntryWithPhonetic}

\begin{EntryWithPhonetic}{以来}{yi3lai2}{4,7}{⼈,⽊}[HSK 3]
  \definition{s.}{desde (em termos de tempo); indica um período de tempo desde um determinado momento no passado até o presente}
  \synonymref{从此}{cong2ci3}
  \synonymref{今后}{jin1hou4}
  \synonymref{凭借}{ping2jie4}
  \synonymref{往后}{wang3hou4}
  \synonymref{依靠}{yi1kao4}
  \synonymref{以后}{yi3hou4}
  \synonymref{自从}{zi4cong2}
\end{EntryWithPhonetic}

\begin{EntryWithPhonetic}{以免}{yi3mian3}{4,7}{⼈,⼉}[HSK 7-9]
  \definition{conj.}{para que não; a fim de não tolar; para evitar; usado no início da segunda metade de uma frase, indica que o objetivo é impedir que a situação descrita a seguir ocorra}
  \seealsoref{免得}{mian3de5}
  \synonymref{免得}{mian3de5}
\end{EntryWithPhonetic}

\begin{EntryWithPhonetic}{以内}{yi3nei4}{4,4}{⼈,⼌}[HSK 4]
  \definition{adv.}{dentro de; menos que; não mais que; dentro de certos limites de tempo, premissas, quantidade e escopo}
  \synonymref{以下}{yi3xia4}
  \synonymref{之内}{zhi1nei4}
  \antonymref{以外}{yi3wai4}
\end{EntryWithPhonetic}

\begin{EntryWithPhonetic}{以期}{yi3qi1}{4,12}{⼈,⽉}
  \definition{conj.}{na esperança de | esperando por  | a fim de | aguardando por}
\end{EntryWithPhonetic}

\begin{EntryWithPhonetic}{以前}{yi3qian2}{4,9}{⼈,⼑}[HSK 2]
  \definition{s.}{antes; antigamente; anteriormente (no tempo); agora ou o período anterior ao tempo indicado}
  \synonymref{此前}{ci3qian2}
  \synonymref{从前}{cong2qian2}
  \synonymref{当年}{dang1nian2}
  \synonymref{当年}{dang4nian2}
  \synonymref{过去}{guo4 qu5}
  \synonymref{往常}{wang3chang2}
  \synonymref{往日}{wang3ri4}
  \synonymref{昔日}{xi1ri4}
  \synonymref{之前}{zhi1qian2}
  \antonymref{此刻}{ci3ke4}
  \antonymref{当前}{dang1qian2}
  \antonymref{将来}{jiang1lai2}
  \antonymref{今后}{jin1hou4}
  \antonymref{往后}{wang3hou4}
  \antonymref{未来}{wei4lai2}
  \antonymref{以后}{yi3hou4}
  \antonymref{正在}{zheng4zai4}
\end{EntryWithPhonetic}

\begin{EntryWithPhonetic}{以求}{yi3qiu2}{4,7}{⼈,⽔}
  \definition{conj.}{a fim de; numa tentativa de | Literário: esperança de; tentativa de}
\end{EntryWithPhonetic}

\begin{EntryWithPhonetic}{以色列}{yi3se4lie4}{4,6,6}{⼈,⾊,⼑}
  \definition*{s.}{Israel}
\end{EntryWithPhonetic}

\begin{EntryWithPhonetic}{以上}{yi3shang4}{4,3}{⼈,⼀}[HSK 2]
  \definition[本]{s.}{mais do que; sobre; acima; indica posição, ordem ou número acima de um determinado ponto | o acima; o precedente; o acima mencionado; refere"-se às palavras anteriores}
  \antonymref{以下}{yi3xia4}
\end{EntryWithPhonetic}

\begin{EntryWithPhonetic}{以身作则}{yi3shen1-zuo4ze2}{4,7,7,6}{⼈,⾝,⼈,⼑}[HSK 7-9]
  \definition{expr.}{liderar pelo exemplo; dar o exemplo; praticar o que se prega; dê o exemplo com suas próprias ações}
\end{EntryWithPhonetic}

\begin{EntryWithPhonetic}{以外}{yi3wai4}{4,5}{⼈,⼣}[HSK 2]
  \definition{s.}{além; exceto; fora; diferente de; fora dos limites de um determinado tempo, quantidade ou lugar}
  \synonymref{除外}{chu2wai4}
  \synonymref{之外}{zhi1wai4}
  \antonymref{以内}{yi3nei4}
\end{EntryWithPhonetic}

\begin{EntryWithPhonetic}{以往}{yi3wang3}{4,8}{⼈,⼻}[HSK 5]
  \definition{s.}{antes; anterior; no passado}
  \seealsoref{以前}{yi3qian2}
  \synonymref{曾经}{ceng2jing1}
  \synonymref{此前}{ci3qian2}
  \synonymref{过往}{guo4wang3}
  \synonymref{先前}{xian1qian2}
  \synonymref{现在}{xian4zai4}
  \antonymref{此刻}{ci3ke4}
  \antonymref{如今}{ru2jin1}
\end{EntryWithPhonetic}

\begin{EntryWithPhonetic}{以为}{yi3wei2}{4,4}{⼈,⼂}[HSK 2]
  \definition{v.}{pensar; acreditar; considerar (geralmente erroneamente); expressa opiniões e atitudes em relação às coisas, geralmente erradas}
  \synonymref{合计}{he2ji4}
  \synonymref{合计}{he2ji5}
  \synonymref{觉得}{jue2de5}
  \synonymref{认为}{ren4wei2}
  \antonymref{怀疑}{huai2yi2}
\end{EntryWithPhonetic}

\begin{EntryWithPhonetic}{以下}{yi3xia4}{4,3}{⼈,⼀}[HSK 2]
  \definition[所]{s.}{abaixo; sob; indica posição, ordem ou número abaixo de um certo ponto | seguinte; refere"-se às seguintes palavras}
  \synonymref{保护}{bao3hu4}
  \synonymref{履行}{lv3xing2}
  \synonymref{实行}{shi2xing2}
  \synonymref{羞愧}{xiu1kui4}
  \synonymref{以内}{yi3nei4}
  \synonymref{之下}{zhi1xia4}
  \antonymref{以上}{yi3shang4}
\end{EntryWithPhonetic}

\begin{EntryWithPhonetic}{以至}{yi3zhi4}{4,6}{⼈,⾄}
  \definition{conj.}{até; usado para indicar uma extensão em grau, escopo, quantidade, tempo |  então\dots isso\dots; a tal ponto que\dots; uma cláusula de conexão é usada no início da cláusula seguinte para expressar o resultado da situação anterior}
  \antonymref{乃至}{nai3zhi4}
  \antonymref{甚至}{shen4zhi4}
  \antonymref{以致}{yi3zhi4}
\end{EntryWithPhonetic}

\begin{EntryWithPhonetic}{以至于}{yi3zhi4yu2}{4,6,3}{⼈,⾄,⼆}[HSK 7-9]
  \definition{conj.}{na medida em que\dots}
  \definition{conj.}{então\dots que\dots; na medida em que\dots; usado no início da última cláusula de uma frase complexa, significando "como resultado"}
  \seealsoref{以至}{yi3zhi4}
\end{EntryWithPhonetic}

\begin{EntryWithPhonetic}{以致}{yi3zhi4}{4,10}{⼈,⾄}[HSK 7-9]
  \definition{conj.}{de modo que; consequentemente; como resultado; com a consequência de que; utilizada antes da segunda frase em uma estrutura de duas frases, indica que o resultado negativo na segunda frase é causado pela situação na primeira frase}
  \synonymref{乃至}{nai3zhi4}
  \synonymref{使得}{shi3de5}
  \synonymref{以至}{yi3zhi4}
\end{EntryWithPhonetic}

%%%%%%%%%% 尾 %%%%%%%%%%
\subsection*{尾}\addcontentsline{loh}{figure}{尾 \dpy{yi3}}

\begin{EntryWithPhonetic}{尾}{yi3}{7}{⼫}
  \definition{s.}{rabo do cavalo | parte posterior pontiaguda de um gafanhoto etc.}
  \seeref{wei3}
\end{EntryWithPhonetic}

%%%%%%%%%% 矣 %%%%%%%%%%
\subsection*{矣}\addcontentsline{loh}{figure}{矣 \dpy{yi3}}

\begin{EntryWithPhonetic}{矣}{yi3}{7}{⽮}[HSK 7-9]
  \definition{part.}{partícula final clássica, indica um tom declarativo e é equivalente a 了 | em frases exclamativas}
  \seealsoref{了}{le5}
\end{EntryWithPhonetic}

%%%%%%%%%% 倚 %%%%%%%%%%
\subsection*{倚}\addcontentsline{loh}{figure}{倚 \dpy{yi3}}

\begin{EntryWithPhonetic}{倚}{yi3}{10}{⼈}[HSK 7-9]
  \definition*{s.}{Sonbrenome: Yi}
  \definition{adj.}{Literário: tendencioso; parcial; distorcido; torto}
  \definition{v.}{apoiar"-se em ou encostar"-se em; descansar em ou encostar"-se em | confiar em; contar com | depender de}
\end{EntryWithPhonetic}

%%%%%%%%%% 椅 %%%%%%%%%%
\subsection*{椅}\addcontentsline{loh}{figure}{椅 \dpy{yi3}}

\begin{EntryWithPhonetic}{椅}{yi3}{12}{⽊}
  \definition*{s.}{Sobrenome: Yi}
  \definition{s.}{cadeira}
\end{EntryWithPhonetic}

\begin{EntryWithPhonetic}{椅子}{yi3zi5}{12,3}{⽊,⼦}[HSK 2]
  \definition[把,套,排]{s.}{cadeira; assentos com encosto, feitos principalmente de madeira, bambu, rattan, etc.; móveis com pernas, mas sem encosto para as pessoas se sentarem}
\end{EntryWithPhonetic}

%%%%%%%%%% 一 %%%%%%%%%%
\subsection*{一}\addcontentsline{loh}{figure}{一 \dpy{yi4}}

\begin{EntryWithPhonetic}{一}{yi4}[(一 $+$ 1º, 2º e 3º tom)]{1}{⼀}[Kangxi 1]
  \definition{adv.}{uma vez | assim que | ao}
  \definition{num.}{um; 1 | um (artigo)}
  \seeref{yi1}
  \seeref{yi2}
\end{EntryWithPhonetic}

\begin{EntryWithPhonetic}{一般}{yi4ban1}{1,10}{⼀,⾈}[HSK 2]
  \definition{adj.}{o mesmo que; exatamente como | geral; ordinário; comum | médio; medíocre; o grau ou nível não é muito alto}
  \definition{adv.}{frequentemente; geralmente}
  \synonymref{凡是}{fan2shi4}
  \synonymref{平常}{ping2chang2}
  \synonymref{普遍}{pu3bian4}
  \synonymref{普通}{pu3tong1}
  \synonymref{日常}{ri4chang2}
  \synonymref{通常}{tong1chang2}
  \synonymref{一些}{yi4xie1}
  \antonymref{出色}{chu1se4}
  \antonymref{非常}{fei1chang2}
  \antonymref{格外}{ge2wai4}
  \antonymref{个别}{ge4bie2}
  \antonymref{特别}{te4bie2}
  \antonymref{特殊}{te4shu1}
  \antonymref{突出}{tu1/chu1}
  \antonymref{无比}{wu2bi3}
  \antonymref{新颖}{xin1ying3}
  \antonymref{尤其}{you2qi2}
\end{EntryWithPhonetic}

\begin{EntryWithPhonetic}{一般来说}{yi4ban1lai2shuo1}{1,10,7,9}{⼀,⾈,⽊,⾔}[HSK 4]
  \definition{expr.}{de modo geral; na média; no caso usual; a declaração usual}
\end{EntryWithPhonetic}

\begin{EntryWithPhonetic}{一边}{yi4bian1}{1,5}{⼀,⾡}[HSK 1]
  \definition{adj.}{igual; idêntico; da mesma forma}
  \definition{adv.}{enquanto; ao mesmo tempo; simultaneamente; indica que uma ação ocorre simultaneamente a outra ação}
  \definition{s.}{lado; um lado; um aspecto | ambos os lados; ao lado de}
  \synonymref{侧面}{ce4mian4}
  \synonymref{一面}{yi2mian4}
  \antonymref{四处}{si4chu4}
\end{EntryWithPhonetic}

\begin{EntryWithPhonetic}{一长一短}{yi4 chang2 yi4 duan3}{1,4,1,12}{⼀,⾧,⼀,⽮}[HSK 7-9]
  \definition{expr.}{um longo e um curto}
\end{EntryWithPhonetic}

\begin{EntryWithPhonetic}{一成不变}{yi4cheng2-bu2bian4}{1,6,4,8}{⼀,⼽,⼀,⼜}[HSK 7-9]
  \definition{expr.}{imutável e fixo; ser imutável; ser inflexível; inalterável; definido e definitivo; cair na rotina; rígido e inflexível (fórmula, etc.); inflexível; invariável; monótono; imutável}
  \synonymref{刻舟求剑}{ke4zhou1-qiu2jian4}
  \synonymref{一模一样}{yi4mu2-yi2yang4}
  \antonymref{沧海桑田}{cang1 hai3 sang1 tian2}
  \antonymref{翻天覆地}{fan1tian1-fu4di4}
  \antonymref{面目全非}{mian4mu4-quan2fei1}
  \antonymref{日新月异}{ri4xin1-yue4yi4}
  \antonymref{水涨船高}{shui3zhang3-chuan2gao1}
  \antonymref{五花八门}{wu3hua1-ba1men2}
  \antonymref{朝三暮四}{zhao1san1-mu4si4}
\end{EntryWithPhonetic}

\begin{EntryWithPhonetic}{一筹莫展}{yi4chou2-mo4zhan3}{1,13,10,10}{⼀,⽵,⾋,⼫}[HSK 7-9]
  \definition{expr.}{estar sem saber o que fazer; estar completamente perdido; não conseguir apresentar nenhum plano ou solução}
  \antonymref{千方百计}{qian1fang1-bai3ji4}
\end{EntryWithPhonetic}

\begin{EntryWithPhonetic}{一点点}{yi4dian3dian3}{1,9,9}{⼀,⽕,⽕}[HSK 2]
  \definition{adj.}{um pouco; muito pouco ou um pouquinho}
\end{EntryWithPhonetic}

\begin{EntryWithPhonetic}{一点儿}{yi4dian3r5}{1,9,2}{⼀,⽕,⼉}[HSK 1]
  \definition{adv.}{um pouco; uma pitada; uma gota; uma amostra; uma pequena quantidade; ({adj.} + (一)点儿, 一点儿 + {s.} ou 有 + (一)点儿 + {s.})}
  \synonymref{一点点}{yi4dian3dian3}
\end{EntryWithPhonetic}

\begin{EntryWithPhonetic}{一帆风顺}{yi4fan1-feng1shun4}{1,6,4,9}{⼀,⼱,⾵,⾴}[HSK 7-9]
  \definition{expr.}{``Vento favorável durante toda a viagem.''; navegação tranquila; navegação sem problemas; ter uma boa viagem; tudo correr bem}
  \synonymref{心想事成}{xin1xiang3-shi4cheng2}
  \synonymref{一路平安}{yi2lu4-ping2'an1}
  \synonymref{一路顺风}{yi2lu4shun4feng1}
  \antonymref{碰钉子}{peng4 ding1zi5}
\end{EntryWithPhonetic}

\begin{EntryWithPhonetic}{一番}{yi4 fan1}{1,12}{⼀,⽥}[HSK 6]
  \definition{adv.}{uma demonstração de, uma dose de, um pedaço de (conversa, investigação, pensamento)}
  \synonymref{一阵}{yi2zhen4}
  \antonymref{多次}{duo1ci4}
\end{EntryWithPhonetic}

\begin{EntryWithPhonetic}{一方面}{yi4fang1mian4}{1,4,9}{⼀,⽅,⾯}[HSK 3]
  \definition{s.}{um lado; um dos dois aspectos opostos ou um lado de algo que está relacionado a outro}
  \seealsoref{一方面……，一方面……}{yi4fang1mian4 yi4fang1mian4}
  \synonymref{一部分}{yi2bu4fen5}
\end{EntryWithPhonetic}

\begin{EntryWithPhonetic}{一方面……，一方面……}{yi4fang1mian4 yi4fang1mian4}{1,4,9,1,4,9}{⼀,⽅,⾯,⼀,⽅,⾯}
  \definition{conj.}{por um lado\dots, por outro lado\dots; conecta duas orações paralelas (devem ser usadas juntas)}[\underline{一方面}觉得兴奋，\underline{一方面}又害怕。===Por um lado, sinto"-me entusiasmado, mas, por outro, também sinto medo.]
\end{EntryWithPhonetic}

\begin{EntryWithPhonetic}{一干二净}{yi4gan1-er4jing4}{1,3,2,8}{⼀,⼲,⼆,⼎}[HSK 7-9]
  \definition{expr.}{limpo e impecável; completamente; totalmente; é uma expressão idiomática chinesa que significa ``completamente limpo'' ou ``totalmente desprovido de qualquer coisa''; tem origem no romance da Dinastia Qing, Flores no Espelho, de Li Ruzhen, 李汝珍}
\end{EntryWithPhonetic}

\begin{EntryWithPhonetic}{一鼓作气}{yi4gu3-zuo4qi4}{1,13,7,4}{⼀,⿎,⼈,⽓}[HSK 7-9]
  \definition{expr.}{realizar algo com um esforço vigoroso; concluir com um único esforço; de uma só vez; ao primeiro toque do tambor, a coragem se desperta; ser estimulado ao som do primeiro tambor; preparar os nervos para um esforço; fazer algo sem trégua; de uma só vez; fazer um esforço vigoroso para terminar algo de uma só vez; prosseguir até o fim sem trégua; despertar o espírito com a batida do tambor; tomar coragem com ambas as mãos; o Zuo Zhuan, Duque Zhuang, 庄公, Ano 10, afirma: ``Na guerra, a coragem é fundamental. Um toque de tambor desperta o espírito, um segundo o enfraquece e um terceiro o exaure.''}
\end{EntryWithPhonetic}

\begin{EntryWithPhonetic}{一锅粥}{yi4guo1zhou1}{1,12,12}{⼀,⾦,⽶}[HSK 7-9]
  \definition{expr.}{uma panela de mingau; uma completa bagunça; tudo em turbilhão}
\end{EntryWithPhonetic}

\begin{EntryWithPhonetic}{一晃}{yi4huang3}{1,10}{⼀,⽇}
  \definition{v.}{piscar; cintilar; piscar muito rapidamente}
  \seeref{yi2huang4}
  \synonymref{毕竟}{bi4jing4}
\end{EntryWithPhonetic}

\begin{EntryWithPhonetic}{一回事}{yi4hui2shi4}{1,6,8}{⼀,⼞,⼅}[HSK 7-9]
  \definition{s.}{uma e a mesma (coisa); a mesma coisa}
\end{EntryWithPhonetic}

\begin{EntryWithPhonetic}{一家人}{yi4jia1ren2}{1,10,2}{⼀,⼧,⼈}[HSK 7-9]
  \definition{s.}{todos da mesma família; uma família}
\end{EntryWithPhonetic}

\begin{EntryWithPhonetic}{一经}{yi4jing1}{1,8}{⼀,⽷}[HSK 7-9]
  \definition{adv.}{uma vez; assim que; significa que, desde que uma determinada etapa ou ação seja tomada, um determinado resultado será alcançado imediatamente}
  \synonymref{已经}{yi3jing1}
\end{EntryWithPhonetic}

\begin{EntryWithPhonetic}{一举}{yi4ju3}{1,9}{⼀,⼂}[HSK 7-9]
  \definition{adv.}{em uma única ação; de uma só vez; com um único esforço; de repente}
  \definition{s.}{uma ação; um golpe; um único movimento brusco}
  \seealsoref{一下子}{yi2xia4zi5}
\end{EntryWithPhonetic}

\begin{EntryWithPhonetic}{一举一动}{yi4ju3-yi2dong4}{1,9,1,6}{⼀,⼂,⼀,⼒}[HSK 7-9]
  \definition{expr.}{cada movimento; cada ação}
  \synonymref{所作所为}{suo3zuo4-suo3wei2}
  \synonymref{一言一行}{yi4yan2-yi4xing2}
\end{EntryWithPhonetic}

\begin{EntryWithPhonetic}{一卡通}{yi4ka3tong1}{1,5,10}{⼀,⼘,⾡}[HSK 7-9]
  \definition*{s.}{Yikatong (cartão inteligente de transporte público de Pequim)}
  \definition{s.}{cartão de crédito universal; cartão multifuncional}
\end{EntryWithPhonetic}

\begin{EntryWithPhonetic}{一口气}{yi4kou3qi4}{1,3,4}{⼀,⼝,⽓}[HSK 5]
  \definition{adv.}{em um só fôlego; sem pausa; fazer algo continuamente}
  \synonymref{连续}{lian2xu4}
\end{EntryWithPhonetic}

\begin{EntryWithPhonetic}{一连}{yi4lian2}{1,7}{⼀,⾡}[HSK 7-9]
  \definition{adv.}{em sequência; em fila; em sucessão; isso indica uma série contínua de ações ou situações, enfatizando uma grande quantidade ou uma longa duração}
  \seealsoref{接连}{jie1lian2}
  \seealsoref{连续}{lian2xu4}
  \seealsoref{一口气}{yi4kou3qi4}
  \synonymref{持续}{chi2xu4}
  \synonymref{继续}{ji4xu4}
  \synonymref{接连}{jie1lian2}
  \synonymref{连续}{lian2xu4}
  \synonymref{陆续}{lu4xu4}
  \synonymref{延续}{yan2xu4}
\end{EntryWithPhonetic}

\begin{EntryWithPhonetic}{一连串}{yi4lian2chuan4}{1,7,7}{⼀,⾡,⼁}[HSK 7-9]
  \definition{adj.}{uma cadeia de; uma série de; uma sequência de; uma sucessão de; (ações, eventos, etc.) um após o outro}
  \synonymref{一系列}{yi2xi4lie4}
\end{EntryWithPhonetic}

\begin{EntryWithPhonetic}{一流}{yi4liu2}{1,10}{⼀,⽔}[HSK 5]
  \definition{adj.}{clássico; de primeira linha; de primeira classe; o melhor}
  \definition[些]{s.}{tipo; mesmo tipo; da mesma classe; da mesma categoria; uma categoria}
  \synonymref{顶尖}{ding3jian1}
  \synonymref{优秀}{you1xiu4}
  \synonymref{优异}{you1yi4}
\end{EntryWithPhonetic}

\begin{EntryWithPhonetic}{一毛不拔}{yi4mao2-bu4ba2}{1,4,4,8}{⼀,⽑,⼀,⼿}[HSK 7-9]
  \definition{expr.}{avarento demais para dar um fio de cabelo; mesquinho como um avarento; ser pão"-duro; extremamente mesquinho; não dar um centavo a; não ceder nem um fio de cabelo; não mover um dedo para ajudar; recusar"-se a contribuir com um único centavo; relutante em ceder nem um fio de cabelo; muito econômico com dinheiro; muito avarento | Mencius, Livro 7, Capítulo 1: ``Yang Zi não arrancaria um único fio de cabelo para o benefício do mundo, mesmo que fosse para o meu próprio benefício.''}
\end{EntryWithPhonetic}

\begin{EntryWithPhonetic}{一模一样}{yi4mu2-yi2yang4}{1,14,1,10}{⼀,⽊,⼀,⽊}[HSK 6]
  \definition{expr.}{tão parecidos quanto duas ervilhas; ser exatamente iguais; muito parecido, a mesma aparência}
  \synonymref{大同小异}{da4tong2-xiao3yi4}
  \synonymref{一成不变}{yi4cheng2-bu2bian4}
  \antonymref{截然不同}{jie2ran2-bu4tong2}
\end{EntryWithPhonetic}

\begin{EntryWithPhonetic}{一年到头}{yi4nian2-dao4tou2}{1,6,8,5}{⼀,⼲,⼑,⼤}[HSK 7-9]
  \definition{expr.}{ao longo do ano; durante todo o ano; na temporada e fora de temporada; ano após ano; o ano todo; o ano inteiro}
  \synonymref{一天到晚}{yi4tian1-dao4wan3}
\end{EntryWithPhonetic}

\begin{EntryWithPhonetic}{一旁}{yi4pang2}{1,10}{⼀,⽅}[HSK 7-9]
  \definition{s.}{de um lado; ao lado de; na lateral do campo}
\end{EntryWithPhonetic}

\begin{EntryWithPhonetic}{一齐}{yi4qi2}{1,6}{⼀,⿑}[HSK 6]
  \definition{adv.}{juntos; em uníssono; simultaneamente; ao mesmo tempo; indica que diferentes sujeitos emitem simultaneamente o mesmo comportamento ou o mesmo sujeito emite vários comportamentos diferentes ao mesmo tempo}
  \synonymref{全部}{quan2bu4}
  \synonymref{所有}{suo3you3}
  \synonymref{一块}{yi2kuai4}
  \synonymref{一路}{yi2lu4}
  \synonymref{一切}{yi2qie4}
  \synonymref{一起}{yi4qi3}
  \synonymref{一同}{yi4tong2}
  \antonymref{分散}{fen1san4}
  \antonymref{先后}{xian1hou4}
\end{EntryWithPhonetic}

\begin{EntryWithPhonetic}{一起}{yi4qi3}{1,10}{⼀,⾛}[HSK 1]
  \definition{adv.}{juntos; em companhia; indica o mesmo local, ao mesmo tempo que se faz algo | no total; em todos; no conjunto}
  \definition{s.}{no mesmo lugar}
  \synonymref{全部}{quan2bu4}
  \synonymref{所有}{suo3you3}
  \synonymref{同步}{tong2bu4}
  \synonymref{一道}{yi2dao4}
  \synonymref{一块}{yi2kuai4}
  \synonymref{一路}{yi2lu4}
  \synonymref{一切}{yi2qie4}
  \synonymref{一齐}{yi4qi2}
  \synonymref{一同}{yi4tong2}
  \antonymref{独自}{du2zi4}
\end{EntryWithPhonetic}

\begin{EntryWithPhonetic}{一如既往}{yi4ru2-ji4wang3}{1,6,9,8}{⼀,⼥,⽆,⼻}[HSK 7-9]
  \definition{expr.}{como no passado; como sempre; como antes; exatamente como antes; significa que a atitude não mudou e continua exatamente a mesma de antes}
  \antonymref{翻天覆地}{fan1tian1-fu4di4}
  \antonymref{面目全非}{mian4mu4-quan2fei1}
  \antonymref{取而代之}{qu3'er2dai4zhi1}
\end{EntryWithPhonetic}

\begin{EntryWithPhonetic}{一身}{yi4shen1}{1,7}{⼀,⾝}[HSK 5]
  \definition{s.}{o corpo inteiro; em todo o corpo | um terno; (um conjunto completo de) roupas | sozinho; uma única pessoa; relativo a uma única pessoa}
\end{EntryWithPhonetic}

\begin{EntryWithPhonetic}{一生}{yi4sheng1}{1,5}{⼀,⽣}[HSK 2]
  \definition{s.}{vida inteira; toda a vida; ao longo da vida; todo o tempo desde o nascimento até a morte; às vezes exagerado para indicar um longo período de tempo no curso da vida}
  \synonymref{一辈子}{yi2bei4zi5}
  \synonymref{终身}{zhong1shen1}
  \antonymref{瞬间}{shun4jian1}
\end{EntryWithPhonetic}

\begin{EntryWithPhonetic}{一声不吭}{yi4sheng1-bu4keng1}{1,7,4,7}{⼀,⼠,⼀,⼝}[HSK 7-9]
  \definition{expr.}{``Sem dizer uma palavra.''}
  \synonymref{一言不发}{yi4yan2-bu4fa1}
\end{EntryWithPhonetic}

\begin{EntryWithPhonetic}{一时}{yi4shi2}{1,7}{⼀,⽇}[HSK 6]
  \definition{adv.}{por um curto período; temporário | (usado em pares) agora\dots, agora\dots; este momento\dots, e o próximo\dots; o mesmo que 时而}
  \definition{s.}{um período de tempo | um momento; um breve momento; um tempo muito curto}
  \seealsoref{时而}{shi2'er2}
  \seealsoref{一时……，一时……}{yi4shi2 yi4shi2}
  \synonymref{间或}{jian4huo4}
  \synonymref{临时}{lin2shi2}
  \synonymref{偶尔}{ou3'er3}
  \synonymref{有时}{you3shi2}
  \synonymref{暂时}{zan4shi2}
  \antonymref{常常}{chang2chang2}
  \antonymref{持久}{chi2jiu3}
  \antonymref{时常}{shi2chang2}
\end{EntryWithPhonetic}

\begin{EntryWithPhonetic}{一时……，一时……}{yi4shi2 yi4shi2}{1,7,1,7}{⼀,⽇,⼀,⽇}
  \definition{adv.}{por um tempo\dots, por um tempo\dots}
  \seealsoref{一时}{yi4shi2}
\end{EntryWithPhonetic}

\begin{EntryWithPhonetic}{一手}{yi4shou3}{1,4}{⼀,⼿}[HSK 7-9]
  \definition{adv.}{sozinho; completamente por si só; totalmente sozinho}
  \definition{s.}{proficiência; habilidade | truque; movimento}
\end{EntryWithPhonetic}

\begin{EntryWithPhonetic}{一塌糊涂}{yi4ta1hu2tu2}{1,13,15,10}{⼀,⼟,⽶,⽔}[HSK 7-9]
  \definition{expr.}{em uma completa bagunça; em um estado terrível; uma grande confusão; fazer uma grande besteira; uma baita confusão; uma situação bastante complicada; em total desordem; fazer uma grande bagunça; aproveitar ao máximo}
  \synonymref{乱七八糟}{luan4qi1ba1zao1}
  \synonymref{杂乱无章}{za2luan4-wu2zhang1}
\end{EntryWithPhonetic}

\begin{EntryWithPhonetic}{一体}{yi4ti3}{1,7}{⼀,⼈}[HSK 7-9]
  \definition{s.}{todo orgânico (ou integral); eles estão intimamente relacionados, como um todo | todas as pessoas envolvidas; sem exceção}
\end{EntryWithPhonetic}

\begin{EntryWithPhonetic}{一天到晚}{yi4tian1-dao4wan3}{1,4,8,11}{⼀,⼤,⼑,⽇}[HSK 7-9]
  \definition{expr.}{o dia todo}
  \synonymref{从早到晚}{cong2zao3-dao4wan3}
  \synonymref{日复一日}{ri4fu4yi2ri4}
  \synonymref{一年到头}{yi4nian2-dao4tou2}
\end{EntryWithPhonetic}

\begin{EntryWithPhonetic}{一同}{yi4tong2}{1,6}{⼀,⼝}[HSK 6]
  \definition{adv.}{juntos; ao mesmo tempo e lugar}
  \synonymref{一块}{yi2kuai4}
  \synonymref{一路}{yi2lu4}
  \synonymref{一齐}{yi4qi2}
  \synonymref{一起}{yi4qi3}
  \antonymref{分开}{fen1/kai1}
  \antonymref{混合}{hun4he2}
\end{EntryWithPhonetic}

\begin{EntryWithPhonetic}{一头}{yi4tou2}{1,5}{⼀,⼤}[HSK 7-9]
  \definition{adv.}{simultaneamente; indica fazer várias coisas ao mesmo tempo | diretamente; de cabeça; de repente | de repente; de uma só vez; rapidamente}
  \definition{clas.}{uma cabeça (de gado, porco, mula, etc.)}
  \definition{s.}{final; extremidade (ferramentas, armas, etc.) | lado; aspecto | um grupo de; mesmo lado; mesmo aspecto | uma cabeça cheia de algo (coberta de cabelos grisalhos, poeira, pele, doenças, etc.)}
\end{EntryWithPhonetic}

\begin{EntryWithPhonetic}{一无所知}{yi4wu2suo3zhi1}{1,4,8,8}{⼀,⽆,⼾,⽮}[HSK 7-9]
  \definition[本]{expr.}{``Não saber absolutamente nada.''; não saber nada sobre; não ter a mínima ideia de; ser completamente ignorante de; sem a menor ideia}
  \synonymref{不得而知}{bu4de2'er2zhi1}
\end{EntryWithPhonetic}

\begin{EntryWithPhonetic}{一些}{yi4xie1}{1,8}{⼀,⼆}[HSK 1]
  \definition{clas.}{alguns; um número de; quantidade indeterminada | um pouco; uma pequena quantidade | mais de um; mais de uma vez; indica mais de um ou mais de uma vez, etc. | uma ligeira mudança no grau, intensidade; usado após certos verbos, adjetivos, etc., para indicar uma quantidade muito pequena}
  \definition{pron.}{uns; alguns}
  \synonymref{一般}{yi4ban1}
  \synonymref{有些}{you3xie1}
\end{EntryWithPhonetic}

\begin{EntryWithPhonetic}{一心}{yi4xin1}{1,4}{⼀,⼼}[HSK 7-9]
  \definition{adv.}{de todo o coração; com toda a alma | de uma só mente; em uníssono}
  \synonymref{用心}{yong4 xin1}
  \synonymref{专心}{zhuan1xin1}
  \antonymref{多心}{duo1xin1}
\end{EntryWithPhonetic}

\begin{EntryWithPhonetic}{一行}{yi4xing2}{1,6}{⼀,⾏}[HSK 6]
  \definition{s.}{delegação; um grupo viajando junto; festa}
\end{EntryWithPhonetic}

\begin{EntryWithPhonetic}{一言不发}{yi4yan2-bu4fa1}{1,7,4,5}{⼀,⾔,⼀,⼜}[HSK 7-9]
  \definition{expr.}{``Sem dizer uma palavra.''; não dizer (ou proferir) uma palavra; manter a boca fechada}
  \synonymref{一声不吭}{yi4sheng1-bu4keng1}
  \antonymref{滔滔不绝}{tao1tao1-bu4jue2}
  \antonymref{脱口而出}{tuo1kou3'er2chu1}
\end{EntryWithPhonetic}

\begin{EntryWithPhonetic}{一言一行}{yi4yan2-yi4xing2}{1,7,1,6}{⼀,⾔,⼀,⾏}[HSK 7-9]
  \definition{expr.}{``Cada palavra e ação.''; o que se diz e o que se faz}
  \synonymref{一举一动}{yi4ju3-yi2dong4}
\end{EntryWithPhonetic}

\begin{EntryWithPhonetic}{一眼}{yi4yan3}{1,11}{⼀,⽬}[HSK 7-9]
  \definition{expr.}{um olhar rápido; uma olhada rápida; um vislumbre; frequentemente usado para descrever uma percepção veloz ou imediata}
\end{EntryWithPhonetic}

\begin{EntryWithPhonetic}{一应俱全}{yi4ying1-ju4quan2}{1,7,10,6}{⼀,⼴,⼈,⼊}[HSK 7-9]
  \definition{expr.}{tudo o que é necessário está lá; os produtos estão disponíveis em todas as variedades; disponíveis em todas as variedades; completo em todas as linhas; completo com tudo; tudo o que é necessário é fornecido; do início ao fim; há uma linha completa de suprimentos; não falta nada; fornecer tudo; disponibilizar todas as variedades}
  \synonymref{面面俱到}{mian4mian4-ju4dao4}
  \synonymref{无微不至}{wu2wei1-bu2zhi4}
  \synonymref{应有尽有}{ying1you3-jin4you3}
  \antonymref{一无所有}{yi1wu2suo3you3}
\end{EntryWithPhonetic}

\begin{EntryWithPhonetic}{一早}{yi4zao3}{1,6}{⼀,⽇}[HSK 7-9]
  \definition{s.}{Coloquial: de manhã cedo}
\end{EntryWithPhonetic}

\begin{EntryWithPhonetic}{一直}{yi4zhi2}{1,8}{⼀,⽬}[HSK 2]
  \definition{adv.}{direto; indica que permanece inalterado em uma direção | sempre; continuamente; o tempo todo; o tempo todo; indica que a ação é sempre ininterrupta ou o estado é sempre inalterado | de um ponto a outro sem enfatizar nenhuma exceção}
  \synonymref{不断}{bu2duan4}
  \synonymref{不停}{bu4ting2}
  \synonymref{继续}{ji4xu4}
  \synonymref{连续}{lian2xu4}
  \synonymref{始终}{shi3zhong1}
  \synonymref{向来}{xiang4lai2}
  \synonymref{一贯}{yi2guan4}
  \synonymref{一向}{yi2xiang4}
  \antonymref{间断}{jian4duan4}
  \antonymref{中断}{zhong1duan4}
\end{EntryWithPhonetic}

%%%%%%%%%% 义 %%%%%%%%%%
\subsection*{义}\addcontentsline{loh}{figure}{义 \dpy{yi4}}

\begin{EntryWithPhonetic}{义}{yi4}{3}{⼂}
  \definition*{s.}{Sobrenome: Yi}
  \definition{adj.}{justo; equitativo | adotado; adotivo | juramentado | artificial; falso}
  \definition[个]{s.}{justiça; retidão | laços humanos; relacionamento | significado; importância}
\end{EntryWithPhonetic}

\begin{EntryWithPhonetic}{义工}{yi4gong1}{3,3}{⼂,⼯}[HSK 7-9]
  \definition{s.}{voluntário; isso se refere a indivíduos e grupos que, sem levar em conta a recompensa material, contribuem com seu tempo, energia e conhecimento técnico pessoal para melhorar a sociedade com base na moralidade, crenças, consciência, compaixão e responsabilidade}
  \synonymref{志愿者}{zhi4yuan4zhe3}
\end{EntryWithPhonetic}

\begin{EntryWithPhonetic}{义务}{yi4wu4}{3,5}{⼂,⼒}[HSK 4]
  \definition{adj.}{voluntário; fornecer serviços ou ajuda a outros gratuitamente}
  \definition[项]{s.}{dever; obrigação; responsabilidades perante a lei | obrigação moral; responsabilidade moral}
  \synonymref{负担}{fu4dan1}
  \synonymref{任务}{ren4wu5}
  \synonymref{无偿}{wu2chang2}
  \synonymref{仔肩}{zi1jian1}
  \antonymref{权利}{quan2li4}
  \antonymref{权力}{quan2li4}
  \antonymref{权益}{quan2yi4}
  \antonymref{责任}{ze2ren4}
\end{EntryWithPhonetic}

%%%%%%%%%% 亿 %%%%%%%%%%
\subsection*{亿}\addcontentsline{loh}{figure}{亿 \dpy{yi4}}

\begin{EntryWithPhonetic}{亿}{yi4}{3}{⼈}[HSK 2]
  \definition*{s.}{Sobrenome: Yi}
  \definition{num.}{cem milhões; 100.000.000; 1.0000.0000}
\end{EntryWithPhonetic}

%%%%%%%%%% 艺 %%%%%%%%%%
\subsection*{艺}\addcontentsline{loh}{figure}{艺 \dpy{yi4}}

\begin{EntryWithPhonetic}{艺}{yi4}{4}{⾋}
  \definition*{s.}{Sobrenome: Yi}
  \definition[个,种]{s.}{habilidade | arte | regra; norma | padrão; critério; diretrizes | limite}
  \definition{v.}{plantar; crescer}
\end{EntryWithPhonetic}

\begin{EntryWithPhonetic}{艺人}{yi4ren2}{4,2}{⾋,⼈}[HSK 6]
  \definition[位]{s.}{artista performático; ator profissional; ator ou artista (em teatro local, narradores, acrobacia ou outro show business); atores de ópera, arte popular, acrobacia, cinema e televisão, etc. | artesão; artífice}
\end{EntryWithPhonetic}

\begin{EntryWithPhonetic}{艺术}{yi4shu4}{4,5}{⾋,⽊}[HSK 3]
  \definition{adj.}{artístico,único; elegante}
  \definition[个,种,门,项,类]{s.}{arte; literatura e arte | habilidade; arte; ofício; métodos criativos}
\end{EntryWithPhonetic}

%%%%%%%%%% 艾 %%%%%%%%%%
\subsection*{艾}\addcontentsline{loh}{figure}{艾 \dpy{yi4}}

\begin{EntryWithPhonetic}{艾}{yi4}{5}{⾋}
  \definition{adj.}{estável}
  \definition{v.}{ser corrigido; estar corrigido}
  \seeref{ai4}
\end{EntryWithPhonetic}

%%%%%%%%%% 议 %%%%%%%%%%
\subsection*{议}\addcontentsline{loh}{figure}{议 \dpy{yi4}}

\begin{EntryWithPhonetic}{议}{yi4}{5}{⾔}[HSK 7-9]
  \definition[个,则,条]{s.}{opinião; visão}
  \definition{v.}{discutir; trocar pontos de vista sobre; conversar sobre | comentar; observar | fofocar; comentar}
\end{EntryWithPhonetic}

\begin{EntryWithPhonetic}{议程}{yi4cheng2}{5,12}{⾔,⽲}[HSK 7-9]
  \definition[项,个]{s.}{agenda; procedimentos para discussão de propostas na reunião}
\end{EntryWithPhonetic}

\begin{EntryWithPhonetic}{议会}{yi4hui4}{5,6}{⾔,⼈}[HSK 7-9]
  \definition{s.}{congresso; parlamento; assembleia legislativa; em alguns países, o parlamento é a autoridade máxima ou o órgão responsável por promulgar e alterar leis; os membros do parlamento são eleitos}
  \synonymref{国会}{guo2hui4}
\end{EntryWithPhonetic}

\begin{EntryWithPhonetic}{议论}{yi4lun4}{5,6}{⾔,⾔}[HSK 4]
  \definition[个]{s.}{comentário; discussão; opiniões ou pontos de vista sobre o que é bom ou ruim, certo ou errado em relação a pessoas ou coisas}
  \definition{v.}{discutir; comentar; falar sobre; expressar opiniões e trocar pontos de vista sobre o bom, o ruim, o certo e o errado de pessoas ou coisas}
\end{EntryWithPhonetic}

\begin{EntryWithPhonetic}{议题}{yi4ti2}{5,15}{⾔,⾴}[HSK 6]
  \definition[项,个]{s.}{assunto; assunto em discussão; tópico para discussão}
\end{EntryWithPhonetic}

\begin{EntryWithPhonetic}{议员}{yi4yuan2}{5,7}{⾔,⼝}[HSK 7-9]
  \definition{s.}{membro de uma assembleia legislativa; deputado; (Reino Unido) membro do parlamento; (EUA) congressista; (Japão) membro da assembleia parlamentar; membros que possuem representação formal e direito a voto no parlamento}
\end{EntryWithPhonetic}

%%%%%%%%%% 亦 %%%%%%%%%%
\subsection*{亦}\addcontentsline{loh}{figure}{亦 \dpy{yi4}}

\begin{EntryWithPhonetic}{亦}{yi4}{6}{⼇}[HSK 7-9]
  \definition*{s.}{Sobrenome: Yi}
  \definition{adv.}{também; também (que significa o mesmo)}
\end{EntryWithPhonetic}

%%%%%%%%%% 屹 %%%%%%%%%%
\subsection*{屹}\addcontentsline{loh}{figure}{屹 \dpy{yi4}}

\begin{EntryWithPhonetic}{屹}{yi4}{6}{⼭}
  \definition{adj.}{Literário: elevando como um pico de montanha; representando um pico de montanha imponente)}
  \seeref{ge1}
\end{EntryWithPhonetic}

\begin{EntryWithPhonetic}{屹立}{yi4li4}{6,5}{⼭,⽴}[HSK 7-9]
  \definition{v.}{erguer"-se imponente como um gigante; ficar ereto; erguer"-se alta e firme como o pico de uma montanha, é frequentemente usada como metáfora para algo constante e inabalável}
  \synonymref{耸立}{song3li4}
  \synonymref{挺拔}{ting3ba2}
  \synonymref{挺立}{ting3li4}
\end{EntryWithPhonetic}

%%%%%%%%%% 异 %%%%%%%%%%
\subsection*{异}\addcontentsline{loh}{figure}{异 \dpy{yi4}}

\begin{EntryWithPhonetic}{异}{yi4}{6}{⼶}
  \definition{adj.}{diferente | estranho; incomum; extraordinário; especial | outro}
  \definition{v.}{surpreender | separar; divorciar"-se}
\end{EntryWithPhonetic}

\begin{EntryWithPhonetic}{异常}{yi4chang2}{6,11}{⼶,⼱}[HSK 6]
  \definition{adj.}{incomum; anormal; descreve uma situação diferente do normal}
  \definition{adv.}{extremamente; particularmente; excepcionalmente; descreve uma situação que atingiu um nível extremamente alto}
\end{EntryWithPhonetic}

\begin{EntryWithPhonetic}{异口同声}{yi4kou3-tong2sheng1}{6,3,6,7}{⼶,⼝,⼝,⼠}[HSK 7-9]
  \definition{expr.}{a uma só voz; em uníssono; pessoas diferentes dizem a mesma coisa; isso descreve uma situação em que a opinião de todos é completamente idêntica; falar em uníssono; concordar em comum acordo; todos concordam; estar todos na mesma história; por consenso comum; gritar em uníssono; juntar"-se ao mesmo coro; uma voz vinda de bocas diferentes; falar em conjunto (coro; uníssono); unanimemente}
  \synonymref{不约而同}{bu4yue1'er2tong2}
\end{EntryWithPhonetic}

\begin{EntryWithPhonetic}{异想天开}{yi4xiang3-tian1kai1}{6,13,4,4}{⼶,⼼,⼤,⼶}[HSK 7-9]
  \definition{expr.}{dar asas à fantasia mais desvairada; pedir (clamar) pela lua; ser uma ilusão fantasiosa; estar completamente deslumbrado; ter ideias fantásticas; esperar maravilhas; ideias fantasiosas; fantástico; bizarro; ter ideias mirabolantes na cabeça; dar asas à imaginação; ter uma ideia muito fantástica; ter morcegos no campanário; ter abelhas na cabeça; ter a cabeça cheia de abelhas; ter ideias estranhas e fantásticas; ter ideias mirabolantes; (tais) ideias são fantasiosas e falaciosas; imaginar algo impossível; dar asas à imaginação; inventivo; deixar a imaginação correr solta; esticar a imaginação; caprichoso; pensamento ilusório; esperanças desmedidas}
  \synonymref{胡思乱想}{hu2si1-luan4xiang3}
  \antonymref{实事求是}{shi2shi4-qiu2shi4}
\end{EntryWithPhonetic}

\begin{EntryWithPhonetic}{异性}{yi4xing4}{6,8}{⼶,⼼}[HSK 7-9]
  \definition{adj.}{diferente por natureza}
  \definition{s.}{sexo oposto; pessoas de diferentes gêneros; coisas de diferentes naturezas}
\end{EntryWithPhonetic}

\begin{EntryWithPhonetic}{异议}{yi4yi4}{6,5}{⼶,⾔}[HSK 7-9]
  \definition{s.}{discordância; objeção; opiniões diferentes}
  \synonymref{反驳}{fan3bo2}
  \synonymref{反对}{fan3dui4}
  \antonymref{赞同}{zan4tong2}
\end{EntryWithPhonetic}

%%%%%%%%%% 衣 %%%%%%%%%%
\subsection*{衣}\addcontentsline{loh}{figure}{衣 \dpy{yi4}}

\begin{EntryWithPhonetic}{衣}{yi4}{6}{⾐}
  \definition{v.}{vestir"-se; vestir alguém}
  \seeref{yi1}
\end{EntryWithPhonetic}

%%%%%%%%%% 抑 %%%%%%%%%%
\subsection*{抑}\addcontentsline{loh}{figure}{抑 \dpy{yi4}}

\begin{EntryWithPhonetic}{抑}{yi4}{7}{⼿}
  \definition{conj.}{ou; indica uma escolha, equivalente a 或是 ou 还是 | mas; ligação de orações para indicar um contraste, equivalente a 可是 | mas também; indica progressão, equivalente a 而且 | no entanto; indicando uma concessão | então; indica aceitação}
  \definition{v.}{pressionar; conter; reprimir; refrear}
  \seealsoref{而且}{er2qie3}
  \seealsoref{还是}{hai2shi5}
  \seealsoref{或是}{huo4shi4}
  \seealsoref{可是}{ke3shi4}
  \antonymref{扶}{fu2}
  \antonymref{扬}{yang2}
\end{EntryWithPhonetic}

\begin{EntryWithPhonetic}{抑扬顿挫}{yi4yang2-dun4cuo4}{7,6,10,10}{⼿,⼿,⾴,⼿}[HSK 7-9]
  \definition{expr.}{entonação; cadência; modulação de tom; falar em tons medidos; com cadência ascendente e descendente; agudos e graves (como em notas musicais); entonação; ritmo; subindo e descendo em cadência; rítmico}
\end{EntryWithPhonetic}

\begin{EntryWithPhonetic}{抑郁}{yi4yu4}{7,8}{⼿,⾢}[HSK 7-9]
  \definition{adj.}{melancólico; deprimido; abatido; um transtorno psicológico}
  \synonymref{烦闷}{fan2men4}
  \synonymref{忧郁}{you1yu4}
  \antonymref{开朗}{kai1lang3}
  \antonymref{舒畅}{shu1chang4}
\end{EntryWithPhonetic}

\begin{EntryWithPhonetic}{抑郁症}{yi4yu4zheng4}{7,8,10}{⼿,⾢,⽧}[HSK 7-9]
  \definition[种]{s.}{depressão; um transtorno mental comum que envolve períodos prolongados de humor deprimido ou perda de prazer ou interesse em atividades}
\end{EntryWithPhonetic}

\begin{EntryWithPhonetic}{抑制}{yi4zhi4}{7,8}{⼿,⼑}[HSK 7-9]
  \definition{v.}{restringir; controlar; suprimir; inibir | inibir; um dos dois processos neurais básicos do córtex cerebral; é gerado sob estimulação externa ou interna; sua função é inibir a excitação cortical e enfraquecer a função dos órgãos; o sono é um fenômeno no qual o córtex cerebral é completamente inibido}
  \synonymref{抵制}{di3zhi4}
  \synonymref{禁止}{jin4zhi3}
  \synonymref{克服}{ke4fu2}
  \synonymref{克制}{ke4zhi4}
  \synonymref{控制}{kong4zhi4}
  \synonymref{强迫}{qiang3po4}
  \synonymref{收敛}{shou1lian3}
  \synonymref{压抑}{ya1yi4}
  \synonymref{约束}{yue1shu4}
  \antonymref{不禁}{bu4jin1}
  \antonymref{冲动}{chong1dong4}
  \antonymref{促成}{cu4cheng2}
  \antonymref{发扬}{fa1yang2}
  \antonymref{放肆}{fang4si4}
  \antonymref{放纵}{fang4zong4}
  \antonymref{激活}{ji1huo2}
  \antonymref{激励}{ji1li4}
  \antonymref{兴奋}{xing1fen4}
\end{EntryWithPhonetic}

%%%%%%%%%% 译 %%%%%%%%%%
\subsection*{译}\addcontentsline{loh}{figure}{译 \dpy{yi4}}

\begin{EntryWithPhonetic}{译}{yi4}{7}{⾔}[HSK 7-9]
  \definition{v.}{traduzir; interpretar; decodificar}
  \seealsoref{翻译}{fan1yi4}
\end{EntryWithPhonetic}

%%%%%%%%%% 易 %%%%%%%%%%
\subsection*{易}\addcontentsline{loh}{figure}{易 \dpy{yi4}}

\begin{EntryWithPhonetic}{易}{yi4}{8}{⽇}
  \definition*{s.}{Sobrenome: Yi}
  \definition{adj.}{fácil | amigável; pacífico}
  \definition{v.}{modificar; transformar | trocar | subestimar; desprezar}
\end{EntryWithPhonetic}

\begin{EntryWithPhonetic}{易拉罐}{yi4la1guan4}{8,8,23}{⽇,⼿,⽸}[HSK 7-9]
  \definition{s.}{lata de alumínio; lata com tampa de pressão (ou com anel de abertura, tampa de puxar, tampa articulada) | lata de fácil abertura (com anel de puxar) | lata com tampa de abrir}
\end{EntryWithPhonetic}

%%%%%%%%%% 疫 %%%%%%%%%%
\subsection*{疫}\addcontentsline{loh}{figure}{疫 \dpy{yi4}}

\begin{EntryWithPhonetic}{疫}{yi4}{9}{⽧}
  \definition[场]{s.}{epidemia; doença epidêmica; pestilência; praga}
\end{EntryWithPhonetic}

\begin{EntryWithPhonetic}{疫苗}{yi4miao2}{9,8}{⽧,⾋}[HSK 7-9]
  \definition[种]{s.}{vacina; preparações que podem induzir imunidade no organismo}
  \antonymref{病毒}{bing4du2}
\end{EntryWithPhonetic}

%%%%%%%%%% 益 %%%%%%%%%%
\subsection*{益}\addcontentsline{loh}{figure}{益 \dpy{yi4}}

\begin{EntryWithPhonetic}{益}{yi4}{10}{⽫}
  \definition*{s.}{Sobrenome: Yi}
  \definition{adj.}{benéfico}
  \definition{adv.}{Literário: mais; cada vez mais}[空气污染问题日益严重。===O problema da poluição do ar está se tornando cada vez mais sério.]
  \definition{s.}{benefício; lucro; vantagem}
  \definition{v.}{aumentar}
\end{EntryWithPhonetic}

\begin{EntryWithPhonetic}{益虫}{yi4chong2}{10,6}{⽫,⾍}
  \definition{s.}{inseto benéfico}
  \antonymref{害虫}{hai4chong2}
\end{EntryWithPhonetic}

\begin{EntryWithPhonetic}{益处}{yi4chu4}{10,5}{⽫,⼡}[HSK 7-9]
  \definition{s.}{lucro; benefício; vantagem; fatores que são benéficos para uma pessoa ou coisa}
  \synonymref{长处}{chang2chu4}
  \synonymref{好处}{hao3chu5}
  \synonymref{利益}{li4yi4}
  \synonymref{便宜}{pian2yi5}
  \synonymref{甜头}{tian2tou5}
  \synonymref{优点}{you1dian3}
  \antonymref{弊端}{bi4duan1}
  \antonymref{坏处}{huai4chu5}
\end{EntryWithPhonetic}

%%%%%%%%%% 意 %%%%%%%%%%
\subsection*{意}\addcontentsline{loh}{figure}{意 \dpy{yi4}}

\begin{EntryWithPhonetic}{意}{yi4}{13}{⼼}
  \definition*{s.}{Itália, abreviação de 意大利}
  \definition{s.}{ideia; significado; pensamento | desejo; vontade; intenção | significância}
  \seealsoref{意大利}{yi4da4li4}
\end{EntryWithPhonetic}

\begin{EntryWithPhonetic}{意大利}{yi4da4li4}{13,3,7}{⼼,⼤,⼑}
  \definition*{s.}{Itália}
\end{EntryWithPhonetic}

\begin{EntryWithPhonetic}{意见}{yi4jian5}{13,4}{⼼,⾒}[HSK 2]
  \definition[种,点,条]{s.}{ideia; visão; opinião; sugestão; uma certa visão ou ideia sobre algo | objeção; reclamação; opinião divergente; (em relação a uma pessoa ou coisa) o sentimento de estar insatisfeito com algo porque está errado}
\end{EntryWithPhonetic}

\begin{EntryWithPhonetic}{意料}{yi4liao4}{13,10}{⼼,⽃}[HSK 7-9]
  \definition{v.}{esperar; antecipar; estimar; estimativa prévia de circunstâncias, resultados, etc.}
  \synonymref{预见}{yu4jian4}
  \synonymref{预料}{yu4liao4}
  \antonymref{不料}{bu2liao4}
  \antonymref{意外}{yi4wai4}
\end{EntryWithPhonetic}

\begin{EntryWithPhonetic}{意料之外}{yi4liao4 zhi1 wai4}{13,10,3,5}{⼼,⽃,⼂,⼣}[HSK 7-9]
  \definition{expr.}{contrário ao esperado; inesperado}
  \synonymref{出人意料}{chu1ren2yi4liao4}
  \synonymref{意想不到}{yi4xiang3bu2dao4}
\end{EntryWithPhonetic}

\begin{EntryWithPhonetic}{意识}{yi4shi5}{13,7}{⼼,⾔}[HSK 5]
  \definition{s.}{consciência; percepção; grau de reconhecimento e importância atribuído a uma determinada questão}
  \definition{s.}{consciência; percepção; o reflexo da mente humana no mundo material objetivo é a soma de vários processos psicológicos, como sensação e pensamento | consciência; conscientização; o grau de conscientização e atenção dada a um problema}
  \definition{v.}{perceber; despertar para; estar ciente de; sentir, descobrir o que antes não se sentia ou não se descobria; geralmente é usado junto com 到}
  \seealsoref{到}{dao4}
\end{EntryWithPhonetic}

\begin{EntryWithPhonetic}{意思}{yi4si5}{13,9}{⼼,⼼}[HSK 2]
  \definition[个]{s.}{ideia; significado; o significado da linguagem e das palavras; conteúdo ideológico | desejo; vontade; opiniões | um símbolo de afeto, apreciação, gratidão, etc. | dica; traço; sugestão; refere"-se principalmente ao afeto entre homens e mulheres | diversão; interesse}
  \definition{v.}{dar uma dica; demonstrar sua gratidão com presentes ou outros meios}
\end{EntryWithPhonetic}

\begin{EntryWithPhonetic}{意图}{yi4tu2}{13,8}{⼼,⼞}[HSK 7-9]
  \definition[个,种]{s.}{intenção; propósito; objetivo; a intenção de alcançar um determinado objetivo}
  \definition{v.}{tentar; pretender}
  \seealsoref{意向}{yi4xiang4}
  \synonymref{企图}{qi3tu2}
  \synonymref{妄想}{wang4xiang3}
  \synonymref{意愿}{yi4yuan4}
  \antonymref{无意}{wu2yi4}
\end{EntryWithPhonetic}

\begin{EntryWithPhonetic}{意外}{yi4wai4}{13,5}{⼼,⼣}[HSK 3]
  \definition{adj.}{inesperado; imprevisto}
  \definition{adv.}{acidentalmente}
  \definition[个,种]{s.}{acidente; infortúnio; um infortúnio inesperado}
\end{EntryWithPhonetic}

\begin{EntryWithPhonetic}{意味着}{yi4wei4zhe5}{13,8,11}{⼼,⼝,⽬}[HSK 5]
  \definition{v.}{significar; subentender; implicar; entender que tem vários significados}
\end{EntryWithPhonetic}

\begin{EntryWithPhonetic}{意想不到}{yi4xiang3bu2dao4}{13,13,4,8}{⼼,⼼,⼀,⼑}[HSK 6]
  \definition{expr.}{anteriormente inimaginável | inesperado}
\end{EntryWithPhonetic}

\begin{EntryWithPhonetic}{意向}{yi4xiang4}{13,6}{⼼,⼝}[HSK 7-9]
  \definition{s.}{intenção; propósito; intenções e planos}
  \seealsoref{意图}{yi4tu2}
  \synonymref{抱负}{bao4fu4}
  \synonymref{动向}{dong4xiang4}
  \synonymref{理想}{li3xiang3}
  \synonymref{希望}{xi1wang4}
  \synonymref{愿望}{yuan4wang4}
  \synonymref{志愿}{zhi4yuan4}
\end{EntryWithPhonetic}

\begin{EntryWithPhonetic}{意义}{yi4yi4}{13,3}{⼼,⼂}[HSK 3]
  \definition[个,种,层,重,点]{s.}{sentido; significado; o significado expresso por meio de linguagem escrita ou outros sinais; o significado identificado por meio de ações ou obtenção | valor; efeito; significado; influência; impacto}
\end{EntryWithPhonetic}

\begin{EntryWithPhonetic}{意译}{yi4yi4}{13,7}{⼼,⾔}
  \definition{s.}{tradução livre | significado (de expressão estrangeira) | paráfrase | tradução do significado}
  \antonymref{直译}{zhi2yi4}
\end{EntryWithPhonetic}

\begin{EntryWithPhonetic}{意愿}{yi4yuan4}{13,14}{⼼,⽕}[HSK 6]
  \definition{s.}{desejo; aspiração; vontade}
\end{EntryWithPhonetic}

\begin{EntryWithPhonetic}{意指}{yi4zhi3}{13,9}{⼼,⼿}
  \definition{v.}{implicar | significar}
\end{EntryWithPhonetic}

\begin{EntryWithPhonetic}{意志}{yi4zhi4}{13,7}{⼼,⼼}[HSK 5]
  \definition[个,股]{s.}{vontade; determinação; desejo; força de vontade; o estado psicológico produzido pela determinação de atingir um determinado objetivo, muitas vezes expresso por meio de linguagem e ações}
\end{EntryWithPhonetic}

%%%%%%%%%% 溢 %%%%%%%%%%
\subsection*{溢}\addcontentsline{loh}{figure}{溢 \dpy{yi4}}

\begin{EntryWithPhonetic}{溢}{yi4}{13}{⽔}[HSK 7-9]
  \definition{adj.}{excessivo}
  \definition{v.}{transbordar; derramar}
\end{EntryWithPhonetic}

%%%%%%%%%% 毅 %%%%%%%%%%
\subsection*{毅}\addcontentsline{loh}{figure}{毅 \dpy{yi4}}

\begin{EntryWithPhonetic}{毅}{yi4}{15}{⽎}
  \definition{adj.}{firme; resoluto; inabalável}
\end{EntryWithPhonetic}

\begin{EntryWithPhonetic}{毅力}{yi4li4}{15,2}{⽎,⼒}[HSK 7-9]
  \definition{s.}{coragem; vontade; resistência; força de vontade; perseverança; o espírito de nunca desistir e perseverar diante das dificuldades}
  \synonymref{坚韧}{jian1ren4}
  \synonymref{顽强}{wan2qiang2}
  \synonymref{意志}{yi4zhi4}
  \antonymref{动摇}{dong4yao2}
  \antonymref{软弱}{ruan3ruo4}
  \antonymref{颓废}{tui2fei4}
\end{EntryWithPhonetic}

\begin{EntryWithPhonetic}{毅然}{yi4ran2}{15,12}{⽎,⽕}[HSK 7-9]
  \definition{adv.}{firmemente; resolutamente; com determinação}
  \seealsoref{坚决}{jian1jue2}
  \synonymref{果断}{guo3duan4}
  \synonymref{坚决}{jian1jue2}
  \antonymref{犹豫}{you2yu4}
\end{EntryWithPhonetic}

%%%%%%%%%% 因 %%%%%%%%%%
\subsection*{因}\addcontentsline{loh}{figure}{因 \dpy{yin1}}

\begin{EntryWithPhonetic}{因}{yin1}{6}{⼞}[HSK 6]
  \definition*{s.}{Sobrenome: Yin}
  \definition{conj.}{porque; orações de conexão, indicando relações de causa e efeito}
  \definition{prep.}{com base em; à luz de; de acordo com; a introdução da ação comportamental equivale a 按照 ou 根据}
  \definition{s.}{causa; motivo; condições em que algo ocorre ou causa um determinado resultado}
  \definition{v.}{seguir; continuar; fazer como sempre fez | estar em conformidade com; estar de acordo com; depender; contar com}
  \seealsoref{按照}{an4zhao4}
  \seealsoref{根据}{gen1ju4}
  \antonymref{果}{guo3}
\end{EntryWithPhonetic}

\begin{EntryWithPhonetic}{因此}{yin1ci3}{6,6}{⼞,⽌}[HSK 3]
  \definition{conj.}{assim; portanto; consequentemente}
\end{EntryWithPhonetic}

\begin{EntryWithPhonetic}{因此就}{yin1ci3 jiu4}{6,6,12}{⼞,⽌,⼪}
  \definition{conj.}{portanto}
\end{EntryWithPhonetic}

\begin{EntryWithPhonetic}{因而}{yin1'er2}{6,6}{⼞,⽽}[HSK 5]
  \definition{conj.}{assim; como resultado; com o resultado que; conecta frases, indicando relação de causa e efeito}
\end{EntryWithPhonetic}

\begin{EntryWithPhonetic}{因人而异}{yin1ren2'er2yi4}{6,2,6,6}{⼞,⼈,⽽,⼶}[HSK 7-9]
  \definition{expr.}{``Isso varia de pessoa para pessoa.''; varia de pessoa para pessoa; pessoas diferentes se comportam de maneira diferente; métodos diferentes devem ser usados dependendo do objeto}
\end{EntryWithPhonetic}

\begin{EntryWithPhonetic}{因素}{yin1su4}{6,10}{⼞,⽷}[HSK 6]
  \definition[个,种]{s.}{fator; elemento; os componentes que constituem a essência das coisas | fator; as razões ou condições que determinam o sucesso ou o fracasso de algo}
\end{EntryWithPhonetic}

\begin{EntryWithPhonetic}{因为}{yin1wei5}{6,4}{⼞,⼂}[HSK 2]
  \definition{conj.}{porque; indica o motivo e a frase seguinte indica o resultado}
  \definition{prep.}{por causa de; por conta de; indica razão ou justificativa}
\end{EntryWithPhonetic}

\begin{EntryWithPhonetic}{因为……所以……}{yin1wei5 suo3yi3}{6,4,8,4}{⼞,⼂,⼾,⼈}
  \definition{conj.}{porque\dots portanto\dots}
\end{EntryWithPhonetic}

%%%%%%%%%% 阴 %%%%%%%%%%
\subsection*{阴}\addcontentsline{loh}{figure}{阴 \dpy{yin1}}

\begin{EntryWithPhonetic}{阴}{yin1}{6}{⾩}[HSK 2]
  \definition*{s.}{Yin, o princípio negativo de Yin e Yang | A Lua; refere"-se a Taiyin | Sobrenome: Yin}
  \definition{adj.}{nublado; opaco; sombrio | escondido; secreto; não exposto | sinistro | do mundo inferior; dos fantasmas | Física: negativo; cátodo | nublado; mais de 80\% do céu estão cobertos por nuvens | em talhe"-doce; rebaixado | (matéria) carregada negativamente}
  \definition[片]{s.}{sombra; lugar sombrio | partes íntimas (especialmente da mulher) | ao norte de uma colina ou ao sul de um rio | verso | entalhe}
  \seealsoref{阳}{yang2}
  \seealsoref{阴阳}{yin1yang2}
\end{EntryWithPhonetic}

\begin{EntryWithPhonetic}{阴暗}{yin1'an4}{6,13}{⾩,⽇}[HSK 7-9]
  \definition{adj.}{escuro; sombrio; metaforicamente, não confiável ou honroso | escuro; sombrio; deprimido; essa metáfora descreve o mau humor de uma pessoa; uma mentalidade negativa; pensamentos negativos sobre pessoas e coisas; ou um ambiente social que carece de justiça ou equidade}
  \synonymref{黑暗}{hei1'an4}
  \antonymref{明朗}{ming2lang3}
  \antonymref{明亮}{ming2liang4}
  \antonymref{晴朗}{qing2lang3}
\end{EntryWithPhonetic}

\begin{EntryWithPhonetic}{阴谋}{yin1mou2}{6,11}{⾩,⾔}[HSK 6]
  \definition[个,场,起]{s.}{trama; conspiração; um esquema para fazer o mal em segredo}
  \definition{v.}{tramar; conspirar secretamente (fazer algo ruim)}
\end{EntryWithPhonetic}

\begin{EntryWithPhonetic}{阴天}{yin1tian1}{6,4}{⾩,⼤}[HSK 2]
  \definition[个]{s.}{nublado; céu nublado; dia nublado; uma condição climática em que 80\% do céu está coberto por nuvens e apenas um pouco de sol pode ser visto}
\end{EntryWithPhonetic}

\begin{EntryWithPhonetic}{阴性}{yin1xing4}{6,8}{⾩,⼼}[HSK 7-9]
  \definition{s.}{Medicina: negativo | Linguística: gênero feminino | feminino}[链球菌本质是阴性的。===A cultura para estreptococos deu negativa. | ``门''在德语里是阴性名词。===A palavra ``porta'' é um substantivo feminino em alemão. | 她的气质明显偏阴性。===Seu temperamento é visivelmente mais feminino.]
  \antonymref{阳性}{yang2xing4}
\end{EntryWithPhonetic}

\begin{EntryWithPhonetic}{阴阳}{yin1yang2}{6,6}{⾩,⾩}
  \definition*{s.}{Yin e Yang}
  \seealsoref{阳}{yang2}
  \seealsoref{阴}{yin1}
\end{EntryWithPhonetic}

\begin{EntryWithPhonetic}{阴影}{yin1ying3}{6,15}{⾩,⼺}[HSK 6]
  \definition{s.}{sombra; sombra escura | uma analogia de elementos negativos em negócios, relacionamentos, estado mental, etc.}
\end{EntryWithPhonetic}

%%%%%%%%%% 音 %%%%%%%%%%
\subsection*{音}\addcontentsline{loh}{figure}{音 \dpy{yin1}}

\begin{EntryWithPhonetic}{音}{yin1}{9}{⾳}[Kangxi 180]
  \definition[个,种]{s.}{som; som musical | notícias; novidades; informação | tom; refere"-se especificamente a uma sílaba ou fonética | sílaba; refere"-se a sílabas (um caractere chinês é uma sílaba)}
  \definition{v.}{vocalizar}
\end{EntryWithPhonetic}

\begin{EntryWithPhonetic}{音节}{yin1jie2}{9,5}{⾳,⾋}[HSK 2]
  \definition{s.}{sílaba}
\end{EntryWithPhonetic}

\begin{EntryWithPhonetic}{音量}{yin1liang4}{9,12}{⾳,⾥}[HSK 6]
  \definition[把]{s.}{volume; volume do som; a força de um som}
\end{EntryWithPhonetic}

\begin{EntryWithPhonetic}{音响}{yin1xiang3}{9,9}{⾳,⼝}[HSK 7-9]
  \definition[套,款]{s.}{áudio; som (referindo"-se principalmente aos efeitos produzidos pelo som) | alto"-falante; caixa de som; sistema de som}[音响设备有点过时了。===O equipamento de áudio está um pouco desatualizado.]
  \synonymref{喇叭}{la3ba1}
\end{EntryWithPhonetic}

\begin{EntryWithPhonetic}{音像}{yin1xiang4}{9,13}{⾳,⼈}[HSK 6]
  \definition{s.}{audiovisual; produtos audiovisuais; o nome coletivo para gravações de áudio e vídeo}
\end{EntryWithPhonetic}

\begin{EntryWithPhonetic}{音乐}{yin1yue4}{9,5}{⾳,⼃}[HSK 2]
  \definition[种,段,张,曲]{s.}{música; ramo da arte que cria imagens artísticas, expressa pensamentos e sentimentos e reflete a vida real por meio da melodia e do ritmo da música; geralmente é dividido em duas categorias: música vocal e música instrumental}
\end{EntryWithPhonetic}

\begin{EntryWithPhonetic}{音乐光碟}{yin1yue4 guang1die2}{9,5,6,14}{⾳,⼃,⼉,⽯}
  \definition{s.}{CD de música}
\end{EntryWithPhonetic}

\begin{EntryWithPhonetic}{音乐会}{yin1yue4hui4}{9,5,6}{⾳,⼃,⼈}[HSK 2]
  \definition[场]{s.}{concerto; atividades de execução de obras musicais}
\end{EntryWithPhonetic}

\begin{EntryWithPhonetic}{音乐家}{yin1yue4jia1}{9,5,10}{⾳,⼃,⼧}
  \definition{s.}{músico}
\end{EntryWithPhonetic}

\begin{EntryWithPhonetic}{音乐节}{yin1yue4jie2}{9,5,5}{⾳,⼃,⾋}
  \definition{s.}{festival de música}
\end{EntryWithPhonetic}

\begin{EntryWithPhonetic}{音乐厅}{yin1yue4ting1}{9,5,4}{⾳,⼃,⼚}
  \definition{s.}{auditório | teatro | \emph{concert hall}}
\end{EntryWithPhonetic}

\begin{EntryWithPhonetic}{音乐学}{yin1yue4xue2}{9,5,8}{⾳,⼃,⼦}
  \definition{s.}{musicologia}
\end{EntryWithPhonetic}

\begin{EntryWithPhonetic}{音乐学院}{yin1yue4xue2yuan4}{9,5,8,9}{⾳,⼃,⼦,⾩}
  \definition{s.}{conservatório | academia de música}
\end{EntryWithPhonetic}

\begin{EntryWithPhonetic}{音乐院}{yin1yue4yuan4}{9,5,9}{⾳,⼃,⾩}
  \definition{s.}{conservatório | instituto de música}
\end{EntryWithPhonetic}

%%%%%%%%%% 殷 %%%%%%%%%%
\subsection*{殷}\addcontentsline{loh}{figure}{殷 \dpy{yin1}}

\begin{EntryWithPhonetic}{殷}{yin1}{10}{⽎}
  \definition*{s.}{Dinastia Yin (1300--1046 a.C.; nome dado à última parte da Dinastia Shang) | Sobrenome: Yin}
  \definition{adj.}{Literário: abundante; rico; farto | Literário: ansioso; ardente; profundo | Literário: hospitaleiro; atento | grandioso; magnífico; grande | numeroso}
  \seeref{yan1}
  \seeref{yin3}
\end{EntryWithPhonetic}

\begin{EntryWithPhonetic}{殷勤}{yin1qin2}{10,13}{⽎,⼒}[HSK 7-9]
  \definition{adj.}{solícito; ansiosamente atencioso; desejoso de agradar; caloroso e atencioso}
  \synonymref{热情}{re4qing2}
  \antonymref{怠慢}{dai4man4}
  \antonymref{冷淡}{leng3dan4}
\end{EntryWithPhonetic}

%%%%%%%%%% 吟 %%%%%%%%%%
\subsection*{吟}\addcontentsline{loh}{figure}{吟 \dpy{yin2}}

\begin{EntryWithPhonetic}{吟}{yin2}{7}{⼝}
  \definition*{s.}{Sobrenome: Yin}
  \definition{s.}{canção (como um tipo de poesia clássica) | grito de certos animais ou insetos}
  \definition{v.}{cantar; recitar | gemer; lamentar}
\end{EntryWithPhonetic}

\begin{EntryWithPhonetic}{吟诗}{yin2shi1}{7,8}{⼝,⾔}
  \definition{v.}{recitar poesia}
\end{EntryWithPhonetic}

%%%%%%%%%% 银 %%%%%%%%%%
\subsection*{银}\addcontentsline{loh}{figure}{银 \dpy{yin2}}

\begin{EntryWithPhonetic}{银}{yin2}{11}{⾦}[HSK 3]
  \definition*{s.}{Sobrenome: Yin}
  \definition{adj.}{prateado; como a cor da prata}
  \definition[锭]{s.}{Ag, prata | refere"-se a moeda ou a coisas relacionadas com moeda}
\end{EntryWithPhonetic}

\begin{EntryWithPhonetic}{银行}{yin2hang2}{11,6}{⾦,⾏}[HSK 2]
  \definition[个,家,所]{s.}{banco; instituições financeiras que operam depósitos, empréstimos, câmbio, poupança e outros negócios}
\end{EntryWithPhonetic}

\begin{EntryWithPhonetic}{银行卡}{yin2hang2ka3}{11,6,5}{⾦,⾏,⼘}[HSK 2]
  \definition{s.}{cartão bancário; cartão ATM}
\end{EntryWithPhonetic}

\begin{EntryWithPhonetic}{银河}{yin2he2}{11,8}{⾦,⽔}
  \definition*{s.}{Via Láctea}
  \seealsoref{银河系}{yin2he2xi4}
\end{EntryWithPhonetic}

\begin{EntryWithPhonetic}{银河系}{yin2he2xi4}{11,8,7}{⾦,⽔,⽷}
  \definition*{s.}{Galáxia Via Láctea}
  \seealsoref{银河}{yin2he2}
\end{EntryWithPhonetic}

\begin{EntryWithPhonetic}{银幕}{yin2mu4}{11,13}{⾦,⼱}[HSK 7-9]
  \definition[块]{s.}{tela; tela de projeção; uma tela branca usada para exibir vídeos, slides ou imagens projetadas}
  \synonymref{屏幕}{ping2mu4}
\end{EntryWithPhonetic}

\begin{EntryWithPhonetic}{银牌}{yin2pai2}{11,12}{⾦,⽚}[HSK 3]
  \definition[枚]{s.}{medalha de prata; um tipo de medalha, concedida ao segundo colocado}
\end{EntryWithPhonetic}

\begin{EntryWithPhonetic}{银色}{yin2 se4}{11,6}{⾦,⾊}
  \definition{s.}{cor prata; prateado}
\end{EntryWithPhonetic}

%%%%%%%%%% 引 %%%%%%%%%%
\subsection*{引}\addcontentsline{loh}{figure}{引 \dpy{yin3}}

\begin{EntryWithPhonetic}{引}{yin3}{4}{⼸}[HSK 4]
  \definition*{s.}{Sobrenome: Yin}
  \definition{clas.}{uma unidade de comprimento equivalente a =33,333\dots metros}
  \definition{v.}{puxar; esticar | liderar; conduzir; guiar | sair; deixar | sobressair | atrair; provocar; trazer à existência | causar; provocar | citar; ser usado como evidência ou justificativa}
\end{EntryWithPhonetic}

\begin{EntryWithPhonetic}{引导}{yin3dao3}{4,6}{⼸,⼨}[HSK 4]
  \definition{v.}{conduzir; guiar; liderar; andar na frente e deixar que os outros sigam atrás para ver ou andar; usar imagens ou sinais para mostrar às pessoas para onde ir | esclarecer; fornecer orientação em termos de ideias, métodos, conceitos, etc.}
\end{EntryWithPhonetic}

\begin{EntryWithPhonetic}{引发}{yin3fa1}{4,5}{⼸,⼜}[HSK 7-9]
  \definition{v.}{despertar; desencadear; acionar; iniciar (ou acender, provocar)}
  \synonymref{激发}{ji1fa1}
  \synonymref{激励}{ji1li4}
\end{EntryWithPhonetic}

\begin{EntryWithPhonetic}{引进}{yin3jin4}{4,7}{⼸,⾡}[HSK 4]
  \definition{v.}{importar; trazer de fora | recomendar; dar uma indicação}
\end{EntryWithPhonetic}

\begin{EntryWithPhonetic}{引经据典}{yin3jing1-ju4dian3}{4,8,11,8}{⼸,⽷,⼿,⼋}[HSK 7-9]
  \definition{expr.}{``Citar capítulo e versículo.''; citar os clássicos; ciare obras de referência}
\end{EntryWithPhonetic}

\begin{EntryWithPhonetic}{引领}{yin3ling3}{4,11}{⼸,⾴}[HSK 7-9]
  \definition{v.}{liderar; guiar | aguardar ansiosamente por algo; esperar ansiosamente}
\end{EntryWithPhonetic}

\begin{EntryWithPhonetic}{引起}{yin3qi3}{4,10}{⼸,⾛}[HSK 4]
  \definition{v.}{causar; despertar; levar a; desencadear; dar origem a}
\end{EntryWithPhonetic}

\begin{EntryWithPhonetic}{引擎}{yin3qing2}{4,16}{⼸,⼿}[HSK 7-9]
  \definition[个,台]{s.}{motor, referindo"-se especificamente a motores térmicos como motores a vapor e motores de combustão interna | Empréstimo Linguístico: \emph{engine}}
\end{EntryWithPhonetic}

\begin{EntryWithPhonetic}{引人入胜}{yin3ren2-ru4sheng4}{4,2,2,9}{⼸,⼈,⼊,⾁}[HSK 7-9]
  \definition{expr.}{conduzir alguém à parte interessante de algo; sedutor; atraente; cativante; emocionante; encantador; fascinante; que atraia muito alguém; interessante e absorvente (intrigante); levar alguém a uma nova perspectiva; abrir uma nova visão; deslumbrante; envolver o leitor em uma cena bela (referindo"-se à paisagem, à escrita, etc.)}
  \synonymref{扣人心弦}{kou4ren2xin1xian2}
\end{EntryWithPhonetic}

\begin{EntryWithPhonetic}{引人注目}{yin3ren2-zhu4mu4}{4,2,8,5}{⼸,⼈,⽔,⽬}[HSK 7-9]
  \definition{expr.}{impressionante; chamativo; que prende a atenção; significa descrever uma pessoa ou coisa que é muito peculiar e atrai a atenção das pessoas}
  \synonymref{举世瞩目}{ju3shi4-zhu3mu4}
  \synonymref{怡然自得}{yi2ran2-zi4de2}
  \antonymref{默默无闻}{mo4mo4-wu2wen2}
\end{EntryWithPhonetic}

\begin{EntryWithPhonetic}{引入}{yin3ru4}{4,2}{⼸,⼊}[HSK 7-9]
  \definition{v.}{introduzir; atrair; trazer de outro lugar | introduzir (pessoal, capital, tecnologia, equipamentos, etc.) de outras regiões ou países estrangeiros}
  \synonymref{引进}{yin3jin4}
  \synonymref{引用}{yin3yong4}
\end{EntryWithPhonetic}

\begin{EntryWithPhonetic}{引用}{yin3yong4}{4,5}{⼸,⽤}[HSK 7-9]
  \definition{v.}{citar; transcrever; utilizar as palavras de outras pessoas ou eventos passados como evidência para apoiar e comprovar o próprio ponto de vista | nomear; recomendar; nomear alguém para um cargo ou para realizar um trabalho}
  \synonymref{引入}{yin3ru4}
\end{EntryWithPhonetic}

\begin{EntryWithPhonetic}{引诱}{yin3you4}{4,9}{⼸,⾔}[HSK 7-9]
  \definition{v.}{atrair; seduzir; tentar; induzir e orientar pessoas a fazer coisas ruins}
  \synonymref{勾结}{gou1jie2}
  \synonymref{迷惑}{mi2huo4}
  \synonymref{引导}{yin3dao3}
  \synonymref{诱惑}{you4huo4}
  \antonymref{导致}{dao3zhi4}
\end{EntryWithPhonetic}

%%%%%%%%%% 听 %%%%%%%%%%
\subsection*{听}\addcontentsline{loh}{figure}{听 \dpy{yin3}}

\begin{EntryWithPhonetic}{听}{yin3}{7}{⼝}
  \definition[个]{s.}{lata; embalagem metálica}
  \seeref{ting1}
\end{EntryWithPhonetic}

%%%%%%%%%% 饮 %%%%%%%%%%
\subsection*{饮}\addcontentsline{loh}{figure}{饮 \dpy{yin3}}

\begin{EntryWithPhonetic}{饮}{yin3}{7}{⾷}
  \definition{s.}{bebidas; \emph{drinks}; algo para beber | uma decocção da medicina chinesa para ser tomada fria | fluido retido}
  \definition{v.}{beber | cuidar; engolir a pílula amarga}
  \seeref{yin4}
\end{EntryWithPhonetic}

\begin{EntryWithPhonetic}{饮料}{yin3liao4}{7,10}{⾷,⽃}[HSK 5]
  \definition[杯,瓶,种]{s.}{bebida; drinque; líquidos processados e fabricados para consumo, como vinho, chá, refrigerantes, suco de laranja, etc.}
\end{EntryWithPhonetic}

\begin{EntryWithPhonetic}{饮食}{yin3shi2}{7,9}{⾷,⾷}[HSK 5]
  \definition{s.}{dieta; alimentos e bebidas}
  \definition{v.}{comer; beber}
\end{EntryWithPhonetic}

\begin{EntryWithPhonetic}{饮水}{yin3 shui3}{7,4}{⾷,⽔}[HSK 7-9]
  \definition{v.}{beber água}[多饮水对身体有利。===Beber bastante água faz bem para a saúde.]
\end{EntryWithPhonetic}

\begin{EntryWithPhonetic}{饮用水}{yin3yong4shui3}{7,5,4}{⾷,⽤,⽔}[HSK 7-9]
  \definition{s.}{água potável}
\end{EntryWithPhonetic}

%%%%%%%%%% 殷 %%%%%%%%%%
\subsection*{殷}\addcontentsline{loh}{figure}{殷 \dpy{yin3}}

\begin{EntryWithPhonetic}{殷}{yin3}{10}{⽎}
  \definition{s.}{Onomatopeia: estrondo do trovão}
\end{EntryWithPhonetic}

%%%%%%%%%% 隐 %%%%%%%%%%
\subsection*{隐}\addcontentsline{loh}{figure}{隐 \dpy{yin3}}

\begin{EntryWithPhonetic}{隐}{yin3}{11}{⾩}
  \definition*{s.}{Sobrenome: Yin}
  \definition{adj.}{escondido; escondido profundamente | latente; adormecido; à espreita}
  \definition{pref.}{cripto-}
  \definition{s.}{segredo; assuntos ocultos}
  \definition{v.}{esconder; esconder da vista; ocultar}
\end{EntryWithPhonetic}

\begin{EntryWithPhonetic}{隐蔽}{yin3bi4}{11,14}{⾩,⾋}[HSK 7-9]
  \definition{adj.}{escondido; oculto; oculto e difícil de detectar}
  \definition{v.}{ocultar; proteger"-se; usar outra coisa para cobrir e assim não ser descoberto}
  \seealsoref{隐藏}{yin3cang2}
  \synonymref{藏匿}{cang2ni4}
  \synonymref{隐藏}{yin3cang2}
  \synonymref{隐患}{yin3huan4}
  \synonymref{隐瞒}{yin3man2}
  \antonymref{暴露}{bao4lu4}
  \antonymref{查明}{cha2ming2}
  \antonymref{公开}{gong1kai1}
  \antonymref{明显}{ming2xian3}
  \antonymref{显出}{xian3 chu1}
  \antonymref{显现}{xian3xian4}
  \antonymref{显眼}{xian3yan3}
\end{EntryWithPhonetic}

\begin{EntryWithPhonetic}{隐藏}{yin3cang2}{11,17}{⾩,⾋}[HSK 6]
  \definition{v.}{esconder; ocultar}
\end{EntryWithPhonetic}

\begin{EntryWithPhonetic}{隐患}{yin3huan4}{11,11}{⾩,⼼}[HSK 7-9]
  \definition{s.}{perigo oculto; dano oculto; perigo escondido em algo; infortúnio não visível na superfície}
  \synonymref{危险}{wei1xian3}
  \synonymref{问题}{wen4ti2}
  \synonymref{隐蔽}{yin3bi4}
  \antonymref{安全}{an1quan2}
\end{EntryWithPhonetic}

\begin{EntryWithPhonetic}{隐瞒}{yin3man2}{11,15}{⾩,⽬}[HSK 7-9]
  \definition{v.}{esconder; ocultar; reter; encobrir a verdade; impedir que outros saibam}
  \synonymref{保密}{bao3/mi4}
  \synonymref{秘密}{mi4mi4}
  \synonymref{掩盖}{yan3gai4}
  \synonymref{掩饰}{yan3shi4}
  \synonymref{隐蔽}{yin3bi4}
  \synonymref{遮盖}{zhe1gai4}
  \antonymref{报告}{bao4gao4}
  \antonymref{暴露}{bao4lu4}
  \antonymref{承认}{cheng2ren4}
  \antonymref{打听}{da3ting5}
  \antonymref{告诉}{gao4su4}
  \antonymref{坦白}{tan3bai2}
  \antonymref{坦率}{tan3shuai4}
  \antonymref{提醒}{ti2/xing3}
  \antonymref{提示}{ti2shi4}
  \antonymref{通知}{tong1zhi1}
\end{EntryWithPhonetic}

\begin{EntryWithPhonetic}{隐情}{yin3qing2}{11,11}{⾩,⼼}[HSK 7-9]
  \definition{s.}{fatos ou circunstâncias que se deseja ocultar; segredos | um assunto que é melhor evitar | algo que se deseja manter em segredo | razões posteriores}
\end{EntryWithPhonetic}

\begin{EntryWithPhonetic}{隐身}{yin3shen1}{11,7}{⾩,⾝}[HSK 7-9]
  \definition{v.}{ocultar"-se; esconder"-se; tornar"-se invisível}
\end{EntryWithPhonetic}

\begin{EntryWithPhonetic}{隐私}{yin3si1}{11,7}{⾩,⽲}[HSK 6]
  \definition[点,些]{s.}{privacidade; segredos de alguém; assuntos pessoais que você não quer contar ou tornar públicos}
\end{EntryWithPhonetic}

\begin{EntryWithPhonetic}{隐形}{yin3xing2}{11,7}{⾩,⼺}[HSK 7-9]
  \definition{adj.}{invisível; indetectável}
  \synonymref{躲避}{duo3bi4}
\end{EntryWithPhonetic}

\begin{EntryWithPhonetic}{隐性}{yin3xing4}{11,8}{⾩,⼼}[HSK 7-9]
  \definition{adj.}{Biologia: recessivo | oculto; invisível}
\end{EntryWithPhonetic}

\begin{EntryWithPhonetic}{隐约}{yin3yue1}{11,6}{⾩,⽷}[HSK 7-9]
  \definition{adj.}{fraco; vago; indistinto; descreve algo como parecendo ou soando confuso; não claramente compreendido}
  \synonymref{朦胧}{meng2long2}
  \synonymref{模糊}{mo2hu5}
  \antonymref{分明}{fen1ming2}
  \antonymref{明显}{ming2xian3}
  \antonymref{清晰}{qing1xi1}
  \antonymref{清楚}{qing1chu5}
\end{EntryWithPhonetic}

%%%%%%%%%% 瘾 %%%%%%%%%%
\subsection*{瘾}\addcontentsline{loh}{figure}{瘾 \dpy{yin3}}

\begin{EntryWithPhonetic}{瘾}{yin3}{16}{⽧}[HSK 7-9]
  \definition{s.}{vício; desejo habitual | forte interesse (em um esporte ou passatempo)}
\end{EntryWithPhonetic}

%%%%%%%%%% 印 %%%%%%%%%%
\subsection*{印}\addcontentsline{loh}{figure}{印 \dpy{yin4}}

\begin{EntryWithPhonetic}{印}{yin4}{5}{⼙}[HSK 6]
  \definition*{s.}{Sobrenome: Yin}
  \definition[个,枚,道,条]{s.}{selo; lacre; carimbo; estampilha | marca; estampa; impressão}
  \definition{v.}{imprimir; gravar | corresponder; conformar; estar em conformidade com}
\end{EntryWithPhonetic}

\begin{EntryWithPhonetic}{印刷}{yin4shua1}{5,8}{⼙,⼑}[HSK 5]
  \definition{v.}{imprimir; imprimir textos, imagens, etc. em papel}
\end{EntryWithPhonetic}

\begin{EntryWithPhonetic}{印刷术}{yin4shua1shu4}{5,8,5}{⼙,⼑,⽊}[HSK 7-9]
  \definition{s.}{arte da impressão; impressão}[印刷术改变了世界。===A impressão mudou o mundo.]
\end{EntryWithPhonetic}

\begin{EntryWithPhonetic}{印象}{yin4xiang4}{5,11}{⼙,⾗}[HSK 3]
  \definition[种]{s.}{impressão; marca; ideia; os vestígios deixados por coisas objetivas na mente das pessoas}
\end{EntryWithPhonetic}

\begin{EntryWithPhonetic}{印章}{yin4zhang1}{5,11}{⼙,⾳}[HSK 7-9]
  \definition[个]{s.}{selo; carimbo; sinete}
\end{EntryWithPhonetic}

\begin{EntryWithPhonetic}{印证}{yin4zheng4}{5,7}{⼙,⾔}[HSK 7-9]
  \definition{s.}{prova; evidência; coisas usadas para verificar}
  \definition{v.}{verificar; confirmar; corroborar; comprovar; a prova é consistente com os fatos}
  \synonymref{表明}{biao3ming2}
  \synonymref{见证}{jian4zheng4}
  \synonymref{证明}{zheng4ming2}
  \antonymref{编造}{bian1zao4}
\end{EntryWithPhonetic}

%%%%%%%%%% 饮 %%%%%%%%%%
\subsection*{饮}\addcontentsline{loh}{figure}{饮 \dpy{yin4}}

\begin{EntryWithPhonetic}{饮}{yin4}{7}{⾷}
  \definition{v.}{dar (aos animais) água para beber}
  \seeref{yin3}
\end{EntryWithPhonetic}

%%%%%%%%%% 应 %%%%%%%%%%
\subsection*{应}\addcontentsline{loh}{figure}{应 \dpy{ying1}}

\begin{EntryWithPhonetic}{应}{ying1}{7}{⼴}[HSK 4,5]
  \definition{v.}{ecoar; responder; responder a; responder às chamadas, saudações, perguntas, etc. de outras pessoas | conceder; cumprir | adequar; adaptar; responder a | lidar com; enfrentar; abordar | tornar"-se realidade; ser cumprido}
\end{EntryWithPhonetic}

\begin{EntryWithPhonetic}{应当}{ying1dang1}{7,6}{⼴,⼹}[HSK 3]
  \definition{v.}{dever}[学生们应当努力学习。===Os alunos devem se esforçar nos estudos.]
\end{EntryWithPhonetic}

\begin{EntryWithPhonetic}{应该}{ying1gai1}{7,8}{⼴,⾔}[HSK 2]
  \definition{v.}{deveria; deve ser assim | deveria; acho que deve ser esse o caso}
\end{EntryWithPhonetic}

\begin{EntryWithPhonetic}{应有尽有}{ying1you3-jin4you3}{7,6,6,6}{⼴,⽉,⼫,⽉}[HSK 7-9]
  \definition{expr.}{``Tudo o que você poderia desejar.''; ter tudo o que se espera encontrar; ter tudo o que é necessário; ter todas as coisas que se deseja; tudo o que deveria estar lá, está lá, ou seja, tudo está pronto}
  \synonymref{层出不穷}{ceng2chu1-bu4qiong2}
  \synonymref{面面俱到}{mian4mian4-ju4dao4}
  \synonymref{形形色色}{xing2xing2se4se4}
  \synonymref{一应俱全}{yi4ying1-ju4quan2}
  \antonymref{一无所有}{yi1wu2suo3you3}
\end{EntryWithPhonetic}

%%%%%%%%%% 英 %%%%%%%%%%
\subsection*{英}\addcontentsline{loh}{figure}{英 \dpy{ying1}}

\begin{EntryWithPhonetic}{英}{ying1}{8}{⾋}
  \definition*{s.}{Reino Unido, abreviação de 英国 | Sobrenome: Ying}
  \definition{s.}{flor | herói; pessoa excepcional | uma pessoa de talento ou sabedoria extraordinários}
  \seealsoref{英国}{ying1guo2}
\end{EntryWithPhonetic}

\begin{EntryWithPhonetic}{英镑}{ying1bang4}{8,15}{⾋,⾦}[HSK 7-9]
  \definition{s.}{libra; libra esterlina; a moeda padrão do Reino Unido}
\end{EntryWithPhonetic}

\begin{EntryWithPhonetic}{英国}{ying1guo2}{8,8}{⾋,⼞}
  \definition*{s.}{Reino Unido; Grã-Bretanha; Inglaterra}
\end{EntryWithPhonetic}

\begin{EntryWithPhonetic}{英国人}{ying1guo2ren2}{8,8,2}{⾋,⼞,⼈}
  \definition{s.}{inglês | pessoa ou povo do Reino Unido}
\end{EntryWithPhonetic}

\begin{EntryWithPhonetic}{英俊}{ying1jun4}{8,9}{⾋,⼈}[HSK 7-9]
  \definition{adj.}{brilhante; extremamente talentoso; excepcionalmente talentoso | bonito e espirituoso}
  \synonymref{俊俏}{jun4qiao4}
  \synonymref{美丽}{mei3li4}
  \synonymref{潇洒}{xiao1sa3}
  \antonymref{丑陋}{chou3lou4}
  \antonymref{难看}{nan2kan4}
\end{EntryWithPhonetic}

\begin{EntryWithPhonetic}{英明}{ying1ming2}{8,8}{⾋,⽇}
  \definition{adj.}{sábio; brilhante; excelente e sábio}
\end{EntryWithPhonetic}

\begin{EntryWithPhonetic}{英文}{ying1wen2}{8,4}{⾋,⽂}[HSK 2]
  \definition{s.}{inglês, língua inglesa; a forma escrita do inglês}
\end{EntryWithPhonetic}

\begin{EntryWithPhonetic}{英雄}{ying1xiong2}{8,12}{⾋,⾫}[HSK 6]
  \definition{adj.}{heróico}
  \definition[名,个,位]{s.}{herói; uma pessoa cujas habilidades e coragem superam as das pessoas comuns | herói; aqueles que não têm medo das dificuldades, dos perigos ou da morte e que lutam bravamente pelos interesses do povo, mesmo ao custo das suas próprias vidas}
\end{EntryWithPhonetic}

\begin{EntryWithPhonetic}{英勇}{ying1yong3}{8,9}{⾋,⼒}[HSK 4]
  \definition{adj.}{heroico; valente; bravo; corajoso; extraordinariamente corajoso}
\end{EntryWithPhonetic}

\begin{EntryWithPhonetic}{英语}{ying1yu3}{8,9}{⾋,⾔}[HSK 2]
  \definition{s.}{inglês, língua inglesa}
\end{EntryWithPhonetic}

%%%%%%%%%% 婴 %%%%%%%%%%
\subsection*{婴}\addcontentsline{loh}{figure}{婴 \dpy{ying1}}

\begin{EntryWithPhonetic}{婴}{ying1}{11}{⼥}
  \definition*{s.}{Sobrenome: Ying}
  \definition{adj.}{infantil}
  \definition[个,名,位]{s.}{bebê; criança}
\end{EntryWithPhonetic}

\begin{EntryWithPhonetic}{婴儿}{ying1'er2}{11,2}{⼥,⼉}[HSK 7-9]
  \definition[个,名,位]{s.}{bebê; criança; crianças desde o nascimento até um ano de idade}
  \synonymref{宝宝}{bao3bao5}
\end{EntryWithPhonetic}

%%%%%%%%%% 樱 %%%%%%%%%%
\subsection*{樱}\addcontentsline{loh}{figure}{樱 \dpy{ying1}}

\begin{EntryWithPhonetic}{樱}{ying1}{15}{⽊}
  \definition[个,棵,朵]{s.}{cereja | cerejeira oriental; flores de cerejeira}
\end{EntryWithPhonetic}

\begin{EntryWithPhonetic}{樱桃}{ying1tao2}{15,10}{⽊,⽊}
  \definition{s.}{cereja}
\end{EntryWithPhonetic}

%%%%%%%%%% 鹦 %%%%%%%%%%
\subsection*{鹦}\addcontentsline{loh}{figure}{鹦 \dpy{ying1}}

\begin{EntryWithPhonetic}{鹦}{ying1}{16}{⿃}
  \definition[只]{s.}{papagaio}
\end{EntryWithPhonetic}

\begin{EntryWithPhonetic}{鹦鹉}{ying1wu3}{16,13}{⿃,⿃}
  \definition{s.}{papagaio (ave)}
\end{EntryWithPhonetic}

%%%%%%%%%% 鹰 %%%%%%%%%%
\subsection*{鹰}\addcontentsline{loh}{figure}{鹰 \dpy{ying1}}

\begin{EntryWithPhonetic}{鹰}{ying1}{18}{⿃}[HSK 7-9]
  \definition[只,群]{s.}{falcão; águia}
\end{EntryWithPhonetic}

%%%%%%%%%% 迎 %%%%%%%%%%
\subsection*{迎}\addcontentsline{loh}{figure}{迎 \dpy{ying2}}

\begin{EntryWithPhonetic}{迎}{ying2}{7}{⾡}[HSK 7-9]
  \definition*{s.}{Sobrenome: Ying}
  \definition{v.}{ir ao encontro; cumprimentar; acolher; receber | mover"-se em direção a; encontrar"-se cara a cara}
  \seealsoref{欢迎}{huan1ying2}
  \synonymref{接}{jie1}
\end{EntryWithPhonetic}

\begin{EntryWithPhonetic}{迎合}{ying2he2}{7,6}{⾡,⼝}[HSK 7-9]
  \definition{v.}{atender a; bajular; agradar a; concordar com o outro lado; fazer com que as palavras ou ações de alguém se conformem intencionalmente aos desejos de outros}
\end{EntryWithPhonetic}

\begin{EntryWithPhonetic}{迎接}{ying2jie1}{7,11}{⾡,⼿}[HSK 3]
  \definition{v.}{conhecer; cumprimentar; felicitar; dar as boas"-vindas | cumprimentar; felicitar; dar as boas"-vindas; preparar"-se; aguardar a chegada de um determinado momento ou evento}
\end{EntryWithPhonetic}

\begin{EntryWithPhonetic}{迎来}{ying2lai2}{7,7}{⾡,⽊}[HSK 6]
  \definition{v.}{dar boas"-vindas; cumprimentar | introduzir}
\end{EntryWithPhonetic}

%%%%%%%%%% 盈 %%%%%%%%%%
\subsection*{盈}\addcontentsline{loh}{figure}{盈 \dpy{ying2}}

\begin{EntryWithPhonetic}{盈}{ying2}{9}{⽫}
  \definition{v.}{estar cheio de; estar repleto de | ter um excedente de  | aumentar; expandir; crescer}
  \synonymref{满}{man3}
  \antonymref{亏}{kui1}
\end{EntryWithPhonetic}

\begin{EntryWithPhonetic}{盈利}{ying2/li4}{9,7}{⽫,⼑}[HSK 7-9]
  \definition[笔,份,些,点]{s.}{ganho; lucro; rendimento}
  \definition{v.+compl.}{lucrar; obter lucro; ganhar dinheiro}
  \synonymref{利润}{li4run4}
  \synonymref{剩余}{sheng4yu2}
  \antonymref{亏本}{kui1/ben3}
  \antonymref{亏损}{kui1sun3}
\end{EntryWithPhonetic}

%%%%%%%%%% 荧 %%%%%%%%%%
\subsection*{荧}\addcontentsline{loh}{figure}{荧 \dpy{ying2}}

\begin{EntryWithPhonetic}{荧}{ying2}{9}{⾋}
  \definition*{s.}{Marte (planeta)}
  \definition{adj.}{cintilante; brilhante; reluzente | perplexo; deslumbrado}
  \definition{s.}{um vislumbre | um brilho | fluorescência | fosforescência}
\end{EntryWithPhonetic}

\begin{EntryWithPhonetic}{荧光}{ying2guang1}{9,6}{⾋,⼉}[HSK 7-9]
  \definition{s.}{fluorescência; luz fluorescente; a fluorescência é a luz visível emitida por certas substâncias quando expostas à luz, elétrons ou outras formas de irradiação}
\end{EntryWithPhonetic}

%%%%%%%%%% 营 %%%%%%%%%%
\subsection*{营}\addcontentsline{loh}{figure}{营 \dpy{ying2}}

\begin{EntryWithPhonetic}{营}{ying2}{11}{⾋}
  \definition*{s.}{Sobrenome: Ying}
  \definition{s.}{acampamento; quartel; onde o exército está estacionado | batalhão; unidades militares}
  \definition{v.}{procurar | operar; executar; gerenciar}
\end{EntryWithPhonetic}

\begin{EntryWithPhonetic}{营救}{ying2jiu4}{11,11}{⾋,⽁}[HSK 7-9]
  \definition{v.}{socorrer; resgatar | salvar; tentar resgatar}
  \synonymref{救济}{jiu4ji4}
  \synonymref{救援}{jiu4yuan2}
  \synonymref{救助}{jiu4zhu4}
  \synonymref{抢救}{qiang3jiu4}
  \synonymref{援助}{yuan2zhu4}
  \synonymref{拯救}{zheng3jiu4}
\end{EntryWithPhonetic}

\begin{EntryWithPhonetic}{营养}{ying2yang3}{11,9}{⾋,⼋}[HSK 3]
  \definition[种]{s.}{nutrição; alimentação; a função do organismo de absorver as substâncias necessárias do meio externo para manter atividades vitais, como crescimento e desenvolvimento | nutrição; alimentação; ato ou processo de fornecer nutrição}
\end{EntryWithPhonetic}

\begin{EntryWithPhonetic}{营业}{ying2ye4}{11,5}{⾋,⼀}[HSK 4]
  \definition{v.}{fazer negócios; estar aberto para negócios}
\end{EntryWithPhonetic}

\begin{EntryWithPhonetic}{营造}{ying2zao4}{11,10}{⾋,⾡}[HSK 7-9]
  \definition{v.}{construir; edificar; produzir | fazer/construir algo com um propósito/sob planejamento; criar; construir}
  \antonymref{摧毁}{cui1hui3}
  \antonymref{毁灭}{hui3mie4}
  \antonymref{破坏}{po4huai4}
\end{EntryWithPhonetic}

%%%%%%%%%% 赢 %%%%%%%%%%
\subsection*{赢}\addcontentsline{loh}{figure}{赢 \dpy{ying2}}

\begin{EntryWithPhonetic}{赢}{ying2}{17}{⾙}[HSK 3]
  \definition{v.}{vencer; derrotar | ganhar; lucrar}
\end{EntryWithPhonetic}

\begin{EntryWithPhonetic}{赢得}{ying2de2}{17,11}{⾙,⼻}[HSK 4]
  \definition{v.}{ganhar; obter; conquistar; assegurar; garantir}
\end{EntryWithPhonetic}

\begin{EntryWithPhonetic}{赢家}{ying2jia1}{17,10}{⾙,⼧}[HSK 7-9]
  \definition[个]{s.}{vencedor; o vencedor de uma aposta ou competição}
  \antonymref{输家}{shu1jia1}
\end{EntryWithPhonetic}

%%%%%%%%%% 影 %%%%%%%%%%
\subsection*{影}\addcontentsline{loh}{figure}{影 \dpy{ying3}}

\begin{EntryWithPhonetic}{影}{ying3}{15}{⼺}
  \definition*{s.}{Sobrenome: Ying}
  \definition{s.}{sombra | reflexão; imagem | traço; sinal; impressão vaga | fotografia; imagem | filme | jogo de sombras; pantomima de sombra}
  \definition{v.}{(dialeto) esconder; ocultar | copiar; rastrear | fotocopiar}
\end{EntryWithPhonetic}

\begin{EntryWithPhonetic}{影迷}{ying3mi2}{15,9}{⼺,⾡}[HSK 6]
  \definition[个,名,位]{s.}{fã de cinema; entusiasta de cinema; pessoas viciadas em assistir filmes}
\end{EntryWithPhonetic}

\begin{EntryWithPhonetic}{影片}{ying3pian4}{15,4}{⼺,⽚}[HSK 2]
  \definition[部,盘,盒,卷]{s.}{filme; imagem | filme; película usada para reproduzir filmes}
\end{EntryWithPhonetic}

\begin{EntryWithPhonetic}{影视}{ying3shi4}{15,8}{⼺,⾒}[HSK 3]
  \definition{s.}{cinema e televisão combinados; denominação conjunta para cinema e TV}
\end{EntryWithPhonetic}

\begin{EntryWithPhonetic}{影响}{ying3xiang3}{15,9}{⼺,⼝}[HSK 2]
  \definition{s.}{efeito; influência; efeitos sobre pessoas ou coisas}
  \definition{v.}{afetar; influenciar; influência sobre os pensamentos ou ações dos outros}
\end{EntryWithPhonetic}

\begin{EntryWithPhonetic}{影响力}{ying3 xiang3 li4}{15,9,2}{⼺,⼝,⼒}
  \definition{s.}{impacto | influência}
\end{EntryWithPhonetic}

\begin{EntryWithPhonetic}{影像}{ying3xiang4}{15,13}{⼺,⼈}[HSK 7-9]
  \definition{s.}{retrato | imagem; forma; figura; impressão | ícone; silhueta; vídeo; fantasma; a forma de um objeto conforme percebida por dispositivos ópticos, dispositivos eletrônicos, etc.}
  \synonymref{镜子}{jing4zi5}
  \synonymref{印象}{yin4xiang4}
\end{EntryWithPhonetic}

\begin{EntryWithPhonetic}{影星}{ying3xing1}{15,9}{⼺,⽇}[HSK 6]
  \definition{s.}{estrela de cinema}
\end{EntryWithPhonetic}

\begin{EntryWithPhonetic}{影子}{ying3zi5}{15,3}{⼺,⼦}[HSK 4]
  \definition[个,片]{s.}{sombra; imagem projetada por um objeto, etc., que bloqueia a luz | reflexão; reflexo; imagem de um objeto, etc., conforme aparece em um refletor, como um espelho, uma superfície de água, etc. | sinal; vestígio; vaga impressão}
\end{EntryWithPhonetic}

%%%%%%%%%% 应 %%%%%%%%%%
\subsection*{应}\addcontentsline{loh}{figure}{应 \dpy{ying4}}

\begin{EntryWithPhonetic}{应酬}{ying4chou5}{7,13}{⼴,⾣}[HSK 7-9]
  \definition[次,场]{s.}{festa; evento social; atividades sociais como encontros e refeições em conjunto}
  \definition{v.}{socializar; tratar com cortesia; participar de atividades sociais}
  \synonymref{应付}{ying4fu5}
\end{EntryWithPhonetic}

\begin{EntryWithPhonetic}{应对}{ying4dui4}{7,5}{⼴,⼨}[HSK 6]
  \definition{v.}{reagir; responder; lidar com; dar uma resposta; tomar medidas e contramedidas para lidar com a situação}
\end{EntryWithPhonetic}

\begin{EntryWithPhonetic}{应付}{ying4fu5}{7,5}{⼴,⼈}[HSK 7-9]
  \definition{v.}{encontrar; lidar com; enfrentar; superar; tomar medidas ou usar métodos para lidar com (uma pessoa ou uma situação) | improvizar; fazer algo descuidadamente; fazer algo de forma superficial}
  \seealsoref{处理}{chu3li3}
  \synonymref{对待}{dui4dai4}
  \synonymref{对付}{dui4fu5}
  \synonymref{应对}{ying4dui4}
  \synonymref{应酬}{ying4chou5}
  \antonymref{竭力}{jie2li4}
  \antonymref{探索}{tan4suo3}
  \antonymref{悉心}{xi1xin1}
\end{EntryWithPhonetic}

\begin{EntryWithPhonetic}{应急}{ying4 ji2}{7,9}{⼴,⼼}[HSK 6]
  \definition{v.}{atender a uma necessidade urgente (emergência, contingência, etc.)}
\end{EntryWithPhonetic}

\begin{EntryWithPhonetic}{应聘}{ying4pin4}{7,13}{⼴,⽿}[HSK 7-9]
  \definition{v.}{candidatar"-se a uma vaga; candidatar"-se a um emprego enviando um currículo para empresas de recrutamento | aceitar uma oferta de emprego ou ser convidado a trabalhar para uma empresa}
  \antonymref{招聘}{zhao1pin4}
\end{EntryWithPhonetic}

\begin{EntryWithPhonetic}{应邀}{ying4yao1}{7,16}{⼴,⾡}[HSK 7-9]
  \definition{v.}{aceitar o convite de alguém}
  \synonymref{聘请}{pin4qing3}
  \antonymref{邀请}{yao1qing3}
\end{EntryWithPhonetic}

\begin{EntryWithPhonetic}{应用}{ying4yong4}{7,5}{⼴,⽤}[HSK 3]
  \definition{adj.}{aplicado (na vida ou na produção); usado diretamente na vida ou na produção}
  \definition{v.}{usar; aplicar}
\end{EntryWithPhonetic}

\begin{EntryWithPhonetic}{应用程序}{ying4yong4 cheng2xu4}{7,5,12,7}{⼴,⽤,⽲,⼴}
  \definition{s.}{programa aplicativo; principais categorias de \emph{software}}
\end{EntryWithPhonetic}

\begin{EntryWithPhonetic}{应用程序编程接口}{ying4yong4 cheng2xu4 bian1cheng2 jie1kou3}{7,5,12,7,12,12,11,3}{⼴,⽤,⽲,⼴,⽷,⽲,⼿,⼝}
  \definition{s.}{API (\emph{application programming interface})}
  \seealsoref{应用程序接口}{ying4yong4 cheng2xu4 jie1kou3}
\end{EntryWithPhonetic}

\begin{EntryWithPhonetic}{应用程序接口}{ying4yong4 cheng2xu4 jie1kou3}{7,5,12,7,11,3}{⼴,⽤,⽲,⼴,⼿,⼝}
  \definition{s.}{API (\emph{application programming interface})}
  \seealsoref{应用程序编程接口}{ying4yong4 cheng2xu4 bian1cheng2 jie1kou3}
\end{EntryWithPhonetic}

%%%%%%%%%% 映 %%%%%%%%%%
\subsection*{映}\addcontentsline{loh}{figure}{映 \dpy{ying4}}

\begin{EntryWithPhonetic}{映}{ying4}{9}{⽇}[HSK 7-9]
  \definition{v.}{refletir; espelhar; brilhar; a imagem de um objeto é revelada pela luz | projetar um filme}
\end{EntryWithPhonetic}

%%%%%%%%%% 硬 %%%%%%%%%%
\subsection*{硬}\addcontentsline{loh}{figure}{硬 \dpy{ying4}}

\begin{EntryWithPhonetic}{硬}{ying4}{12}{⽯}[HSK 5]
  \definition{adj.}{duro; rígido; resistente;  objeto resistente e não se deforma facilmente quando submetido a forças externas | firme; forte; resistente; obstinado; (vontade, atitude, etc.) inabalável, forte e poderoso | capaz (pessoa); boa (qualidade) | rígido; severo; sem flexibilidade | duro; rígido; rigoroso; imutável}
  \definition{adv.}{conseguir fazer algo com dificuldade; indica fazer algo à força, independentemente das circunstâncias}
  \antonymref{软}{ruan3}
\end{EntryWithPhonetic}

\begin{EntryWithPhonetic}{硬币}{ying4bi4}{12,4}{⽯,⼱}[HSK 7-9]
  \definition[枚]{s.}{moeda; moeda metálica; moeda forte}
  \synonymref{钞票}{chao1piao4}
  \synonymref{纸币}{zhi3bi4}
\end{EntryWithPhonetic}

\begin{EntryWithPhonetic}{硬件}{ying4jian4}{12,6}{⽯,⼈}[HSK 5]
  \definition[种]{s.}{\emph{hardware}; nome genérico dado aos vários elementos, componentes e dispositivos que constituem um computador | máquina, materiais; equipamento; referência a máquinas, equipamentos, materiais físicos, etc., utilizados nos processos de produção, pesquisa científica, gestão, etc.}
\end{EntryWithPhonetic}

\begin{EntryWithPhonetic}{硬朗}{ying4lang5}{12,10}{⽯,⽉}[HSK 7-9]
  \definition{adj.}{saudável; com boa saúde; robusto; vigoroso; significa que os idosos estão com ótima saúde}
  \synonymref{健壮}{jian4zhuang4}
  \antonymref{软弱}{ruan3ruo4}
\end{EntryWithPhonetic}

\begin{EntryWithPhonetic}{硬盘}{ying4pan2}{12,11}{⽯,⽫}[HSK 7-9]
  \definition{s.}{disco rígido; abreviação de 硬磁盘}[硬盘扩展，10块9硬盘。===Expansão de disco rígido, 10 discos rígidos de 9 polegadas.]
\end{EntryWithPhonetic}

%%%%%%%%%% 拥 %%%%%%%%%%
\subsection*{拥}\addcontentsline{loh}{figure}{拥 \dpy{yong1}}

\begin{EntryWithPhonetic}{拥}{yong1}{8}{⼿}
  \definition{v.}{segurar nos braços; abraçar | reunir em volta; envolver em volta | aglomerar"-se; enxamear | para apoiar | (literário) ter; possuir}
\end{EntryWithPhonetic}

\begin{EntryWithPhonetic}{拥抱}{yong1bao4}{8,8}{⼿,⼿}[HSK 5]
  \definition[个,次]{s.}{abraço}
  \definition{v.}{abraçar; segurar em seus braços; abraçar para demonstrar afeto}
\end{EntryWithPhonetic}

\begin{EntryWithPhonetic}{拥护}{yong1hu4}{8,7}{⼿,⼿}[HSK 7-9]
  \definition{v.}{apoiar (usado para partidos políticos, líderes, políticas, rotas, etc.); defender; endossar}
  \synonymref{称赞}{cheng1zan4}
  \synonymref{附和}{fu4he4}
  \synonymref{赞成}{zan4cheng2}
  \synonymref{赞同}{zan4tong2}
  \synonymref{支持}{zhi1chi2}
  \antonymref{反驳}{fan3bo2}
  \antonymref{反对}{fan3dui4}
  \antonymref{叛逆}{pan4ni4}
  \antonymref{嫌弃}{xian2qi4}
\end{EntryWithPhonetic}

\begin{EntryWithPhonetic}{拥挤}{yong1ji3}{8,9}{⼿,⼿}[HSK 7-9]
  \definition{adj.}{lotado; cheio; descreve um grande número ou uma concentração densa de pessoas ou veículos}
  \definition{v.}{aglomerar; empurrar e espremer; (pessoas, veículos, navios, etc.) estar comprimidos uns contra os outros}
  \seealsoref{挤}{ji3}
  \synonymref{堵塞}{du3se4}
  \antonymref{宽广}{kuan1guang3}
  \antonymref{宽松}{kuan1song1}
\end{EntryWithPhonetic}

\begin{EntryWithPhonetic}{拥有}{yong1you3}{8,6}{⼿,⽉}[HSK 5]
  \definition{v.}{possuir; deter; ter (grande quantidade de terras, população, bens, etc.)}
\end{EntryWithPhonetic}

%%%%%%%%%% 庸 %%%%%%%%%%
\subsection*{庸}\addcontentsline{loh}{figure}{庸 \dpy{yong1}}

\begin{EntryWithPhonetic}{庸}{yong1}{11}{⼴}
  \definition{adj.}{comum; medíocre | sem inteligência; ineficaz}
  \definition{adv.}{como; por que; como poderia (usado em perguntas); palavras interrogativas, indicando perguntas retóricas}[庸知对错？===Como poderia distinguir o certo do errado?]
  \definition{v.}{precisar (usado em frases negativas); usar (para a forma negativa)}
\end{EntryWithPhonetic}

\begin{EntryWithPhonetic}{庸俗}{yong1su2}{11,9}{⼴,⼈}[HSK 7-9]
  \definition{adj.}{vulgar; burguês; baixo; medíocre e vulgar; nada nobre}
  \synonymref{低下}{di1xia4}
  \synonymref{平凡}{ping2fan2}
  \synonymref{平方}{ping2fang1}
  \antonymref{高明}{gao1ming2}
  \antonymref{高尚}{gao1shang4}
  \antonymref{高雅}{gao1ya3}
  \antonymref{名贵}{ming2gui4}
  \antonymref{文雅}{wen2ya3}
\end{EntryWithPhonetic}

%%%%%%%%%% 永 %%%%%%%%%%
\subsection*{永}\addcontentsline{loh}{figure}{永 \dpy{yong3}}

\begin{EntryWithPhonetic}{永}{yong3}{5}{⽔}
  \definition*{s.}{Sobrenome: Yong}
  \definition{adj.}{sempre; para sempre; perpetuamente}
  \definition{adv.}{para sempre; significa um tempo muito longo sem fim, o que equivale a 永远}
  \seealsoref{永远}{yong3yuan3}
\end{EntryWithPhonetic}

\begin{EntryWithPhonetic}{永不}{yong3 bu4}{5,4}{⽔,⼀}[HSK 7-9]
  \definition{adv.}{nunca; jamais}
\end{EntryWithPhonetic}

\begin{EntryWithPhonetic}{永恒}{yong3heng2}{5,9}{⽔,⼼}[HSK 7-9]
  \definition{adj.}{eterno; perpétuo; para sempre}
  \synonymref{长久}{chang2jiu3}
  \synonymref{长期}{chang2qi1}
  \synonymref{永久}{yong3jiu3}
  \synonymref{永远}{yong3yuan3}
  \antonymref{短暂}{duan3zan4}
  \antonymref{瞬间}{shun4jian1}
  \antonymref{暂时}{zan4shi2}
\end{EntryWithPhonetic}

\begin{EntryWithPhonetic}{永久}{yong3jiu3}{5,3}{⽔,⼃}[HSK 7-9]
  \definition{adj.}{permanente; perpétuo; eterno; para sempre; para o bem (e tudo)}
  \seealsoref{永远}{yong3yuan3}
  \synonymref{长久}{chang2jiu3}
  \synonymref{长期}{chang2qi1}
  \synonymref{长远}{chang2yuan3}
  \synonymref{持久}{chi2jiu3}
  \synonymref{好久}{hao3jiu3}
  \synonymref{永恒}{yong3heng2}
  \synonymref{永远}{yong3yuan3}
  \synonymref{悠久}{you1jiu3}
  \antonymref{临时}{lin2shi2}
  \antonymref{暂时}{zan4shi2}
\end{EntryWithPhonetic}

\begin{EntryWithPhonetic}{永远}{yong3yuan3}{5,7}{⽔,⾡}[HSK 2]
  \definition{adv.}{sempre; para sempre; Indica um longo período de tempo sem fim}
  \definition{s.}{eternidade; um futuro que nunca acaba}
\end{EntryWithPhonetic}

%%%%%%%%%% 泳 %%%%%%%%%%
\subsection*{泳}\addcontentsline{loh}{figure}{泳 \dpy{yong3}}

\begin{EntryWithPhonetic}{泳}{yong3}{8}{⽔}
  \definition{v.}{nadar}
\end{EntryWithPhonetic}

\begin{EntryWithPhonetic}{泳池}{yong3chi2}{8,6}{⽔,⽔}
  \definition{s.}{piscina}
  \seealsoref{游泳池}{you2yong3chi2}
  \seealsoref{游泳馆}{you2yong3guan3}
\end{EntryWithPhonetic}

\begin{EntryWithPhonetic}{泳衣}{yong3yi1}{8,6}{⽔,⾐}
  \definition{s.}{roupa de banho | maiô}
  \seealsoref{游泳衣}{you2yong3yi1}
\end{EntryWithPhonetic}

%%%%%%%%%% 勇 %%%%%%%%%%
\subsection*{勇}\addcontentsline{loh}{figure}{勇 \dpy{yong3}}

\begin{EntryWithPhonetic}{勇}{yong3}{9}{⼒}
  \definition*{s.}{Sobrenome: Yong}
  \definition{adj.}{bravo; valente; corajoso}
  \definition{s.}{recrutas temporários em tempos de guerra na Dinastia Qing}
\end{EntryWithPhonetic}

\begin{EntryWithPhonetic}{勇敢}{yong3gan3}{9,11}{⼒,⽁}[HSK 4]
  \definition{adj.}{bravo; valente; galante; corajoso}
\end{EntryWithPhonetic}

\begin{EntryWithPhonetic}{勇气}{yong3qi4}{9,4}{⼒,⽓}[HSK 4]
  \definition[种,股]{s.}{coragem; arrojo; nervos; coragem para agir sem medo}
\end{EntryWithPhonetic}

\begin{EntryWithPhonetic}{勇士}{yong3shi4}{9,3}{⼒,⼠}
  \definition{s.}{um guerreiro | uma pessoa corajosa}
\end{EntryWithPhonetic}

\begin{EntryWithPhonetic}{勇往直前}{yong3wang3-zhi2qian2}{9,8,8,9}{⼒,⼻,⽬,⼑}[HSK 7-9]
  \definition{expr.}{``Siga em frente com coragem.''; marchar para a frente com coragem; avançar bravamente; tomar a coragem com as duas mãos; avance com coragem; siga em frente destemidamente}
\end{EntryWithPhonetic}

\begin{EntryWithPhonetic}{勇于}{yong3yu2}{9,3}{⼒,⼆}[HSK 7-9]
  \definition{v.}{ser corajoso em; ser ousado em; ter a coragem de; não recuar}
  \synonymref{敢于}{gan3yu2}
  \antonymref{胆小}{dan3xiao3}
\end{EntryWithPhonetic}

%%%%%%%%%% 涌 %%%%%%%%%%
\subsection*{涌}\addcontentsline{loh}{figure}{涌 \dpy{yong3}}

\begin{EntryWithPhonetic}{涌}{yong3}{10}{⽔}[HSK 7-9]
  \definition*{s.}{Sobrenome: Yong}
  \definition{v.}{jorrar; derramar; transbordar | surgir; brotar; emergir}
  \seeref{chong1}
\end{EntryWithPhonetic}

\begin{EntryWithPhonetic}{涌入}{yong3ru4}{10,2}{⽔,⼊}[HSK 7-9]
  \definition{v.}{afluir; entrar em grande quantidade/escala}
  \synonymref{流向}{liu2xiang4}
\end{EntryWithPhonetic}

\begin{EntryWithPhonetic}{涌现}{yong3xian4}{10,8}{⽔,⾒}[HSK 7-9]
  \definition{v.}{surgir; vir à tona; emergir em grande número; um grande número de casos surgiu ao mesmo tempo}
  \synonymref{表现}{biao3xian4}
  \synonymref{呈现}{cheng2xian4}
  \synonymref{出现}{chu1xian4}
  \synonymref{发现}{fa1xian4}
  \synonymref{浮现}{fu2xian4}
  \synonymref{显示}{xian3shi4}
  \synonymref{展示}{zhan3shi4}
  \synonymref{展现}{zhan3xian4}
  \antonymref{埋没}{mai2mo4}
\end{EntryWithPhonetic}

%%%%%%%%%% 踊 %%%%%%%%%%
\subsection*{踊}\addcontentsline{loh}{figure}{踊 \dpy{yong3}}

\begin{EntryWithPhonetic}{踊}{yong3}{14}{⾜}
  \definition{v.}{saltar; pular}[她兴奋地踊跃起来。===Ela pulou de alegria.]
\end{EntryWithPhonetic}

\begin{EntryWithPhonetic}{踊跃}{yong3yue4}{14,11}{⾜,⾜}[HSK 7-9]
  \definition{adj.}{positivo; entusiasmado; descreve um estado de grande entusiasmo e desejo de ser o primeiro}
  \definition{v.}{pular; saltar; dar um salto}
  \synonymref{奋勇}{fen4yong3}
  \synonymref{积极}{ji1ji2}
  \synonymref{主动}{zhu3dong4}
  \antonymref{消极}{xiao1ji2}
\end{EntryWithPhonetic}

%%%%%%%%%% 用 %%%%%%%%%%
\subsection*{用}\addcontentsline{loh}{figure}{用 \dpy{yong4}}

\begin{EntryWithPhonetic}{用}{yong4}{5}{⽤}[HSK 1][Kangxi 101]
  \definition*{s.}{Sobrenome: Yong}
  \definition{conj.}{portanto; por isso; assim sendo; razões para a introdução, equivalentes a 因}
  \definition{prep.}{com; ação de introduzir ferramentas, meios, etc. utilizados ou empregados}
  \definition{s.}{despesas; gastos; custos | uso; utilidade; eficácia}
  \definition{v.}{usar; aplicar; empregar | necessitar (normalmente na forma negativa) | respeitosamente: comer; beber}
  \seealsoref{因}{yin1}
\end{EntryWithPhonetic}

\begin{EntryWithPhonetic}{用不着}{yong4bu4zhao2}{5,4,11}{⽤,⼀,⽬}[HSK 5]
  \definition{v.}{não precisar; não ter utilidade para; não haver necessidade de}
\end{EntryWithPhonetic}

\begin{EntryWithPhonetic}{用餐}{yong4/can1}{5,16}{⽤,⾷}[HSK 7-9]
  \definition{v.+compl.}{jantar; fazer uma refeição; comer (com tom solene)}
\end{EntryWithPhonetic}

\begin{EntryWithPhonetic}{用处}{yong4chu3}{5,5}{⽤,⼡}[HSK 6]
  \definition[个]{s.}{uso; usabilidade; utilidade}
\end{EntryWithPhonetic}

\begin{EntryWithPhonetic}{用得着}{yong4de5zhao2}{5,11,11}{⽤,⼻,⽬}[HSK 6]
  \definition{adj.}{útil; necessário}
  \definition{v.}{precisar; achar algo útil | ter necessidade de;  ser necessário; valer a pena}
\end{EntryWithPhonetic}

\begin{EntryWithPhonetic}{用法}{yong4fa3}{5,8}{⽤,⽔}[HSK 6]
  \definition[种,个]{s.}{uso; emprego; a maneira de usar}
\end{EntryWithPhonetic}

\begin{EntryWithPhonetic}{用功}{yong4gong1}{5,5}{⽤,⼒}[HSK 7-9]
  \definition{adj.}{diligente; estudioso; trabalhador}
  \synonymref{勤奋}{qin2fen4}
  \synonymref{勤劳}{qin2lao2}
  \synonymref{辛勤}{xin1qin2}
  \antonymref{懒惰}{lan3duo4}
\end{EntryWithPhonetic}

\begin{EntryWithPhonetic}{用户}{yong4hu4}{5,4}{⽤,⼾}[HSK 5]
  \definition[个,位,名]{s.}{usuário; consumidor; entidades e indivíduos que utilizam determinados equipamentos públicos ou bens de consumo}
\end{EntryWithPhonetic}

\begin{EntryWithPhonetic}{用来}{yong4lai2}{5,7}{⽤,⽊}[HSK 5]
  \definition{v.}{ser usado para; depender (dele) ou usar (ele) para atingir algum objetivo}
\end{EntryWithPhonetic}

\begin{EntryWithPhonetic}{用力}{yong4/li4}{5,2}{⽤,⼒}[HSK 7-9]
  \definition{v.+compl.}{esforçar"-se; empregar as próprias forças; usar a força; exercer poder}
  \synonymref{使劲}{shi3/jin4}
\end{EntryWithPhonetic}

\begin{EntryWithPhonetic}{用料}{yong4liao4}{5,10}{⽤,⽃}
  \definition{s.}{ingredientes | materiais}
\end{EntryWithPhonetic}

\begin{EntryWithPhonetic}{用品}{yong4pin3}{5,9}{⽤,⼝}[HSK 6]
  \definition[批,件,种]{s.}{suprimentos; artigos para uso; itens para usar}
\end{EntryWithPhonetic}

\begin{EntryWithPhonetic}{用人}{yong4/ren2}{5,2}{⽤,⼈}[HSK 7-9]
  \definition{s.}{criado; servo; empregado; funcionário}
  \definition{v.+compl.}{escolher uma pessoa para um trabalho; utilizar recursos humanos; gerenciar pessoas | contratar alguém para um emprego; precisar de funcionários}
\end{EntryWithPhonetic}

\begin{EntryWithPhonetic}{用途}{yong4tu2}{5,10}{⽤,⾡}[HSK 4]
  \definition[个,种]{s.}{uso; aplicação; aspectos ou escopo da aplicação}
\end{EntryWithPhonetic}

\begin{EntryWithPhonetic}{用心}{yong4 xin1}{5,4}{⽤,⼼}[HSK 6]
  \definition{adj.}{diligente; atento; com atenção concentrada}
  \definition{s.}{motivo; intenção; o verdadeiro propósito ou razão para fazer algo}
\end{EntryWithPhonetic}

\begin{EntryWithPhonetic}{用意}{yong4yi4}{5,13}{⽤,⼼}[HSK 7-9]
  \definition[个]{s.}{intenção; propósito | significado}
  \synonymref{存心}{cun2xin1}
  \synonymref{故意}{gu4yi4}
  \synonymref{有意}{you3yi4}
  \synonymref{有益}{you3yi4}
  \synonymref{作用}{zuo4yong4}
  \antonymref{无意}{wu2yi4}
\end{EntryWithPhonetic}

\begin{EntryWithPhonetic}{用于}{yong4yu2}{5,3}{⽤,⼆}[HSK 5]
  \definition{v.}{usar para; ser usado para; usar em}
\end{EntryWithPhonetic}

%%%%%%%%%% 优 %%%%%%%%%%
\subsection*{优}\addcontentsline{loh}{figure}{优 \dpy{you1}}

\begin{EntryWithPhonetic}{优}{you1}{6}{⼈}[HSK 7-9]
  \definition*{s.}{Sobrenome: You}
  \definition{adj.}{excelente; ótimo; bom; excepcional | amplo; abundante}
  \definition{s.}{ator ou atriz de ópera chinesa; na antiguidade, referia"-se a pessoas que ganhavam a vida apresentando música, dança ou espetáculos de variedades; mais tarde, passou a se referir, de forma geral, a atores de ópera}
  \definition{v.}{dar tratamento preferencial; tratar bem; tratar com consideração especial}
  \synonymref{好}{hao3}
  \synonymref{良}{liang2}
  \antonymref{差}{cha4}
  \antonymref{劣}{lie4}
\end{EntryWithPhonetic}

\begin{EntryWithPhonetic}{优等}{you1deng3}{6,12}{⼈,⽵}
  \definition{adj.}{excelente | de primeira linha | alta classe | da mais alta ordem, superior}
  \antonymref{劣质}{lie4zhi4}
\end{EntryWithPhonetic}

\begin{EntryWithPhonetic}{优点}{you1dian3}{6,9}{⼈,⽕}[HSK 3]
  \definition[个,项,种,些]{s.}{mérito; virtude; ponto forte; vantagem (em oposição a 缺点)}
  \seealsoref{缺点}{que1dian3}
  \synonymref{长处}{chang2chu4}
  \synonymref{好处}{hao3chu5}
  \synonymref{便宜}{pian2yi5}
  \synonymref{所长}{suo3zhang3}
  \synonymref{甜头}{tian2tou5}
  \synonymref{益处}{yi4chu4}
  \antonymref{弊端}{bi4duan1}
  \antonymref{短处}{duan3chu5}
  \antonymref{毛病}{mao2bing5}
  \antonymref{缺点}{que1dian3}
  \antonymref{缺陷}{que1xian4}
  \antonymref{弱点}{ruo4dian3}
\end{EntryWithPhonetic}

\begin{EntryWithPhonetic}{优格}{you1ge2}{6,10}{⼈,⽊}
  \definition{s.}{Empréstimo linguístico: iogurte}
\end{EntryWithPhonetic}

\begin{EntryWithPhonetic}{优厚}{you1hou4}{6,9}{⼈,⼚}
  \definition{adj.}{munificente; generoso; abundante; favorável}
  \synonymref{丰厚}{feng1hou4}
  \synonymref{优越}{you1yue4}
  \antonymref{苛刻}{ke1ke4}
\end{EntryWithPhonetic}

\begin{EntryWithPhonetic}{优化}{you1hua4}{6,4}{⼈,⼔}[HSK 7-9]
  \definition{v.}{otimizar; fazer alterações ou seleções para tornar superior}[我们需要优化网站。===Precisamos otimizar o site.]
  \synonymref{改进}{gai3jin4}
  \synonymref{完善}{wan2shan4}
  \synonymref{整合}{zheng3he2}
\end{EntryWithPhonetic}

\begin{EntryWithPhonetic}{优惠}{you1hui4}{6,12}{⼈,⼼}[HSK 5]
  \definition{adj.}{especial; pechincha; reduzido; com desconto | favorável; preferencial; melhores condições ou tratamento do que o normal, permitindo que as pessoas obtenham mais benefícios}
  \synonymref{打折}{da3/zhe2}
  \synonymref{实惠}{shi2hui4}
  \synonymref{特价}{te4jia4}
  \synonymref{折扣}{zhe2kou4}
\end{EntryWithPhonetic}

\begin{EntryWithPhonetic}{优良}{you1liang2}{6,7}{⼈,⾉}[HSK 4]
  \definition{adj.}{ótimo; bom; excelente; (variedade, qualidade, desempenho, estilo, etc.) muito bom}
  \synonymref{崇高}{chong2gao1}
  \synonymref{杰出}{jie2chu1}
  \synonymref{良好}{liang2hao3}
  \synonymref{优秀}{you1xiu4}
  \synonymref{优异}{you1yi4}
  \synonymref{优越}{you1yue4}
  \antonymref{恶劣}{e4lie4}
  \antonymref{劣质}{lie4zhi4}
\end{EntryWithPhonetic}

\begin{EntryWithPhonetic}{优伶}{you1ling2}{6,7}{⼈,⼈}
  \definition{s.}{Obsoleto: ator; atriz; artista performático}
  \synonymref{演员}{yan3yuan2}
  \synonymref{艺人}{yi4ren2}
\end{EntryWithPhonetic}

\begin{EntryWithPhonetic}{优美}{you1mei3}{6,9}{⼈,⽺}[HSK 4]
  \definition{adj.}{fino; elegante; gracioso; bonito}
  \synonymref{精美}{jing1mei3}
  \synonymref{美好}{mei3hao3}
  \synonymref{美丽}{mei3li4}
  \synonymref{美妙}{mei3miao4}
  \synonymref{优雅}{you1ya3}
  \antonymref{丑陋}{chou3lou4}
\end{EntryWithPhonetic}

\begin{EntryWithPhonetic}{优盘}{you1pan2}{6,11}{⼈,⽫}
  \definition[个]{s.}{unidade de memória USB}
  \seealsoref{闪存盘}{shan3cun2pan2}
\end{EntryWithPhonetic}

\begin{EntryWithPhonetic}{优势}{you1shi4}{6,8}{⼈,⼒}[HSK 3]
  \definition[种,个]{s.}{vantagem; superioridade; preponderância; posição dominante; uma situação favorável que permite superar o adversário}
  \antonymref{劣势}{lie4shi4}
  \antonymref{弱点}{ruo4dian3}
\end{EntryWithPhonetic}

\begin{EntryWithPhonetic}{优先}{you1xian1}{6,6}{⼈,⼉}[HSK 5]
  \definition{adj.}{anterior; sênior; subjacente}
  \definition{v.}{ter prioridade; ter precedência; colocar"-se à frente de outras pessoas ou assuntos}
  \antonymref{稍后}{shao1hou4}
\end{EntryWithPhonetic}

\begin{EntryWithPhonetic}{优秀}{you1xiu4}{6,7}{⼈,⽲}[HSK 4]
  \definition{adj.}{esplêndido; excelente; extraordinário; excepcional; notável; descreve moral, qualidades, realizações, aprendizado, etc. muito bons.}
  \seealsoref{优美}{you1mei3}
  \synonymref{出色}{chu1se4}
  \synonymref{杰出}{jie2chu1}
  \synonymref{良好}{liang2hao3}
  \synonymref{突出}{tu1/chu1}
  \synonymref{先进}{xian1jin4}
  \synonymref{一流}{yi4liu2}
  \synonymref{优异}{you1yi4}
  \synonymref{优越}{you1yue4}
  \synonymref{优质}{you1zhi4}
  \antonymref{恶劣}{e4lie4}
\end{EntryWithPhonetic}

\begin{EntryWithPhonetic}{优选}{you1xuan3}{6,9}{⼈,⾡}
  \definition{adj.}{preferido; selecionado a dedo; escolhido a dedo}
  \definition{v.}{otimizar; escolher a melhor opção}
\end{EntryWithPhonetic}

\begin{EntryWithPhonetic}{优雅}{you1ya3}{6,12}{⼈,⾫}[HSK 7-9]
  \definition{adj.}{elegante; gracioso}
  \synonymref{高雅}{gao1ya3}
  \synonymref{温柔}{wen1rou2}
  \synonymref{文雅}{wen2ya3}
  \synonymref{优美}{you1mei3}
  \antonymref{粗暴}{cu1bao4}
  \antonymref{粗鲁}{cu1lu3}
\end{EntryWithPhonetic}

\begin{EntryWithPhonetic}{优异}{you1yi4}{6,6}{⼈,⼶}[HSK 7-9]
  \definition{adj.}{excelente; excepcional; excepcionalmente bom}
  \synonymref{崇高}{chong2gao1}
  \synonymref{杰出}{jie2chu1}
  \synonymref{良好}{liang2hao3}
  \synonymref{一流}{yi4liu2}
  \synonymref{优良}{you1liang2}
  \synonymref{优秀}{you1xiu4}
  \synonymref{优越}{you1yue4}
  \synonymref{优质}{you1zhi4}
\end{EntryWithPhonetic}

\begin{EntryWithPhonetic}{优于}{you1yu2}{6,3}{⼈,⼆}
  \definition{v.}{superar}
\end{EntryWithPhonetic}

\begin{EntryWithPhonetic}{优裕}{you1yu4}{6,12}{⼈,⾐}
  \definition{adj.}{abundante | bastante}
  \definition{s.}{abundância}
  \synonymref{充足}{chong1zu2}
  \synonymref{富足}{fu4zu2}
  \antonymref{拮据}{jie2ju1}
  \antonymref{窘迫}{jiong3po4}
  \antonymref{贫穷}{pin2qiong2}
\end{EntryWithPhonetic}

\begin{EntryWithPhonetic}{优越}{you1yue4}{6,12}{⼈,⾛}[HSK 7-9]
  \definition{adj.}{superior; vantajoso; excelente}
  \synonymref{出色}{chu1se4}
  \synonymref{杰出}{jie2chu1}
  \synonymref{良好}{liang2hao3}
  \synonymref{优厚}{you1hou4}
  \synonymref{优良}{you1liang2}
  \synonymref{优秀}{you1xiu4}
  \synonymref{优异}{you1yi4}
  \antonymref{平凡}{ping2fan2}
\end{EntryWithPhonetic}

\begin{EntryWithPhonetic}{优质}{you1zhi4}{6,8}{⼈,⾙}[HSK 6]
  \definition{adj.}{excelente qualidade; alta qualidade; qualidade superior; alto grau}
  \synonymref{优秀}{you1xiu4}
  \synonymref{优异}{you1yi4}
  \antonymref{变质}{bian4/zhi4}
  \antonymref{劣质}{lie4zhi4}
\end{EntryWithPhonetic}

%%%%%%%%%% 忧 %%%%%%%%%%
\subsection*{忧}\addcontentsline{loh}{figure}{忧 \dpy{you1}}

\begin{EntryWithPhonetic}{忧}{you1}{7}{⼼}
  \definition{s.}{tristeza; ansiedade; preocupação; cuidado; coisas que causam tristeza}
  \definition{v.}{preocupar"-se; estar preocupado; estar ansioso; estar triste}
\end{EntryWithPhonetic}

\begin{EntryWithPhonetic}{忧愁}{you1chou2}{7,13}{⼼,⼼}[HSK 7-9]
  \definition{adj.}{triste; preocupado; deprimido; melancólico}
  \synonymref{担心}{dan1/xin1}
  \synonymref{担忧}{dan1you1}
  \synonymref{发愁}{fa1/chou2}
  \synonymref{烦闷}{fan2men4}
  \synonymref{烦恼}{fan2nao3}
  \synonymref{忧虑}{you1lv4}
  \synonymref{忧郁}{you1yu4}
  \antonymref{高兴}{gao1xing4}
  \antonymref{欢快}{huan1kuai4}
  \antonymref{欢乐}{huan1le4}
  \antonymref{开朗}{kai1lang3}
  \antonymref{快乐}{kuai4le4}
  \antonymref{快活}{kuai4huo5}
  \antonymref{痛快}{tong4kuai5}
  \antonymref{喜悦}{xi3yue4}
  \antonymref{欣喜}{xin1xi3}
\end{EntryWithPhonetic}

\begin{EntryWithPhonetic}{忧虑}{you1lv4}{7,10}{⼼,⾌}[HSK 7-9]
  \definition{adj.}{preocupado; ansioso; apreensivo}
  \definition[个,丝]{s.}{preocupação; apreensão; um estado de espírito ou sentimento de ansiedade}
  \synonymref{担心}{dan1/xin1}
  \synonymref{担忧}{dan1you1}
  \synonymref{顾虑}{gu4lv4}
  \synonymref{焦虑}{jiao1lv4}
  \synonymref{困扰}{kun4rao3}
  \synonymref{忧愁}{you1chou2}
  \synonymref{忧郁}{you1yu4}
  \synonymref{着急}{zhao2/ji2}
  \antonymref{安心}{an1xin1}
  \antonymref{放心}{fang4/xin1}
  \antonymref{快活}{kuai4huo5}
  \antonymref{期望}{qi1wang4}
  \antonymref{舒畅}{shu1chang4}
\end{EntryWithPhonetic}

\begin{EntryWithPhonetic}{忧郁}{you1yu4}{7,8}{⼼,⾢}[HSK 7-9]
  \definition{adj.}{abatido; melancólico; desanimado; triste}
  \definition{s.}{depressão | melancolia}
  \synonymref{担忧}{dan1you1}
  \synonymref{难过}{nan2guo4}
  \synonymref{忧愁}{you1chou2}
  \synonymref{忧虑}{you1lv4}
  \antonymref{高兴}{gao1xing4}
  \antonymref{开朗}{kai1lang3}
  \antonymref{愉快}{yu2kuai4}
\end{EntryWithPhonetic}

%%%%%%%%%% 幽 %%%%%%%%%%
\subsection*{幽}\addcontentsline{loh}{figure}{幽 \dpy{you1}}

\begin{EntryWithPhonetic}{幽}{you1}{9}{⼳}
  \definition*{s.}{Sobrenome: You}
  \definition{adj.}{profundo e remoto; isolado; escuro | secreto; escondido; oculto; não público | quieto; tranquilo; sereno | do mundo inferior}
  \definition{s.}{mundo inferior}
\end{EntryWithPhonetic}

\begin{EntryWithPhonetic}{幽默}{you1mo4}{9,16}{⼳,⿊}[HSK 5]
  \definition{adj.}{humorístico; interessante ou engraçado, mas com um significado profundo}
  \definition{s.}{humor; lado engraçado; graça; características, temperamento, palavras ou comportamentos interessantes, engraçados ou significativos}
\end{EntryWithPhonetic}

%%%%%%%%%% 悠 %%%%%%%%%%
\subsection*{悠}\addcontentsline{loh}{figure}{悠 \dpy{you1}}

\begin{EntryWithPhonetic}{悠}{you1}{11}{⼼}
  \definition{adj.}{velho; demorado; remoto (no tempo ou no espaço) | tranquilo; relaxado; pensativo}
  \definition{v.}{balançar}
\end{EntryWithPhonetic}

\begin{EntryWithPhonetic}{悠久}{you1jiu3}{11,3}{⼼,⼃}[HSK 7-9]
  \definition{adj.}{longo; ancestral; de longa data; descreve uma história, cultura ou outro período histórico que tenha durado muito tempo}
  \synonymref{长久}{chang2jiu3}
  \synonymref{长远}{chang2yuan3}
  \synonymref{持久}{chi2jiu3}
  \synonymref{好久}{hao3jiu3}
  \synonymref{漫长}{man4chang2}
  \synonymref{深远}{shen1yuan3}
  \synonymref{永久}{yong3jiu3}
  \synonymref{永远}{yong3yuan3}
  \antonymref{短促}{duan3cu4}
  \antonymref{短暂}{duan3zan4}
\end{EntryWithPhonetic}

\begin{EntryWithPhonetic}{悠闲}{you1xian2}{11,7}{⼼,⾨}[HSK 7-9]
  \definition{adj.}{descontraído e despreocupado; um estado de lazer e contentamento}
  \antonymref{匆忙}{cong1mang2}
  \antonymref{繁忙}{fan2mang2}
  \antonymref{忙碌}{mang2lu4}
\end{EntryWithPhonetic}

%%%%%%%%%% 尤 %%%%%%%%%%
\subsection*{尤}\addcontentsline{loh}{figure}{尤 \dpy{you2}}

\begin{EntryWithPhonetic}{尤}{you2}{4}{⼪}
  \definition*{s.}{Sobrenome: You}
  \definition{adj.}{excelente; peculiar; notável}
  \definition{adv.}{particularmente; especialmente}
  \definition{s.}{falha; erro | irregularidade}
  \definition{v.}{ter rancor contra; culpar}
\end{EntryWithPhonetic}

\begin{EntryWithPhonetic}{尤其}{you2qi2}{4,8}{⼪,⼋}[HSK 5]
  \definition{adv.}{especialmente; particularmente; indica um grau mais avançado, equivalente a 更加}
  \seealsoref{更加}{geng4jia1}
\end{EntryWithPhonetic}

\begin{EntryWithPhonetic}{尤为}{you2wei2}{4,4}{⼪,⼂}[HSK 7-9]
  \definition{adv.}{especialmente; particularmente; usado antes de adjetivos ou verbos dissilábicos para indicar que algo se destaca no todo ou em comparação com outras coisas}[这件事尤为重要。===Este assunto é particularmente importante.]
  \synonymref{非常}{fei1chang2}
  \synonymref{极其}{ji2qi2}
  \synonymref{特别}{te4bie2}
  \synonymref{尤其}{you2qi2}
\end{EntryWithPhonetic}

%%%%%%%%%% 由 %%%%%%%%%%
\subsection*{由}\addcontentsline{loh}{figure}{由 \dpy{you2}}

\begin{EntryWithPhonetic}{由}{you2}{5}{⽥}[HSK 3]
  \definition*{s.}{Sobrenome: You}
  \definition{prep.}{por causa de; devido a | por; indica que algo deve ser feito por alguém | indica confiança em; indica dependência em | de; indica o ponto de partida | por; através de}
  \definition[个]{s.}{causa; razão; motivo}
  \definition{v.}{atravessar; passar por; seguir o caminho de | obedecer; seguir}
\end{EntryWithPhonetic}

\begin{EntryWithPhonetic}{由此}{you2ci3}{5,6}{⽥,⽌}[HSK 5]
  \definition{adv.}{assim; por meio disto; disto; daí; por causa disto; portanto; daqui; de agora em diante}
\end{EntryWithPhonetic}

\begin{EntryWithPhonetic}{由此看来}{you2ci3-kan4lai2}{5,6,9,7}{⽥,⽌,⽬,⽊}[HSK 7-9]
  \definition{expr.}{portanto, pode"-se observar que; deste ponto de vista; julgando por isso; tendo em vista isso}
\end{EntryWithPhonetic}

\begin{EntryWithPhonetic}{由此可见}{you2ci3-ke3jian4}{5,6,5,4}{⽥,⽌,⼝,⾒}[HSK 7-9]
  \definition{expr.}{portanto, pode"-se observar que\dots; disso podemos ver que\dots; isto demonstra que\dots; portanto, pode"-se observar que isso representa a conclusão}
\end{EntryWithPhonetic}

\begin{EntryWithPhonetic}{由来}{you2lai2}{5,7}{⽥,⽊}[HSK 7-9]
  \definition{s.}{razão; causa; fonte; motivo | origem; tempo desde o início até agora}
  \synonymref{来历}{lai2li4}
  \synonymref{来源}{lai2yuan2}
  \synonymref{理由}{li3you2}
  \synonymref{缘故}{yuan2gu4}
  \synonymref{原因}{yuan2yin1}
\end{EntryWithPhonetic}

\begin{EntryWithPhonetic}{由于}{you2yu2}{5,3}{⽥,⼆}[HSK 3]
  \definition{conj.}{porque; uma vez que; visto que;  usado no início da frase anterior, indica a razão, e a frase seguinte indica o resultado}
  \definition{prep.}{devido a; graças a; por causa de; em virtude de; como resultado de; introduzir a causa da ocorrência de eventos, ações, etc.}
\end{EntryWithPhonetic}

\begin{EntryWithPhonetic}{由衷}{you2zhong1}{5,10}{⽥,⾐}[HSK 7-9]
  \definition{adj.}{sincero; genuíno; autêntico; do fundo do coração}
  \synonymref{忠心}{zhong1xin1}
  \antonymref{虚假}{xu1jia3}
\end{EntryWithPhonetic}

%%%%%%%%%% 犹 %%%%%%%%%%
\subsection*{犹}\addcontentsline{loh}{figure}{犹 \dpy{you2}}

\begin{EntryWithPhonetic}{犹}{you2}{7}{⽝}
  \definition*{s.}{Sobrenome: You}
  \definition{adv.}{ainda | assim como; exatamente como; como se}
  \definition{v.}{ser exatamente como; ser como}
\end{EntryWithPhonetic}

\begin{EntryWithPhonetic}{犹如}{you2ru2}{7,6}{⽝,⼥}[HSK 7-9]
  \definition{conj.}{como; igualzinho a; exatamente como; tipo; como se fosse}
  \synonymref{好似}{hao3si4}
  \synonymref{好像}{hao3xiang4}
  \synonymref{如同}{ru2tong2}
  \synonymref{似乎}{si4hu1}
  \synonymref{相似}{xiang1si4}
  \synonymref{正如}{zheng4ru2}
  \antonymref{果断}{guo3duan4}
  \antonymref{绝对}{jue2dui4}
\end{EntryWithPhonetic}

\begin{EntryWithPhonetic}{犹豫}{you2yu4}{7,15}{⽝,⾗}[HSK 5]
  \definition{adj.}{hesitante; indeciso, incapaz de decidir ou agir}
  \definition{v.}{hesitar; ser indeciso}
\end{EntryWithPhonetic}

\begin{EntryWithPhonetic}{犹豫不决}{you2yu4-bu4jue2}{7,15,4,6}{⽝,⾗,⼀,⼎}[HSK 7-9]
  \definition{expr.}{hesitar; permanecer indeciso; ser irresoluto; hesitação; duvidoso; indeciso; irresoluto; vacilar; incapaz de se decidir}
  \synonymref{迟疑}{chi2yi2}
  \antonymref{不假思索}{bu4jia3-si1suo3}
  \antonymref{毫不犹豫}{hao2 bu4 you2yu4}
  \antonymref{死心塌地}{si3xin1-ta1di4}
\end{EntryWithPhonetic}

%%%%%%%%%% 邮 %%%%%%%%%%
\subsection*{邮}\addcontentsline{loh}{figure}{邮 \dpy{you2}}

\begin{EntryWithPhonetic}{邮}{you2}{7}{⾢}
  \definition*{s.}{Sobrenome: You}
  \definition{s.}{postal; correio; refere"-se a serviços postais | agência dos correios}
  \definition{v.}{postar; enviar pelo correio}
\end{EntryWithPhonetic}

\begin{EntryWithPhonetic}{邮包}{you2bao1}{7,5}{⾢,⼓}
  \definition{s.}{encomenda postal}
\end{EntryWithPhonetic}

\begin{EntryWithPhonetic}{邮编}{you2bian1}{7,12}{⾢,⽷}[HSK 7-9]
  \definition{s.}{código postal; código de área; abreviação de 邮政编码}
\end{EntryWithPhonetic}

\begin{EntryWithPhonetic}{邮递}{you2di4}{7,10}{⾢,⾡}
  \definition{v.}{enviar por correio}
\end{EntryWithPhonetic}

\begin{EntryWithPhonetic}{邮电}{you2dian4}{7,5}{⾢,⽥}
  \definition*{s.}{Correios e Telecomunicações}
\end{EntryWithPhonetic}

\begin{EntryWithPhonetic}{邮费}{you2fei4}{7,9}{⾢,⾙}
  \definition{s.}{postagem}
  \definition{v.}{postar}
\end{EntryWithPhonetic}

\begin{EntryWithPhonetic}{邮件}{you2jian4}{7,6}{⾢,⼈}[HSK 3]
  \definition[封,份,个,条]{s.}{correspondência; correio; assunto postal; termo que se refere a cartas, encomendas, etc., recebidos, transportados e entregues pelos correios | \emph{e"-mail}; refere"-se a \emph{e"-mails}, informações recebidas e enviadas através de caixas de correio eletrônico na \emph{Internet}, etc.}
\end{EntryWithPhonetic}

\begin{EntryWithPhonetic}{邮局}{you2ju2}{7,7}{⾢,⼫}[HSK 4]
  \definition[家,个]{s.}{correio; agência dos correios; organizações que lidam com serviços postais}
\end{EntryWithPhonetic}

\begin{EntryWithPhonetic}{邮迷}{you2mi2}{7,9}{⾢,⾡}
  \definition{s.}{filatelista | colecionador de selos}
\end{EntryWithPhonetic}

\begin{EntryWithPhonetic}{邮票}{you2piao4}{7,11}{⾢,⽰}[HSK 3]
  \definition[枚,张,套,版]{s.}{selo; selo postal; comprovante vendido pelos correios, usado para colar nas correspondências para indicar que o porte foi pago}
\end{EntryWithPhonetic}

\begin{EntryWithPhonetic}{邮市}{you2shi4}{7,5}{⾢,⼱}
  \definition{s.}{mercado postal}
\end{EntryWithPhonetic}

\begin{EntryWithPhonetic}{邮箱}{you2xiang1}{7,15}{⾢,⾋}[HSK 3]
  \definition{s.}{caixa de correio | \emph{mailbox}; refere"-se ao endereço de \emph{e"-mail}}
\end{EntryWithPhonetic}

\begin{EntryWithPhonetic}{邮政}{you2zheng4}{7,9}{⾢,⽁}[HSK 7-9]
  \definition{s.}{serviço postal; departamentos que gerenciam ou operam serviços de correio, telecomunicações, vale postal e distribuição de jornais/revistas}
  \synonymref{快递}{kuai4di4}
\end{EntryWithPhonetic}

\begin{EntryWithPhonetic}{邮资}{you2zi1}{7,10}{⾢,⾙}
  \definition{s.}{postagem}
\end{EntryWithPhonetic}

%%%%%%%%%% 油 %%%%%%%%%%
\subsection*{油}\addcontentsline{loh}{figure}{油 \dpy{you2}}

\begin{EntryWithPhonetic}{油}{you2}{8}{⽔}[HSK 2]
  \definition*{s.}{Sobrenome: You}
  \definition{adj.}{oleoso; gorduroso}
  \definition[瓶,滴,层]{s.}{óleo; gordura; graxa; petróleo}
  \definition{v.}{aplicar óleo de tungue, verniz ou tinta | estar manchado ou sujo com óleo ou graxa | aplicar óleo de tungue ou tinta}
\end{EntryWithPhonetic}

\begin{EntryWithPhonetic}{油画}{you2hua4}{8,8}{⽔,⽥}[HSK 7-9]
  \definition[幅]{s.}{pintura a óleo; um tipo de pintura ocidental, que utiliza tintas à base de óleo sobre tela ou painéis de madeira}[她在欣赏那幅油画。===Ela está admirando a pintura a óleo.]
\end{EntryWithPhonetic}

%%%%%%%%%% 游 %%%%%%%%%%
\subsection*{游}\addcontentsline{loh}{figure}{游 \dpy{you2}}

\begin{EntryWithPhonetic}{游}{you2}{12}{⽔}[HSK 3]
  \definition*{s.}{Sobrenome: You}
  \definition{adj.}{itinerante; não fixo; que se move frequentemente}
  \definition{s.}{parte de um rio; uma seção do rio; trecho; bacia; curso}
  \definition{v.}{nadar | vagar por aí; caminhar; viajar; fazer turismo | associar com (comunicação) | vagar; passear; andar tranquilamente por todos os lugares}
\end{EntryWithPhonetic}

\begin{EntryWithPhonetic}{游船}{you2chuan2}{12,11}{⽔,⾈}[HSK 7-9]
  \definition{s.}{navio de cruzeiro; barco de recreio; barcos para passeios turísticos}
\end{EntryWithPhonetic}

\begin{EntryWithPhonetic}{游客}{you2ke4}{12,9}{⽔,⼧}[HSK 2]
  \definition[个,位,名,群]{s.}{visitante; turista | (jogo \emph{on"-line}) jogador convidado}
\end{EntryWithPhonetic}

\begin{EntryWithPhonetic}{游览}{you2lan3}{12,9}{⽔,⾒}[HSK 7-9]
  \definition{v.}{fazer um \emph{tour}; visitar; passear; durante a caminhada, você poderá admirar locais históricos e paisagens}
  \seealsoref{观光}{guan1guang1}
  \synonymref{参观}{can1guan1}
  \synonymref{观察}{guan1cha2}
  \synonymref{观光}{guan1guang1}
  \synonymref{旅行}{lv3xing2}
  \synonymref{旅游}{lv3you2}
  \synonymref{视察}{shi4cha2}
  \synonymref{瞻仰}{zhan1yang3}
\end{EntryWithPhonetic}

\begin{EntryWithPhonetic}{游人}{you2ren2}{12,2}{⽔,⼈}[HSK 6]
  \definition[个,名,位,批]{s.}{visitante (de um parque, etc.); turista}
\end{EntryWithPhonetic}

\begin{EntryWithPhonetic}{游艇}{you2ting3}{12,12}{⽔,⾈}
  \definition[只]{s.}{barcaça | iate}
\end{EntryWithPhonetic}

\begin{EntryWithPhonetic}{游玩}{you2wan2}{12,8}{⽔,⽟}[HSK 6]
  \definition{v.}{brincar; jogar; divertir"-se | passear; vagar; fazer turismo}
\end{EntryWithPhonetic}

\begin{EntryWithPhonetic}{游戏}{you2xi4}{12,6}{⽔,⼽}[HSK 3]
  \definition[场]{s.}{jogo; recreação; atividades recreativas, como esconde"-esconde, adivinhar charadas, etc.; certas atividades esportivas não competitivas; jogos recreativos}
  \definition{v.}{jogar; fazer atividades divertidas e agradáveis, sozinho ou com outras pessoas}
\end{EntryWithPhonetic}

\begin{EntryWithPhonetic}{游戏机}{you2xi4ji1}{12,6,6}{⽔,⼽,⽊}[HSK 6]
  \definition[台]{s.}{jogador de videogame | console | videogame}
\end{EntryWithPhonetic}

\begin{EntryWithPhonetic}{游行}{you2xing2}{12,6}{⽔,⾏}[HSK 6]
  \definition{s.}{desfilar; marchar; manifestar"-se; marchar em grupos nas ruas para celebrar, comemorar, manifestar"-se, etc.}
\end{EntryWithPhonetic}

\begin{EntryWithPhonetic}{游泳}{you2/yong3}{12,8}{⽔,⽔}[HSK 3]
  \definition[次]{s.}{natação; refere"-se ao esporte ou atividade de natação}
  \definition{v.+compl.}{nadar; pessoas ou animais nadando na água}
\end{EntryWithPhonetic}

\begin{EntryWithPhonetic}{游泳池}{you2yong3chi2}{12,8,6}{⽔,⽔,⽔}[HSK 5]
  \definition[场,个]{s.}{piscina; piscinas artificiais para natação, divididas em duas categorias: internas e externas}
  \seealsoref{泳池}{yong3chi2}
  \seealsoref{游泳馆}{you2yong3guan3}
\end{EntryWithPhonetic}

\begin{EntryWithPhonetic}{游泳馆}{you2yong3guan3}{12,8,11}{⽔,⽔,⾷}
  \definition{s.}{natatório; piscina coberta; edifícios esportivos usados principalmente para esportes aquáticos, como natação, mergulho e polo aquático}
  \seealsoref{泳池}{yong3chi2}
  \seealsoref{游泳池}{you2yong3chi2}
\end{EntryWithPhonetic}

\begin{EntryWithPhonetic}{游泳镜}{you2yong3jing4}{12,8,16}{⽔,⽔,⾦}
  \definition{s.}{óculos de natação}
\end{EntryWithPhonetic}

\begin{EntryWithPhonetic}{游泳衣}{you2yong3yi1}{12,8,6}{⽔,⽔,⾐}
  \definition{s.}{roupa de banho}
  \seealsoref{泳衣}{yong3yi1}
\end{EntryWithPhonetic}

%%%%%%%%%% 友 %%%%%%%%%%
\subsection*{友}\addcontentsline{loh}{figure}{友 \dpy{you3}}

\begin{EntryWithPhonetic}{友}{you3}{4}{⼜}
  \definition{adj.}{amigável; bom relacionamento; próximo | de relações amigáveis}
  \definition{s.}{amigo; pessoas intimamente relacionadas}
\end{EntryWithPhonetic}

\begin{EntryWithPhonetic}{友好}{you3hao3}{4,6}{⼜,⼥}[HSK 2]
  \definition{adj.}{amigável; amistoso; muito próximo, relacionamento muito bom; como bons amigos}
  \definition{s.}{amigo próximo, íntimo; em ocasiões formais referem"-se a bons amigos}
\end{EntryWithPhonetic}

\begin{EntryWithPhonetic}{友情}{you3qing2}{4,11}{⼜,⼼}[HSK 7-9]
  \definition[份]{s.}{amizade; sentimentos amistosos; sentimentos de amizade entre amigos}
  \synonymref{交情}{jiao1qing5}
  \synonymref{情谊}{qing2yi4}
  \synonymref{友好}{you3hao3}
  \synonymref{友谊}{you3yi4}
\end{EntryWithPhonetic}

\begin{EntryWithPhonetic}{友人}{you3ren2}{4,2}{⼜,⼈}[HSK 7-9]
  \definition{s.}{amigo}
  \synonymref{伙伴}{huo3ban4}
  \synonymref{朋友}{peng2you5}
  \synonymref{亲人}{qin1ren2}
  \synonymref{同伴}{tong2ban4}
  \antonymref{仇人}{chou2ren2}
  \antonymref{敌人}{di2ren2}
  \antonymref{恩人}{en1ren2}
\end{EntryWithPhonetic}

\begin{EntryWithPhonetic}{友善}{you3shan4}{4,12}{⼜,⼝}[HSK 7-9]
  \definition{adj.}{amigável; cordial; afável; íntimos e harmoniosos entre amigos}
  \synonymref{和睦}{he2mu4}
  \synonymref{善良}{shan4liang2}
  \synonymref{友好}{you3hao3}
  \antonymref{冷漠}{leng3mo4}
  \antonymref{凶恶}{xiong1'e4}
\end{EntryWithPhonetic}

\begin{EntryWithPhonetic}{友谊}{you3yi4}{4,10}{⼜,⾔}[HSK 5]
  \definition[段,份]{s.}{amizade; amizade entre amigos}
\end{EntryWithPhonetic}

%%%%%%%%%% 有 %%%%%%%%%%
\subsection*{有}\addcontentsline{loh}{figure}{有 \dpy{you3}}

\begin{EntryWithPhonetic}{有}{you3}{6}{⽉}[HSK 1]
  \definition*{s.}{Sobrenome: You}
  \definition{pref.}{usado antes do nome de certas dinastias ou etnias}
  \definition{v.}{ter; possuir; indica posse ou propriedade | existe; há; indica que certas coisas existem em certos lugares | fazer uma estimativa ou uma comparação; expressar estimativa ou comparação | indicar ação; indica que algo aconteceu ou ocorreu | usado antes de substantivos abstratos, indica quantidade ou grandeza | em termos gerais, semelhante a 某; refere"-se de maneira geral a algo semelhante | usado antes de pessoa, hora e lugar, indica a existência parcial | usado antes de certos verbos para formar uma expressão idiomática, indicando cortesia, polidez}
  \seealsoref{某}{mou3}
\end{EntryWithPhonetic}

\begin{EntryWithPhonetic}{有待}{you3dai4}{6,9}{⽉,⼻}[HSK 7-9]
  \definition{v.}{esperar; aguardar; precisar}
\end{EntryWithPhonetic}

\begin{EntryWithPhonetic}{有道理}{you3dao4li5}{6,12,11}{⽉,⾡,⽟}
  \definition{v.}{fazer sentido; ser bem fundamentado; haver verdade em}
\end{EntryWithPhonetic}

\begin{EntryWithPhonetic}{有的}{you3de5}{6,8}{⽉,⽩}[HSK 1]
  \definition{pron.}{algum, alguns}
\end{EntryWithPhonetic}

\begin{EntryWithPhonetic}{有的时候}{you3de5shi2hou4}{6,8,7,10}{⽉,⽩,⽇,⼈}
  \definition{adv.}{às vezes; ocasionalmente}
  \seealsoref{有时}{you3shi2}
  \seealsoref{有时候}{you3shi2hou5}
\end{EntryWithPhonetic}

\begin{EntryWithPhonetic}{有的是}{you3de5shi4}{6,8,9}{⽉,⽩,⽇}[HSK 3]
  \definition{expr.}{ter em abundância; não faltar; enfatizar que há muitos}
\end{EntryWithPhonetic}

\begin{EntryWithPhonetic}{有的放矢}{you3di4-fang4shi3}{6,8,8,5}{⽉,⽩,⽅,⽮}[HSK 7-9]
  \definition{expr.}{direcionado; atirar uma flecha em um alvo é uma metáfora para ter um objetivo claro em palavras ou ações}
\end{EntryWithPhonetic}

\begin{EntryWithPhonetic}{有点儿}{you3dian3r5}{6,9,2}{⽉,⽕,⼉}[HSK 2]
  \definition{adv.}{um pouco; indica um grau inferior, equivalente a 稍微 (usado principalmente para coisas que são insatisfatórias)}
  \definition{v.}{há um pouco; tem (ou ser de) algum; existem alguns}
  \seealsoref{稍微}{shao1wei1}
  \seealsoref{有（一）点儿}{you3 yi4 dian3r5}
\end{EntryWithPhonetic}

\begin{EntryWithPhonetic}{有毒}{you3du2}{6,9}{⽉,⽏}[HSK 5]
  \definition{adj.}{venenoso; tóxico; nocivo; geralmente é usada para descrever as propriedades nocivas à saúde de produtos químicos, plantas ou animais.}
\end{EntryWithPhonetic}

\begin{EntryWithPhonetic}{有关}{you3guan1}{6,6}{⽉,⼋}[HSK 6]
  \definition{prep.}{no caminho de; sobre}
  \definition{v.}{preocupar"-se com; relacionar"-se com; ter algo a ver com; existir algum tipo de relacionamento}
\end{EntryWithPhonetic}

\begin{EntryWithPhonetic}{有害}{you3hai4}{6,10}{⽉,⼧}[HSK 5]
  \definition{adj.}{prejudicial; nocivo; danoso; que pode causar danos ou prejuízos a algo}
\end{EntryWithPhonetic}

\begin{EntryWithPhonetic}{有机}{you3ji1}{6,6}{⽉,⽊}[HSK 7-9]
  \definition{adj.}{orgânico; vegetais e frutas cultivados e produzidos sem o uso de ingredientes sintetizados artificialmente | orgânico; compostos que contêm átomos de carbono (excluindo algumas substâncias simples que contêm carbono) | orgânico; intrínseco; isso descreve algo cujas partes estão todas relacionadas e não podem ser separadas}
\end{EntryWithPhonetic}

\begin{EntryWithPhonetic}{有劲儿}{you3jin4er5}{6,7,2}{⽉,⼒,⼉}[HSK 4]
  \definition{adj.}{interessante; divertido; estimulante | energético}
  \definition{v.}{ter força}
\end{EntryWithPhonetic}

\begin{EntryWithPhonetic}{有空儿}{you3kong4r5}{6,8,2}{⽉,⽳,⼉}[HSK 2]
  \definition{v.}{estar livre; ter tempo livre}
\end{EntryWithPhonetic}

\begin{EntryWithPhonetic}{有口无心}{you3kou3-wu2xin1}{6,3,4,4}{⽉,⼝,⽆,⼼}[HSK 7-9]
  \definition{expr.}{falar asperamente, mas sem má intenção; ser mordaz, mas não malicioso; ou: ``O latido de alguém é pior que a mordida.''}
\end{EntryWithPhonetic}

\begin{EntryWithPhonetic}{有劳}{you3lao2}{6,7}{⽉,⼒}
  \definition{v.}{posso incomodá-lo; desculpe incomodá-lo | (educado) obrigado pelo seu trabalho (usado ao pedir um favor ou após ter recebido um)}
\end{EntryWithPhonetic}

\begin{EntryWithPhonetic}{有力}{you3li4}{6,2}{⽉,⼒}[HSK 5]
  \definition{adj.}{forte; vigoroso; poderoso; energético}
\end{EntryWithPhonetic}

\begin{EntryWithPhonetic}{有利}{you3li4}{6,7}{⽉,⼑}[HSK 3]
  \definition{adj.}{benéfico; favorável; vantajoso}
\end{EntryWithPhonetic}

\begin{EntryWithPhonetic}{有利于}{you3li4yu2}{6,7,3}{⽉,⼑,⼆}[HSK 5]
  \definition{prep.}{disponível; é benéfico para alguém ou alguma coisa e pode ajudar e promovê-lo}
\end{EntryWithPhonetic}

\begin{EntryWithPhonetic}{有两下子}{you3 liang3xia4zi5}{6,7,3,3}{⽉,⼀,⼀,⼦}[HSK 7-9]
  \definition{v.}{Coloquial: ter habilidade real; conhecer bem o assunto}
\end{EntryWithPhonetic}

\begin{EntryWithPhonetic}{有没有}{you3mei2you3}{6,7,6}{⽉,⽔,⽉}[HSK 6]
  \definition{adv.}{Você tem\dots?; Você já\dots? ; Existe algum\dots?}
\end{EntryWithPhonetic}

\begin{EntryWithPhonetic}{有名}{you3 ming2}{6,6}{⽉,⼝}[HSK 1]
  \definition{adj.}{conhecido; famoso; célebre; nome conhecido por todos}
\end{EntryWithPhonetic}

\begin{EntryWithPhonetic}{有名无实}{you3ming2wu2shi2}{6,6,4,8}{⽉,⼝,⽆,⼧}
  \definition{v.}{(literal) tem um nome, mas não tem realidade | existe apenas no nome}
\end{EntryWithPhonetic}

\begin{EntryWithPhonetic}{有趣}{you3qu4}{6,15}{⽉,⾛}[HSK 4]
  \definition{adj.}{interessante; fascinante; divertido; excitante; estimulante}
\end{EntryWithPhonetic}

\begin{EntryWithPhonetic}{有人}{you3ren2}{6,2}{⽉,⼈}[HSK 2]
  \definition{adj.}{ocupado (como no banheiro)}
  \definition{pron.}{qualquer um; alguém}
  \definition[所]{s.}{pessoas}
  \definition{v.}{ter alguém ali}
\end{EntryWithPhonetic}

\begin{EntryWithPhonetic}{有声有色}{you3sheng1-you3se4}{6,7,6,6}{⽉,⼠,⽉,⾊}[HSK 7-9]
  \definition{expr.}{``Que tem som e cor.''; repleto de som e cor; vívido e dramático; deslumbrante; impressionante}
  \synonymref{绘声绘色}{hui4sheng1-hui4se4}
\end{EntryWithPhonetic}

\begin{EntryWithPhonetic}{有时}{you3shi2}{6,7}{⽉,⽇}
  \definition{expr.}{às vezes; ocasionalmente; de vez em quando}
  \seealsoref{有的时候}{you3de5shi2hou4}
  \seealsoref{有时候}{you3shi2hou5}
\end{EntryWithPhonetic}

\begin{EntryWithPhonetic}{有时候}{you3shi2hou5}{6,7,10}{⽉,⽇,⼈}[HSK 1]
  \definition{adv.}{às vezes; indica um momento incerto, mas não frequente}
  \seealsoref{有的时候}{you3de5shi2hou4}
  \seealsoref{有时}{you3shi2}
\end{EntryWithPhonetic}

\begin{EntryWithPhonetic}{有时……有时……}{you3shi2 you3shi2}{6,7,6,7}{⽉,⽇,⽉,⽇}
  \definition{adv.}{às vezes\dots às vezes\dots}
\end{EntryWithPhonetic}

\begin{EntryWithPhonetic}{有事}{you3shi4}{6,8}{⽉,⼅}[HSK 6]
  \definition{v.}{estar ocupado; estar envolvido | ter algo acontecendo; sofrer um acidente; se meter em encrenca | (com 心里) ter algo em mente; estar ansioso; preocupar"-se | ter um emprego; estar empregado}[看他这几天愁眉苦脸的, 心里一定有事。===Vendo como ele parece triste ultimamente, deve haver algo em sua mente.]
  \seealsoref{心里}{xin1li3}
\end{EntryWithPhonetic}

\begin{EntryWithPhonetic}{有所}{you3suo3}{6,8}{⽉,⼾}[HSK 7-9]
  \definition{adv.}{de certa forma; em certa medida; até certo ponto; frequentemente seguido por verbos dissilábicos}
\end{EntryWithPhonetic}

\begin{EntryWithPhonetic}{有所不同}{you3suo3 bu4 tong2}{6,8,4,6}{⽉,⼾,⼀,⼝}[HSK 7-9]
  \definition{v.}{diferir em certa medida (expressão idiomática)}
\end{EntryWithPhonetic}

\begin{EntryWithPhonetic}{有望}{you3wang4}{6,11}{⽉,⽉}[HSK 7-9]
  \definition{adj.}{esperançoso; promissor}
  \definition{v.}{Literário: ter esperança; ser promissor}
  \synonymref{希望}{xi1wang4}
  \antonymref{绝望}{jue2/wang4}
\end{EntryWithPhonetic}

\begin{EntryWithPhonetic}{有限}{you3xian4}{6,8}{⽉,⾩}[HSK 4]
  \definition{adj.}{finito; limitado; restrito | indica baixo grau; indica pouco número; número baixo; nível baixo}
\end{EntryWithPhonetic}

\begin{EntryWithPhonetic}{有限公司}{you3xian4gong1si1}{6,8,4,5}{⽉,⾩,⼋,⼝}
  \definition{s.}{companhia limitada | corporação}
\end{EntryWithPhonetic}

\begin{EntryWithPhonetic}{有效}{you3xiao4}{6,10}{⽉,⽁}[HSK 3]
  \definition{adj.}{válido; eficiente; eficaz; capaz de alcançar os objetivos esperados}
\end{EntryWithPhonetic}

\begin{EntryWithPhonetic}{有效期}{you3xiao4qi1}{6,10,12}{⽉,⽁,⽉}[HSK 7-9]
  \definition{s.}{prazo de validade; período de validade}
  \synonymref{保质期}{bao3zhi4qi1}
\end{EntryWithPhonetic}

\begin{EntryWithPhonetic}{有些}{you3xie1}{6,8}{⽉,⼆}[HSK 1]
  \definition{adv.}{um pouco; bastante; ligeiramente}
  \definition{pron.}{uma parte; alguns}
  \definition{v.}{usado para indicar que há alguns, mas não muitos;}
  \seealsoref{有（一）些}{you3 (yi4) xie1}
\end{EntryWithPhonetic}

\begin{EntryWithPhonetic}{有幸}{you3xing4}{6,8}{⽉,⼲}[HSK 7-9]
  \definition{adv.}{felizmente; afortunadamente; por sorte}
  \synonymref{庆幸}{qing4xing4}
\end{EntryWithPhonetic}

\begin{EntryWithPhonetic}{有序}{you3xu4}{6,7}{⽉,⼴}[HSK 7-9]
  \definition{adj.}{ordenado; metódico; sistemático; em boa ordem; regular; sucessivo}
  \antonymref{混乱}{hun4luan4}
\end{EntryWithPhonetic}

\begin{EntryWithPhonetic}{有（一）点儿}{you3 yi4 dian3r5}{6,1,9,2}{⽉,⼀,⽕,⼉}
  \definition{adv.}{um pouco (有点儿 + {s.} ou {v. mental})}
  \seealsoref{有点儿}{you3dian3r5}
\end{EntryWithPhonetic}

\begin{EntryWithPhonetic}{有（一）些}{you3 (yi4) xie1}{6,1,8}{⽉,⼀,⼆}
  \definition{adv.}{em vez disso; em vez de; de certa forma}
  \definition{pron.}{de certa forma}
  \seealsoref{有些}{you3xie1}
\end{EntryWithPhonetic}

\begin{EntryWithPhonetic}{有益}{you3yi4}{6,10}{⽉,⽫}[HSK 7-9]
  \definition{adj.}{útil; benéfico}
  \definition{v.}{beneficiar; ser benéfico}
  \synonymref{存心}{cun2xin1}
  \synonymref{用意}{yong4yi4}
  \synonymref{有利}{you3li4}
  \antonymref{有害}{you3hai4}
\end{EntryWithPhonetic}

\begin{EntryWithPhonetic}{有意}{you3yi4}{6,13}{⽉,⼼}[HSK 7-9]
  \definition{adv.}{deliberadamente; propositalmente}
  \definition{v.}{ter a intenção; pretender | ter sentimentos; estar interessado em alguém; refere"-se ao afeto entre duas pessoas}
  \seealsoref{故意}{gu4yi4}
  \synonymref{存心}{cun2xin1}
  \synonymref{故意}{gu4yi4}
  \synonymref{用意}{yong4yi4}
  \antonymref{妨害}{fang2hai4}
  \antonymref{无故}{wu2gu4}
  \antonymref{无意}{wu2yi4}
  \antonymref{意外}{yi4wai4}
\end{EntryWithPhonetic}

\begin{EntryWithPhonetic}{有意思}{you3yi4si5}{6,13,9}{⽉,⼼,⼼}[HSK 2]
  \definition{adj.}{significativo; significativo e intrigante | interessante; agradável}
  \definition{v.}{ter interesse por; ser atraído sexualmente}
\end{EntryWithPhonetic}

\begin{EntryWithPhonetic}{有用}{you3yong4}{6,5}{⽉,⽤}[HSK 1]
  \definition{adj.}{útil; prático; funcional}
\end{EntryWithPhonetic}

\begin{EntryWithPhonetic}{有朝一日}{you3zhao1-yi1ri4}{6,12,1,4}{⽉,⽉,⼀,⽇}[HSK 7-9]
  \definition{expr.}{algum dia; se chegar o dia em que\dots; se por acaso\dots; um dia; em algum momento no futuro}
\end{EntryWithPhonetic}

\begin{EntryWithPhonetic}{有着}{you3zhe5}{6,11}{⽉,⽬}[HSK 5]
  \definition{v.}{ter; possuir; haver; existir}
\end{EntryWithPhonetic}

\begin{EntryWithPhonetic}{有助于}{you3zhu4yu2}{6,7,3}{⽉,⼒,⼆}[HSK 7-9]
  \definition{v.}{contribuir para; conduzir a; ser propício a; refere"-se ao processo de causar ou produzir um resultado específico sob circunstâncias específicas}
\end{EntryWithPhonetic}

%%%%%%%%%% 又 %%%%%%%%%%
\subsection*{又}\addcontentsline{loh}{figure}{又 \dpy{you4}}

\begin{EntryWithPhonetic}{又}{you4}{2}{⼜}[HSK 2][Kangxi 29]
  \definition{adv.}{indica repetição ou continuação | indica que várias situações ou propriedades existem simultaneamente | indica um nível mais profundo de significado | indica adicionar zero a números inteiros | indica duas coisas contraditórias | indica um ponto de virada, significando 可是 | usado em frases negativas ou perguntas retóricas para fortalecer o tom | além disso; indica informações adicionais ou suplementares}
  \seealsoref{可是}{ke3shi4}
\end{EntryWithPhonetic}

\begin{EntryWithPhonetic}{又称}{you4cheng1}{2,10}{⼜,⽲}
  \definition{s.}{também conhecido como}
\end{EntryWithPhonetic}

\begin{EntryWithPhonetic}{又及}{you4ji2}{2,3}{⼜,⼃}
  \definition{s.}{P.S., \emph{postscript}}
\end{EntryWithPhonetic}

\begin{EntryWithPhonetic}{又名}{you4ming2}{2,6}{⼜,⼝}
  \definition{s.}{também conhecido como | nome alternativo}
\end{EntryWithPhonetic}

\begin{EntryWithPhonetic}{又一次}{you4yi2ci4}{2,1,6}{⼜,⼀,⽋}
  \definition{adv.}{outra vez | mais uma vez | de novo}
\end{EntryWithPhonetic}

\begin{EntryWithPhonetic}{又……又……}{you4 you4}{2,2}{⼜,⼜}
  \definition{conj.}{\dots e\dots; tanto\dots como\dots; as duas palavras usadas depois de 又 não devem ter nenhuma conotação de contraste, ambas devem ser positivas ou negativas}[这件毛衣挺不错的，\underline{又}便宜\underline{又}漂亮。===Este suéter é muito bom, barato e bonito.]
\end{EntryWithPhonetic}

%%%%%%%%%% 右 %%%%%%%%%%
\subsection*{右}\addcontentsline{loh}{figure}{右 \dpy{you4}}

\begin{EntryWithPhonetic}{右}{you4}{5}{⼝}[HSK 1]
  \definition*{s.}{Sobrenome: You}
  \definition{adj.}{conservador; reacionário}
  \definition{s.}{a direita; o lado direito | oeste; na antiguidade, referia"-se especificamente à direção oeste (com base na orientação para o sul) | o lado direito como o lado de precedência; posição ou nível mais elevado (os antigos costumavam considerar a direita como mais respeitável)}
  \definition{v.}{favorecer; apoiar; reverenciar}
\end{EntryWithPhonetic}

\begin{EntryWithPhonetic}{右边}{you4bian5}{5,5}{⼝,⾡}[HSK 1]
  \definition{s.}{a direita; o lado direito; do lado direito}
\end{EntryWithPhonetic}

\begin{EntryWithPhonetic}{右侧}{you4ce4}{5,8}{⼝,⼈}
  \definition{s.}{lateral direita | lado direito}
\end{EntryWithPhonetic}

\begin{EntryWithPhonetic}{右面}{you4mian4}{5,9}{⼝,⾯}
  \definition{s.}{lado direito}
\end{EntryWithPhonetic}

\begin{EntryWithPhonetic}{右倾}{you4qing1}{5,10}{⼝,⼈}
  \definition{adj.}{conservador | reacionário}
\end{EntryWithPhonetic}

\begin{EntryWithPhonetic}{右手}{you4shou3}{5,4}{⼝,⼿}
  \definition{s.}{mão direita | lado direito}
\end{EntryWithPhonetic}

\begin{EntryWithPhonetic}{右袒}{you4tan3}{5,10}{⼝,⾐}
  \definition{v.}{ser tendencioso | ser parcial | favorecer um lado | tomar partido}
\end{EntryWithPhonetic}

\begin{EntryWithPhonetic}{右转}{you4zhuan3}{5,8}{⼝,⾞}
  \definition{v.}{virar à direita}
\end{EntryWithPhonetic}

%%%%%%%%%% 幼 %%%%%%%%%%
\subsection*{幼}\addcontentsline{loh}{figure}{幼 \dpy{you4}}

\begin{EntryWithPhonetic}{幼}{you4}{5}{⼳}
  \definition{adj.}{jovem; menor de idade}
  \definition{s.}{crianças; os jovens}
  \antonymref{老}{lao3}
\end{EntryWithPhonetic}

\begin{EntryWithPhonetic}{幼儿园}{you4'er2yuan2}{5,2,7}{⼳,⼉,⼞}[HSK 4]
  \definition[家,所]{s.}{jardim de infância; escola maternal; escola infantil; instituição para a educação de crianças pequenas}
\end{EntryWithPhonetic}

\begin{EntryWithPhonetic}{幼稚}{you4zhi4}{5,13}{⼳,⽲}[HSK 7-9]
  \definition{adj.}{infantil; jovem}
  \antonymref{成熟}{cheng2shu2}
\end{EntryWithPhonetic}

%%%%%%%%%% 诱 %%%%%%%%%%
\subsection*{诱}\addcontentsline{loh}{figure}{诱 \dpy{you4}}

\begin{EntryWithPhonetic}{诱}{you4}{9}{⾔}
  \definition{v.}{guiar; liderar; dirigir | atrair; seduzir; aliciar | induzir; causar; resultar de; levar a}
\end{EntryWithPhonetic}

\begin{EntryWithPhonetic}{诱饵}{you4'er3}{9,9}{⾔,⾷}[HSK 7-9]
  \definition{s.}{isca; chamariz}
\end{EntryWithPhonetic}

\begin{EntryWithPhonetic}{诱发}{you4fa1}{9,5}{⾔,⼜}[HSK 7-9]
  \definition{v.}{induzir; causar; trazer à tona (algo potencial ou latente) | provocar (uma doença); levar a (geralmente referindo"-se a doenças)}
  \synonymref{导致}{dao3zhi4}
  \synonymref{发作}{fa1zuo4}
  \synonymref{引起}{yin3qi3}
  \synonymref{造成}{zao4cheng2}
  \antonymref{杜绝}{du4jue2}
\end{EntryWithPhonetic}

\begin{EntryWithPhonetic}{诱惑}{you4huo4}{9,12}{⾔,⼼}[HSK 7-9]
  \definition[股,种,个]{s.}{tentação; refere"-se à atração que certas pessoas ou coisas exercem}
  \definition{v.}{atrair; tentar; seduzir; utilizar métodos para confundir as pessoas e induzi"-las a fazer coisas ruins | atrair; seduzir}
  \synonymref{引诱}{yin3you4}
\end{EntryWithPhonetic}

\begin{EntryWithPhonetic}{诱人}{you4ren2}{9,2}{⾔,⼈}[HSK 7-9]
  \definition{adj.}{tentador; sedutor; atraente para as pessoas}
  \synonymref{迷人}{mi2ren2}
  \synonymref{诱惑}{you4huo4}
  \antonymref{恶心}{e3xin5}
\end{EntryWithPhonetic}

%%%%%%%%%% 淤 %%%%%%%%%%
\subsection*{淤}\addcontentsline{loh}{figure}{淤 \dpy{yu1}}

\begin{EntryWithPhonetic}{淤}{yu1}{11}{⽔}
  \definition{adj.}{assoreado}
  \definition{s.}{lodo}
  \definition[出]{s.}{(medicina chinesa) estase de sangue}
  \definition{v.}{ficar assoreado; ficar sufocado com lodo | derramar; transbordar}
\end{EntryWithPhonetic}

\begin{EntryWithPhonetic}{淤泥}{yu1ni2}{11,8}{⽔,⽔}
  \definition[出]{s.}{gosma | limo | lodo}
\end{EntryWithPhonetic}

%%%%%%%%%% 与 %%%%%%%%%%
\subsection*{与}\addcontentsline{loh}{figure}{与 \dpy{yu2}}

\begin{EntryWithPhonetic}{与}{yu2}{3}{⼀}
  \definition{part.}{não; o quê; hein}
  \variantof{欤}
  \seeref{yu3}
  \seeref{yu4}
\end{EntryWithPhonetic}

%%%%%%%%%% 于 %%%%%%%%%%
\subsection*{于}\addcontentsline{loh}{figure}{于 \dpy{yu2}}

\begin{EntryWithPhonetic}{于}{yu2}{3}{⼆}[HSK 6]
  \definition*{s.}{Sobrenome: Yu}
  \definition{prep.}{indica hora, lugar, alcance, etc. | indica a direção da ação | usado depois de um verbo para indicar dar, entregar, etc. | apresentar a relação do objeto ou entidade introduzida | indica o ponto de início ou de partida | indica comparação | indica passividade}
\end{EntryWithPhonetic}

\begin{EntryWithPhonetic}{于是}{yu2shi4}{3,9}{⼆,⽇}[HSK 4]
  \definition{conj.}{então; portanto; consequentemente; como resultado; indica que o último segue o primeiro e que o último é frequentemente causado pelo primeiro}
  \synonymref{所以}{suo3yi3}
  \synonymref{因而}{yin1'er2}
  \synonymref{因此}{yin1ci3}
  \antonymref{但是}{dan4shi4}
  \antonymref{然而}{ran2'er2}
\end{EntryWithPhonetic}

%%%%%%%%%% 予 %%%%%%%%%%
\subsection*{予}\addcontentsline{loh}{figure}{予 \dpy{yu2}}

\begin{EntryWithPhonetic}{予}{yu2}{4}{⼅}
  \definition{pron.}{eu; mim}
  \seeref{yu3}
\end{EntryWithPhonetic}

%%%%%%%%%% 余 %%%%%%%%%%
\subsection*{余}\addcontentsline{loh}{figure}{余 \dpy{yu2}}

\begin{EntryWithPhonetic}{余}{yu2}{7}{⼈}[HSK 7-9]
  \definition*{s.}{Sobrenome: Yu}
  \definition{adj.}{reserva; excedente; restante; extra}
  \definition{num.}{ímpar; acima de; mais que; representa uma fração de um número inteiro, equivalente a 多}
  \definition{pron.}{eu; mim; quando o falante se refere a si mesmo, é equivalente a 我}
  \definition{s.}{após (um evento); (de tempo) além; fora ou depois de (um determinado evento ou situação) | excedente; excesso; resto; parte extra}
  \definition{v.}{permanecer; sair (após)}
  \seealsoref{多}{duo1}
  \seealsoref{剩}{sheng4}
  \seealsoref{我}{wo3}
  \antonymref{亏}{kui1}
  \antonymref{缺}{que1}
\end{EntryWithPhonetic}

\begin{EntryWithPhonetic}{余地}{yu2di4}{7,6}{⼈,⼟}[HSK 7-9]
  \definition{s.}{margem; folga; espaço; refere"-se a uma situação em que há margem de manobra em palavras ou ações}
\end{EntryWithPhonetic}

\begin{EntryWithPhonetic}{余额}{yu2'e2}{7,15}{⼈,⾴}[HSK 7-9]
  \definition{s.}{excedente; saldo; restante; quantia remanescente; saldo de fundos nas contas | vagas ainda não preenchidas; cotas excedentes; cotas vagas}
  \synonymref{剩余}{sheng4yu2}
\end{EntryWithPhonetic}

%%%%%%%%%% 欤 %%%%%%%%%%
\subsection*{欤}\addcontentsline{loh}{figure}{欤 \dpy{yu2}}

\begin{EntryWithPhonetic}{欤}{yu2}{7}{⽋}
  \definition{part.}{Literário: partícula final semelhante a 吗, 呢 ou 啊}
  \seealsoref{啊}{a1}
  \seealsoref{吗}{ma5}
  \seealsoref{呢}{ne5}
\end{EntryWithPhonetic}

%%%%%%%%%% 鱼 %%%%%%%%%%
\subsection*{鱼}\addcontentsline{loh}{figure}{鱼 \dpy{yu2}}

\begin{EntryWithPhonetic}{鱼}{yu2}{8}{⿂}[HSK 2][Kangxi 195]
  \definition*{s.}{Sobrenome: Yu}
  \definition[条,种,尾]{s.}{peixe; um vertebrado que vive na água; geralmente possui um corpo achatado lateralmente, fusiforme e com muitas escamas; nada com as nadadeiras e respira com as brânquias; sua temperatura corporal varia de acordo com a temperatura externa; existem muitas espécies, a maioria das quais comestíveis | carne de peixe; peixe (como alimento)}
\end{EntryWithPhonetic}

\begin{EntryWithPhonetic}{鱼船}{yu2chuan2}{8,11}{⿂,⾈}
  \definition{s.}{barco de pesca}
  \seealsoref{渔船}{yu2chuan2}
\end{EntryWithPhonetic}

\begin{EntryWithPhonetic}{鱼具}{yu2ju4}{8,8}{⿂,⼋}
  \variantof{渔具}
\end{EntryWithPhonetic}

\begin{EntryWithPhonetic}{鱼片}{yu2pian4}{8,4}{⿂,⽚}
  \definition{s.}{fatia de peixe | filé de peixe}
\end{EntryWithPhonetic}

\begin{EntryWithPhonetic}{鱼网}{yu2wang3}{8,6}{⿂,⽹}
  \variantof{渔网}
\end{EntryWithPhonetic}

\begin{EntryWithPhonetic}{鱼香}{yu2xiang1}{8,9}{⿂,⾹}
  \definition{s.}{um tempero da culinária chinesa que normalmente contém alho, cebolinha, gengibre, açúcar, sal, pimenta, etc.; embora 鱼香 signifique literalmente ``fragrância de peixe'', não contém frutos do mar}
\end{EntryWithPhonetic}

\begin{EntryWithPhonetic}{鱼香肉丝}{yu2xiang1rou4si1}{8,9,6,5}{⿂,⾹,⾁,⼀}
  \definition{s.}{Prato: tiras de carne de porco salteadas com molho picante}
  \seealsoref{鱼香}{yu2xiang1}
\end{EntryWithPhonetic}

\begin{EntryWithPhonetic}{鱼汛}{yu2xun4}{8,6}{⿂,⽔}
  \variantof{渔汛}
\end{EntryWithPhonetic}

%%%%%%%%%% 舁 %%%%%%%%%%
\subsection*{舁}\addcontentsline{loh}{figure}{舁 \dpy{yu2}}

\begin{EntryWithPhonetic}{舁}{yu2}{9}{⾅}
  \definition{v.}{levantar; elevar | aumentar}
\end{EntryWithPhonetic}

%%%%%%%%%% 娱 %%%%%%%%%%
\subsection*{娱}\addcontentsline{loh}{figure}{娱 \dpy{yu2}}

\begin{EntryWithPhonetic}{娱}{yu2}{10}{⼥}
  \definition{s.}{alegria; prazer; diversão; felicidade}
  \definition{v.}{dar prazer a; divertir; fazer feliz}
\end{EntryWithPhonetic}

\begin{EntryWithPhonetic}{娱乐}{yu2le4}{10,5}{⼥,⼃}[HSK 6]
  \definition[项]{s.}{entretenimento; diversão; recreação; passa"-tempo; atividades recreativas, prazerosas e divertidas}
  \definition{v.}{recrear; divertir; distrair; entreter; passar o tempo}
\end{EntryWithPhonetic}

%%%%%%%%%% 渔 %%%%%%%%%%
\subsection*{渔}\addcontentsline{loh}{figure}{渔 \dpy{yu2}}

\begin{EntryWithPhonetic}{渔}{yu2}{11}{⽔}
  \definition[条]{s.}{pescador}
  \definition{v.}{pescar}
\end{EntryWithPhonetic}

\begin{EntryWithPhonetic}{渔场}{yu2chang3}{11,6}{⽔,⼟}
  \definition{s.}{área de pesca}
\end{EntryWithPhonetic}

\begin{EntryWithPhonetic}{渔船}{yu2chuan2}{11,11}{⽔,⾈}[HSK 7-9]
  \definition[条]{s.}{barco de pesca; embarcação de pesca}
  \seealsoref{鱼船}{yu2chuan2}
\end{EntryWithPhonetic}

\begin{EntryWithPhonetic}{渔船队}{yu2chuan2 dui4}{11,11,4}{⽔,⾈,⾩}
  \definition{s.}{frota pesqueira}
\end{EntryWithPhonetic}

\begin{EntryWithPhonetic}{渔夫}{yu2fu1}{11,4}{⽔,⼤}
  \definition{s.}{pescador}
\end{EntryWithPhonetic}

\begin{EntryWithPhonetic}{渔具}{yu2ju4}{11,8}{⽔,⼋}
  \definition{s.}{equipamento de pesca}
\end{EntryWithPhonetic}

\begin{EntryWithPhonetic}{渔捞}{yu2lao1}{11,10}{⽔,⼿}
  \definition{s.}{pesca (como atividade comercial); pesca em massa; operações de pesca em grande escala}
\end{EntryWithPhonetic}

\begin{EntryWithPhonetic}{渔猎}{yu2lie4}{11,11}{⽔,⽝}
  \definition{s.}{pesca e caça}
  \definition{v.}{saquear | pilhar}
\end{EntryWithPhonetic}

\begin{EntryWithPhonetic}{渔笼}{yu2long2}{11,11}{⽔,⽵}
  \definition{s.}{gaiola de pesca | armadilha de pesca}
\end{EntryWithPhonetic}

\begin{EntryWithPhonetic}{渔轮}{yu2lun2}{11,8}{⽔,⾞}
  \definition{s.}{navio de pesca}
\end{EntryWithPhonetic}

\begin{EntryWithPhonetic}{渔民}{yu2min2}{11,5}{⽔,⽒}[HSK 7-9]
  \definition[个]{s.}{pescadores; pessoas que ganham a vida pescando}
\end{EntryWithPhonetic}

\begin{EntryWithPhonetic}{渔网}{yu2wang3}{11,6}{⽔,⽹}
  \definition{s.}{rede de pesca | tresmalho}
\end{EntryWithPhonetic}

\begin{EntryWithPhonetic}{渔汛}{yu2xun4}{11,6}{⽔,⽔}
  \definition{s.}{temporada de pesca}
\end{EntryWithPhonetic}

%%%%%%%%%% 愉 %%%%%%%%%%
\subsection*{愉}\addcontentsline{loh}{figure}{愉 \dpy{yu2}}

\begin{EntryWithPhonetic}{愉}{yu2}{12}{⼼}
  \definition{adj.}{satisfeito; feliz; alegre}
\end{EntryWithPhonetic}

\begin{EntryWithPhonetic}{愉快}{yu2kuai4}{12,7}{⼼,⼼}[HSK 6]
  \definition{adj.}{feliz; alegre; de bom humor, muito feliz}
\end{EntryWithPhonetic}

%%%%%%%%%% 逾 %%%%%%%%%%
\subsection*{逾}\addcontentsline{loh}{figure}{逾 \dpy{yu2}}

\begin{EntryWithPhonetic}{逾}{yu2}{12}{⾡}
  \definition{adv.}{Literário: ainda mais}
  \definition{v.}{exceder; cruzar; ir além}
\end{EntryWithPhonetic}

\begin{EntryWithPhonetic}{逾期}{yu2/qi1}{12,12}{⾡,⽉}[HSK 7-9]
  \definition{v.+compl.}{estar em atraso; não cumprir o prazo; ultrapassar o período prescrito}
  \synonymref{过期}{guo4/qi1}
\end{EntryWithPhonetic}

%%%%%%%%%% 愚 %%%%%%%%%%
\subsection*{愚}\addcontentsline{loh}{figure}{愚 \dpy{yu2}}

\begin{EntryWithPhonetic}{愚}{yu2}{13}{⼼}
  \definition{adj.}{tolo; estúpido}
  \definition{pron.}{eu; termos humildes usados para se referir a si mesmo}
  \definition{v.}{fazer de bobo; enganar}
  \synonymref{笨}{ben4}
  \synonymref{蠢}{chun3}
  \synonymref{傻}{sha3}
  \antonymref{智}{zhi4}
\end{EntryWithPhonetic}

\begin{EntryWithPhonetic}{愚蠢}{yu2chun3}{13,21}{⼼,⾍}[HSK 7-9]
  \definition{adj.}{estúpido; tolo; bobo; não inteligente}
  \synonymref{痴呆}{chi1dai1}
  \synonymref{无知}{wu2zhi1}
  \antonymref{聪慧}{cong1hui4}
  \antonymref{聪明}{cong1ming5}
  \antonymref{明智}{ming2zhi4}
  \antonymref{巧妙}{qiao3miao4}
  \antonymref{智慧}{zhi4hui4}
\end{EntryWithPhonetic}

\begin{EntryWithPhonetic}{愚公移山}{yu2gong1-yi2shan1}{13,4,11,3}{⼼,⼋,⽲,⼭}[HSK 7-9]
  \definition{expr.}{``O velho insensato move a montanha.''; figurativamente, querer é poder; faça coisas aparentemente impossíveis com perseverança obstinada e, eventualmente, tenha sucesso; conta a lenda que, na antiguidade, havia um velho chamado Yu Gong, de Beishan; duas grandes montanhas bloqueavam seu caminho em frente à sua casa, e ele resolveu arrasá"-las; outro velho, Zhi Sou, de Hequ, riu de sua tolice, achando impossível. Yu Gong respondeu: ``Quando eu morrer, terei filhos; quando meus filhos morrerem, terei netos; nossa descendência é infinita. Essas duas montanhas não crescerão mais. Cada pedaço que removermos será um pouco a menos; um dia elas serão arrasadas.'' (De Liezi, ``Tang Wen''《列子·汤问》); esta história é uma analogia para a perseverança e a coragem diante das dificuldades}
\end{EntryWithPhonetic}

%%%%%%%%%% 瑜 %%%%%%%%%%
\subsection*{瑜}\addcontentsline{loh}{figure}{瑜 \dpy{yu2}}

\begin{EntryWithPhonetic}{瑜}{yu2}{13}{⽟}
  \definition{s.}{Arcaico: jade fino; gema | (literário) brilho das gemas; virtudes; pontos positivos | excelência}
\end{EntryWithPhonetic}

\begin{EntryWithPhonetic}{瑜伽}{yu2jia1}{13,7}{⽟,⼈}
  \definition*{s.}{Ioga}
\end{EntryWithPhonetic}

\begin{EntryWithPhonetic}{瑜珈}{yu2jia1}{13,9}{⽟,⽟}
  \variantof{瑜伽}
\end{EntryWithPhonetic}

%%%%%%%%%% 虞 %%%%%%%%%%
\subsection*{虞}\addcontentsline{loh}{figure}{虞 \dpy{yu2}}

\begin{EntryWithPhonetic}{虞}{yu2}{13}{⾌}
  \definition*{s.}{Reino Yu, uma dinastia lendária fundada por Shun 舜 |Yu (um estado da Dinastia Zhou 周) | Sobrenome: Yu}
  \definition{s.}{Literário: suposição; previsão | Literário: ansiedade; preocupação}
  \definition{v.}{Literário: enganar; trapacear; fazer de bobo}
  \seealsoref{舜}{shun4}
  \seealsoref{周}{zhou1}
\end{EntryWithPhonetic}

\begin{EntryWithPhonetic}{虞世南}{yu2 shi4nan2}{13,5,9}{⾌,⼀,⼗}
  \definition*{s.}{Yu Shi'nan (558--638), político dos períodos Sui e Tang inicial, poeta e calígrafo, um dos Quatro Grandes Calígrafos do início da Dinastia Tang, 唐初四大家}
  \seealsoref{唐初四大家}{tang2 chu1 si4 da4jia1}
\end{EntryWithPhonetic}

%%%%%%%%%% 舆 %%%%%%%%%%
\subsection*{舆}\addcontentsline{loh}{figure}{舆 \dpy{yu2}}

\begin{EntryWithPhonetic}{舆}{yu2}{14}{⾅}
  \definition*{s.}{Sobrenome: Yu}
  \definition{adj.}{Obsoleto: público; popular}
  \definition{s.}{Arcaico: carruagem; carroça | liteira; cadeira de arruar | Obsoleto: área; território}
\end{EntryWithPhonetic}

\begin{EntryWithPhonetic}{舆论}{yu2lun4}{14,6}{⾅,⾔}[HSK 7-9]
  \definition{s.}{opinião pública}
  \synonymref{言论}{yan2lun4}
  \synonymref{议论}{yi4lun4}
\end{EntryWithPhonetic}

%%%%%%%%%% 与 %%%%%%%%%%
\subsection*{与}\addcontentsline{loh}{figure}{与 \dpy{yu3}}

\begin{EntryWithPhonetic}{与}{yu3}{3}{⼀}[HSK 6]
  \definition*{s.}{Sobrenome: Yu}
  \definition{conj.}{e; junto com}
  \definition{prep.}{com}
  \definition{v.}{dar; oferecer; conceder | conviver com; estar em bons termos com; socializar; ser amigável | ajudar; apoiar; patrocinar | Literário: esperar}
  \seeref{yu2}
  \seeref{yu4}
  \synonymref{给}{gei3}
  \synonymref{跟}{gen1}
  \synonymref{和}{he2}
  \antonymref{取}{qu3}
\end{EntryWithPhonetic}

\begin{EntryWithPhonetic}{与此同时}{yu3ci3-tong2shi2}{3,6,6,7}{⼀,⽌,⼝,⽇}[HSK 7-9]
  \definition{expr.}{ao mesmo tempo; entretanto; enquanto isso; simultaneamente; simultaneamente a esse evento; frequentemente usado para conectar duas frases ou orações que estão em relação paralela}[他在学习，与此同时，我在做饭。===Ele estava estudando, e ao mesmo tempo, eu estava cozinhando.]
  \synonymref{不相上下}{bu4xiang1-shang4xia4}
\end{EntryWithPhonetic}

\begin{EntryWithPhonetic}{与否}{yu3 fou3}{3,7}{⼀,⼝}[HSK 7-9]
  \definition{adv./part.}{independentemente de ser ou não; se ou não (no final de uma frase)}
  \synonymref{是否}{shi4fou3}
\end{EntryWithPhonetic}

\begin{EntryWithPhonetic}{与其}{yu3qi2}{3,8}{⼀,⼋}[HSK 7-9]
  \definition{conj.}{em vez de\dots; orações de ligação, indicando uma decisão tomada após comparação; 与其 é usado para o aspecto a ser abandonado, enquanto o aspecto a ser escolhido é frequentemente expresso por frases como 不如 ou 宁可}
  \seealsoref{不如}{bu4ru2}
  \seealsoref{宁可}{ning4ke3}
  \synonymref{如果}{ru2guo3}
  \synonymref{因为}{yin1wei5}
\end{EntryWithPhonetic}

\begin{EntryWithPhonetic}{与其……不如……}{yu3qi2 bu4ru2}{3,8,4,6}{⼀,⼋,⼀,⼥}
  \definition{conj.}{em vez de\dots é melhor\dots}
\end{EntryWithPhonetic}

\begin{EntryWithPhonetic}{与其……宁可……}{yu3qi2 ning4ke3}{3,8,5,5}{⼀,⼋,⼧,⼝}
  \definition{conj.}{em vez de\dots eu prefereria\dots}
\end{EntryWithPhonetic}

\begin{EntryWithPhonetic}{与日俱增}{yu3ri4-ju4zeng1}{3,4,10,15}{⼀,⽇,⼈,⼟}[HSK 7-9]
  \definition{expr.}{aumentando a cada dia; crescer a cada dia que passa; crescendo ao longo do tempo}
\end{EntryWithPhonetic}

\begin{EntryWithPhonetic}{与时俱进}{yu3shi2-ju4jin4}{3,7,10,7}{⼀,⽇,⼈,⾡}[HSK 7-9]
  \definition{expr.}{continuar a se desenvolver e avançar com o passar do tempo; manter"-se atualizado; estar a par dos desenvolvimentos modernos; progressivo; oportuno; acompanhar os tempos}
\end{EntryWithPhonetic}

\begin{EntryWithPhonetic}{与众不同}{yu3zhong4-bu4tong2}{3,6,4,6}{⼀,⼈,⼀,⼝}[HSK 7-9]
  \definition{expr.}{fora do comum; diferente dos demais; peculiar; incomum; significa ser diferente dos outros}
  \synonymref{鹤立鸡群}{he4li4ji1qun2}
  \antonymref{司空见惯}{si1kong1-jian4guan4}
\end{EntryWithPhonetic}

%%%%%%%%%% 予 %%%%%%%%%%
\subsection*{予}\addcontentsline{loh}{figure}{予 \dpy{yu3}}

\begin{EntryWithPhonetic}{予}{yu3}{4}{⼅}
  \definition{v.}{dar; conceder; outorgar}
  \seeref{yu2}
  \antonymref{夺}{duo2}
  \antonymref{取}{qu3}
\end{EntryWithPhonetic}

\begin{EntryWithPhonetic}{予以}{yu3yi3}{4,4}{⼅,⼈}[HSK 7-9]
  \definition{v.}{dar; conceder}
  \synonymref{给予}{ji3yu3}
  \antonymref{索取}{suo3qu3}
\end{EntryWithPhonetic}

%%%%%%%%%% 宇 %%%%%%%%%%
\subsection*{宇}\addcontentsline{loh}{figure}{宇 \dpy{yu3}}

\begin{EntryWithPhonetic}{宇}{yu3}{6}{⼧}
  \definition*{s.}{Sobrenome: Yu}
  \definition[座,栋]{s.}{beirais; calha; casa | espaço; universo; mundo | postura; porte}
\end{EntryWithPhonetic}

\begin{EntryWithPhonetic}{宇航员}{yu3hang2yuan2}{6,10,7}{⼧,⾈,⼝}[HSK 6]
  \definition[位,名,个,些]{s.}{astronauta; cosmonauta;}
\end{EntryWithPhonetic}

\begin{EntryWithPhonetic}{宇宙}{yu3zhou4}{6,8}{⼧,⼧}[HSK 7-9]
  \definition[个]{s.}{cosmos; universo; todo o espaço, incluindo a Terra}
  \synonymref{世界}{shi4jie4}
  \synonymref{天地}{tian1di4}
\end{EntryWithPhonetic}

%%%%%%%%%% 羽 %%%%%%%%%%
\subsection*{羽}\addcontentsline{loh}{figure}{羽 \dpy{yu3}}

\begin{EntryWithPhonetic}{羽}{yu3}{6}{⽻}[Kangxi 124]
  \definition*{s.}{Sobrenome: Yu}
  \definition{s.}{pena; pluma | asas (de pássaros ou insetos) | uma nota da antiga escala chinesa de cinco tons, correspondente a 6 na notação musical numerada}
\end{EntryWithPhonetic}

\begin{EntryWithPhonetic}{羽冠}{yu3guan1}{6,9}{⽻,⼍}
  \definition{s.}{crista emplumada (de pássaro)}
\end{EntryWithPhonetic}

\begin{EntryWithPhonetic}{羽林}{yu3lin2}{6,8}{⽻,⽊}
  \definition{s.}{escolta armada}
\end{EntryWithPhonetic}

\begin{EntryWithPhonetic}{羽流}{yu3liu2}{6,10}{⽻,⽔}
  \definition{s.}{pluma}
\end{EntryWithPhonetic}

\begin{EntryWithPhonetic}{羽毛}{yu3mao2}{6,4}{⽻,⽑}
  \definition{s.}{pena | plumagem | pluma}
\end{EntryWithPhonetic}

\begin{EntryWithPhonetic}{羽毛笔}{yu3mao2bi3}{6,4,10}{⽻,⽑,⽵}
  \definition{s.}{caneta de pena}
\end{EntryWithPhonetic}

\begin{EntryWithPhonetic}{羽毛球}{yu3mao2qiu2}{6,4,11}{⽻,⽑,⽟}[HSK 5]
  \definition[只,个]{s.}{\emph{badminton}; esporte com bola, as regras e equipamentos são bastante semelhantes ao tênis | peteca}
\end{EntryWithPhonetic}

\begin{EntryWithPhonetic}{羽绒服}{yu3rong2fu2}{6,9,8}{⽻,⽷,⽉}[HSK 5]
  \definition[件,个]{s.}{jaqueta de plumas; peça de vestuário com enchimento de plumas; casaco cujo interior é preenchido com penas de pato ou ganso}
\end{EntryWithPhonetic}

%%%%%%%%%% 雨 %%%%%%%%%%
\subsection*{雨}\addcontentsline{loh}{figure}{雨 \dpy{yu3}}

\begin{EntryWithPhonetic}{雨}{yu3}{8}{⾬}[HSK 1][Kangxi 173]
  \definition*{s.}{Sobrenome: Yu}
  \definition[场,阵,滴]{s.}{chuva; água que cai das nuvens para o solo}
  \seeref{yu4}
\end{EntryWithPhonetic}

\begin{EntryWithPhonetic}{雨伞}{yu3san3}{8,6}{⾬,⼈}
  \definition[把]{s.}{guarda"-chuva}
\end{EntryWithPhonetic}

\begin{EntryWithPhonetic}{雨蚀}{yu3shi2}{8,9}{⾬,⾷}
  \definition{s.}{erosão da chuva}
\end{EntryWithPhonetic}

\begin{EntryWithPhonetic}{雨水}{yu3shui3}{8,4}{⾬,⽔}[HSK 5]
  \definition{s.}{água da chuva; precipitação; chuva; água proveniente da chuva}
\end{EntryWithPhonetic}

\begin{EntryWithPhonetic}{雨靴}{yu3xue1}{8,13}{⾬,⾰}
  \definition[双]{s.}{botas de chuva}
\end{EntryWithPhonetic}

\begin{EntryWithPhonetic}{雨衣}{yu3yi1}{8,6}{⾬,⾐}[HSK 6]
  \definition[件,个]{s.}{capa de chuva; jaqueta impermeável; roupas impermeáveis}
\end{EntryWithPhonetic}

%%%%%%%%%% 语 %%%%%%%%%%
\subsection*{语}\addcontentsline{loh}{figure}{语 \dpy{yu3}}

\begin{EntryWithPhonetic}{语}{yu3}{9}{⾔}
  \definition{s.}{língua; linguagem | dito; provérbio; refere"-se especialmente a coloquialismos, provérbios, expressões idiomáticas ou palavras de livros antigos | sinal; meio não linguístico de comunicar ideias ; ações ou sinais que substituem palavras para expressar significado | palavras; expressão; refere"-se a uma palavra, frase ou sentença}
  \definition{v.}{dizer; falar | (pássaros, insetos, etc.) gorjear; pipilar}
  \seeref{yu4}
\end{EntryWithPhonetic}

\begin{EntryWithPhonetic}{语调}{yu3diao4}{9,10}{⾔,⾔}
  \definition[个]{s.}{entonação}
\end{EntryWithPhonetic}

\begin{EntryWithPhonetic}{语法}{yu3fa3}{9,8}{⾔,⽔}[HSK 4]
  \definition[个]{s.}{gramática; maneira como o idioma é estruturado, incluindo a formação e as variações de palavras, a organização de frases e sentenças | estudo da gramática; estudo das regras de estrutura linguística}
\end{EntryWithPhonetic}

\begin{EntryWithPhonetic}{语法术语}{yu3fa3 shu4yu3}{9,8,5,9}{⾔,⽔,⽊,⾔}
  \definition{s.}{termo gramatical}
\end{EntryWithPhonetic}

\begin{EntryWithPhonetic}{语气}{yu3qi4}{9,4}{⾔,⽓}[HSK 7-9]
  \definition[种,个]{s.}{tom; maneira de falar; as emoções expressas no tom de voz | humor; o modo verbal é uma categoria gramatical que indica a atitude do falante em relação ao conteúdo expresso; geralmente, divide"-se em quatro modos: declarativo, interrogativo, imperativo e exclamativo; o modo verbal costuma ser expresso por meio da entonação e de partículas modais}
  \synonymref{口气}{kou3qi4}
\end{EntryWithPhonetic}

\begin{EntryWithPhonetic}{语言}{yu3yan2}{9,7}{⾔,⾔}[HSK 2]
  \definition[种,门]{s.}{linguagem; é uma ferramenta exclusiva dos humanos para expressar ideias e comunicar pensamentos; é um fenômeno social especial e consiste em um sistema específico de pronúncia, vocabulário e gramática | linguagem falada}
\end{EntryWithPhonetic}

\begin{EntryWithPhonetic}{语言实验室}{yu3yan2shi2yan4shi4}{9,7,8,10,9}{⾔,⾔,⼧,⾺,⼧}
  \definition{s.}{laboratório de línguas}
\end{EntryWithPhonetic}

\begin{EntryWithPhonetic}{语音}{yu3yin1}{9,9}{⾔,⾳}[HSK 4]
  \definition{s.}{voz; pronúncia; sons da fala; som de alguém falando | pronúncia; som do idioma}
\end{EntryWithPhonetic}

%%%%%%%%%% 与 %%%%%%%%%%
\subsection*{与}\addcontentsline{loh}{figure}{与 \dpy{yu4}}

\begin{EntryWithPhonetic}{与}{yu4}{3}{⼀}
  \definition{v.}{participar de; tomar parte em}
  \seeref{yu2}
  \seeref{yu3}
  \synonymref{参加}{can1jia1}
  \synonymref{出席}{chu1/xi2}
  \synonymref{加入}{jia1ru4}
  \synonymref{介入}{jie4ru4}
  \antonymref{旁观}{pang2guan1}
\end{EntryWithPhonetic}

%%%%%%%%%% 玉 %%%%%%%%%%
\subsection*{玉}\addcontentsline{loh}{figure}{玉 \dpy{yu4}}

\begin{EntryWithPhonetic}{玉}{yu4}{5}{⽟}[HSK 4][Kangxi 96]
  \definition*{s.}{Sobrenome: Yu}
  \definition{adj.}{(pessoa, especialmente uma mulher) pura; justa; bonita; bela | cristalino, branco e belo como o jade | (vida) rica; luxuosa}
  \definition{pron.}{seu; um termo de respeito, usado para honrar o corpo, as ações ou as coisas associadas à outra pessoa}
  \definition[块,种]{s.}{jade}
\end{EntryWithPhonetic}

\begin{EntryWithPhonetic}{玉帛}{yu4bo2}{5,8}{⽟,⼱}
  \definition{s.}{objetos de jade e tecidos de seda, usados como presentes de estado; paz e harmonia}[化干戈为玉帛。===Transforme guerra em paz.]
  \antonymref{干戈}{gan1ge1}
\end{EntryWithPhonetic}

\begin{EntryWithPhonetic}{玉米}{yu4mi3}{5,6}{⽟,⽶}[HSK 4]
  \definition[根,粒,棵,片]{s.}{milho}
\end{EntryWithPhonetic}

\begin{EntryWithPhonetic}{玉米饼}{yu4mi3bing3}{5,6,9}{⽟,⽶,⾷}
  \definition{s.}{tortilha mexicana | bolo de milho}
\end{EntryWithPhonetic}

\begin{EntryWithPhonetic}{玉米粉}{yu4mi3fen3}{5,6,10}{⽟,⽶,⽶}
  \definition{s.}{amido de milho | farinha de milho}
\end{EntryWithPhonetic}

\begin{EntryWithPhonetic}{玉米糕}{yu4mi3gao1}{5,6,16}{⽟,⽶,⽶}
  \definition{s.}{bolo de milho | polenta}
\end{EntryWithPhonetic}

\begin{EntryWithPhonetic}{玉米花}{yu4mi3hua1}{5,6,7}{⽟,⽶,⾋}
  \definition{s.}{pipoca}
\end{EntryWithPhonetic}

\begin{EntryWithPhonetic}{玉米面}{yu4mi3mian4}{5,6,9}{⽟,⽶,⾯}
  \definition{s.}{fubá | farinha de milho}
\end{EntryWithPhonetic}

\begin{EntryWithPhonetic}{玉米片}{yu4mi3pian4}{5,6,4}{⽟,⽶,⽚}
  \definition{s.}{flocos de milho | chips de tortilha}
\end{EntryWithPhonetic}

\begin{EntryWithPhonetic}{玉米糁}{yu4mi3 san3}{5,6,14}{⽟,⽶,⽶}
  \definition{s.}{grãos de milho}
\end{EntryWithPhonetic}

\begin{EntryWithPhonetic}{玉米笋}{yu4mi3 sun3}{5,6,10}{⽟,⽶,⽵}
  \definition{s.}{broto de milho}
\end{EntryWithPhonetic}

%%%%%%%%%% 芋 %%%%%%%%%%
\subsection*{芋}\addcontentsline{loh}{figure}{芋 \dpy{yu4}}

\begin{EntryWithPhonetic}{芋}{yu4}{6}{⾋}
  \definition*{s.}{Sobrenome: Yu}
  \definition{s.}{taro; erva perene | tubérculos; geralmente se refere a batatas, etc.}
\end{EntryWithPhonetic}

\begin{EntryWithPhonetic}{芋头}{yu4tou5}{6,5}{⾋,⼤}
  \definition{s.}{taro, similar ao inhame e batata doce}
\end{EntryWithPhonetic}

\begin{EntryWithPhonetic}{芋头色}{yu4tou5se4}{6,5,6}{⾋,⼤,⾊}
  \definition{s.}{cor lilás}
\end{EntryWithPhonetic}

%%%%%%%%%% 郁 %%%%%%%%%%
\subsection*{郁}\addcontentsline{loh}{figure}{郁 \dpy{yu4}}

\begin{EntryWithPhonetic}{郁}{yu4}{8}{⾢}
  \definition*{s.}{Sobrenome: Yu}
  \definition{adj.}{fortemente perfumado | luxuriante; exuberante | sombrio; deprimido}
\end{EntryWithPhonetic}

\begin{EntryWithPhonetic}{郁郁葱葱}{yu4yu4cong1cong1}{8,8,12,12}{⾢,⾢,⾋,⾋}
  \definition{adj.}{exuberante e verde}
  \definition{expr.}{verdejante e exuberante; uma profusão selvagem de vegetação; luxuriantemente verde; ela cresce mais verde e mais fresca}
\end{EntryWithPhonetic}

%%%%%%%%%% 雨 %%%%%%%%%%
\subsection*{雨}\addcontentsline{loh}{figure}{雨 \dpy{yu4}}

\begin{EntryWithPhonetic}{雨}{yu4}{8}{⾬}[Kangxi 173]
  \definition{v.}{cair (chuva, neve, etc.) | precipitar | chover | molhar}
  \seeref{yu3}
\end{EntryWithPhonetic}

%%%%%%%%%% 语 %%%%%%%%%%
\subsection*{语}\addcontentsline{loh}{figure}{语 \dpy{yu4}}

\begin{EntryWithPhonetic}{语}{yu4}{9}{⾔}
  \definition{v.}{contar; informar}
  \seeref{yu3}
\end{EntryWithPhonetic}

%%%%%%%%%% 浴 %%%%%%%%%%
\subsection*{浴}\addcontentsline{loh}{figure}{浴 \dpy{yu4}}

\begin{EntryWithPhonetic}{浴}{yu4}{10}{⽔}
  \definition{v.}{tomar banho; banhar"-se}
\end{EntryWithPhonetic}

\begin{EntryWithPhonetic}{浴室}{yu4shi4}{10,9}{⽔,⼧}[HSK 7-9]
  \definition[间]{s.}{banho; banheiro; lavabo; sala de banho}
\end{EntryWithPhonetic}

%%%%%%%%%% 预 %%%%%%%%%%
\subsection*{预}\addcontentsline{loh}{figure}{预 \dpy{yu4}}

\begin{EntryWithPhonetic}{预}{yu4}{10}{⾴}
  \definition{adv.}{antecipadamente}
  \definition{v.}{avançar | preparar}
\end{EntryWithPhonetic}

\begin{EntryWithPhonetic}{预报}{yu4bao4}{10,7}{⾴,⼿}[HSK 3]
  \definition[个,项]{s.}{boletim meteorológico; previsões meteorológicas antecipadas}
  \definition{v.}{prever (o tempo); relatar antes que algo aconteça, usado principalmente em relação ao clima, astronomia, desastres naturais, etc.}
\end{EntryWithPhonetic}

\begin{EntryWithPhonetic}{预备}{yu4bei4}{10,8}{⾴,⼡}[HSK 5]
  \definition{v.}{preparar"-se; ficar pronto}
\end{EntryWithPhonetic}

\begin{EntryWithPhonetic}{预测}{yu4ce4}{10,9}{⾴,⽔}[HSK 4]
  \definition{v.}{prever; prognosticar; predizer}
\end{EntryWithPhonetic}

\begin{EntryWithPhonetic}{预订}{yu4ding4}{10,4}{⾴,⾔}[HSK 4]
  \definition{v.}{reservar; fazer uma reserva}
\end{EntryWithPhonetic}

\begin{EntryWithPhonetic}{预定}{yu4ding4}{10,8}{⾴,⼧}[HSK 7-9]
  \definition{v.}{definir com antecedência; predeterminar; agendar; agendar com antecedência}
  \synonymref{预订}{yu4ding4}
  \synonymref{预留}{yu4liu2}
  \synonymref{预约}{yu4yue1}
\end{EntryWithPhonetic}

\begin{EntryWithPhonetic}{预防}{yu4fang2}{10,6}{⾴,⾩}[HSK 3]
  \definition{v.}{prevenir; proteger"-se contra; tomar precauções contra; preparar"-se com antecedência para evitar que algo ruim aconteça}
\end{EntryWithPhonetic}

\begin{EntryWithPhonetic}{预付}{yu4fu4}{10,5}{⾴,⼈}
  \definition{s.}{pré-pago}
  \definition{v.}{pagar antecipadamente}
\end{EntryWithPhonetic}

\begin{EntryWithPhonetic}{预感}{yu4gan3}{10,13}{⾴,⼼}[HSK 7-9]
  \definition[种,个]{s.}{premonição; pressentimento}
  \definition{v.}{ter uma premonição; ter pressentimentos; pressagiar}
  \synonymref{预见}{yu4jian4}
  \synonymref{预料}{yu4liao4}
\end{EntryWithPhonetic}

\begin{EntryWithPhonetic}{预告}{yu4gao4}{10,7}{⾴,⼝}[HSK 7-9]
  \definition{s.}{anúncio prévio; pré"-visualização}
  \definition{v.}{prever; anunciar; anunciar com antecedência}
  \synonymref{预报}{yu4bao4}
  \synonymref{预言}{yu4yan2}
\end{EntryWithPhonetic}

\begin{EntryWithPhonetic}{预购}{yu4gou4}{10,8}{⾴,⾙}
  \definition{s.}{compra antecipada}
  \definition{v.}{comprar antecipadamente}
\end{EntryWithPhonetic}

\begin{EntryWithPhonetic}{预计}{yu4ji4}{10,4}{⾴,⾔}[HSK 3]
  \definition{v.}{estimar; calcular com antecedência}
\end{EntryWithPhonetic}

\begin{EntryWithPhonetic}{预见}{yu4jian4}{10,4}{⾴,⾒}[HSK 7-9]
  \definition{s.}{previsão; perspicácia; análise ou julgamento feito antecipadamente sobre situações futuras}
  \definition{v.}{prever; antecipar; analisar ou prever situações futuras com base nas leis que regem o desenvolvimento das coisas}
  \synonymref{意料}{yi4liao4}
  \synonymref{预感}{yu4gan3}
  \synonymref{预料}{yu4liao4}
\end{EntryWithPhonetic}

\begin{EntryWithPhonetic}{预警}{yu4jing3}{10,19}{⾴,⾔}
  \definition{s.}{aviso | aviso antecipado}
\end{EntryWithPhonetic}

\begin{EntryWithPhonetic}{预览}{yu4lan3}{10,9}{⾴,⾒}
  \definition{s.}{visualização}
  \definition{v.}{visualizar}
\end{EntryWithPhonetic}

\begin{EntryWithPhonetic}{预料}{yu4liao4}{10,10}{⾴,⽃}[HSK 7-9]
  \definition{s.}{expectativa; previsão; antecipação; especulação anterior}
  \definition{v.}{esperar; prever; antecipar}
  \synonymref{意料}{yi4liao4}
  \synonymref{预测}{yu4ce4}
  \synonymref{预感}{yu4gan3}
  \synonymref{预见}{yu4jian4}
  \synonymref{预期}{yu4qi1}
  \antonymref{意外}{yi4wai4}
\end{EntryWithPhonetic}

\begin{EntryWithPhonetic}{预留}{yu4liu2}{10,10}{⾴,⽥}
  \definition{v.}{separar | reservar}
\end{EntryWithPhonetic}

\begin{EntryWithPhonetic}{预谋}{yu4mou2}{10,11}{⾴,⾔}
  \definition{adj.}{premeditado}
  \definition{v.}{planejar algo com antecedência (especialmente um crime)}
\end{EntryWithPhonetic}

\begin{EntryWithPhonetic}{预判}{yu4pan4}{10,7}{⾴,⼑}
  \definition{v.}{prever | antecipar}
\end{EntryWithPhonetic}

\begin{EntryWithPhonetic}{预配}{yu4pei4}{10,10}{⾴,⾣}
  \definition{s.}{pré-alocado | pré-cabeado}
  \definition{v.}{pré-alocar | pré-cabear}
\end{EntryWithPhonetic}

\begin{EntryWithPhonetic}{预期}{yu4qi1}{10,12}{⾴,⽉}[HSK 5]
  \definition{v.}{esperar; antecipar; imaginar; antecipar com expectativa}
\end{EntryWithPhonetic}

\begin{EntryWithPhonetic}{预赛}{yu4sai4}{10,14}{⾴,⾙}[HSK 7-9]
  \definition{s.}{Esporte: competição preliminar; eliminatórias; preliminar; partida de teste}
  \antonymref{决赛}{jue2sai4}
\end{EntryWithPhonetic}

\begin{EntryWithPhonetic}{预示}{yu4shi4}{10,5}{⾴,⽰}[HSK 7-9]
  \definition{v.}{pressagiar; prenunciar | prognosticar; augurar}
  \synonymref{预警}{yu4jing3}
  \synonymref{预兆}{yu4zhao4}
  \antonymref{保密}{bao3/mi4}
\end{EntryWithPhonetic}

\begin{EntryWithPhonetic}{预售}{yu4shou4}{10,11}{⾴,⼝}[HSK 7-9]
  \definition{v.}{vender antecipadamente; reservar com antecedência; abrir para reservas; fazer reservas antecipadas}
\end{EntryWithPhonetic}

\begin{EntryWithPhonetic}{预算}{yu4suan4}{10,14}{⾴,⽵}[HSK 7-9]
  \definition[笔]{s.}{orçamento; \emph{budget}; planos de receitas e despesas para um determinado período de tempo, elaborados por agências governamentais, organizações e instituições públicas}
  \definition{v.}{elaborar um orçamento; preparar um orçamento}
  \synonymref{估算}{gu1suan4}
\end{EntryWithPhonetic}

\begin{EntryWithPhonetic}{预提}{yu4ti2}{10,12}{⾴,⼿}
  \definition{s.}{retenção}
  \definition{v.}{reter (imposto)}
\end{EntryWithPhonetic}

\begin{EntryWithPhonetic}{预习}{yu4xi2}{10,3}{⾴,⼄}[HSK 3]
  \definition{v.}{pré-visualizar; preparar uma lição; estudar antecipadamente as matérias que serão abordadas nas aulas}
\end{EntryWithPhonetic}

\begin{EntryWithPhonetic}{预先}{yu4xian1}{10,6}{⾴,⼉}[HSK 7-9]
  \definition{adv.}{antecipadamente; com antecedência; antes que as coisas aconteçam ou prossigam}
  \seealsoref{事先}{shi4xian1}
  \synonymref{事先}{shi4xian1}
  \antonymref{事后}{shi4hou4}
\end{EntryWithPhonetic}

\begin{EntryWithPhonetic}{预言}{yu4yan2}{10,7}{⾴,⾔}[HSK 7-9]
  \definition[个]{s.}{previsão; profecia; declarações pré"-combinadas sobre o que acontecerá no futuro}
  \definition{v.}{profetizar; predizer; pressagiar; prever (o que acontecerá no futuro)}
\end{EntryWithPhonetic}

\begin{EntryWithPhonetic}{预约}{yu4yue1}{10,6}{⾴,⽷}[HSK 6]
  \definition[个]{s.}{reserva}
  \definition{v.}{reservar; agendar; marcar compromisso; marcar uma consulta}
\end{EntryWithPhonetic}

\begin{EntryWithPhonetic}{预兆}{yu4zhao4}{10,6}{⾴,⼉}[HSK 7-9]
  \definition[个,种]{s.}{sinal; presságio; prenúncio; arauto; sinais pré"-existentes}
  \definition{v.}{pressagiar; ser um presságio; (algum sinal) prenunciar algo que está prestes a acontecer}
\end{EntryWithPhonetic}

\begin{EntryWithPhonetic}{预祝}{yu4zhu4}{10,9}{⾴,⽰}
  \definition{v.}{parabenizar de antemão | oferecer os melhores votos para}
\end{EntryWithPhonetic}

%%%%%%%%%% 欲 %%%%%%%%%%
\subsection*{欲}\addcontentsline{loh}{figure}{欲 \dpy{yu4}}

\begin{EntryWithPhonetic}{欲}{yu4}{11}{⽋}
  \definition{adj.}{desejo; anseio; vontade}
  \definition{adv.}{em breve; prestes a}
  \definition{v.}{desejar; querer; ansiar | precisar; querer; exigir}
\end{EntryWithPhonetic}

\begin{EntryWithPhonetic}{欲望}{yu4wang4}{11,11}{⽋,⽉}[HSK 7-9]
  \definition{v.}{desejar; cobiçar; ansiar; o desejo de obter algo ou alcançar um determinado objetivo}
  \synonymref{渴望}{ke3wang4}
  \synonymref{愿望}{yuan4wang4}
\end{EntryWithPhonetic}

%%%%%%%%%% 喻 %%%%%%%%%%
\subsection*{喻}\addcontentsline{loh}{figure}{喻 \dpy{yu4}}

\begin{EntryWithPhonetic}{喻}{yu4}{12}{⼝}
  \definition{s.}{analogia | símile | metáfora | alegoria}
  \definition{v.}{descrever algo como}
\end{EntryWithPhonetic}

%%%%%%%%%% 寓 %%%%%%%%%%
\subsection*{寓}\addcontentsline{loh}{figure}{寓 \dpy{yu4}}

\begin{EntryWithPhonetic}{寓}{yu4}{12}{⼧}
  \definition[座,间,栋]{s.}{residência; morada}
  \definition{v.}{(literário) residir; viver | implicar; conter}
\end{EntryWithPhonetic}

\begin{EntryWithPhonetic}{寓言}{yu4yan2}{12,7}{⼧,⾔}[HSK 7-9]
  \definition[则,个,篇]{s.}{fábula; alegoria; parábola; uma forma literária que utiliza histórias fictícias para ilustrar um princípio, servindo a um propósito educativo ou satírico}
  \synonymref{故事}{gu4shi5}
  \synonymref{童话}{tong2hua4}
\end{EntryWithPhonetic}

\begin{EntryWithPhonetic}{寓意}{yu4yi4}{12,13}{⼧,⼼}[HSK 7-9]
  \definition{s.}{moral; importância; mensagem; significado implícito ou subentendido}
\end{EntryWithPhonetic}

%%%%%%%%%% 粥 %%%%%%%%%%
\subsection*{粥}\addcontentsline{loh}{figure}{粥 \dpy{yu4}}

\begin{EntryWithPhonetic}{粥}{yu4}{12}{⽶}
  \definition{v.}{dar a luz; ter filhos}
  \seeref{zhou1}
\end{EntryWithPhonetic}

%%%%%%%%%% 遇 %%%%%%%%%%
\subsection*{遇}\addcontentsline{loh}{figure}{遇 \dpy{yu4}}

\begin{EntryWithPhonetic}{遇}{yu4}{12}{⾡}[HSK 4]
  \definition*{s.}{Sobrenome: Yu}
  \definition{s.}{chance; oportunidade}
  \definition{v.}{encontrar; deparar"-se com; encontrar"-se | tratar; receber}
\end{EntryWithPhonetic}

\begin{EntryWithPhonetic}{遇到}{yu4/dao4}{12,8}{⾡,⼑}[HSK 4]
  \definition{v.+compl.}{esbarrar em; encontrar; deparar"-se com; conhecer alguém ou algo (inesperado)}
\end{EntryWithPhonetic}

\begin{EntryWithPhonetic}{遇见}{yu4/jian5}{12,4}{⾡,⾒}[HSK 4]
  \definition{v.+compl.}{encontrar; deparar"-se com}
\end{EntryWithPhonetic}

\begin{EntryWithPhonetic}{遇难}{yu4/nan4}{12,10}{⾡,⾫}[HSK 7-9]
  \definition{v.+compl.}{morrer em um acidente | ser assassinado | enfrentar perigo; estar em apuros | perecer | ser morto}
  \synonymref{遇险}{yu4/xian3}
  \synonymref{遭殃}{zao1/yang1}
  \antonymref{幸存}{xing4cun2}
\end{EntryWithPhonetic}

\begin{EntryWithPhonetic}{遇上}{yu4shang4}{12,3}{⾡,⼀}[HSK 7-9]
  \definition{v.}{encontrar alguém; esbarrar em alguém}
\end{EntryWithPhonetic}

\begin{EntryWithPhonetic}{遇险}{yu4/xian3}{12,9}{⾡,⾩}[HSK 7-9]
  \definition{v.+compl.}{encontrar"-se com um contratempo; encontrar"-se com perigo; estar em perigo (ou em apuros)}
  \synonymref{不幸}{bu2xing4}
  \synonymref{探险}{tan4/xian3}
  \synonymref{遇难}{yu4/nan4}
  \antonymref{脱险}{tuo1xian3}
\end{EntryWithPhonetic}

%%%%%%%%%% 愈 %%%%%%%%%%
\subsection*{愈}\addcontentsline{loh}{figure}{愈 \dpy{yu4}}

\begin{EntryWithPhonetic}{愈}{yu4}{13}{⼼}
  \definition{v.}{curar; recuperar; ficar bem | Literário: ser melhor que; destacar"-se}
  \definition{v.}{(em 愈\dots 愈\dots) quanto mais\dots quanto mais\dots | mais; cada vez mais; a repetição indica que o grau se desenvolve à medida que as condições se desenvolvem}
  \seealsoref{愈……愈……}{yu4 yu4}
\end{EntryWithPhonetic}

\begin{EntryWithPhonetic}{愈合}{yu4he2}{13,6}{⼼,⼝}[HSK 7-9]
  \definition{v.}{curar (feridas ou lesões cicatrizadas)}
  \synonymref{治愈}{zhi4yu4}
  \antonymref{创伤}{chuang1shang1}
  \antonymref{分裂}{fen1lie4}
  \antonymref{伤痕}{shang1hen2}
\end{EntryWithPhonetic}

\begin{EntryWithPhonetic}{愈来愈}{yu4 lai2 yu4}{13,7,13}{⼼,⽊,⼼}[HSK 7-9]
  \definition{expr.}{cada vez mais; mais e mais}
\end{EntryWithPhonetic}

\begin{EntryWithPhonetic}{愈演愈烈}{yu4yan3-yu4lie4}{13,14,13,10}{⼼,⽔,⼼,⽕}[HSK 7-9]
  \definition{expr.}{escalando; aumentar em intensidade; tornar"-se cada vez mais intenso; piorar; cada vez mais crítico; os problemas se tornam cada vez mais intensos}
\end{EntryWithPhonetic}

\begin{EntryWithPhonetic}{愈……愈……}{yu4 yu4}{13,13}{⼼,⼼}
  \definition{conj.}{quanto mais\dots quanto mais\dots}
  \seealsoref{愈}{yu4}
\end{EntryWithPhonetic}

%%%%%%%%%% 豫 %%%%%%%%%%
\subsection*{豫}\addcontentsline{loh}{figure}{豫 \dpy{yu4}}

\begin{EntryWithPhonetic}{豫}{yu4}{15}{⾗}
  \definition*{s.}{Província de Henan, abreviatura de 河南}
  \definition{adj.}{satisfeito; encantado | anterior; preliminar; preparatório}
  \definition{adv.}{com antecedência; antecipadamente}
  \definition{v.}{viver com facilidade e conforto | participar de}
  \seealsoref{河南}{he2nan2}
  \seealsoref{预}{yu4}
\end{EntryWithPhonetic}

%%%%%%%%%% 冤 %%%%%%%%%%
\subsection*{冤}\addcontentsline{loh}{figure}{冤 \dpy{yuan1}}

\begin{EntryWithPhonetic}{冤}{yuan1}{10}{⼍}[HSK 7-9]
  \definition{adj.}{que não vale a pena; inútil; um desperdício (de tempo/dinheiro/esforço) ; custou tempo e dinheiro, mas não atingiu o objetivo; não valeu a pena | sentir"-se injustiçado; sentir ressentimento; isso descreve a sensação de ser criticado ou punido por algo que você não fez; sentir que é injusto}
  \definition{s.}{tratamento injusto; acusação injusta; injustiça; erro judiciário; condenação injusta | rancor; ressentimento; ódio; rixa; inimizade}
  \definition{v.}{enganar; ludibriar; iludir | ser falsamente acusado; ser injustiçado; ser tratado injustamente ou causar injustiça a alguém}
\end{EntryWithPhonetic}

\begin{EntryWithPhonetic}{冤枉}{yuan1wang5}{10,8}{⼍,⽊}[HSK 7-9]
  \definition{adj.}{injusto; injustificado; descreve a sensação de ser criticado injustamente quando você não fez nada de errado | indigno; não merecido; descreve o uso de tempo, dinheiro, etc., como algo que não vale a pena}
  \definition{v.}{prejudicar; tratar injustamente; ser acusado falsamente}
  \synonymref{勉强}{mian3qiang3}
  \synonymref{曲折}{qu1zhe2}
  \synonymref{委屈}{wei3qu5}
  \synonymref{无辜}{wu2gu1}
  \antonymref{美化}{mei3hua4}
\end{EntryWithPhonetic}

%%%%%%%%%% 渊 %%%%%%%%%%
\subsection*{渊}\addcontentsline{loh}{figure}{渊 \dpy{yuan1}}

\begin{EntryWithPhonetic}{渊}{yuan1}{11}{⽔}
  \definition*{s.}{Sobrenome: Yuan}
  \definition{adj.}{profundo}
  \definition{s.}{piscina profunda}
\end{EntryWithPhonetic}

\begin{EntryWithPhonetic}{渊源}{yuan1yuan2}{11,13}{⽔,⽔}[HSK 7-9]
  \definition{s.}{uma origem; uma fonte; a nascente de um riacho; uma metáfora para a origem ou fonte de algo}
\end{EntryWithPhonetic}

%%%%%%%%%% 元 %%%%%%%%%%
\subsection*{元}\addcontentsline{loh}{figure}{元 \dpy{yuan2}}

\begin{EntryWithPhonetic}{元}{yuan2}{4}{⼉}[HSK 1]
  \definition*{s.}{Dinastia Yuan (1271--1368) | Sobrenome: Yuan}
  \definition{adj.}{primeiro; principal; primário | chefe; diretor; líder | básico; fundamental; principal | unidade; componente; formando um todo}
  \definition{clas.}{yuan, a unidade monetária da China}
  \definition{s.}{moeda com valor e peso fixos | origem; elemento}
\end{EntryWithPhonetic}

\begin{EntryWithPhonetic}{元旦}{yuan2dan4}{4,5}{⼉,⽇}[HSK 5]
  \definition*{s.}{Dia de Ano Novo (1 de janeiro)}
\end{EntryWithPhonetic}

\begin{EntryWithPhonetic}{元来}{yuan2lai2}{4,7}{⼉,⽊}
  \variantof{原来}
\end{EntryWithPhonetic}

\begin{EntryWithPhonetic}{元老}{yuan2lao3}{4,6}{⼉,⽼}[HSK 7-9]
  \definition{s.}{estadista sênior; membro fundador | membro fundador (de uma organização política, etc.); patriarca; veterano; pessoas com grande experiência e prestígio na arena política}
\end{EntryWithPhonetic}

\begin{EntryWithPhonetic}{元气}{yuan2qi4}{4,4}{⼉,⽓}
  \definition{s.}{força | vigor | vitalidade | energial vital}
\end{EntryWithPhonetic}

\begin{EntryWithPhonetic}{元首}{yuan2shou3}{4,9}{⼉,⾸}[HSK 7-9]
  \definition[位,名]{s.}{chefe de estado; refere"-se ao líder máximo de um país}
  \synonymref{领导}{ling3dao3}
  \synonymref{领袖}{ling3xiu4}
  \synonymref{首脑}{shou3nao3}
  \synonymref{率领}{shuai4ling3}
  \synonymref{指导}{zhi3dao3}
  \synonymref{指挥}{zhi3hui1}
  \synonymref{总统}{zong3tong3}
\end{EntryWithPhonetic}

\begin{EntryWithPhonetic}{元素}{yuan2su4}{4,10}{⼉,⽷}[HSK 6]
  \definition{s.}{elemento; fator essencial | elemento químico | Geometria: as partes que compõem uma figura, como os lados e ângulos de um triângulo | Algebra: os elementos algébricos incluem números e símbolos}
\end{EntryWithPhonetic}

\begin{EntryWithPhonetic}{元宵}{yuan2xiao1}{4,10}{⼉,⼧}
  \definition*{s.}{Festival das Lanternas}
  \seealsoref{元宵节}{yuan2xiao1jie2}
  \seealsoref{元夜}{yuan2ye4}
\end{EntryWithPhonetic}

\begin{EntryWithPhonetic}{元宵节}{yuan2xiao1jie2}{4,10,5}{⼉,⼧,⾋}[HSK 7-9]
  \definition*{s.}{O Festival das Lanternas, também conhecido como Festival Shangyuan, 上元节, é um festival tradicional chinês celebrado no 15º~dia do primeiro mês lunar; desde a Dinastia Tang, existe o costume de contemplar lanternas nesta noite}
  \seealsoref{元宵}{yuan2xiao1}
  \seealsoref{元夜}{yuan2ye4}
  \synonymref{上元节}{shang4yuan2jie2}
\end{EntryWithPhonetic}

\begin{EntryWithPhonetic}{元夜}{yuan2ye4}{4,8}{⼉,⼣}
  \definition*{s.}{Festival das Lanternas}
  \seealsoref{元宵}{yuan2xiao1}
  \seealsoref{元宵节}{yuan2xiao1jie2}
\end{EntryWithPhonetic}

%%%%%%%%%% 员 %%%%%%%%%%
\subsection*{员}\addcontentsline{loh}{figure}{员 \dpy{yuan2}}

\begin{EntryWithPhonetic}{员}{yuan2}{7}{⼝}[HSK 3]
  \definition{clas.}{para comandantes militares}
  \definition{s.}{uma pessoa envolvida em algum campo de atividade; refere"-se a pessoas que trabalham ou estudam | membro; refere"-se aos membros de um grupo ou organização | vizinhança}
\end{EntryWithPhonetic}

\begin{EntryWithPhonetic}{员工}{yuan2gong1}{7,3}{⼝,⼯}[HSK 3]
  \definition[位,名,个]{s.}{equipe; funcionário; trabalhador; pessoal}
\end{EntryWithPhonetic}

%%%%%%%%%% 园 %%%%%%%%%%
\subsection*{园}\addcontentsline{loh}{figure}{园 \dpy{yuan2}}

\begin{EntryWithPhonetic}{园}{yuan2}{7}{⼞}[HSK 6]
  \definition*{s.}{Sobrenome: Yuan}
  \definition[个]{s.}{jardim; terreno; plantação; terra para cultivar plantas | local para recreação pública; locais para passeios turísticos e entretenimento | área para fins especiais | uma área de terra para o cultivo de plantas; um lugar onde vegetais, flores, frutas e árvores são cultivados}
\end{EntryWithPhonetic}

\begin{EntryWithPhonetic}{园地}{yuan2di4}{7,6}{⼞,⼟}[HSK 6]
  \definition{s.}{jardim; campo | Figurativo: campo; escopo}
\end{EntryWithPhonetic}

\begin{EntryWithPhonetic}{园林}{yuan2lin2}{7,8}{⼞,⽊}[HSK 5]
  \definition[处,座,个]{s.}{parque; jardim; área paisagística com plantas e árvores para as pessoas apreciarem e descansarem.}
\end{EntryWithPhonetic}

%%%%%%%%%% 原 %%%%%%%%%%
\subsection*{原}\addcontentsline{loh}{figure}{原 \dpy{yuan2}}

\begin{EntryWithPhonetic}{原}{yuan2}{10}{⼚}[HSK 6]
  \definition*{s.}{Sobrenome: Yuan}
  \definition{adj.}{inicial; básico; primitivo | cru; bruto; não processado | virgem; primário; original; antigo; inalterado}
  \definition{adv.}{originalmente}
  \definition[项,条,片]{s.}{planície; país aberto; terreno plano e amplo | início; fonte; origem; aparência original | origem; a raiz ou o começo das coisas}
  \definition{v.}{desculpar; perdoar; tolerar; compreender | rastrear; sondar; investigar (a origem das coisas)}
\end{EntryWithPhonetic}

\begin{EntryWithPhonetic}{原本}{yuan2ben3}{10,5}{⼚,⽊}[HSK 7-9]
  \definition{adv.}{originalmente; anteriormente}
  \definition{s.}{cópia mestra; manuscrito original; manuscrito | a primeira edição | o original; o livro original no qual a tradução se baseia}
  \synonymref{本来}{ben3lai2}
  \synonymref{来自}{lai2zi4}
  \synonymref{其实}{qi2shi2}
  \synonymref{原来}{yuan2lai2}
  \synonymref{原先}{yuan2xian1}
  \synonymref{自从}{zi4cong2}
\end{EntryWithPhonetic}

\begin{EntryWithPhonetic}{原材料}{yuan2cai2liao4}{10,7,10}{⼚,⽊,⽃}[HSK 7-9]
  \definition[种]{s.}{matéria"-prima e matéria"-prima processada}
\end{EntryWithPhonetic}

\begin{EntryWithPhonetic}{原创}{yuan2chuang4}{10,6}{⼚,⼑}[HSK 7-9]
  \definition{s.}{criador; originador; refere"-se ao criador original}
  \definition{v.}{criar; iniciar; ser o original}
  \antonymref{抄袭}{chao1xi2}
  \antonymref{借鉴}{jie4jian4}
  \antonymref{模仿}{mo2fang3}
\end{EntryWithPhonetic}

\begin{EntryWithPhonetic}{原地}{yuan2di4}{10,6}{⼚,⼟}[HSK 7-9]
  \definition{s.}{(no) lugar original; o local original; o ponto de partida}
\end{EntryWithPhonetic}

\begin{EntryWithPhonetic}{原告}{yuan2gao4}{10,7}{⼚,⼝}[HSK 6]
  \definition{s.}{(em casos civis) autor; solicitante | (em casos criminais) promotor; acusador; reclamante}
  \antonymref{被告}{bei4gao4}
\end{EntryWithPhonetic}

\begin{EntryWithPhonetic}{原来}{yuan2lai2}{10,7}{⼚,⽊}[HSK 2]
  \definition{adj.}{original; anterior; em primeiro lugar; inicialmente; inalterado}
  \definition{adv.}{na verdade; de fato; como se vê; expressar compreensão repentina}
  \definition{s.}{a princípio; no passado; antigamente}
\end{EntryWithPhonetic}

\begin{EntryWithPhonetic}{原理}{yuan2li3}{10,11}{⼚,⽟}[HSK 5]
  \definition[个,条]{s.}{princípio; axioma; teoria; teoria básica ou princípio científico de significado universal}
\end{EntryWithPhonetic}

\begin{EntryWithPhonetic}{原谅}{yuan2liang4}{10,10}{⼚,⾔}[HSK 6]
  \definition{v.}{perdoar; perdoar a negligência, os erros ou as falhas das pessoas sem culpá"-las ou puni"-las}
\end{EntryWithPhonetic}

\begin{EntryWithPhonetic}{原料}{yuan2liao4}{10,10}{⼚,⽃}[HSK 4]
  \definition[种,个]{s.}{matéria"-prima; refere"-se a materiais que não foram processados e fabricados, como minérios para metalurgia e algodão para têxteis}
\end{EntryWithPhonetic}

\begin{EntryWithPhonetic}{原木}{yuan2mu4}{10,4}{⼚,⽊}
  \definition{s.}{registro | \emph{logs}}
\end{EntryWithPhonetic}

\begin{EntryWithPhonetic}{原色}{yuan2 se4}{10,6}{⼚,⾊}
  \definition{s.}{cor primária}
\end{EntryWithPhonetic}

\begin{EntryWithPhonetic}{原始}{yuan2shi3}{10,8}{⼚,⼥}[HSK 5]
  \definition{s.}{original; de primeira mão | primitivo; mais antigo; não desenvolvido; não civilizado}
\end{EntryWithPhonetic}

\begin{EntryWithPhonetic}{原先}{yuan2xian1}{10,6}{⼚,⼉}[HSK 5]
  \definition{adj.}{antigo; original}
  \definition{s.}{antigamente; no início; no passado; no começo}
\end{EntryWithPhonetic}

\begin{EntryWithPhonetic}{原型}{yuan2xing2}{10,9}{⼚,⼟}[HSK 7-9]
  \definition{s.}{modelo; protótipo; arquétipo; molde original | proterótipo; protomodelo; o tipo ou modelo original refere"-se especificamente às pessoas da vida real nas quais os personagens das obras literárias narrativas são baseados}
  \synonymref{事实}{shi4shi2}
  \synonymref{真相}{zhen1xiang4}
  \antonymref{模型}{mo2xing2}
\end{EntryWithPhonetic}

\begin{EntryWithPhonetic}{原因}{yuan2yin1}{10,6}{⼚,⼞}[HSK 2]
  \definition[个,条,种,些]{s.}{causa; razão; motivo; as condições que fazem com que algo aconteça ou produzam um certo resultado}
\end{EntryWithPhonetic}

\begin{EntryWithPhonetic}{原有}{yuan2you3}{10,6}{⼚,⽉}[HSK 5]
  \definition{v.}{já estar pronto, não é necessário fazer ou procurar nada; ser o original}
\end{EntryWithPhonetic}

\begin{EntryWithPhonetic}{原则}{yuan2ze2}{10,6}{⼚,⼑}[HSK 4]
  \definition{adv.}{em geral; em princípio; refere"-se a um aspecto geral; geralmente}
  \definition[个,条,项,点]{s.}{princípios; leis ou padrões pelos quais alguém fala ou age}
\end{EntryWithPhonetic}

\begin{EntryWithPhonetic}{原汁原味}{yuan2zhi1-yuan2wei4}{10,5,10,8}{⼚,⽔,⼚,⼝}[HSK 7-9]
  \definition{expr.}{(alimentos) suco e sabor naturais; autêntico; genuíno; original}
  \synonymref{地道}{di4dao5}
\end{EntryWithPhonetic}

\begin{EntryWithPhonetic}{原装}{yuan2zhuang1}{10,12}{⼚,⾐}[HSK 7-9]
  \definition{adj.}{genuíno}
  \definition{s.}{embalado na fábrica (no país de origem) | embalagem original (ou encadernação); produto genuíno | intacto na embalagem original (não montado e embalado localmente)}
  \antonymref{改装}{gai3zhuang1}
\end{EntryWithPhonetic}

%%%%%%%%%% 圆 %%%%%%%%%%
\subsection*{圆}\addcontentsline{loh}{figure}{圆 \dpy{yuan2}}

\begin{EntryWithPhonetic}{圆}{yuan2}{10}{⼞}[HSK 4]
  \definition*{s.}{Sobrenome: Yuan}
  \definition{adj.}{redondo; circular; esférico; arredondado | diplomático; satisfatório}
  \definition[个,轮]{s.}{círculo; circunferência | uma moeda de valor e peso fixos}
  \definition{v.}{tornar plausível; justificar; tornar completo; completar}
\end{EntryWithPhonetic}

\begin{EntryWithPhonetic}{圆满}{yuan2man3}{10,13}{⼞,⽔}[HSK 4]
  \definition{adj.}{perfeito; satisfatório; sem defeitos}
\end{EntryWithPhonetic}

\begin{EntryWithPhonetic}{圆形}{yuan2xing2}{10,7}{⼞,⼺}[HSK 7-9]
  \definition{adj.}{circular; redondo}
  \definition{s.}{círculo; a figura formada pela circunferência de um círculo em um plano | com formato de bola ou hemisfério}
\end{EntryWithPhonetic}

\begin{EntryWithPhonetic}{圆珠笔}{yuan2zhu1bi3}{10,10,10}{⼞,⽟,⽵}[HSK 6]
  \definition[支,枝]{s.}{caneta esferográfica}
\end{EntryWithPhonetic}

%%%%%%%%%% 援 %%%%%%%%%%
\subsection*{援}\addcontentsline{loh}{figure}{援 \dpy{yuan2}}

\begin{EntryWithPhonetic}{援}{yuan2}{12}{⼿}
  \definition*{s.}{Sobrenome: Yuan}
  \definition{v.}{puxar com a mão; segurar | citar; referenciar | ajudar; auxiliar; resgatar}
\end{EntryWithPhonetic}

\begin{EntryWithPhonetic}{援助}{yuan2zhu4}{12,7}{⼿,⼒}[HSK 6]
  \definition{s.}{ajuda; assistência; auxílio}
  \definition{v.}{ajudar; apoiar; auxiliar}
\end{EntryWithPhonetic}

%%%%%%%%%% 缘 %%%%%%%%%%
\subsection*{缘}\addcontentsline{loh}{figure}{缘 \dpy{yuan2}}

\begin{EntryWithPhonetic}{缘}{yuan2}{12}{⽷}
  \definition{s.}{causa | razão | karma | destino | predestinação}
\end{EntryWithPhonetic}

\begin{EntryWithPhonetic}{缘分}{yuan2fen4}{12,4}{⽷,⼑}[HSK 7-9]
  \definition[段,份,次,场,种]{s.}{afinidade ou relacionamento predestinado; pessoas supersticiosas acreditam em encontros predestinados entre pessoas; em termos gerais, isso se refere à possibilidade de conexões entre pessoas ou entre pessoas e coisas}
  \antonymref{悲伤}{bei1shang1}
\end{EntryWithPhonetic}

\begin{EntryWithPhonetic}{缘故}{yuan2gu4}{12,9}{⽷,⽁}[HSK 6]
  \definition{s.}{causa; razão}
\end{EntryWithPhonetic}

%%%%%%%%%% 源 %%%%%%%%%%
\subsection*{源}\addcontentsline{loh}{figure}{源 \dpy{yuan2}}

\begin{EntryWithPhonetic}{源}{yuan2}{13}{⽔}
  \definition*{s.}{Sobrenome: Yuan}
  \definition{s.}{nascente (de um rio); fonte | fonte; origem; causa}
  \definition{v.}{originar"-se; provir de}
\end{EntryWithPhonetic}

\begin{EntryWithPhonetic}{源泉}{yuan2quan2}{13,9}{⽔,⽔}[HSK 7-9]
  \definition{s.}{fonte; nascente; o lugar onde um curso de água começa; metaforicamente, a fonte ou causa de força, conhecimento, emoções, etc}
  \synonymref{来源}{lai2yuan2}
\end{EntryWithPhonetic}

\begin{EntryWithPhonetic}{源头}{yuan2tou2}{13,5}{⽔,⼤}[HSK 7-9]
  \definition{s.}{cabeça; fonte; nascente; o lugar onde a água nasce}
\end{EntryWithPhonetic}

\begin{EntryWithPhonetic}{源于}{yuan2yu2}{13,3}{⽔,⼆}[HSK 7-9]
  \definition{v.}{vir de; resultar de; ter origem em; derivar de; ter sua origem em}
  \synonymref{出自}{chu1zi4}
  \synonymref{基于}{ji1yu2}
  \synonymref{来自}{lai2zi4}
  \synonymref{起源}{qi3yuan2}
  \synonymref{由于}{you2yu2}
  \synonymref{最初}{zui4chu1}
\end{EntryWithPhonetic}

\begin{EntryWithPhonetic}{源源不断}{yuan2yuan2-bu2duan4}{13,13,4,11}{⽔,⽔,⼀,⽄}[HSK 7-9]
  \definition{expr.}{um fluxo constante; uma correnteza interminável}
  \synonymref{川流不息}{chuan1liu2-bu4xi1}
  \synonymref{接二连三}{jie1'er4-lian2san1}
  \synonymref{络绎不绝}{luo4yi4-bu4jue2}
  \synonymref{滔滔不绝}{tao1tao1-bu4jue2}
  \antonymref{断断续续}{duan4duan4xu4xu4}
\end{EntryWithPhonetic}

%%%%%%%%%% 远 %%%%%%%%%%
\subsection*{远}\addcontentsline{loh}{figure}{远 \dpy{yuan3}}

\begin{EntryWithPhonetic}{远}{yuan3}{7}{⾡}[HSK 1]
  \definition*{s.}{Sobrenome: Yuan}
  \definition{adj.}{distante (no tempo ou no espaço); longe; remoto; Longa distância espacial ou temporal | (relações de parentesco) distante | com grande diferença}
  \definition{v.}{manter"-se afastado de; não se aproximar}
  \antonymref{近}{jin4}
\end{EntryWithPhonetic}

\begin{EntryWithPhonetic}{远程}{yuan3cheng2}{7,12}{⾡,⽲}[HSK 7-9]
  \definition{adj.}{de longo alcance; de longa distância | remoto; isso se refere a realizar uma determinada função ou atividade a longa distância}
  \definition{pref.}{tele-}
  \synonymref{长途}{chang2tu2}
  \synonymref{资料}{zi1liao4}
\end{EntryWithPhonetic}

\begin{EntryWithPhonetic}{远处}{yuan3chu4}{7,5}{⾡,⼡}[HSK 5]
  \definition{s.}{distância; lugar distante}
\end{EntryWithPhonetic}

\begin{EntryWithPhonetic}{远方}{yuan3fang1}{7,4}{⾡,⽅}[HSK 6]
  \definition{s.}{distância; de longe; lugar distante}
\end{EntryWithPhonetic}

\begin{EntryWithPhonetic}{远见}{yuan3jian4}{7,4}{⾡,⾒}[HSK 7-9]
  \definition{s.}{previsão; visão | visão de longo alcance}
\end{EntryWithPhonetic}

\begin{EntryWithPhonetic}{远近闻名}{yuan3jin4-wen2ming2}{7,7,9,6}{⾡,⾡,⾨,⼝}[HSK 7-9]
  \definition{expr.}{ser conhecido por toda parte; ser conhecido por todos, de perto e de longe; ser amplamente conhecido}
  \synonymref{大名鼎鼎}{da4ming2-ding3ding3}
  \synonymref{家喻户晓}{jia1yu4-hu4xiao3}
  \synonymref{举世闻名}{ju3shi4-wen2ming2}
  \antonymref{默默无闻}{mo4mo4-wu2wen2}
\end{EntryWithPhonetic}

\begin{EntryWithPhonetic}{远离}{yuan3li2}{7,10}{⾡,⼇}[HSK 6]
  \definition{adj.}{afastado; distante}
  \definition{v.}{partir para; ficar longe}
\end{EntryWithPhonetic}

\begin{EntryWithPhonetic}{远天}{yuan3tian1}{7,4}{⾡,⼤}
  \definition{s.}{paraíso | o céu distante}
\end{EntryWithPhonetic}

\begin{EntryWithPhonetic}{远远}{yuan3yuan3}{7,7}{⾡,⾡}[HSK 6]
  \definition{adv.}{de longe; em grande medida; para descrever um alto grau ou uma grande quantidade}
\end{EntryWithPhonetic}

\begin{EntryWithPhonetic}{远征}{yuan3zheng1}{7,8}{⾡,⼻}
  \definition{s.}{expedição; expedições de longa distância ou longas marchas}
  \definition{v.}{fazer uma expedição}
  \synonymref{长征}{chang2zheng1}
\end{EntryWithPhonetic}

%%%%%%%%%% 怨 %%%%%%%%%%
\subsection*{怨}\addcontentsline{loh}{figure}{怨 \dpy{yuan4}}

\begin{EntryWithPhonetic}{怨}{yuan4}{9}{⼼}[HSK 5]
  \definition{s.}{ressentimento; inimizade; rancor}
  \definition{v.}{culpar; reclamar}
\end{EntryWithPhonetic}

\begin{EntryWithPhonetic}{怨恨}{yuan4hen4}{9,9}{⼼,⼼}[HSK 7-9]
  \definition{s.}{ressentimento; forte insatisfação ou ódio; intensa insatisfação ou ódio}
  \definition{v.}{sentir ressentimento; guardar rancor de alguém; forte insatisfação ou ódio em relação a uma pessoa ou coisa}
  \synonymref{抱怨}{bao4yuan5}
  \synonymref{仇恨}{chou2hen4}
  \synonymref{后悔}{hou4hui3}
  \synonymref{悔恨}{hui3hen4}
  \synonymref{埋怨}{man2yuan4}
  \antonymref{恩惠}{en1hui4}
  \antonymref{感激}{gan3ji1}
  \antonymref{宽恕}{kuan1shu4}
\end{EntryWithPhonetic}

\begin{EntryWithPhonetic}{怨气}{yuan4qi4}{9,4}{⼼,⽓}[HSK 7-9]
  \definition{s.}{queixa; reclamação; ressentimento; uma expressão ou emoção de ressentimento}
  \antonymref{和气}{he2qi5}
\end{EntryWithPhonetic}

\begin{EntryWithPhonetic}{怨言}{yuan4yan2}{9,7}{⼼,⾔}[HSK 7-9]
  \definition{s.}{reclamação; resmungo}
  \synonymref{抱怨}{bao4yuan5}
  \antonymref{赞美}{zan4mei3}
\end{EntryWithPhonetic}

%%%%%%%%%% 院 %%%%%%%%%%
\subsection*{院}\addcontentsline{loh}{figure}{院 \dpy{yuan4}}

\begin{EntryWithPhonetic}{院}{yuan4}{9}{⾩}[HSK 2]
  \definition*{s.}{Sobrenome: Yuan}
  \definition[个]{s.}{pátio; quintal; complexo | designação para certos escritórios governamentais e locais públicos | faculdade; academia; instituto de ensino superior | hospital}
\end{EntryWithPhonetic}

\begin{EntryWithPhonetic}{院士}{yuan4shi4}{9,3}{⾩,⼠}[HSK 7-9]
  \definition{s.}{acadêmico | membro (de uma academia)}
\end{EntryWithPhonetic}

\begin{EntryWithPhonetic}{院长}{yuan4zhang3}{9,4}{⾩,⾧}[HSK 2]
  \definition[个,位,名]{s.}{reitor; diretor; o mais alto funcionário de qualquer instituição ou escola pública ou privada}
\end{EntryWithPhonetic}

\begin{EntryWithPhonetic}{院子}{yuan4zi5}{9,3}{⾩,⼦}[HSK 2]
  \definition[个,座,处]{s.}{quintal; pátio; o espaço aberto na frente ou atrás de uma casa cercado por muros ou cercas}
\end{EntryWithPhonetic}

%%%%%%%%%% 愿 %%%%%%%%%%
\subsection*{愿}\addcontentsline{loh}{figure}{愿 \dpy{yuan4}}

\begin{EntryWithPhonetic}{愿}{yuan4}{14}{⽕}[HSK 5]
  \definition{adj.}{honesto e prudente}
  \definition{s.}{esperança; desejo; vontade; a ideia de alcançar algum objetivo no futuro | voto (feito perante o Buda ou um deus); o desejo de retribuição feito ao rezar para os deuses e Buda}
  \definition{v.}{estar disposto; estar pronto; de bom grado, concordar porque está de acordo com seus desejos | ter esperança; desejar; qerer alcançar algum desejo}
\end{EntryWithPhonetic}

\begin{EntryWithPhonetic}{愿望}{yuan4wang4}{14,11}{⽕,⽉}[HSK 3]
  \definition[个,种]{s.}{desejo; aspiração; a ideia de alcançar algum objetivo no futuro.}
\end{EntryWithPhonetic}

\begin{EntryWithPhonetic}{愿意}{yuan4yi5}{14,13}{⽕,⼼}[HSK 2]
  \definition{v.}{estar disposto; estar pronto | desejar; ter esperança}
\end{EntryWithPhonetic}

%%%%%%%%%% 曰 %%%%%%%%%%
\subsection*{曰}\addcontentsline{loh}{figure}{曰 \dpy{yue1}}

\begin{EntryWithPhonetic}{曰}{yue1}{4}{⽈}[HSK 7-9]
  \definition{v.}{dizer | chamar; nomear | falar}
\end{EntryWithPhonetic}

%%%%%%%%%% 约 %%%%%%%%%%
\subsection*{约}\addcontentsline{loh}{figure}{约 \dpy{yue1}}

\begin{EntryWithPhonetic}{约}{yue1}{6}{⽷}[HSK 3]
  \definition*{s.}{Sobrenome: Yue}
  \definition{adj.}{econômico; frugal | simples; breve; resumido | indistinto; confuso}
  \definition{adv.}{cerca de; ao redor; aproximadamente}
  \definition{s.}{pacto; acordo; nomeação; o que foi combinado}
  \definition{v.}{combinar; propor ou discutir antecipadamente (o que deve ser respeitado por todos) | convidar com antecedência | restringir; conter | reduzir (fração aproximada)}
  \seeref{yao1}
\end{EntryWithPhonetic}

\begin{EntryWithPhonetic}{约定}{yue1ding4}{6,8}{⽷,⼧}[HSK 6]
  \definition{s.}{acordo; compromisso}
  \definition{v.}{determinar; chegar a um acordo; concordar com}
\end{EntryWithPhonetic}

\begin{EntryWithPhonetic}{约定俗成}{yue1ding4-su2cheng2}{6,8,9,6}{⽷,⼧,⼈,⼽}[HSK 7-9]
  \definition{expr.}{(nome ou hábito social) estabelecido (ou sancionado) pelo uso popular; aceito pela prática comum; convenção costumeira; estabelecido pelo uso popular; acordo de uso comum}
\end{EntryWithPhonetic}

\begin{EntryWithPhonetic}{约会}{yue1hui5}{6,6}{⽷,⼈}[HSK 4]
  \definition[个,次]{s.}{data; compromisso; engajamento; reunião pré-agendada}
  \definition{v.}{marcar uma reunião; marcar um encontro}
\end{EntryWithPhonetic}

\begin{EntryWithPhonetic}{约束}{yue1shu4}{6,7}{⽷,⽊}[HSK 5]
  \definition{adj.}{amarrado}
  \definition{s.}{restrição; constrangimento; engajamento}
  \definition{v.}{amarrar; prender; reprimir; restringir; manter dentro de si}
\end{EntryWithPhonetic}

%%%%%%%%%% 月 %%%%%%%%%%
\subsection*{月}\addcontentsline{loh}{figure}{月 \dpy{yue4}}

\begin{EntryWithPhonetic}{月}{yue4}{4}{⽉}[HSK 1][Kangxi 74]
  \definition*{s.}{Sobrenome: Yue}
  \definition[个]{s.}{mês | por mês | lua | redondo; em forma de lua cheia}
\end{EntryWithPhonetic}

\begin{EntryWithPhonetic}{月饼}{yue4bing5}{4,9}{⽉,⾷}[HSK 5]
  \definition[个,块,盒,筒]{s.}{bolinho da lua; comida típica do Festival do Meio Outono; redonda e recheada; simboliza a reunião familiar}
\end{EntryWithPhonetic}

\begin{EntryWithPhonetic}{月初}{yue4chu1}{4,7}{⽉,⾐}[HSK 7-9]
  \definition[个]{s.}{início do mês; começo do mês; os primeiros dias do mês}
  \antonymref{月底}{yue4di3}
\end{EntryWithPhonetic}

\begin{EntryWithPhonetic}{月底}{yue4di3}{4,8}{⽉,⼴}[HSK 4]
  \definition[个]{s.}{final do mês; últimos dias do mês}
\end{EntryWithPhonetic}

\begin{EntryWithPhonetic}{月份}{yue4fen4}{4,6}{⽉,⼈}[HSK 2]
  \definition[个]{s.}{mês; refere"-se a um determinado mês}
\end{EntryWithPhonetic}

\begin{EntryWithPhonetic}{月径}{yue4jing4}{4,8}{⽉,⼻}
  \definition{s.}{diâmetro da lua | diâmetro da órbita da lua | caminho iluminado pela lua}
\end{EntryWithPhonetic}

\begin{EntryWithPhonetic}{月亮}{yue4liang5}{4,9}{⽉,⼇}[HSK 2]
  \definition[个,轮,挂,快,页]{s.}{Lua; Lua é o nome comum do satélite da Terra}
\end{EntryWithPhonetic}

\begin{EntryWithPhonetic}{月票}{yue4piao4}{4,11}{⽉,⽰}[HSK 7-9]
  \definition{s.}{bilhete mensal; passe mensal; (britânico) bilhete de transporte; (inglês) passe de temporada; bilhetes mensais para uso em ônibus, bondes, parques, etc.}
\end{EntryWithPhonetic}

\begin{EntryWithPhonetic}{月球}{yue4qiu2}{4,11}{⽉,⽟}[HSK 5]
  \definition*[颗,个]{s.}{Lua}
  \definition{pref.}{seleno-; seleni-}
\end{EntryWithPhonetic}

\begin{EntryWithPhonetic}{月壤}{yue4rang3}{4,20}{⽉,⼟}
  \definition{s.}{solo lunar}
\end{EntryWithPhonetic}

\begin{EntryWithPhonetic}{月相}{yue4xiang4}{4,9}{⽉,⽬}
  \definition{s.}{fases da lua, a saber: lua nova 朔, lua crescente 上弦, lua cheia 望 e lua minguante 下弦}
\end{EntryWithPhonetic}

\begin{EntryWithPhonetic}{月月}{yue4yue4}{4,4}{⽉,⽉}
  \definition{adv.}{todo mês}
\end{EntryWithPhonetic}

%%%%%%%%%% 乐 %%%%%%%%%%
\subsection*{乐}\addcontentsline{loh}{figure}{乐 \dpy{yue4}}

\begin{EntryWithPhonetic}{乐}{yue4}{5}{⼃}
  \definition*{s.}{Sobrenome: Yue}
  \definition{s.}{música}
  \seeref{le4}
\end{EntryWithPhonetic}

\begin{EntryWithPhonetic}{乐队}{yue4dui4}{5,4}{⼃,⾩}[HSK 3]
  \definition[支,个]{s.}{orquestra; banda; um grupo composto por muitas pessoas que tocam diferentes instrumentos musicais}
  \synonymref{歌曲}{ge1qu3}
  \synonymref{音乐}{yin1yue4}
\end{EntryWithPhonetic}

\begin{EntryWithPhonetic}{乐器}{yue4qi4}{5,16}{⼃,⼝}[HSK 7-9]
  \definition[件,种,个]{s.}{instrumento musical; instrumentos que podem produzir sons musicais e são usados para tocar música}
\end{EntryWithPhonetic}

\begin{EntryWithPhonetic}{乐曲}{yue4qu3}{5,6}{⼃,⽈}[HSK 6]
  \definition[支,首,段]{s.}{música; composição musical}
  \synonymref{歌曲}{ge1qu3}
  \synonymref{音乐}{yin1yue4}
\end{EntryWithPhonetic}

%%%%%%%%%% 岳 %%%%%%%%%%
\subsection*{岳}\addcontentsline{loh}{figure}{岳 \dpy{yue4}}

\begin{EntryWithPhonetic}{岳}{yue4}{8}{⼭}
  \definition*{s.}{As Cinco Grandes Montanhas; as cinco montanhas mais importantes na cultura tradicional chinesa simbolizam a unidade e a santidade nacional, e são frequentemente usadas em contextos religiosos, de sacrifício e culturais | Sobrenome: Yue}
  \definition{s.}{montanha alta | pais da esposa | sogro/sogra; as formas de tratamento para os pais ou tios da esposa}
\end{EntryWithPhonetic}

\begin{EntryWithPhonetic}{岳父}{yue4fu4}{8,4}{⼭,⽗}[HSK 7-9]
  \definition{s.}{pai da esposa; sogro}
\end{EntryWithPhonetic}

\begin{EntryWithPhonetic}{岳母}{yue4mu3}{8,5}{⼭,⽏}[HSK 7-9]
  \definition[位,名,个]{s.}{mãe da esposa, sogra}
\end{EntryWithPhonetic}

%%%%%%%%%% 说 %%%%%%%%%%
\subsection*{说}\addcontentsline{loh}{figure}{说 \dpy{yue4}}

\begin{EntryWithPhonetic}{说}{yue4}{9}{⾔}
  \variantof{悦}
  \seeref{shui4}
  \seeref{shuo1}
\end{EntryWithPhonetic}

%%%%%%%%%% 悦 %%%%%%%%%%
\subsection*{悦}\addcontentsline{loh}{figure}{悦 \dpy{yue4}}

\begin{EntryWithPhonetic}{悦}{yue4}{10}{⼼}
  \definition*{s.}{Sobrenome: Yue}
  \definition{adj.}{Linterário: feliz; satisfeito; encantado}
  \definition{v.}{agradar; deliciar}
\end{EntryWithPhonetic}

\begin{EntryWithPhonetic}{悦耳}{yue4'er3}{10,6}{⼼,⽿}[HSK 7-9]
  \definition{adj.}{agradável ao ouvido; de som doce; melodioso}
  \synonymref{动听}{dong4ting1}
  \synonymref{好听}{hao3ting1}
  \synonymref{顺耳}{shun4'er3}
  \antonymref{刺耳}{ci4'er3}
  \antonymref{难听}{nan2ting1}
  \antonymref{跑调}{pao3diao4}
\end{EntryWithPhonetic}

%%%%%%%%%% 阅 %%%%%%%%%%
\subsection*{阅}\addcontentsline{loh}{figure}{阅 \dpy{yue4}}

\begin{EntryWithPhonetic}{阅}{yue4}{10}{⾨}
  \definition{v.}{ler; repassar; examinar | revisar; inspecionar | experimentar; passar por}
\end{EntryWithPhonetic}

\begin{EntryWithPhonetic}{阅兵式}{yue4bing1shi4}{10,7,6}{⾨,⼋,⼷}
  \definition{s.}{parada militar; desfile militar}
\end{EntryWithPhonetic}

\begin{EntryWithPhonetic}{阅读}{yue4du2}{10,10}{⾨,⾔}[HSK 4]
  \definition{v.}{ler; examinar; olhar (livros, jornais, etc.) e entender seu conteúdo}
\end{EntryWithPhonetic}

\begin{EntryWithPhonetic}{阅读广度}{yue4du2guang3du4}{10,10,3,9}{⾨,⾔,⼴,⼴}
  \definition{s.}{intervalo de leitura}
\end{EntryWithPhonetic}

\begin{EntryWithPhonetic}{阅读理解}{yue4du2li3jie3}{10,10,11,13}{⾨,⾔,⽟,⾓}
  \definition{s.}{compreensão de leitura}
\end{EntryWithPhonetic}

\begin{EntryWithPhonetic}{阅读器}{yue4du2qi4}{10,10,16}{⾨,⾔,⼝}
  \definition{s.}{leitor (\emph{software})}
\end{EntryWithPhonetic}

\begin{EntryWithPhonetic}{阅读时间}{yue4 du2 shi2 jian1}{10,10,7,7}{⾨,⾔,⽇,⾨}
  \definition{s.}{tempo de leitura}
\end{EntryWithPhonetic}

\begin{EntryWithPhonetic}{阅读障碍}{yue4du2zhang4ai4}{10,10,13,13}{⾨,⾔,⾩,⽯}
  \definition{s.}{dislexia}
\end{EntryWithPhonetic}

\begin{EntryWithPhonetic}{阅读装置}{yue4du2zhuang1zhi4}{10,10,12,13}{⾨,⾔,⾐,⽹}
  \definition{s.}{dispositivo de leitura (por exemplo, para códigos de barras, etiquetas RFID, etc.)}
\end{EntryWithPhonetic}

\begin{EntryWithPhonetic}{阅览室}{yue4lan3shi4}{10,9,9}{⾨,⾒,⼧}[HSK 5]
  \definition[个,间]{s.}{sala de leitura; a biblioteca dispõe de salas para leitura e pesquisa, equipadas com mesas e cadeiras adequadas, livros, jornais, revistas, etc.}
\end{EntryWithPhonetic}

\begin{EntryWithPhonetic}{阅历}{yue4li4}{10,4}{⾨,⼚}[HSK 7-9]
  \definition{s.}{experiência; conhecimento e experiência adquiridos através de experiências de vida}
  \definition{v.}{ver, ouvir ou fazer por si mesmo}
  \synonymref{经历}{jing1li4}
  \synonymref{经验}{jing1yan4}
  \synonymref{体验}{ti3yan4}
  \synonymref{眼界}{yan3jie4}
\end{EntryWithPhonetic}

%%%%%%%%%% 粤 %%%%%%%%%%
\subsection*{粤}\addcontentsline{loh}{figure}{粤 \dpy{yue4}}

\begin{EntryWithPhonetic}{粤}{yue4}{12}{⾔}
  \definition*{s.}{Outro nome para a Província de Guangdong, 广东}
  \seealsoref{广东}{guang3dong1}
\end{EntryWithPhonetic}

\begin{EntryWithPhonetic}{粤语}{yue4yu3}{12,9}{⾔,⾔}
  \definition{s.}{cantonês | língua cantonesa}
\end{EntryWithPhonetic}

%%%%%%%%%% 越 %%%%%%%%%%
\subsection*{越}\addcontentsline{loh}{figure}{越 \dpy{yue4}}

\begin{EntryWithPhonetic}{越}{yue4}{12}{⾛}[HSK 2]
  \definition{adj.}{superior; excede ou ultrapassa o ordinário}
  \definition{adv.}{quanto mais\dots mais; usados juntos, eles formam o formato de ``越…越…'' para indicar que o grau de uma situação se torna mais sério à medida que se desenvolve; ``成年…'' para indicar que o grau de uma situação se torna mais sério à medida que o tempo passa}
  \definition{v.}{passar por cima; pular; cruzar | exceder; ultrapassar | estar em um tom alto; estar animado | saquear; pilhar; expoliar; apreender; roubar | passar; passar através; atravessar}
  \seealsoref{越……越……}{yue4 yue4}
  \seealsoref{越来越……}{yue4lai2yue4}
\end{EntryWithPhonetic}

\begin{EntryWithPhonetic}{越发}{yue4fa1}{12,5}{⾛,⼜}[HSK 7-9]
  \definition{adv.}{ainda mais; tanto mais; cada vez mais; isso indica um grau adicional, equivalente a 更加 | quanto mais\dots quantomais\dots; ecoa o 越 ou 或 anterior e sua função é a mesma de 越……越…… (usado quando duas ou mais orações se repetem)}
  \seealsoref{更加}{geng4jia1}
  \seealsoref{或}{huo4}
  \seealsoref{越}{yue4}
  \seealsoref{越……越……}{yue4 yue4}
  \synonymref{更加}{geng4jia1}
  \synonymref{尤其}{you2qi2}
  \antonymref{依旧}{yi1jiu4}
\end{EntryWithPhonetic}

\begin{EntryWithPhonetic}{越过}{yue4/guo4}{12,6}{⾛,⾡}[HSK 7-9]
  \definition{v.+compl.}{cruzar; superar; negociar; transpor obstáculos, fronteiras, etc. | exceder; ultrapassar}
  \synonymref{超出}{chao1 chu1}
  \synonymref{超过}{chao1/guo4}
  \synonymref{超越}{chao1yue4}
  \synonymref{穿过}{chuan1guo4}
  \synonymref{跨越}{kua4yue4}
  \synonymref{突出}{tu1/chu1}
\end{EntryWithPhonetic}

\begin{EntryWithPhonetic}{越境}{yue4jing4}{12,14}{⾛,⼟}
  \definition{v.}{cruzar a fronteira ilegalmente; entrar ou sair clandestinamente de um país}
\end{EntryWithPhonetic}

\begin{EntryWithPhonetic}{越来越……}{yue4lai2yue4}{12,7,12}{⾛,⽊,⾛}[HSK 2]
  \definition{adv.}{cada vez mais\dots; isso significa que o grau de algo se aprofunda à medida que o tempo passa}
\end{EntryWithPhonetic}

\begin{EntryWithPhonetic}{越……越……}{yue4 yue4}{12,12}{⾛,⾛}
  \definition{expr.}{quanto mais\dots tanto mais\dots}
\end{EntryWithPhonetic}

\begin{EntryWithPhonetic}{越障}{yue4zhang4}{12,13}{⾛,⾩}
  \definition{s.}{curso com obstáculos para treinamento de tropas}
  \definition{v.}{superar obstáculos}
\end{EntryWithPhonetic}

%%%%%%%%%% 龠 %%%%%%%%%%
\subsection*{龠}\addcontentsline{loh}{figure}{龠 \dpy{yue4}}

\begin{EntryWithPhonetic}{龠}{yue4}{17}{⿕}[Kangxi 214]
  \definition{clas.}{yue, uma unidade de medida seca para grãos (= 0,5 decilitro);}
  \definition{s.}{uma flauta curta antiga}
\end{EntryWithPhonetic}

%%%%%%%%%% 晕 %%%%%%%%%%
\subsection*{晕}\addcontentsline{loh}{figure}{晕 \dpy{yun1}}

\begin{EntryWithPhonetic}{晕}{yun1}{10}{⽇}[HSK 6]
  \definition{adj.}{tonto; vertiginoso; confuso; sensação de que as coisas estão girando ao seu redor e, às vezes, sensação de que você vai cair}
  \definition{v.}{desmaiar; desfalecer}
  \seeref{yun4}
\end{EntryWithPhonetic}

\begin{EntryWithPhonetic}{晕倒}{yun1dao3}{10,10}{⽇,⼈}[HSK 7-9]
  \definition{v.}{desmaiar; perder a consciência; ficar inconsciente; desmaiar; cair em um estado de desmaio; perder a consciência}
  \antonymref{苏醒}{su1xing3}
\end{EntryWithPhonetic}

%%%%%%%%%% 云 %%%%%%%%%%
\subsection*{云}\addcontentsline{loh}{figure}{云 \dpy{yun2}}

\begin{EntryWithPhonetic}{云}{yun2}{4}{⼆}[HSK 2]
  \definition*{s.}{Província de Yunnan, abreviação de 云南 | Sobrenome: Yun}
  \definition[片,朵]{s.}{nuvem}
  \definition{v.}{dizer}
  \seealsoref{云南}{yun2nan2}
\end{EntryWithPhonetic}

\begin{EntryWithPhonetic}{云端}{yun2duan1}{4,14}{⼆,⽴}
  \definition{s.}{alto nas nuvens | (computação) a nuvem}
\end{EntryWithPhonetic}

\begin{EntryWithPhonetic}{云南}{yun2nan2}{4,9}{⼆,⼗}
  \definition*{s.}{Província de Yunnan}
\end{EntryWithPhonetic}

\begin{EntryWithPhonetic}{云云}{yun2yun2}{4,4}{⼆,⼆}
  \definition{adv.}{e assim por diante | assim e assim}
\end{EntryWithPhonetic}

%%%%%%%%%% 允 %%%%%%%%%%
\subsection*{允}\addcontentsline{loh}{figure}{允 \dpy{yun3}}

\begin{EntryWithPhonetic}{允}{yun3}{4}{⼉}
  \definition*{s.}{Sobrenome: Yun}
  \definition{adj.}{justo; imparcial}
  \definition{v.}{permitir; deixar; consentir}
\end{EntryWithPhonetic}

\begin{EntryWithPhonetic}{允许}{yun3xu3}{4,6}{⼉,⾔}[HSK 6]
  \definition{s.}{permitido; permissão}
  \definition{v.}{permitir; deixar; concordar com alguém para fazer algo}
\end{EntryWithPhonetic}

%%%%%%%%%% 陨 %%%%%%%%%%
\subsection*{陨}\addcontentsline{loh}{figure}{陨 \dpy{yun3}}

\begin{EntryWithPhonetic}{陨}{yun3}{9}{⾩}
  \definition{v.}{cair do céu ou do espaço sideral}
\end{EntryWithPhonetic}

\begin{EntryWithPhonetic}{陨石}{yun3shi2}{9,5}{⾩,⽯}[HSK 7-9]
  \definition[块,颗]{s.}{Astronomia: aerólito; meteorito; meteorito rochoso}
  \synonymref{流星}{liu2xing1}
\end{EntryWithPhonetic}

%%%%%%%%%% 孕 %%%%%%%%%%
\subsection*{孕}\addcontentsline{loh}{figure}{孕 \dpy{yun4}}

\begin{EntryWithPhonetic}{孕}{yun4}{5}{⼦}
  \definition[次]{s.}{gravidez}
  \definition{v.}{engravidar | estar grávida}
\end{EntryWithPhonetic}

\begin{EntryWithPhonetic}{孕妇}{yun4fu4}{5,6}{⼦,⼥}[HSK 7-9]
  \definition{s.}{mulher grávida; futura mãe}
\end{EntryWithPhonetic}

\begin{EntryWithPhonetic}{孕育}{yun4yu4}{5,8}{⼦,⾁}[HSK 7-9]
  \definition{v.}{procriar; dar à luz; estar grávida de; gravidez formal e o desenvolvimento de uma criança no corpo | procriar; uma metáfora para criar algo novo a partir de algo que já existe}
  \synonymref{产生}{chan3sheng1}
  \synonymref{出现}{chu1xian4}
  \antonymref{毁灭}{hui3mie4}
  \antonymref{破坏}{po4huai4}
  \antonymref{伤害}{shang1hai4}
\end{EntryWithPhonetic}

%%%%%%%%%% 运 %%%%%%%%%%
\subsection*{运}\addcontentsline{loh}{figure}{运 \dpy{yun4}}

\begin{EntryWithPhonetic}{运}{yun4}{7}{⾡}[HSK 5]
  \definition*{s.}{Sobrenome: Yun}
  \definition{s.}{sorte; destino; fortuna}
  \definition{v.}{mover; deslocar | transportar; levar | usar; empunhar; utilizar}
\end{EntryWithPhonetic}

\begin{EntryWithPhonetic}{运动病}{yun4dong4bing4}{7,6,10}{⾡,⼒,⽧}
  \definition{s.}{enjôo (movimento, carro, etc.)}
\end{EntryWithPhonetic}

\begin{EntryWithPhonetic}{运动场}{yun4dong4chang3}{7,6,6}{⾡,⼒,⼟}
  \definition{s.}{campo desportivo | campo de jogos}
\end{EntryWithPhonetic}

\begin{EntryWithPhonetic}{运动服}{yun4dong4fu2}{7,6,8}{⾡,⼒,⽉}
  \definition{s.}{roupa para prática de esporte}
\end{EntryWithPhonetic}

\begin{EntryWithPhonetic}{运动会}{yun4dong4hui4}{7,6,6}{⾡,⼒,⼈}[HSK 4]
  \definition[届,场,次,个]{s.}{jogos; encontro esportivo; dia de esportes; encontro atlético; competição esportiva abrangente}
\end{EntryWithPhonetic}

\begin{EntryWithPhonetic}{运动家}{yun4dong4jia1}{7,6,10}{⾡,⼒,⼧}
  \definition{s.}{ativista | atleta | esportista}
\end{EntryWithPhonetic}

\begin{EntryWithPhonetic}{运动衫}{yun4dong4shan1}{7,6,8}{⾡,⼒,⾐}
  \definition[件]{s.}{moletom | camisa esportiva}
\end{EntryWithPhonetic}

\begin{EntryWithPhonetic}{运动鞋}{yun4dong4xie2}{7,6,15}{⾡,⼒,⾰}
  \definition{s.}{tênis | sapatos esportivos}
\end{EntryWithPhonetic}

\begin{EntryWithPhonetic}{运动学}{yun4dong4xue2}{7,6,8}{⾡,⼒,⼦}
  \definition{s.}{cinemática; um ramo da ciência do esporte que usa a anatomia e a mecânica humanas para explicar várias atividades esportivas}
\end{EntryWithPhonetic}

\begin{EntryWithPhonetic}{运动员}{yun4dong4yuan2}{7,6,7}{⾡,⼒,⼝}[HSK 4]
  \definition[名,个,班]{s.}{jogador; atleta; esportista; pessoas que participam de competições esportivas}
\end{EntryWithPhonetic}

\begin{EntryWithPhonetic}{运动}{yun4dong5}{7,6}{⾡,⼒}[HSK 2]
  \definition[项,种,场,次]{s.}{esportes; atletismo; exercício; atividades esportivas | movimento; campanha (política); atividades de massa organizadas, intencionais e de alto nível na política, cultura, produção, etc. | movimento; refere"-se a todas as mudanças}
  \definition{v.}{exercitar; fazer atividade física | mover"-se; refere"-se à mudança na posição de um objeto}
\end{EntryWithPhonetic}

\begin{EntryWithPhonetic}{运河}{yun4he2}{7,8}{⾡,⽔}[HSK 7-9]
  \definition[条]{s.}{canal (em um rio)}
\end{EntryWithPhonetic}

\begin{EntryWithPhonetic}{运气}{yun4/qi4}{7,4}{⾡,⽓}
  \definition{v.+compl.}{tentar a sorte | concentrar a energia em uma parte do corpo}[他们地运一口气。===Eles respiraram fundo.]
  \seeref{yun4qi5}
\end{EntryWithPhonetic}

\begin{EntryWithPhonetic}{运气}{yun4qi5}{7,4}{⾡,⽓}[HSK 4]
  \definition{adj.}{sortudo; afortunado}
  \definition{s.}{sorte; fortuna}
  \seeref{yun4/qi4}
\end{EntryWithPhonetic}

\begin{EntryWithPhonetic}{运输}{yun4shu1}{7,13}{⾡,⾞}[HSK 3]
  \definition{v.}{enviar; transportar; transportar pessoas ou coisas de um lugar para outro usando carros, barcos, aviões, etc.}
\end{EntryWithPhonetic}

\begin{EntryWithPhonetic}{运送}{yun4song4}{7,9}{⾡,⾡}[HSK 7-9]
  \definition{v.}{transportar; enviar; conduzir; transportar pessoas ou suprimentos para outros locais}
  \synonymref{输送}{shu1song4}
  \synonymref{运输}{yun4shu1}
\end{EntryWithPhonetic}

\begin{EntryWithPhonetic}{运行}{yun4xing2}{7,6}{⾡,⾏}[HSK 5]
  \definition{v.}{correr; mover; trabalhar; estar em movimento; (veículo, nave, planeta, etc.) mover"-se em um ciclo repetitivo; avançar de maneira regular e direcional}
\end{EntryWithPhonetic}

\begin{EntryWithPhonetic}{运营}{yun4ying2}{7,11}{⾡,⾋}[HSK 7-9]
  \definition{v.}{executar; operar; (veículos, navios, etc.) operação e negócios | operar; a metáfora descreve uma organização que realiza seu trabalho de maneira organizada}
  \synonymref{筹划}{chou2hua4}
  \synonymref{规划}{gui1hua4}
  \synonymref{经营}{jing1ying2}
  \synonymref{运作}{yun4zuo4}
\end{EntryWithPhonetic}

\begin{EntryWithPhonetic}{运用}{yun4yong4}{7,5}{⾡,⽤}[HSK 4]
  \definition{v.}{usar; utilizar; manejar; aplicar; explorar as coisas de acordo com suas características}
\end{EntryWithPhonetic}

\begin{EntryWithPhonetic}{运转}{yun4zhuan3}{7,8}{⾡,⾞}[HSK 7-9]
  \definition{v.}{girar; dar voltas; mover"-se ao longo de uma determinada trilha | (máquinas) funcionar; operar; refere"-se à rotação da máquina | (organização) trabalhar; operar; metaforicamente, refere"-se a (uma instituição, organização, etc.) que exerce poder e realiza trabalho}
  \synonymref{运行}{yun4xing2}
  \synonymref{运作}{yun4zuo4}
  \antonymref{停止}{ting2zhi3}
\end{EntryWithPhonetic}

\begin{EntryWithPhonetic}{运作}{yun4zuo4}{7,7}{⾡,⼈}[HSK 6]
  \definition{v.}{trabalhar; operar; (uma instituição, organização, etc.) realizar trabalho; realizar atividades}
\end{EntryWithPhonetic}

%%%%%%%%%% 晕 %%%%%%%%%%
\subsection*{晕}\addcontentsline{loh}{figure}{晕 \dpy{yun4}}

\begin{EntryWithPhonetic}{晕}{yun4}{10}{⽇}
  \definition{s.}{auréola; o círculo de luz formado pela refração da luz solar ou do luar através dos cristais de gelo nas nuvens | halo em torno de alguma cor ou luz; áreas desfocadas em torno de luz, sombra e cor}
  \definition{v.}{ficar tonto; desmaiar; desfalecer; sensação de tontura, como se os objetos ao seu redor estivessem girando e como se você estivesse prestes a cair}
  \seeref{yun1}
\end{EntryWithPhonetic}

\begin{EntryWithPhonetic}{晕车}{yun4 che1}{10,4}{⽇,⾞}[HSK 6]
  \definition{v.}{ter enjoo no carro; ter tontura e vômito ao andar de carro}
\end{EntryWithPhonetic}

%%%%%%%%%% 酝 %%%%%%%%%%
\subsection*{酝}\addcontentsline{loh}{figure}{酝 \dpy{yun4}}

\begin{EntryWithPhonetic}{酝}{yun4}{11}{⾣}
  \definition{s.}{vinho}
  \definition{v.}{fermentar (vinho); fabricar (cerveja) | fazer vinho}
\end{EntryWithPhonetic}

\begin{EntryWithPhonetic}{酝酿}{yun4niang4}{11,14}{⾣,⾣}[HSK 7-9]
  \definition{v.}{preparar; fermentar; fermentação das matérias"-primas durante a vinificação | preparar; planejar; fazer; essa metáfora descreve um processo no qual uma tarefa é empreendida após um longo período de consideração, discussão, deliberação e preparação, permitindo que ela amadureça gradualmente}
  \synonymref{筹划}{chou2hua4}
  \synonymref{讨论}{tao3lun4}
  \synonymref{预谋}{yu4mou2}
  \synonymref{蕴藏}{yun4cang2}
\end{EntryWithPhonetic}

%%%%%%%%%% 韵 %%%%%%%%%%
\subsection*{韵}\addcontentsline{loh}{figure}{韵 \dpy{yun4}}

\begin{EntryWithPhonetic}{韵}{yun4}{13}{⾳}
  \definition*{s.}{Sobrenome: Yun}
  \definition{s.}{som musical (ou agradável, prazeroso); uma voz agradável | rima | charme}
\end{EntryWithPhonetic}

\begin{EntryWithPhonetic}{韵味}{yun4wei4}{13,8}{⾳,⼝}[HSK 7-9]
  \definition{s.}{riqueza implícita; significado implícito no som e no ritmo; o significado transmitido por sons e rimas | charme; atratividade; interesse}
  \synonymref{风味}{feng1wei4}
  \synonymref{诗意}{shi1yi4}
\end{EntryWithPhonetic}

%%%%%%%%%% 蕴 %%%%%%%%%%
\subsection*{蕴}\addcontentsline{loh}{figure}{蕴 \dpy{yun4}}

\begin{EntryWithPhonetic}{蕴}{yun4}{15}{⾋}
  \definition{s.}{Literário: profundidade | Literário: força interior}
  \definition{v.}{Literário: acumular; coletar; manter em estoque; conter}
\end{EntryWithPhonetic}

\begin{EntryWithPhonetic}{蕴藏}{yun4cang2}{15,17}{⾋,⾋}[HSK 7-9]
  \definition{v.}{conter; manter em estoque; acumular internamente}
  \synonymref{蕴涵}{yun4han2}
  \synonymref{蕴含}{yun4han2}
  \synonymref{酝酿}{yun4niang4}
  \antonymref{展现}{zhan3xian4}
\end{EntryWithPhonetic}

\begin{EntryWithPhonetic}{蕴含}{yun4han2}{15,7}{⾋,⼝}
  \definition{v.}{conter; incluir | acumular}
  \variantof{蕴涵}
  \synonymref{包含}{bao1han2}
  \synonymref{蕴藏}{yun4cang2}
\end{EntryWithPhonetic}

\begin{EntryWithPhonetic}{蕴涵}{yun4han2}{15,11}{⾋,⽔}[HSK 7-9]
  \definition{s.}{Lógica: implicação | vínculo | condição implícita}
  \definition{v.}{Literário: conter | Lógica: implicar; envolver | acumular | abraçar}
\end{EntryWithPhonetic}

%%%%% EOF %%%%%

